\documentclass[11pt]{article}
\ifdefined\pdfgentounicode
\pdfgentounicode=1
\input{glyphtounicode-cmex}
\fi
\usepackage[T1]{fontenc}
\usepackage[utf8]{inputenc}
\usepackage{lmodern}
\usepackage[a4paper,margin=28mm]{geometry}
\usepackage{amsmath,amssymb,amsthm,mathtools}
\usepackage{microtype}
\usepackage{enumitem}
\usepackage{booktabs,longtable,array}
\usepackage{tikz}
\usetikzlibrary{arrows.meta,positioning}
\usepackage{xcolor}
\usepackage[numbers,sort&compress]{natbib}
\usepackage{hyperref}
\hypersetup{colorlinks=true,linkcolor=blue!45!black,citecolor=green!35!black,urlcolor=blue!45!black,
  pdftitle={Uniform Cartier sections for Fano type contractions},pdfauthor={OpenAI}}
\setlength{\emergencystretch}{2em}
\setlist{topsep=5pt,itemsep=3pt,parsep=0pt}
\numberwithin{equation}{section}
\renewcommand{\thesection}{\Alph{section}}
\newtheorem{maintheorem}{Theorem}
\newtheorem{maincorollary}[maintheorem]{Corollary}
\newtheorem{theorem}{Theorem}[section]
\newtheorem{lemma}[theorem]{Lemma}
\newtheorem{proposition}[theorem]{Proposition}
\newtheorem{corollary}[theorem]{Corollary}
\newtheorem{claim}[theorem]{Claim}
\theoremstyle{definition}
\newtheorem{definition}[theorem]{Definition}
\newtheorem{convention}[theorem]{Convention}
\newtheorem{notation}[theorem]{Notation}
\theoremstyle{remark}
\newtheorem{remark}[theorem]{Remark}
\newcommand{\C}{\mathbb C}
\newcommand{\Q}{\mathbb Q}
\newcommand{\R}{\mathbb R}
\newcommand{\Z}{\mathbb Z}
\newcommand{\N}{\mathbb N}
\DeclareMathOperator{\Spec}{Spec}
\DeclareMathOperator{\Proj}{Proj}
\DeclareMathOperator{\ord}{ord}
\DeclareMathOperator{\vol}{vol}
\DeclareMathOperator{\Supp}{Supp}
\DeclareMathOperator{\divisor}{div}
\DeclareMathOperator{\trdeg}{trdeg}
\DeclareMathOperator{\Frac}{Frac}
\DeclareMathOperator{\Cl}{Cl}
\DeclareMathOperator{\Pic}{Pic}
\title{Uniform Cartier sections for Fano type contractions}
\author{OpenAI}
\date{September 25, 2026}
\begin{document}
\maketitle
\begin{abstract}
We prove the Cartier-divisor conjecture of Birkar and Shokurov for rational boundaries in characteristic zero and for real boundaries over $\C$. For an $\epsilon$-lc Fano type contraction with positive-dimensional base and nef negative log canonical class, a Cartier divisor through any prescribed base point can be chosen with pullback log canonical threshold bounded below in terms of the dimension and $\epsilon$ alone.
\end{abstract}
\section*{Introduction}
\addcontentsline{toc}{section}{Introduction}

Through a prescribed point of a singular variety, can one choose a Cartier hypersurface whose log canonical threshold is bounded away from zero? For a Fano type contraction $f:X\to Z$, the corresponding question asks for a hypersurface on $Z$ whose pullback has controlled singularities along the entire fibre. The hypersurface may depend on the contraction and the point. The desired lower bound depends only on the dimension of $X$ and a positive lower bound for its log discrepancies.

We work over an algebraically closed field of characteristic zero. A \emph{contraction} is a projective surjective morphism of normal integral varieties with connected fibres, expressed by $f_*\mathcal O_X=\mathcal O_Z$. It is of \emph{Fano type} if an effective rational boundary $C$ exists such that $(X,C)$ is klt and $-(K_X+C)$ is ample over $Z$. This auxiliary boundary need not equal the boundary in the following theorem.

A rational pair $(X,B)$ has effective rational boundary $B$ and rationally Cartier log canonical divisor $K_X+B$. If $p:W\to X$ is a smooth birational model and $E$ is a prime divisor on $W$, our discrepancy convention is
\[
a(E;X,B)=1+\operatorname{coeff}_E\bigl(K_W-p^*(K_X+B)\bigr).
\]
Thus $\epsilon$-lc means $a(E;X,B)\geq\epsilon$ for every such $E$; lc means that they are nonnegative; and klt means that they are strictly positive.

The obstruction is visible in a single ratio. If $D$ has local equation $g$ on the base, log canonicity after adding $t f^*D$ requires
\[
 t\le\frac{a(E;X,B)}{\ord_E(f^*g)}
\]
for every prime divisor $E$ over the relevant inverse image with positive denominator. A discrepancy bound alone does not control those orders. Even on a smooth curve, the divisor $m[z]$ has threshold $1/m$. The problem is therefore to choose an equation through $z$ for which all these inequalities hold with one uniform positive $t$.

\begin{samepage}
\begin{maintheorem}[Uniform Cartier sections]\label{thm:main}
For every integer $d>0$ and every real number $\epsilon>0$, there is a real number $\tau=\tau(d,\epsilon)>0$ with the following property. Let $(X,B)$ be an $\epsilon$-lc rational pair of dimension $d$, and let $f:X\to Z$ be a Fano type contraction of normal integral quasi-projective varieties with $\dim Z>0$. Suppose that $-(K_X+B)$ is nef over $Z$. For every closed point $z\in Z$, there are an open neighbourhood $U$ of $z$ and a nonzero effective Cartier divisor $D$ on $U$, with $z\in\Supp D$, such that
\[
\bigl(f^{-1}(U),B|_{f^{-1}(U)}+\tau f^*D\bigr)
\quad\text{is log canonical.}
\]
\end{maintheorem}
\end{samepage}

The conclusion of Theorem~\ref{thm:main} concerns the whole inverse image of a neighbourhood of $z$. The constant is independent of Cartier indices, boundary denominators, and any choice of polarization. It is non-effective; the divisor and the neighbourhood may depend on the contraction and the point.

\paragraph{Cartier sections and complements.}
For the identity contraction, the question above concerns arbitrary $\epsilon$-lc singularities. Equivalently, it asks for a uniform positive lower bound for the log canonical threshold of the maximal ideal at a closed point; the passage to a general element of that ideal is recalled in Lemma~\ref{lem:H-field-reduction}. The general contraction requires the corresponding bound for the pulled-back point ideal. This is a different local question from bounding the discriminant coefficients in the canonical bundle formula: those coefficients are defined by thresholds over generic points of divisors on birational models of the base. A Cartier section through a prescribed point must control the whole fibre at once.

The complement formulation asks for a uniformly bounded complement index. Over a neighbourhood of a base point, a \emph{monotonic klt $n$-complement} of $(X,B)$ is an effective rational divisor $\Lambda\ge B$ such that $(X,\Lambda)$ is klt and
\[
 n(K_X+\Lambda)\sim 0.
\]
Here linear equivalence includes the assertion that the displayed divisor is Cartier. In Shokurov's boundedness conjecture, as formulated in \cite[Conjecture~1.1]{Chen2025}, the index $n$ depends only on the total dimension $d$, the discrepancy bound $\epsilon>0$, and a fixed finite set $\mathfrak R\subset[0,1]\cap\Q$ containing the coefficients of $B$. The geometric hypotheses are those of Theorem~\ref{thm:main}, with the zero-dimensional base also allowed.

Complements entered birational geometry through Shokurov's work on threefold log flips, including bounded local lc complements for surfaces \cite[Section~5]{ShokurovFlips}. His surface theorem and bounded-complement conjectures developed the relation between mild singularities and bounded complement indices \cite[Conjecture~1.3 and Main Theorem~1.4]{ShokurovSurfaces}. The stronger expectation that bounded complements preserve a positive lower bound on log discrepancies appears in his 2004 formulation \cite{ShokurovBounded2004}. Birkar proved boundedness of lc complements in arbitrary dimension for Fano type pairs with hyperstandard coefficients \cite[Theorems~1.7--1.8]{BirkarComplements}. Those lc complements provide the bounded initial index used below. A klt complement must have strictly positive discrepancies, and obtaining this additional property uniformly is the issue addressed here.

Several earlier results describe the reach of the Cartier-section problem. Sankaran and Santos bounded the smallest weight of an $\epsilon$-lc weighted blowup of affine space \cite[Theorem~1.3]{SankaranSantos}, proving the weighted-blowup case of Birkar's question. For torus-invariant boundaries with coefficients in a fixed DCC set (one with no infinite strictly decreasing sequence), Ambro constructed invariant Cartier sections through torus-fixed base points in the toric setting \cite[Theorem~3.12, arXiv version~2]{AmbroToric}. Under his hyperstandard coefficient hypothesis, his Theorem~1.2 also gives bounded complements preserving a positive discrepancy bound. Birkar's control of base singularities gives the Cartier assertion over curves, together with bounded klt complements in the curve-base case \cite[Theorem~1.1 and Corollary~1.4]{BirkarCY}. These curve and toric results illustrate two distinct ways to control the orders appearing in the threshold inequalities.

Over $\C$, Chen proves the decisive general reduction. For each bound $\ell$ on the base dimension, bounded klt complements, the Cartier-section assertion, and its identity-contraction case are equivalent; on the complement side, birational contractions with zero boundary suffice \cite[Theorem~1.6]{Chen2025}. He proves the surface local assertion with coefficient $\epsilon^2/24$, and hence both conjectures for bases of dimension at most two \cite[Theorem~1.7 and Corollary~1.8]{Chen2025}. He also removes the coefficient restriction in the toric Cartier-section assertion \cite[Theorem~1.11]{Chen2025}. We use the precise local formulation of the Birkar--Shokurov conjecture in \cite[Conjecture~1.3]{Chen2025}; the relatively ample formulation also appears in \cite[Conjecture~1.27]{BirkarSurvey}.

Theorem~\ref{thm:main} proves the rational-boundary assertion in characteristic zero, and Corollary~\ref{cor:real-boundaries} below gives the real-boundary formulation over $\C$. The conjecture is therefore resolved positively in these settings. The new input to Chen's reduction is the following birational theorem.

\begin{maintheorem}[Birational klt complements]\label{thm:birational}
For each $d>0$ and $\epsilon>0$, there is a positive integer $n=n(d,\epsilon)$ such that the following holds. Let $X$ be an $\epsilon$-lc rationally Gorenstein variety of dimension $d$, and let $f:X\to Z$ be a birational Fano type contraction with $-K_X$ nef over $Z$. For every closed $z\in Z$, there is an effective rational divisor $\Delta$ over a neighbourhood of $z$ such that $(X,\Delta)$ is klt there and $n(K_X+\Delta)$ is Cartier and linearly trivial there, after shrinking the neighbourhood.
\end{maintheorem}

The companion article \cite[Theorem~1.1 and Corollary~1.2]{OpenAIComplements}
treats bounded klt complements for Fano contractions and their
finite-rational-coefficient extension by an independent argument.
The present proof establishes Theorem~\ref{thm:birational} directly,
through its own character and pinning obstruction.

Applying Chen's Theorem~1.6 in arXiv version~4, implication (2)$\Rightarrow$(3), transfers Theorem~\ref{thm:birational} to Theorem~\ref{thm:main}. For a fixed total dimension $d$, use the base-dimension bound $\ell=d$ and the birational theorem in every dimension at most $d$, for every positive discrepancy gap. Implication (2)$\Rightarrow$(1) of the same theorem also supplies the finite-rational-coefficient complement assertion stated above over $\C$. Its Cartier-section conclusion already allows real boundaries, giving the following extension.

\begin{maincorollary}[Real boundaries]\label{cor:real-boundaries}
Over $\C$, the conclusion of Theorem~\ref{thm:main} holds for effective real boundaries $B$ with $K_X+B$ real Cartier, with the same uniform dependence on $d$ and $\epsilon$.
\end{maincorollary}

The reduction requires rationalization only in the zero-boundary birational case: an integral Weil divisor that is real Cartier is rationally Cartier, and a real klt log Fano witness can be chosen rational. Section~\ref{sec:H} supplies these details, the reduction from arbitrary algebraically closed characteristic-zero fields to $\C$ for Theorem~\ref{thm:main}, and the passage between point-ideal thresholds and individual Cartier divisors. No approximation of the arbitrary nef real boundary in Corollary~\ref{cor:real-boundaries} is required.

\paragraph{The argument.}
Suppose that Theorem~\ref{thm:birational} fails. Section~\ref{sec:H} uses bounded lc complements to construct a sequence of complex klt germs $x_i\in S_i$ of bounded dimension, with finite cyclic groups $\Gamma_i$ fixing $x_i$. The quotient maps
\[
 S_i\longrightarrow S_i/\Gamma_i
\]
are quasi-\'etale, meaning \'etale outside subsets of codimension at least two, and the quotient germs have a fixed positive discrepancy gap for primitive prime orders. On $S_i$, a nonzero regular function $h_i$ and a nowhere-vanishing canonical generator have the same character $\chi_i$. The pair $(S_i,\divisor(h_i))$ is lc, and a centred quasi-monomial valuation $T_i$ satisfies
\[
 A_{S_i}(T_i)=T_i(h_i)=1,
\]
where $A_{S_i}$ is the log-discrepancy function with zero boundary. Failure of bounded klt complements gives the stronger condition
\[
 \forall m\in\Z_{>0}\ \exists i(m)\ \forall i\ge i(m):\qquad
 T_i(s)\ge m
 \quad\text{for every nonzero regular $s$ of character $\chi_i^m$.}
\]
Thus the order bound holds for every section in each fixed character degree, rather than for a selected sequence of sections. The reduction uses a canonical cover in the isomorphism case and a signed anti-canonical grading for other birational contractions.

Character relations become concrete after division by powers of $h_i$. For a section $s$ of character $\chi_i^m$, the rational function $Y=s/h_i^m$ is $\Gamma_i$-invariant and
\[
 T_i(Y)=T_i(s)-m.
\]
Requiring nonnegative order for both $Y$ and $Y^{-1}$ forces the exact equality $T_i(Y)=0$. The construction uses pairs of ratios with opposite orders to impose these two inequalities. The induction builds and retains relations of this kind. It optimizes valuations on the constrained $h_i$-lc slice, passes to finitely generated central algebras, and imposes further invariant relations while preserving the algebraic independence of the variables already obtained. Either it produces another independent variable or the remaining character direction yields a quotient prime of arbitrarily small discrepancy. Dimension bounds the first process, while the discrepancy gap excludes the second. Theorem~\ref{thm:G-pinning} formalizes this obstruction, and Section~\ref{sec:H} uses it to prove the birational theorem before applying Chen's equivalence.

The proof requires uniform geometric estimates and exact control of the valuation constraints. Sections~\mbox{\ref{sec:A}--\ref{sec:C}} establish the discrepancy comparisons, phase nonvanishing, and saturation statements for fields obtained by cutting off small discrepancy directions. Their geometric inputs include Birkar's boundedness of weak log Fano varieties \cite[Theorem~1.1]{BirkarBAB}, boundedness of rationally connected generalized Calabi--Yau varieties up to isomorphism in codimension one \cite[Theorem~1.7]{BirkarCY}, and effective birationality for nef and big integral Weil divisors \cite[Theorem~1.1]{BirkarPolarised}. The latter theorem also requires the polarization minus the canonical divisor to be pseudo-effective. These results apply without a prior bound on the nef divisors' Cartier indices. Filipazzi's generalized canonical bundle formula supplies the generalized-pair structures on successive bases \cite[Theorem~1.4]{Filipazzi}; the uniform discrepancy comparisons needed here are established in Section~\ref{sec:A}.

For the local optimization, the normalized-volume theory of Li, Blum, Xu, Li--Xu, and Xu--Zhuang supplies existence, quasi-monomiality, uniqueness, and stable degeneration for ordinary klt-pair minimizers \cite{LiIzumi,Blum,XuQuasimonomial,LiXuStable,XuZhuangUnique,XuZhuangStable}. The order inequalities above impose additional constraints. Section~\ref{sec:E} constructs an effective rational klt pair whose ordinary minimizer lies exactly on the constrained slice; only then is the stable-degeneration theorem applied. Section~\ref{sec:F} transfers the resulting equalities and signs through later degenerations and changes of grading. Together with the character counts of Section~\ref{sec:D}, these exact transfers ensure that each step of the induction in Section~\ref{sec:G} retains the earlier relations.

\paragraph{Conventions on generalized pairs and sequences.}
For a generalized pair, a rationally Cartier nef b-divisor $\mathbf M$ is determined and nef on a model $\pi:Y\to X$. On that model the boundary is defined by
\[
K_Y+B_Y+\mathbf M_Y=\pi^*(K_X+B+\mathbf M_X).
\]
Discrepancies on higher models are the ordinary discrepancies of the resulting, possibly noneffective, boundary. The order of a rationally Cartier b-divisor is computed after pulling it back from a determining model. Unless expressly stated otherwise, rational divisor data are used for MMPs.

Orders and discrepancies extend homogeneously to positive multiples of divisorial valuations. A lower discrepancy bound on genuine singularities is always imposed on primitive orders $\ord_E$, not on all of their positive rescalings. This distinction is essential when valuations on covers, quotients, and subfields are compared.

The auxiliary arguments concern sequences in bounded dimension. The symbols $O(1)$ and $o(1)$ refer to those sequences, with additional fixed parameters indicated at the point of use. Sequence conditions are required on every subsequence. We pass to further subsequences when proving a contradiction and choose auxiliary growth rates slowly enough to satisfy the preceding finite sets of requirements. Each uniform conclusion keeps this order of choices explicit. In Sections~\ref{sec:E}--\ref{sec:H}, germs are complex algebraic germs. An $h$-lc valuation means a valuation of zero discrepancy for the pair with boundary $\divisor(h)$, using a power to descend the rational divisor through a finite quotient when necessary.

\clearpage
\begingroup
\small
\tableofcontents
\endgroup
\clearpage
\section{Discriminant comparison and uniform chambers}\label{sec:A}

We establish three comparison results for generalized Calabi--Yau data.
Their proofs reduce to bounded generic Fano fibres, but impose no bound on
the Cartier indices of the nef b-divisors.
The geometric input is a family of charts with a uniform bound on angular
variation of discrepancies; the contact of a trait with the parameter
space is allowed to be arbitrarily large.

\subsection{Conventions and the canonical bundle formula}

We use the usual generalized pair notation.
A nef part \(\mathbf M\) is a \(\mathbb Q\)-Cartier b-divisor determined and nef on some model (nef over the point unless indicated).
In particular all pullbacks involved in forming generalized discrepancies are ordinary pullbacks on sufficiently high models.
Write \(a(v)\) for log discrepancy, extended to positive multiples of \({\rm ord}_E\) by homogeneity; unless specified, numerical conditions on divisorial valuations are imposed for \emph{primitive} ones (\({\rm ord}_E\)).
Coefficient functions of b-divisors are extended by homogeneity in the same way.
Under a generically finite pullback use log pullback, so discrepancies just compose with restriction (ramification formula, characteristic zero).
Restriction of a divisorial valuation to a finitely generated subfunction
field is a multiple of a divisorial valuation or trivial.
Indeed, let \(L/K\) be the extension, of transcendence degree \(q\), and
let \(v\) be divisorial on \(L\), with nontrivial restriction \(w\).
Both value groups have rational rank one and
\[
 \trdeg_{\kappa(w)}\kappa(v)\le q,
 \qquad \trdeg_k\kappa(w)\le\trdeg_k K-1.
\]
Equality \(\trdeg_k\kappa(v)=\trdeg_k L-1\) forces equality in
both inequalities.  Thus \(w\) is divisorial.  If \(v\) has integral
values, its restricted value group is a subgroup of \(\Z\), so the
multiple of the primitive order is a positive integer.

Here and below an \emph{effective generalized CY over \(S\)} (possibly not lc) on \(X\to S\) means data as above with \(B\ge0\) on \(X\) and \(K_X+B+\mathbf M_X\sim_{\mathbb Q} f^*L\) for some \(\mathbb Q\)-Cartier \(L\) on \(S\).
CY will also indicate its discrepancy function.
Assertions about varying pairs below need only be read on the \emph{rational} parameter values (piecewise linear discrepancy functions can of course be extended on their intervals).
For an affine discrepancy family it suffices to choose common resolutions for finitely many defining members: differences of generalized discrepancies are Cartier b-evaluations.
No denominator bounds on the nef parts are assumed.
Varieties in the statements below are projective normal in characteristic zero; one can work over algebraically closed fields.
The bounds/sequence statements work also with fields allowed to vary: boundedness, complements and families used below are uniform in characteristic zero (or one can work with descents to characteristic zero subfields embeddable in \(\mathbb C\)).
We limit the dimension of the total space by a fixed integer.

We use the klt MMP with big boundary and the extraction and Mori dream
space consequences of \cite{BCHM}; preservation of Fano type under rational
contractions \cite[p.~155 and Lemma~2.8]{ProkhorovShokurov}; rational connectedness of
projective klt log Fano pairs
\cite[Corollary 1.13]{HaconMcKernan}\cite[Theorem 1]{ZhangRC};
BAB \cite[Theorem 1.1]{BirkarBAB}; and ordinary bounded lc complements
\cite[Theorems 1.7 and 1.8]{BirkarComplements}.
For complements, the seed coefficients used below are only \(0\) and
\(1\), the varieties are of Fano type over the relevant base, and the
negative log canonical divisor is nef.  The conversion to ordinary data
will be proved below; no generalized complement theorem with an
unstated nef-index hypothesis is used.
We also use the generalized klt-trivial canonical bundle formula
\cite[Theorem~1.4 and Sections~4--5]{Filipazzi}.
Recall the latter in the projective rational case: for a contraction \(g:X\to W\), data with nef \(\mathbb Q\)-Cartier b-part, \(K_X+B+\mathbf M_X\sim_{\mathbb Q}g^*L_W\), \(B\) effective over the generic point, generalized klt generically there, the discriminant is given on every model by
\begin{equation}\label{eq:A-discriminant-minimum}
   b(P)=\min_{T|_{k(W)}=e\,{\rm ord}_P,\ e>0} a(T)/e
\end{equation}
(log-discrepancy convention) and \(L_W\) minus the sum of canonical and discriminant-boundary terms (that is, \(L_W-(K_W+B_W)\), on higher models using pullback of \(L_W\)) gives a b-nef \(\mathbb Q\)-Cartier moduli part.
This is the b-nefness/b-Cartier theorem for a generalized klt-trivial fibration; the generic rank condition follows by generic klt and horizontal effectivity on \(X\) (on a resolution the rounded-up negative boundary over the generic field is effective exceptional), using \(g_*\mathcal O_X=\mathcal O_W\).
All data here are projective with nef part over the point, NQC by rationality.
We only need the generically klt case as stated here (no global klt condition).
To fit the subboundary convention in the cited theorem when some vertical
trace coefficients exceed one, subtract \(g^*G\) for a sufficiently large
effective Cartier divisor \(G\) on the base containing the images of
those components, and replace \(L_W\) by \(L_W-G\).  Generic klt is
unchanged.  The discriminant boundary changes from \(B_W\) to
\(B_W-G\), so the moduli b-divisor is unchanged.  Thus this harmless
shift permits the stated application to arbitrary vertical coefficients.
Minima are finite and attained (work over the generic DVR of \(P\)).
We apply this to \(g\) over \(S\) with \(L_W\) a pullback of \(L\).
If \(B\) is effective on \(X\), the induced trace \(B_W\) is effective because any prime on \(W\) is covered by a prime in its inverse image.
Discriminants compose by the minimum formula, including through crepant changes and birational contractions preserving the CY discrepancy data.
In fact pushing traces in small modifications and birational contractions over \(S\) (not extracting) preserves effectivity and discrepancies; crepancy follows by writing total log canonical divisor as a pullback from \(S\) plus a rational principal divisor.

\subsection{The comparison statements}

Here are three sequence versions.
Subscripts indexing a sequence will often be suppressed; ``eventually''
allows discarding finitely many indices.  Every assertion about a test
sequence of valuations is required on every subsequence of the original
sequence as well.  Bounds may depend on the fixed dimension and on the
fixed parameters explicitly appearing in a statement.  In particular,
a condition on every test sequence is uniform against a diagonal choice
of counterexamples.  A subsequent diagonal passage may make any
auxiliary divergent parameter grow more slowly.

\begin{lemma}[Comparison along long parameter intervals]\label{lem:A-comparison} Let \(X\to S\) be a sequence of Fano type contractions as above.
Suppose \(a_t\) are effective generalized CYs over \(S\), continuous finite rational piecewise affine families of discrepancies (with a finite subdivision of parameter per example, so the pieces can be treated on a common resolution), on \(0\le t\le L_i\) or \(-L_i\le t\le L_i\) with \(L_i\to\infty\).
Pieces and endpoints can be taken rational.
Values at each valuation are concave in \(t\), and \(a_0\) is klt.
Assume the corresponding one of the following (one-sided for \(0\le t\le L_i\); two-sided for \(-L_i\le t\le L_i\)):
\begin{enumerate}[label=\textup{(\roman*)}]
\item (one-sided) the discrepancy functions are nonincreasing, and whenever \(a_0(T)\to0\), \(a_t(T)/a_0(T)\to1\) for every fixed \(t\);
\item (two-sided) for some fixed \(\delta>0\), whenever \(a_0(T)<\delta\), \(a_t(T)/a_0(T)\to1\) for every fixed signed \(t\).
\end{enumerate}

As throughout these lemmas, tests on sequences of valuations include along subsequences.
Then on every fixed bounded parameter range the data are eventually klt
over the generic point of \(S\).  Write \(b_t\) for their discriminant
discrepancies on \(S\).  In case \textup{(i)}, they are nonincreasing
and satisfy \(b_t(P)/b_0(P)\to1\) for every fixed \(t\) whenever
\(b_0(P)\to0\).  In case \textup{(ii)}, some fixed \(\delta'>0\)
gives the same convergence for every fixed signed \(t\) whenever
\(b_0(P)<\delta'\).  These conclusions include every subsequence and
primitive prime orders on every birational model of \(S\).

\end{lemma}

\begin{lemma}[Phases under pure pullback]\label{lem:A-pullback-phases} Let \(a\) be a sequence of klt effective generalized CYs over \(S\) in the same Fano type setting, and let \(H\) be a sequence of \(\mathbb Q\)-Cartier \(\mathbb Q\)-divisors on the respective bases \(S\).
Suppose for some fixed positive integer \(m\), for every \(a(T)\to0\) the phase \(v_T(m f^*H)\bmod 1\) has lifts to \(\mathbb R\) of size \(O(a(T))\) bounded below by \(-o(a(T))\).
Then with discriminant \(b\), for some fixed positive multiple \(m'\), \(v_P(m'H)\) satisfies the same phase condition wherever \(b(P)\to0\).

\end{lemma}

\begin{lemma}[A uniform linear chamber]\label{lem:A-linear-chamber} In the same setting suppose \(a\) is an lc effective generalized CY over \(S\), with values in \((1/p)\mathbb Z\) for fixed \(p\), and \(a+t d\) for \(0\le t\le t_0\) are effective generalized CYs, klt for positive \(t\), \(d\ge0\) on all divisorials; take \(L_t\) affine, and \(t_0>0\) fixed rational.
Then eventually there is a uniform smaller \(t_1>0\) on which discriminants equal \(a'+t d'\) with \(a'\) of bounded denominator (i.e., the same \(p'\) works on all models).
There is also the following version with extra data:
\begin{enumerate}[label=\textup{(\roman*)}]
\item \(a+\lambda h+s c\) are effective generalized CYs over \(S\) for \(0<\lambda\le\lambda_0,\ 0<s\le\lambda\theta_0\), \(\lambda_0,\theta_0>0\) fixed, \(h\) of bounded denominator, \(c\ge0\);
\item on \(a=0\), \(d\le c\le C d\), \(-\eta_i c\le h\le Cc\) with fixed \(C\) and \(\eta_i\to0^+\).
\end{enumerate}

Then there are uniform \(\lambda_1,\theta_1>0\), with
\(\lambda_1\le\lambda_0\) and \(\theta_1\le\theta_0\), such that
eventually, on the sector
\(0<\lambda\le\lambda_1\), \(0<s\le\lambda\theta_1\), the data
are klt generically over \(S\), their discriminants are exactly
\(a'+\lambda h'+s c'\), and \(h'\) has bounded denominator.
Moreover, eventually there is \emph{no} divisorial \(T\) horizontal
over \(S\) with \(a(T)=0,\ h(T)<0\).
In particular for \(S\) a point exact nonnegativity on \(a=0\) follows even though only the approximate bound \(-\eta_i c\) was assumed.

In this statement the minima are lexicographic: first minimize \(a/e\),
then \(d/e\) on its minimum locus; for the extra data first minimize
\(a/e\), then \(h/e\), then \(c/e\).
Here \(e\ord_P=T|_{k(S)}\).
Thus \(a'\) itself is the discriminant in discrepancy convention even
if the endpoint is not klt generically.
The family \(a'+t d'\), \emph{including its endpoint at zero}, has the generalized effective log data (with nef b-part and total logarithmic divisor \(L_t\)).

\end{lemma}

The endpoint assertion follows from the positive-parameter assertion.
Indeed, for positive parameters use the generalized canonical bundle formula.
The moduli b-parts chosen using \(L_t\) are affine on \emph{all} models by the affine discrepancies; thus the limit b-divisor at zero is a difference of \(\mathbb Q\)-Cartier b-divisors by extrapolation, and nef by passage to the limit on a common determination.
The same applies on any intermediate base using pullback of \(L_t\).

We give details of the geometry and proofs of these lemmas.
Boundedness of denominators here and later does not require that the b-parts of the input have bounded Cartier index.

\subsection{Extraction and the Mori step}\label{subsec:A-mori}
We can pass to a small \(\mathbb Q\)-factorialization on which the
families still have effective traces.
Given an effective generalized klt CY \(A\) over \(S\) and Fano type, we can extract all exceptional divisors with \(A<\delta_*\) (\(0<\delta_*<1\)) to a \(\mathbb Q\)-factorial Fano type model over \(S\), extracting no others (or just select any subset of this list).
There are finitely many: on a log resolution where the nef part descends any further ones must be toroidal for the snc support (the integral reduced-log discrepancy otherwise contributes at least 1: all lc places of a reduced snc pair are monomial at the strata), and all their weights are bounded using positivity on the finitely many components.
For extraction one can pass to an ordinary big klt CY over \(S\) having the selected divisors still strictly positive in log-pulled boundary.
For clarity, mix the given data with small weight of auxiliary crepant sub-data with ample b-part determined on that resolution, horizontal and vertical trace on \(X\) effective and containing a general relatively ample contribution; these auxiliary data with total \(\sim_{\mathbb Q}0/S\) exist by taking the ample b-part sufficiently small and using a klt log Fano pair on \(X/S\).
More explicitly subtract the small pushforward of the ample divisor and a small ample divisor on \(X\) from the log Fano ample class (twist by a pullback if needed).
Sub-klt persists on a common resolution after mixing, as do effectivity on \(X\) and positive coefficients at the selected divisors.
Replace the resulting ample moduli contribution on the high model by general representatives (over \(S\), twisting by pullback), with small coefficients.
Apply klt extraction; the log transform is effective and still big over \(S\) by the general ample contribution on \(X\) avoiding the finitely many centres.
Big klt CY implies Fano type.

Other CY families retain effective traces on the extracted model
provided their discrepancies at the extracted divisors are at most
\(1\).  Each application below verifies this condition.
When the whole list is retained, the resulting \emph{ordinary variety}
is \(\delta_*\)-lc.  For an exceptional divisor on a higher
determination, effectivity of the trace and the negativity lemma for
the difference between the pullback of the nef-part trace and its nef
determination give \(A_X(T)\ge A(T)\).  Every divisor for which
\(A(T)<\delta_*\) is already on the extracted model; its ordinary
discrepancy there is \(1\).  The remaining exceptional divisors have
\(A_X(T)\ge A(T)\ge\delta_*\).
A \(K\)-MMP over \(S\), if \(\dim X>\dim S\), now ends at a \(\delta_*\)-lc Fano Mori step \(g:X_1\to W\) over \(S\) with \(\dim W<\dim X\), preserving the CY discrepancies/effectivity we arranged and Fano type.
Here \(-K_X\) is big over \(S\), and the ordinary discrepancies stay bounded below by monotonicity in the \(K\)-MMP; the geometric generic fibre \(F\) of \(g\) (we also write \(F\) for the fibre over \(k(W)\)) is a bounded Fano by BAB (generic discrepancies via horizontal divisors and separable constant extensions). \(W\) is Fano type over \(S\) (image of Fano type).
If instead the contraction to \(S\) is birational no further step is
needed: the function fields agree and the discriminant is the original
discrepancy function.  This is the terminal case in each induction.

\subsection{Tailored complements and connected minimizing places}

\begin{lemma}[Tools at an lc place]\label{lem:A-lc-place-tools}
Let \(g:X_1\to W\) be a Mori step as above, and let \(A_*\) be an
effective generalized CY discrepancy on \(X_1\) over \(S\).
\begin{enumerate}[label=\textup{(\roman*)}]
\item Suppose that \(A_*\) is lc over the generic point of \(W\),
and that an effective generalized klt CY over \(S\) is available for
perturbation.  Any prescribed horizontal divisorial lc place remains
an lc place of an ordinary boundary \(\Omega\), effective on the
generic fibre \(F\), with \(n(K_F+\Omega_F)\sim0\), for a bounded
positive integer \(n\).
\item For \(P\) on a birational model of \(W\), suppose \(A_*\) is klt generically and \(T\) computes the discriminant \(b_*(P)\).
There is an ordinary \(n\)-CY over the DVR, lc in its b-discrepancies
and effective horizontally on \(F\), with \(n\) bounded, keeping
\(T\) as an lc place.  Here \(n\)-CY means that the ordinary total
log canonical divisor, multiplied by \(n\), is linearly trivial
over the trait.
Moreover for the adjusted discrepancy \(A_*-b_*(P)e\), any two minimizing divisorials can be joined by a chain/interpolation of lc places in the following sense: on a log resolution over the generic DVR their strict divisors (or extractions of them) belong to a connected snc floor consisting of vertical coefficient-one divisors with monomial cones at the intersections.
One can include any additional finite discrepancy data in the resolution.
In particular normalized evaluations \(a/e\) of such other data interpolate continuously along the chain, using rational weights to approximate intermediate values by divisorial evaluations.
\end{enumerate}

\end{lemma}

\begin{proof}
Here everything over a generic DVR can be carried out by shrinking about a generic point of a prime in a smooth model over \(k\).
Only assertions over the generic fibre and that DVR are needed.
\emph{A model retaining the selected fibre components.}
Write $R$ for the trait ring.  Fix a nonempty finite collection of
divisors computing $b_*(P)$ and choose a log resolution $Z$ over
$\Spec R$ on which they occur.
For the adjusted generalized data $A_*-b_*(P)e$, they have boundary
coefficient one.  We first construct an ordinary klt pair on $Z$ whose
minimal model contracts every other special-fibre component and every
horizontal divisor exceptional over $F$.

Subtract a sufficiently small positive multiple of the fibre from the
adjusted boundary.  The resulting generalized subpair is sub-klt, and
its coefficients on the selected components remain positive.  Mix it
with a sufficiently small amount of auxiliary sub-CY data having an
ample b-part on $Z$ and effective nonexceptional horizontal trace.
Such data exist because $F$ is Fano: a small pushed ample contribution
can be subtracted from $-K_F$, and vertical terms can be adjusted freely
over the trait.  The pushforward is rationally Cartier on the
$\Q$-factorial total space $X_1$.
Replace the mixed ample b-part by general relatively free
representatives.  Make these choices on a common resolution and retain
a general big contribution avoiding the components whose coefficients
will be changed.

Increase the coefficients of all horizontal divisors exceptional over
$F$ and all special-fibre components outside the selected collection,
strictly and to values in $[0,1)$.  We obtain an effective ordinary klt
boundary $\Delta$, big over the trait, and
\[
                    K_Z+\Delta\sim_{\Q,R}E,
\]
where $E\ge0$ is supported precisely on the divisors just increased.
In particular at least one special-fibre component is absent from
$\Supp E$.

Here the relative fixed part can be computed directly.  For divisible
$m$, a rational section of $\mathcal O_Z(mE)$ restricts on the generic
fibre to a rational function with poles only along divisors exceptional
over the normal projective variety $F$.  It descends to a regular
function on $F$, and is therefore a constant in the trait field.
The section condition on a retained special-fibre component, whose
coefficient in $E$ is zero, forces that constant to have nonnegative
DVR order.  Conversely every element of $R$ gives such a section.
Thus the relative section module is $R$, and the relative fixed part
of $|mE|$ is exactly $mE$.

The good minimal model theorem for klt pairs with big boundary now
applies, after spreading over an open neighbourhood of the generic
point of $P$.  Every curve in a negative extremal ray lies in $\Supp E$, so
no divisor outside this support is contracted.  The section ring is
preserved by the MMP.  On its semiample output the surviving trace of
$E$ must therefore vanish, since otherwise it would remain a nonzero
fixed part.  Consequently every component of $E$ is contracted.
Call the output $Y$.  Its boundary remains big and klt and its log
canonical divisor is rationally trivial over the trait, so $Y$ is of
Fano type there.  Its generic fibre agrees with $F$ in codimension one.
The original adjusted generalized data are crepant through these
birational operations.  Their trace on $Y$ is effective: the only
special-fibre components are the retained ones, with coefficient one,
and the surviving horizontal coefficients are the effective
nonexceptional coefficients over $F$.

\emph{Connectedness of the minimizing locus.}
Let $\pi:\widetilde Y\to Y$ be a log resolution on which the nef part
of the original adjusted data descends.  Write $B^{\log}$ for its
boundary and $H=(B^{\log})^{=1}$.  The negative boundary is
$\pi$-exceptional because the trace on $Y$ is effective.  Moreover
$-(K_{\widetilde Y}+B^{\log})$ is nef over $Y$, and is big over $Y$
because $\pi$ is birational.  The birational connectedness lemma
therefore makes $H$ connected over each fibre of $\pi$.

One can see the relevant structure-sheaf assertion directly.  Put
$G=\lceil-(B^{\log})^{<1}\rceil$, an effective exceptional divisor.
Relative Kawamata--Viehweg vanishing for
$\mathcal O_{\widetilde Y}(G-H)$ gives a surjection
\[
               \mathcal O_Y\longrightarrow\pi_*\mathcal O_H(G).
\]
The image factors through $\pi_*\mathcal O_H$, which injects into
$\pi_*\mathcal O_H(G)$.  Hence $\pi_*\mathcal O_H$ is the structure
sheaf of the reduced image.  This image is the entire reduced special
fibre of $Y$, since all its components have coefficient one.  The
special fibre is connected by Stein factorization.  Thus $H$ is
connected, and the prescribed divisors lie in one connected snc
coefficient-one locus.  Generic klt ensures that every component of
this locus is vertical.  Monomial valuations along its intersections
join its divisorial orders continuously; the uniformizer has positive
value throughout.  Evaluations of any additional discrepancy data are
linear on these cones after taking a common resolution, so their
ratios to $e$ interpolate continuously.  Rational weights recover the
required divisorial interpolants.

\emph{An ordinary complement preserving the chosen place.}
For a tailored ordinary \(\Omega\), use the model above (or horizontally, extract the horizontal place using a CY slightly perturbed towards klt as funding).
Write \(T_0\) for the resulting prime, which has coefficient one.
Run a \(-(K+T_0)\)-MMP; this divisor is pseudo-effective over the trait (respectively generically) by the effective generalized CY containing \(T_0\).
On every model the pair with strict \(T_0\) is lc by comparison with the generalized data and b-nef negativity.
At the nef model take a bounded ordinary lc complement (coefficients only 0 and 1, Fano type; in the relative case use Birkar's complement over a closed point of the relevant prime/open, after appropriate shrinkings as above, containing the generic trait/field in the resulting neighbourhood).
Indeed over the trait/field the final nef divisor is semiample by Fano type; generation of a multiple spreads out, as does a Fano type witness (relative ampleness and klt) and lc of the indicated pair via resolution, shrinking by projectivity to ensure the conditions for complements over an open.
Its crepant pullback boundary on the model with \(T_0\) still dominates \(T_0\): it adds the pull of the effective complement addition and the effective MMP discrepancy of \(-(K+T_0)\).
In particular contraction of \(T_0\) in a negative step would even contradict lc.
Push generically to \(F\), preserving effective boundary there.
All complement indices use only the absolute dimension bound.
This proves both assertions.

\end{proof}

\subsection{Uniform charts and discrepancy estimates}

\begin{proposition}[Uniform charts]\label{prop:A-uniform-charts}
At a chosen vertical place in Lemma~\ref{lem:A-lc-place-tools}\textup{(ii)},
after a bounded degree extension of \(k(W)\) and a chosen extension of
the DVR and of the place, there is a toroidal cone representing the
place, with nonnegative coordinates
\begin{equation}\label{eq:A-contact-cone}
       (x,e_*),\qquad M x=r e_* .
\end{equation}
Here \(M\) ranges over a finite list of integral matrices of bounded sizes, \(r\) is integral (possibly huge), \(e_*\) is value of the new uniformizer, and original uniformizer value is \(\rho e_*\) with positive bounded \(\rho\).
Nonzero integral solutions on the cone give divisorial valuations (not necessarily primitive after restriction); evaluations of discrepancy data always refer to homogeneous values by restriction.
For horizontal places as in the first bullet one just has an orthant on a bounded generic resolution, without \(e_*\).
For \emph{any} effective generalized CY data \(A\) on \(X_1\) over \(S\), values on these cones are finite continuous homogeneous piecewise linear, and are Lipschitz in \(x\) with uniform constant \emph{at fixed} \(e_*>0\) (uniform homogeneous Lipschitz in the horizontal case).
Constants depend on the boundedness in this Mori step and the complement indices, \emph{not} on \(r\), vertical coefficients, or the nef part index or the chosen data \(A\).
In particular finitely many differences at fixed parameter separations have uniform Lipschitz bounds as well.

\end{proposition}

The proof separates the construction of the charts from the metric
estimate.  The first lemma keeps track of their integral lattices and
logarithmic coordinates; the second converts generic-fibre degree into
an angular bound.  We then replace the nef part at two chosen valuations
to apply that bound to arbitrary generalized CY data.

\begin{lemma}[Bounded toroidal charts]\label{lem:A-toroidal-charts}
For the prescribed place and tailored complement in
Proposition~\ref{prop:A-uniform-charts}, there are charts with its stated
contact equations and integral-valuation property.  The prescribed
place is monomial on one of these charts.  They have rational functions
\(z_1,\ldots,z_f\), \(f=\dim F\), with uniformly bounded pole degrees on
a bounded embedding of \(F\), such that:
\begin{enumerate}[label=\textup{(\roman*)}]
\item their logarithmic differentials form a relative log basis in the
vertical case and an absolute log basis over the generic field in the
horizontal case;
\item their orders are linear functions of the chart coordinates with
uniformly bounded slopes;
\item every rational monomial valuation in the vertical chart with
\(e_*>0\) restricts to a Gauss valuation on
\(K_1(z_1,\ldots,z_f)\), where \(K_1/k(W)\) is the bounded extension
used in constructing the chart.  Horizontally the same assertion holds
with the trivial valuation on the generic field.
\end{enumerate}
Here Gauss means that evaluating a polynomial in the \(z_j\) gives the
minimum of its coefficient-adjusted monomial weights, without
cancellation among terms of equal weight.
\end{lemma}

\begin{proof}
\emph{Bounded marked fibres and descent of an embedding.}
The pairs \((F,{\rm Supp}\,\Omega_F)\) are in bounded embedded families over the geometric field.
Indeed boundedness gives a very ample class \(\mathcal H\) of bounded degree and embedding dimension on each bounded \(F\).
One can use the complete system here (dimension still bounded by degree for the nondegenerate image).
To bound the anticanonical degree, let \(f=\dim F\) and take a
smooth complete-intersection curve \(C\) of \(f-1\) general members
of the very ample system \(\mathcal H\), avoiding the singular
locus.  Adjunction gives
\[
 -K_F\cdot\mathcal H^{f-1}
   =(f-1)\mathcal H^f+2-2g(C)
   \le (f-1)\mathcal H^f+2 .
\]
For \(f=1\), the same formula is the usual degree formula on the
smooth curve \(F\).
The nonzero coefficients of \(\Omega_F\) are at least \(1/n\),
so
\[
 \deg_{\mathcal H}\Supp\Omega_F
 \le n\deg_{\mathcal H}\Omega_F
 =-n K_F\cdot\mathcal H^{f-1}.
\]
The support therefore has bounded degree.
In the DVR case we can secure the embedding over a bounded degree extension of \(k(W)\).
The geometric Picard group is free abelian of bounded rank: it injects by normality into the Picard group of a smooth projective resolution, which is rationally connected (rational connectedness for klt Fano) and has bounded Betti numbers for a suitable resolution, by resolving the bounded families after stratification.
We use here the usual characteristic zero Picard consequences of smooth rational connectedness (in particular trivial connected Picard and no torsion).
The finite image of Galois on the discrete Picard lattice has bounded
order.  One way to see the uniformity is to reduce a finite subgroup of
\(\operatorname{GL}_r(\Z)\) modulo \(3\): the principal congruence
subgroup at \(3\) is torsion free, so this reduction is injective, and
\(\lvert\operatorname{GL}_r(\mathbf F_3)\rvert\) bounds its order.
After bounded extension the chosen \(\mathcal H\) is invariant; its complete linear system descends projectively as a projective-space form of bounded dimension (the ambiguity on section maps is scalar), which splits after another bounded extension (Severi--Brauer index bound).
This suffices.
Bounded-degree embedded varieties and their marked reduced divisors as
above are parametrized by finitely many Chow families
\cite[Definition~3.21, (3.21.1)--(3.21.3)]{KollarChow2017}.
Here we use effective cycles and their reduced supports, with the
universal cycles restricted to smooth characteristic-zero strata.

\emph{Toroidalization of the fixed families.}
Apply marked weak toroidalization
\cite[Theorem 1.1]{AbramovichDenefKaru} to these families.
Its source and base modifications can be taken projective and
birational; allowing a fixed alteration would give the same construction
after a bounded extension.
Here are the uniformity and chart details we use.
By compactification and noetherian stratification, for a dense open of each parameter stratum we can use a single proper toroidalized family over a modification (an alteration would also suffice) of its compactified base.
Over the dense open (which can always be shrunk, treating its complement by noetherian induction) the new fibres give birational models of the original varieties.
Here one works stratum by stratum starting with geometrically integral generic fibres, retaining these and birationality in the spread-out open (in case of base alteration it suffices to modify the base-changed family birationally).
The hypotheses on the fibres parametrized, such as geometric integrality and marked support of lower dimension, can be imposed by ordinary Hilbert stratification.
The compactification of a stratum serves to extend traits whose \emph{field} point is in that dense open; their closed points need not parametrize good pairs at all.
We include horizontal markings as follows: first resolve the generic fibre with its marked support over this stratum, spread out over an open, and mark the strict boundary and the exceptional locus; arrange their preimage in the toroidalization to be in the boundary.
Thus on the resulting generic-fibre models relevant to us any \emph{positive} crepant coefficient of \(\Omega_F\) is supported on the toroidal boundary: outside the old marking one had smooth with crepant subboundary \(\le0\).
We arrange the dense open to be interior on the toroidal base of the family.
We may use smooth strict toroidal models (also by toroidal resolution); all modifications here are of fixed finite type families.
Any base alteration only requires bounded degree field extension to lift the field point on the chosen dense open.
Thus after bounded degree extension of \(k(W)\), say to \(K_1\), and extension of the place, we have a lift of the whole new trait to this base by properness, with its generic point interior.
Write \(R\) for the new trait ring, \(t_*\) for its uniformizer, \(\rho\) the ramification index over the original DVR (bounded).
We always use log pullback over the finite extension.
One may view the trait as the localization at a prime divisor in a smooth model over \(k\), working close to the generic point of that divisor (and spreading the map to such a neighbourhood, shrinking to pull the base boundary only to the closed-point boundary there).
In particular the divisorial log structure of the trait comes from a log-smooth \(k\)-pair locally.

\emph{Full character rank and the log chart.}
We verify log smoothness at the strata, including the case in which the
map on boundary exponent groups has a kernel.
In a completed toroidal description use the full character lattices of
the toric models.  These include invertible torus characters as well as
the characters cutting out the boundary.  Dominance gives an injective
map of full character lattices after tensoring with \(\Q\).  In their
logarithmic differential bases its matrix has constant integral
entries.  It therefore has full target rank at the closed stratum as
well as at the generic point.  Any lattice torsion is tame in
characteristic zero.  The unit characters are essential here: two
target boundary characters can have the same boundary exponents on the
source while their quotient pulls back to a nonconstant unit, whose
differential supplies the missing direction.

To detail the algebraic chart check, at closed points in these smooth strict models write the target boundary equations \(y_j,\ 1\le j\le s\), as pulling to \(u_j x^{\alpha_j}\) with \(x_1,\ldots,x_m\) source boundary parameters and \(u_j\) units.
The map of log differential fibres has full target rank (as above in the completed toroidal models, where boundary equations match up to units).
For \(\alpha:\mathbb Z^s\to\mathbb Z^m\) the exponent map, let
\(K=\ker(\alpha)\).
The quotient \(\Z^s/K\) is the free group \(\operatorname{im}(\alpha)\),
so \(K\) is saturated and admits an integral projection
\(p:\Z^s\to K\).
The unit homomorphism \(u\) from \(\mathbb Z^s\) factors as its kernel-direction contribution via this projection, times a homomorphism on \({\rm im}(\alpha)\).
The latter extends to \(\mathbb Z^m\) after adjoining unit roots \'etale locally, so can be absorbed by rescaling the \(x_\ell\) by units.
We now have the chart \(\mathbb N^s\to \mathbb N^m\oplus K\) using \(\alpha\) and the projection, with the free group summand mapping on the source by unit functions \(u|_K\).
This is injective on groupifications with tame cokernel torsion.
The induced map to the product of the target with the source toric chart scheme over \(\mathbb A^s\) is smooth near the point: besides the rescaled \(x_\ell\), one needs independent ordinary differentials for those kernel units and the pulled target transverse coordinates along the source stratum.
Exactly these are independent there by full log differential rank (their boundary-weight entries are zero).
This proves the smoothness required in Kato's chart criterion
\cite[Theorem~3.5 and Remark~3.6]{KatoLog}; see also the statement and
proof in \cite[Theorem~4.1, Remark~4.2 and Section~6]{FumiharuKatoLog}.
In particular the required rank at a closed stratum follows from the
constant full-character matrix; it does not require specializing an
arbitrary ordinary differential minor.

\emph{The trait and its integral contact cone.}
Take the fine saturated log fibre product with the trait equipped with
its closed-point divisorial log structure.  Log smoothness is preserved
by fs base change and composition \cite[Section~1.1]{IllusieNakayamaTsuji}.
Spread the trait to the log-smooth $k$-pair described above and form the
same fs fibre product there.  Its source is log smooth over $k$ with
the trivial log structure.  It therefore admits smooth local charts
over affine toric varieties
\cite[Theorem~4.7 and Section~7]{FumiharuKatoLog}.
In these charts the monomial stratum quotient is regular, and its
dimension plus the characteristic-monoid rank is the ambient dimension.
Thus the spread-out source is log regular; localization gives this
property for the fs fibre product over the trait.
Its main component is proper over the trait and gives the desired generic model.
Indeed the map from the fs fibre product to the ordinary one is proper here (in charts, integralization and finite toric saturation); the trivial-log open is dense in a log regular model, and maps to the generic point of the trait.
At a stratum with boundary generators from the total family and from the trait, the cone of nonnegative orders is given precisely by \(Mx=r e_*,\ x,e_*\ge0\).
Here rows come from the toroidal base boundary equations, pulled to total boundary monomials times units with multiplicities \(M\), and to the trait with contacts \(r\); thus \(M\), unlike \(r\), comes from fixed chart types.
These are the matching conditions from the pushout chart (the total family and its base before this log base change can have smooth strict toroidal structures, hence boundary-coordinate charts).
After spreading near the trait the base has also ordinary residue/transverse coordinates, with the prime divisor as its single boundary near the generic point in use; only its boundary weight thus enters these equations.
Units have weight zero; at the point used only boundary generators vanishing at its images in the factors are counted.
No further independent zero conditions on weights arise from units of the saturated monoid becoming units at the point: a positive multiple of such an element comes from the original generators and if it is a unit cannot involve any generator vanishing at the point.
An integral match defines an integer-valued homomorphism on the
groupification of the pushout monoid.  Saturation takes place in that
same groupification, so this homomorphism remains integer-valued on
the saturated monoid.
We may work in strict \'etale local charts at a geometric stratum and then restrict valuations.
Thus nonzero rational weights give monomial divisorial evaluations (and the zero vector the trivial evaluation); integral ones restrict to integral multiples of primitive divisorials when nontrivial.
Indeed in the split toroidal chart one is taking the dual cone of the characteristic monoid at the stratum.
A nonzero rational ray even on a face can be introduced in a toroidal subdivision (over the corresponding generization stratum for a face); the primitive vector for its toric lattice then gives a boundary-divisor order.
Matched integral weights as just described are integral in this lattice.
For \(e_*>0\) the restriction to \(K_1\) is \(e_*{\rm ord}_{R}\).
This description is equivalently obtained by fs pulling back the log-smooth monomial charts and subdividing toroidally.
Use the family of monomial valuations for one chosen lift of the stratum in a split local chart (cones/subdivisions there, then restrict to the original function field); integral points give integral orders since they already do so in such a chart.
No degree bound on further strict \'etale local charts is needed here:
restriction of an integer-valued valuation still has value group a
subgroup of \(\Z\), hence is an integral multiple of a primitive
divisorial order.  It is never divided by the chart degree.
Only the initial field extension contributes the separately bounded
ramification index \(\rho\).  The rational functions used in the metric
estimate below are taken on the fixed total family itself.

All special components are log boundary on this model.
The subboundary of the crepantly pulled ordinary \(\Omega\) is bounded above by the reduced log boundary (coefficient bound by lc, horizontal positive support marked as above).
Thus its chosen lc place is a place also of the reduced-log pair, since the difference is effective \(\mathbb Q\)-Cartier.
Such lc places are monomial on the log regular model: pass to a toroidal snc subdivision and use the lc place description for an snc pair.
We take a cone as above containing the extended chosen place.
For a horizontal place as in the first bullet one just takes a bounded stratified log resolution of the generic pairs \((F,{\rm Supp}\,\Omega_F)\); same reasoning yields an ordinary snc monomial cone (over the generic field, or its extension).

\emph{Logarithmic coordinates and Gauss restriction.}
For the later metric estimate, choose rational log coordinates
\(z_1,\ldots,z_f\), \(f=\dim F\), on charts of the \emph{fixed}
toroidalized families, such that \(d\log z_j\) form a basis of relative
log differentials near the stratum in question.
Indeed the latter are locally free by log smoothness, and one can choose among boundary parameters and invertible functions giving generators from the smooth total chart.
Finitely many choices of genuine rational functions on opens of smooth strict charts suffice (Cartier local boundary equations and local units generate logarithmic differentials).
Their weights are bounded linear forms in the cone coordinates (boundary parameters or units), and the basis property persists under fs base change and toroidal resolution.
Moreover the \(z_j\) have bounded pole degrees on the embedded \(F\)'s for these charts: they are chosen on fixed families birational to the embedded ones, so on a sufficiently shrunk dense open of each stratum they give fixed bounded-degree rational maps on the embedded fibres.
This shrinking can be built into the stratification (express each function by fractions of bounded-degree embedded coordinates on the generic fibre of the family and spread).
Use the analogous absolute log bases in the horizontal generic case.

A rational toroidal valuation here with \(e_*>0\) restricts to a Gauss valuation on \(K_1(z_1,\ldots,z_f)\).
To check, rescale to a primitive divisor on a toroidal subdivision.
Write \(\nu_j=v(z_j)\in\mathbb Z\), \(e_*=v(t_*)>0\).
At this divisor the residues of \(z_j^{e_*}t_*^{-\nu_j}\) have independent log differentials.
Indeed relative log cotangent there, modulo the divisor, is identified with residue differentials over the base residue field: \(t_*=\pi^{e_*}\cdot(\text{unit})\), for uniformizer \(\pi\), eliminates \(d\log\pi\), while base unit differentials restrict as usual.
This also follows from the exact sequences for log differentials at the generic smooth divisors (and separability).
The indicated unit log differentials are \(e_*\) times the chosen basis in relative log cotangent.
The residues are thus algebraically independent.
In particular a lattice basis of the weight ties (monomials \(z^\gamma\) with \(\nu\cdot\gamma\) a multiple of \(e_*\), adjusted by that power of \(t_*\)) has algebraically independent residues too, so no cancellation occurs at equal coefficient-adjusted weights.
Horizontally one keeps \(d\log\pi\), without a base uniformizer; for a nontrivial monomial valuation of integral rank one weights, a basis of the kernel characters of the weight gives independent residues by the same log basis calculation.
This again gives Gauss, now for the trivial base valuation.

\end{proof}

\begin{lemma}[Degrees control angular variation]\label{lem:A-function-slopes}
On a chart of Lemma~\ref{lem:A-toroidal-charts}, the order of a nonzero rational
function is continuous, homogeneous and finite piecewise linear,
including on the faces.  If its pole divisor on the embedded generic
fibre has degree at most \(m\), where \(m\ge1\), its Lipschitz constant
in \(x\) at fixed \(e_*>0\) is at most \(Cm\).  Horizontally this bound
holds on the entire cone.  The constant \(C\) depends only on the fixed
families and chart choices.  A function with pole degree less than one
is constant on the generic fibre and has zero angular variation.
\end{lemma}

\begin{proof}
\emph{Piecewise linearity and faces.}
On these cones rational functions have continuous homogeneous finite piecewise linear values, including faces.
For example subdivide a split chart toroidally to snc (subfaces using the corresponding generization strata in the subdivision) and use the monomial order formula at the generic stratum (after splitting a chart if necessary), by expanding at the snc parameters: minima on power series supports are finite polyhedral by monomial noetherianity.
The same minima using weights zero off a subcollection of parameters give the values on its face, i.e. at the more general stratum.
Indeed the monomial ideals in the retained parameters test order there; membership is still tested in the power series ring localized along the prime of those parameters (faithful flatness and the monomial formula; a positive-threshold monomial ideal for the retained parameters only is primary to their prime, as seen already in completion).
This gives continuity and compatibility across subdivisions.
In a local split toroidal chart the subdivisions are toric subdivisions at that stratum; one is using the corresponding monomial order family throughout (weight ideals of boundary monomials), so that subfaces common to two cones use the same such orders.

\emph{A degree bound for angular slopes.}
Suppose a nonzero rational function \(g\) on \(F\) over the extended field has pole degree \(\le m,\ m\ge1\), in the bounded embedding.
It satisfies a polynomial equation over the \(z\)'s of bounded degree in \(g\) and degree \(O(m)\) in the \(z\)'s.
Use the equation of the image hypersurface of \((z,g)\) in a product of projective lines.
Its multidegrees are bounded by intersecting general level divisors of all but one coordinate on \(F\) on the locus of definition: individual divisor degrees are bounded by their pole degrees, so projective B\'ezout bounds isolated general solutions by a constant times the product bound (characteristic zero).
Equivalently intersect the product of those level-divisor supports with diagonal conditions of bounded degree in a product of projective spaces, counting isolated points (the product variety has Segre degree bounded by a fixed multiple of the product of individual degrees).
Write this equation as
\[
 \sum_{j=0}^q c_j(z)g^j=0,
 \qquad q\le C,\qquad \deg_z c_j\le Cm .
\]
By the Gauss property, \(v(c_j)\) is a minimum of affine functions whose
\(x\)-slopes have norm at most \(Cm\).  Their constant terms can involve
arbitrary orders of coefficients in the trait field, but at fixed
\(e_*\) these terms do not vary.
At least two nonzero terms of the equation attain its least valuation.
Consequently \(v(g)\) is one of the finitely many quantities
\[
 \frac{v(c_i)-v(c_j)}{j-i}\qquad(i\ne j).
\]
Each is \(Cm\)-Lipschitz after enlarging the fixed constant.
Along a segment the continuous, finite piecewise linear function
\(v(g)\) selects from this finite list; subdividing at its finitely many
switches gives the same Lipschitz bound on the entire segment.
Rational weights are dense, so the estimate holds also for real
weights.  Horizontally the argument applies on the entire cone.
If the pole degree itself is \(<1\) the function is constant along the generic fibre.

\end{proof}

\begin{proof}[Proof of Proposition~\ref{prop:A-uniform-charts}]
Lemma~\ref{lem:A-toroidal-charts} supplies the chart, the prescribed
monomial place and the integral orders.  We first extend the
piecewise-linearity assertion of Lemma~\ref{lem:A-function-slopes} from
rational functions to the discrepancy data.
Evaluations of \(\mathbb Q\)-Cartier b-divisors are also finite piecewise linear: locally over an affine neighbourhood of the stratum write a determining Cartier multiple on a projective model as a difference of relatively generated divisors, then it uses minima on finite sets of rational functions.
Generalized discrepancies differ on this cone from the reduced-log discrepancy (identically zero) by such Cartier b-evaluations.
Piecewise-linearity here need not have bounded combinatorial complexity.


\emph{Two-point replacement of the nef part.}
Fix two rational valuations on one cone with the same \(e_*\), and
choose a common high model on which their centres and the determination
of the nef part are visible.  On that model let \(L\) be an ample
divisor, relative to the trait, and choose a rational
\(\varepsilon>0\).  After twisting by a pullback from the base, take
general effective representatives
\[
 E_+\sim_{\Q}\mathbf M+\varepsilon L,\qquad
 E_-\sim_{\Q}\varepsilon L
\]
which contain neither centre.  The ample perturbation is selected
\emph{after} the high model: its coefficient can therefore be made so
small that its pushed horizontal degree is arbitrarily small.
Add \(E_+-E_-\) to the crepant generalized boundary and remove the nef
part.  Since both representative divisors have order zero at both
chosen valuations, the resulting ordinary divisor data have exactly the
prescribed discrepancies at these two valuations.

Let \(\Theta\) be its horizontal trace on the bounded fibre \(F\).
Write degrees with respect to the bounded very ample class.
The original trace \(B_F\) is effective, and the trace of the nef part
has nonnegative degree.  Relative triviality gives
\[
 \deg B_F+\deg\mathbf M_F=-\deg K_F .
\]
Negative exceptional crepant coefficients on the high model disappear
when pushed to \(F\).  Hence, if \(L_F\) denotes the pushed ample
contribution, then
\[
 \deg\Theta^+\le -\deg K_F+\varepsilon\deg L_F,
 \qquad
 \deg\Theta^-\le \varepsilon\deg L_F .
\]
Both degrees are uniformly bounded.  This estimate uses no bound on
the height of the determination or the denominator of \(\mathbf M\).
For divisible \(l\), relative rational triviality supplies a rational
\(l\)-pluriform \(\sigma\) for these ordinary data.  Its divisor on the
generic fibre is \(-l\Theta\); the possible pullback term from the trait
has constant value on a slice with fixed \(e_*\).
Divide by the \(l\)-power of the log form \(\psi\) from wedging our relative \(d\log z_j\)'s with a base log volume near this DVR.
At either evaluation the discrepancy is \(v(\sigma/\psi^l)/l\) up to the base term constant at fixed \(e_*\), since \(\psi\) is a log generator at the toroidal divisors.
The generic pole degree of this ratio is \(O(l)\): the log volume of the \(z_j\) has controlled poles (on the regular codimension one locus of \(F\)), hence controlled zeros by the canonical degree (nonpositive here by log CY effectivity horizontally), and the horizontal boundary degree was just bounded.
Lemma~\ref{lem:A-function-slopes} applied to \(\sigma/\psi^l\), followed by division
by \(l\), proves the uniform metric assertion.  The integer \(l\)
may be arbitrarily large.  Continuity extends the estimate from the
two rational endpoint valuations to any two real valuations on the
slice.  Thus neither the nef-part index nor the vertical contact
enters the constant.
Horizontal traces and these canonical-degree computations on the generic fibre can use forms divided by a separating base volume (the absolute log order at a horizontal divisor is the fibre log order, by taking independent base residues).
Horizontally use exactly the same replacement and log volume over the generic field.

\end{proof}

\begin{corollary}[Horizontal variants]\label{cor:A-horizontal-variants}
The chart and metric conclusions have the following two variants.
\begin{enumerate}[label=\textup{(\roman*)}]
\item
The same tailorable bounded complement and homogeneous chart argument works if we are working over an algebraically closed generic field, with the relevant effective CYs on a \(\mathbb Q\)-factorial model of Fano type, and have boundedness only up to isomorphism in codimension one of these models.
\item
A Cartier rational b-order function \(N\) with effective trace of degree
\(O(\xi)\) on the bounded model, b-linearly \(\mathbb Q\)-equivalent
to a difference of b-nef parts both of trace degrees \(O(\xi)\), has
horizontal Lipschitz constant \(O(\xi)\), including when \(\xi\to0^+\).
\end{enumerate}
\end{corollary}

\begin{proof}
For \textup{(i)}, tailor the complement on the Fano type model using
the klt perturbation in Lemma~\ref{lem:A-lc-place-tools}, and then push
the complement and the effective CY data crepantly to the bounded
normal small model.  Their total divisors are \(\Q\)-linearly trivial.
This supplies the bounded embedded marked family and all the degree
estimates in Proposition~\ref{prop:A-uniform-charts}; the
anticanonical degree is bounded by the same curve slicing argument.

For \textup{(ii)}, use the two-point construction for each of the two
b-nef parts, taking the pushed ample errors to be \(o(\xi)\).
As above replace those parts by general effective representatives after the same arbitrarily tiny ample addition on a high determination, avoiding both centres.
The residual rational principal term for \(N\) now uses functions with pole degree divided by denominator \(O(\xi)\) on \(F\); apply Lemma~\ref{lem:A-function-slopes} (using constancy if pole degree \(<1\)).
\end{proof}

\subsection{Integral approximation on the cones}

The charts now give a bound on how discrepancy changes with the angular
coordinates, even when the trait contact is large.  We next approximate
real chart points by integral ones.  Integrality makes the resulting
orders positive integral multiples of prime orders, so the small-prime
hypotheses in the comparison lemmas can be used without changing their
normalization.

The following approximation does not require uniform smoothness over the
DVR.  All norms are fixed Euclidean norms; the matrices range over a
fixed finite list, so their kernel lattices have uniformly equivalent
coordinate norms.

\begin{lemma}[Integral approximation]\label{lem:A-lattice-approximation}
\begin{enumerate}[label=\textup{(\roman*)}]
\item
Let \(u_i\) lie in a fixed nonnegative orthant and
\(\lvert u_i\rvert\to\infty\).  After passing to a subsequence there
are \(s_i\to0^+\), nonzero nonnegative integral vectors \(n_i\), and
a fixed, possibly subsequence-dependent, \(C_*\) such that
\begin{equation}\label{eq:A-horizontal-approximation}
 \lvert n_i-s_i u_i\rvert\le C_*s_i,
 \qquad s_i\lvert u_i\rvert\longrightarrow\infty .
\end{equation}
\item
Let \(M_i x_i=r_i\), where \(x_i\ge0\), \(r_i\) is integral and
\(M_i\) belongs to a fixed finite list of integral matrices.
If \(\alpha_i>0\) tends to zero, then after passing to a subsequence
there are positive integers \(E_i\), nonnegative integral \(n_i\), and
a fixed \(C_*\) such that
\begin{equation}\label{eq:A-vertical-approximation}
 M_i n_i=r_i E_i,\qquad
 s_i=E_i\alpha_i\longrightarrow0^+,\qquad
 \lvert n_i-E_i x_i\rvert\le C_*s_i .
\end{equation}
\item
For every sufficiently small fixed \(\gamma>0\), the same matching
equations admit \(n_i\ge0\) integral and \(1\le E_i\le J(\gamma)\)
with
\begin{equation}\label{eq:A-bounded-dirichlet}
 \lvert n_i-E_i x_i\rvert\le\gamma ,
\end{equation}
where \(J(\gamma)\) is uniform in \(i\) and in the contact vectors
\(r_i\).
\end{enumerate}
\end{lemma}

\begin{proof}
For \textup{(i)}, pass diagonally so that for every integral character
\(\gamma\), the numbers \(\langle\gamma,u_i\rangle\) either converge
in \(\R\) or diverge in absolute value.  The characters with bounded
values form a subgroup \(\Gamma\).  Their limits are additive and
extend to evaluation on a vector \(u_*\).
For a fixed interval \(I=[a,b]\subset(0,\infty)\), the probability
distribution of \(s(u_i-u_*)\) modulo the lattice, with \(s\) uniform
on \(I\), tends to Haar measure on \(\Gamma^\perp\).
Indeed its Fourier coefficient at a character in \(\Gamma\) tends to
one; for a character outside \(\Gamma\), integration of the
exponential bounds the coefficient by
\[
 \frac{2}{2\pi(b-a)
       \lvert\langle\gamma,u_i-u_*\rangle\rvert}\longrightarrow0.
\]
Trigonometric approximation proves weak convergence.
Every neighborhood of zero in \(\Gamma^\perp\) has positive Haar
measure.  Choose intervals tending to zero sufficiently slowly and
neighborhoods of radius \(o(s_i)\), while also requiring
\(s_i\lvert u_i\rvert\to\infty\).
A diagonal choice gives an integral \(n_i\) with
\(\lvert n_i-s_i(u_i-u_*)\rvert=o(s_i)\), proving
\eqref{eq:A-horizontal-approximation}.
The error is eventually less than one in each coordinate.
Since \(s_i u_i\ge0\), an integral coordinate of \(n_i\) cannot be
negative.  Its norm tends to infinity, so \(n_i\ne0\).

For \textup{(ii)}, pass to a subsequence with \(M_i=M\).
Elementary integer linear algebra gives a rational \(x_{0,i}\) with
\(Mx_{0,i}=r_i\) and \(l x_{0,i}\) integral for a positive integer
\(l\) depending only on \(M\).
For example, solve using a fixed maximal nonzero minor and take \(l\)
divisible by its determinant.
Put \(H_0=\ker M\), \(\Lambda=H_0\cap\Z^m\), and
\(U_i=x_i-x_{0,i}\).
On the dual lattice of \(\Lambda\), pass diagonally so that every
normalized distance
\[
 \frac{\operatorname{dist}
   (\langle\gamma,lU_i\rangle,\Z)}{l\alpha_i}
\]
converges to a finite value, with a signed residual, or tends to
infinity.  The characters in the first case form a subgroup
\(\Gamma\).  Their signed residual limits are additive; choose
\(u_*\in H_0\) whose evaluations give these limits.

Sample positive integers \(j\) for which
\(s=jl\alpha_i\in[a,b]\), with \([a,b]\subset(0,\infty)\) fixed.
The distributions of \(jl(U_i-\alpha_i u_*)\) in
\(H_0/\Lambda\) tend to Haar measure on \(\Gamma^\perp\), which may
be disconnected.  For \(\gamma\in\Gamma\) the frequency, reduced
modulo \(\Z\), is \(o(\alpha_i)\), so its Fourier coefficient tends
to one.  Otherwise the geometric-series bound is at most
\[
 \frac{C\alpha_i}
 {\operatorname{dist}
   (\langle\gamma,l(U_i-\alpha_i u_*)\rangle,\Z)}
 \longrightarrow0 .
\]
Use neighborhoods of zero and a slow diagonal choice as before.
There are \(j_i>0\) and \(k_i\in\Lambda\) such that, with
\(s_i=j_i l\alpha_i\to0\),
\[
 \lvert k_i-j_i l(U_i-\alpha_i u_*)\rvert=o(s_i).
\]
Set \(E_i=l j_i\) and \(n_i=E_i x_{0,i}+k_i\).
These are integral matches and
\(\lvert n_i-E_i x_i\rvert\le C_*s_i\).
The same coordinate argument proves \(n_i\ge0\).

For \textup{(iii)}, apply simultaneous Dirichlet approximation to
\(lU_i\) on the fixed kernel torus.  It supplies a positive
\(j_i\le J_0(\gamma)\) with
\(\operatorname{dist}(j_i lU_i,\Lambda)\le\gamma\).
Use the same definition of \(n_i\) and \(E_i=l j_i\).
Matching is exact, positivity follows when \(\gamma<1\), and the
bound on \(E_i\) is independent of \(r_i\).
\end{proof}

In a horizontal discrepancy application take \(u_i=w_i/a_0(w_i)\).
In a vertical application of \eqref{eq:A-contact-cone}, take
\(x_i=w_i/e_{*,i}\) and
\(\alpha_i=a_0(w_i,e_{*,i})/e_{*,i}\).
The constant \(C_*\) in the vanishing-error conclusions is allowed to
depend on the selected subsequence; the bound \(J(\gamma)\) in
\eqref{eq:A-bounded-dirichlet} is uniform.

\subsection{Proofs of the comparison statements}

We apply the two preceding tools in tandem.  The tailored complement
places a minimizing divisor in a bounded chart; integral approximation
then replaces it by a prime order of small discrepancy.  This turns a
failure of comparison on the base into a violation of the corresponding
hypothesis upstairs.

\begin{proof}[Proof of Lemma~\ref{lem:A-comparison}]
We induct on the dimension of the total space; the birational case was
settled in Subsection~\ref{subsec:A-mori}.

\emph{Extraction and effectivity.}
Extract using \(a_0\) with \(\delta_*\) fixed small and take the Mori
step \(X_1\to W\).
Effectivity of all parameters on the extraction holds by monotonicity, or in the two-sided case by the ratio hypothesis taking \(\delta_*<\min(\delta,1)/3\), after trimming \(L_i\to\infty\) more slowly.
Indeed concavity using the opposite endpoint of a fixed interval controls interior upper ratios too.
For example, if \(0\le t\le T\), then concavity at zero gives
\[
 a_t(T')\le (1+t/T)a_0(T')-(t/T)a_{-T}(T').
\]
The hypothesis on all test subsequences makes the right side
\((1+o(1))a_0(T')\), uniformly on the extracted divisors.
The other sign is identical.  Apply this for each fixed rational
\(T\), then choose a slowly increasing interval by diagonalization.
At subsequent Mori steps the discrepancy families are discriminants of
the original family.  They need not have a common global piecewise
affine presentation.  For each fixed prime \(P\), however, resolve the
original family over its generic DVR and include the fibre in the snc
support.  In the trimmed range, original divisorial valuations trivial
on the current function field already have positive discrepancy.
The discriminant is therefore the minimum of the finitely many vertical
component ratios on that resolution.  Minima compose under the nested
field restrictions, so the finite rational parameter transitions used
below persist at each step.  Generic klt over the next base can likewise
be tested on a resolution of the original family over its generic field.

\emph{Generic klt over the next base.}
We claim generic klt over \(W\) in the Mori step for all bounded parameters eventually.
Otherwise take a sign and a first horizontal lc threshold \(t_*\ne0\) with \(|t_*|\) bounded, and a horizontal minimizing divisorial \(T\).
This first threshold is rational by resolution; if tested on the original family, a place giving it restricts nontrivially to the current field (by the previous generic klt), hence gives an lc place there too.
Tailor the chart using \(a_{t_*}\).
Concavity at this place gives
\[
              a_{\pm L}\le (1-L/|t_*|)a_0
\]
for arbitrarily large fixed rational \(L\), sign as chosen.
Let \(w\) be the chart vector representing \(T\).
Put \(u=w/a_0(w)\), and bound \(\lvert t_*\rvert\) by a fixed
\(T_0\).  If \(\lvert u\rvert\) were bounded along a subsequence,
the uniform homogeneous Lipschitz bound would bound
\(\lvert a_{\pm L}(u)\rvert\) independently of fixed \(L\), whereas
\(a_{\pm L}(u)\le1-L/T_0\).  Choosing \(L\) large gives a
contradiction.  Thus, after a further subsequence,
\(\lvert u\rvert\to\infty\).
Lemma~\ref{lem:A-lattice-approximation}\textup{(i)} gives
\(\lvert n_i-s_i u_i\rvert\le C_*s_i\), \(n_i\ne0\), and hence
\[
 0<a_0(n_i)\le (1+CC_*)s_i,\qquad
 a_{\pm L}(n_i)\le(1-L/T_0+CC_*)s_i .
\]
The same \(C\) works for every effective CY discrepancy
\(a_{\pm L}\).
Choose \(L>T_0(1+CC_*)\) after fixing the subsequence.
The restriction of \(n_i\) is \(k_i\ord_{T_i}\), with
\(k_i\ge1\).  Therefore \(a_0(T_i)\to0\), while
\(a_{\pm L}(T_i)<0\), contradicting the appropriate ratio
hypothesis.
This proves the claim; trim again to keep generic klt throughout and induce effective generalized data on \(W\).

\emph{Transfer of the ratio hypothesis.}
We must pass the corresponding ratio hypothesis to \(b_t(P)\) on \(W\).
Suppose there is failure with \(b_0(P)\to0\), respectively with \(b_0(P)<\delta'\) where \(\delta'>0\) is fixed sufficiently small below.
Concavity (also using the other sign in the two-sided case) gives a bounded parameter with \(b_t\le0\).
Move from zero along a sign to the first \(b_t=b_0/2\); we may take \(b_s\le2b_0\) up to that time: if there is a higher peak use the opposite sign, where the drop is then forced.
In the one-sided case this uses monotonicity instead.

There is a rational parameter \(p\) on this segment and a simultaneous minimizer over \(P\) with
\[
        a_0(T)/e(T)\to0\quad
        (\text{respectively }\ a_0(T)/e(T)\le 10\delta'),\qquad
        a_p(T)\le \tfrac12 a_0(T).
\]
Indeed take a cutoff radius \(\sqrt{b_0}\) or \(8\delta'\).
At zero all minimizers have ratio \(a_0/e=b_0\).
At the half-drop any minimizer below the cutoff works.
If minimizers exit the cutoff first, at a transition interpolate where a minimizer below and one at/above tie, to obtain a ratio near the cutoff, where the desired discrepancy drop holds since \(a_p/e=b_p\le 2b_0\).
It suffices to track the finite minimizing candidates from log resolution (or the divisorial restrictions of the finite candidates minimizing in the original family at each parameter; restriction of any minimizing original candidate gives a current minimizer by composition); use the LC-place connected interpolation above at the rational transition parameter.

Tailor the vertical chart to this \(T\).
In its notation \(\alpha=\rho\,a_0(T)/e(T)\) and concavity gives a decrease as above by arbitrarily large factors at a fixed signed \(L\).
If \(\alpha\to0\), kernel approximation and uniform Lipschitz at fixed uniformizer value give again
\[
                   a_0(n,E)=O(s),\qquad a_{\pm L}(n,E)<0
\]
for sufficiently large fixed \(L\), impossible.
Otherwise we are in the two-sided case with \(\alpha\le C\delta'\) but bounded below after passage to a subsequence.
Use the bounded Dirichlet approximation with \(\gamma\) so small that the Lipschitz error for \(a_0\) is \(<\delta/2\).
The constants here (also ramification and \(J\)) are uniform for this Mori step, so choose \(\delta'\) small enough that \(E\alpha<\delta/2\).
A sufficiently large signed \(L\) fixed \emph{on the subsequence} still makes the discrepancy negative by the uniform Lipschitz bound even for \(a_{\pm L}\).
This contradicts the strong ratio hypothesis.
We have thus transferred that hypothesis with fixed \(\delta'>0\) (and the one-sided hypothesis without a fixed cutoff as required).
Iterate using \(\dim W<\dim X\), stopping if birational over \(S\).
This proves Lemma~\ref{lem:A-comparison}.

\end{proof}

\begin{proof}[Proof of Lemma~\ref{lem:A-pullback-phases}] Absorb the input multiple into \(H\).
Use the same Mori and discriminant steps with fixed klt CY \(a\).
For \(b(P)\to0\) at a step choose a minimizer, so in the vertical chart \(\alpha=a(w,e_*)/e_*=\rho b(P)\to0\).
Here write \(h=v_P(H_W)\) where \(H_W\) is the base pullback.
Use the kernel argument on the augmented torus with sampling vector \(l(U,\rho h)\).
Write \((u_*,z_*)\) for an extension of the signed limits divided by \(l\alpha\) on the bounded-frequency subgroup \(\Gamma\).
The same sampling equidistributes modulo the augmented lattice after subtracting \(s(u_*,z_*)\).

Every vector in the Lie space of \(\Gamma^\perp\) must have last coordinate zero.
Otherwise choose its sign with negative last coordinate and target a small multiple \(\xi v\) with \(\xi\to0^+\) on that torus.
By taking intervals \(0<s\ll\xi\to0\) sufficiently slowly and approximating the target to \(o(\xi)\) along a diagonal subsequence, the integral matched valuation has angular error \(O(\xi)\), hence \(a(n,E)=O(\xi)\), but signed phase \(\xi v_{\rm last}+o(\xi)\), contradicting the lower sign bound.
The sign bound applies to integral multiples whenever their total
discrepancy tends to zero.  Indeed, for \(v_i=k_i\ord_{T_i}\) with
\(a(v_i)\to0\), the primitive discrepancies also tend to zero.
If their phase lifts are
\(r_i=O(a(T_i))\), \(r_i\ge-\varepsilon_i a(T_i)\), with
\(\varepsilon_i\to0\), then \(k_i r_i\) are lifts with
\[
 k_i r_i=O(a(v_i)),\qquad
 k_i r_i\ge-\varepsilon_i a(v_i).
\]
They tend to zero, so no phase wrapping affects the sign.
The projection of the compact subgroup \(\Gamma^\perp\) to its
last circle now has zero Lie algebra, hence finite image.
Some \(q>0\) annihilates this image; by character duality
\((0,q)\in\Gamma\).  Consequently
\({\rm dist}(q l\rho h,\mathbb Z)=O(l\alpha)\).
Its normalized residual limit is nonnegative: indeed \(z_*\ge0\) by the same sign test targeting zero on the torus with precision \(o(s)\), angular error \(O(s)\).
Pass to a subsequence with \(l,\rho\) fixed.

For every test sequence \(b(P)\to0\) we thus have a subsequence and a fixed multiple there giving both the \(O(b)\) bound and the lower \(-o(b)\).
There exists a \emph{single} fixed multiple giving the upper size bound
for all such sequences.  Otherwise choose diagonal violations with
\[
 {\rm dist}(k!h,\Z)>k!\,k\,b(P),\qquad b(P)\longrightarrow0.
\]
The subsequence property gives a fixed good multiple \(M\) on a
further subsequence, with
\({\rm dist}(Mh,\Z)\le Cb(P)\).
For large \(k\), \(M\mid k!\), and thus
\[
 {\rm dist}(k!h,\Z)
 \le (k!/M){\rm dist}(Mh,\Z)
 \le (k!/M)C b(P),
\]
a contradiction.
Fix such a multiple \(M_0\).  If its small lift has a negative
ratio to \(b(P)\) bounded away from zero along a subsequence, apply
the subsequence property again and take a common multiple of \(M_0\)
and the resulting good multiple.  Both small lifts, multiplied by
fixed positive integers, tend to zero and represent the same phase,
so eventually they are equal.  Their incompatible lower signs
give a contradiction.  This proves the required lower
\(-o(b(P))\) bound for \(M_0\).
This gives the required induction and proves Lemma~\ref{lem:A-pullback-phases}.

\end{proof}

\begin{proof}[Proof of Lemma~\ref{lem:A-linear-chamber}] Use the Mori step extracting by \(a+t_0d\).
Since \(a\) has bounded denominator and \(d\ge0\), all extracted divisors with this discrepancy tiny have \(a=0,\ d\) tiny.
Thus in the extra-data case too, effectivity on the extraction holds uniformly, by the upper bounds on \(a=0\).

\begin{claim}[Quantized chart cells]\label{claim:A-quantized-cells}
The bounded-denominator functions \(a,h\), when present, have only a
bounded finite set of possible slopes in a horizontal chart, or in
integral kernel coordinates on a vertical slice \(e_*=1\).
Their vertical affine constant terms belong to a rational grid with
bounded denominator, although their magnitudes need not be bounded.
Their common cells are cut out by a fixed finite set of possible
normal vectors, with offsets of bounded denominator.  They are
pointed, their vertices have bounded denominator, and they satisfy a
uniform polyhedral distance bound.
\end{claim}

\begin{proof}[Proof of the claim]
Horizontally, the functions of bounded denominator \(a,h\) have homogeneous linear pieces with slopes in a fixed bounded finite rational grid: their Lipschitz constants are bounded using values of effective CY discrepancies at finitely many fixed parameters, and their values on integral vectors lie in a prescribed rational grid, which controls the linear form on any full-dimensional open conical piece (take lattice differences there).
Thus their cells can be cut by the comparisons of those finitely many linear forms.
Vertically work on \(e_*=1\), writing \(x=x_0+U\) as before in an integral kernel basis.
Linear parts in \(U\) again belong to a bounded finite rational grid, and the possibly unbounded constant terms at \(x_0\) to a rational grid with bounded denominator (test lattice differences on the open \emph{homogeneous} pieces, using integrality after multiplying \(x_0\) by \(l\)).
Coordinates forced zero on the slice can first be dropped (use \(x_0\) for the reduced system), giving full relative dimension in \(x_0+\ker M\).
More explicitly, an interior of linearity there lifts by homogeneity to an open in the linear space of matches.
Take a lattice point interior to this homogeneous piece and translate a sufficiently large positive multiple of it by each of \((l x_0,l)\) and the kernel basis vectors (with last coordinate zero).
Both evaluations in each difference are on integral matches in the same linearity piece, which gives the intercept and kernel-slope quantization regardless of the size of \(x_0\).
Closures of cells can be cut by inequalities from a finite set of normals, with offsets of bounded denominator: compare pairs of the affine forms occurring and include nonnegativity constraints.
Values are affine on these comparison cells by continuity.
Vertices have bounded denominator, and the polyhedra are pointed.
We also use the elementary polyhedral error bound (Hoffman's bound): distance to a nonempty polyhedron with rows in a fixed finite list is bounded by a constant times the maximum violation.
Here is the distance estimate explicitly.  Write a nonempty polyhedron
as \(P=\{y:R_\nu y\le b_\nu\}\), with rows \(R_\nu\) in the
fixed finite list.  If \(y\) is the nearest point of \(P\) to \(x\),
express \(x-y\) in its normal cone using independent active rows:
\[
 x-y=\sum_{\nu\in I}\lambda_\nu R_\nu^{\mathsf T},
 \qquad \lambda_\nu\ge0,\qquad
 \sum_{\nu\in I}\lambda_\nu\le C\lvert x-y\rvert .
\]
The constant is uniform because only finitely many independent row
sets can occur.  Taking the scalar product with \(x-y\) gives
\[
 \lvert x-y\rvert^2
   =\sum_{\nu\in I}\lambda_\nu(R_\nu x-b_\nu)
   \le C\lvert x-y\rvert
      \max_\nu(R_\nu x-b_\nu)_+.
\]
Cancel \(\lvert x-y\rvert\) when it is nonzero.
\end{proof}

\emph{Horizontal signs and generic klt.}
First in the extra-data case, \(h<0\) on a horizontal \(a\)-lc place over \(W\) is impossible eventually.
Tailor using \(a\).  The homogeneous comparison cells for \(a,h\)
come from finitely many rational cones.  On a cell containing a
witness, the face \(a=0\) is rational and \(h\) is linear.
Some extremal ray of that face has \(h<0\).  The possible primitive
integral ray generators form a bounded finite set, so there is a
bounded nonzero integral witness \(a=0,\ h<0\).
But \(c\) on it is bounded by Lipschitz (again by fixed parameters), and \(h\) quantized and \(\ge-\eta_i c\), impossible.
Next generic klt for \(a+\lambda(h+\theta c)\) holds uniformly for small \(\lambda>0\), \(0<\theta\le\theta_0\): otherwise take \(\lambda_i\to0\) equal to a first positive lc threshold with the corresponding fixed \(\theta_i\) in each example.
There is such a first threshold by common resolution, since on horizontal \(a=0\) we have \(h\ge0,\ c>0\).
Tailor using this threshold place.
In the homogeneous chart distance from its vector \(w\) to \(a=0\) (including the zero vector) is bounded by \(C_1 a(w)\) by the slope/polyhedral bound.
On that set \(h+\theta_i c\ge0\), by the divisorial conditions and rational approximation there.
Thus \(h+\theta_i c\ge-C_2 a(w)\) at \(w\), inconsistent with the threshold tending to zero (\(a(w)>0\)).

\emph{Lexicographic minima and their denominators.}
For \(P\) over \(W\) denote the claimed lexicographic minima by \(a',d'\) (and \(a',h',c'\)).
They exist: on a common resolution over the DVR use the vertical components and the formula for discrepancies by monomial orders plus the nonnegative reduced-log discrepancy.
Horizontal \(a\) is nonnegative, and in the extra-data case on horizontal \(a=0\), \(h\) is nonnegative and \(c>0\).
Explicitly include the fibre among the divisors on the resolution, so \(e\) is a nonnegative linear combination of vertical orders only.
At a generically klt parameter a minimizer has reduced-log discrepancy zero (on the resolution) and all its positive orders in those coordinates lie on minimizing vertical components, by subtracting the minimum times \(e\); thus tracking these finitely many components captures departures from the faces in the argument.
For each example and \(P\), sufficiently small strictly positive parameters in lexicographic order thus have only lexicographically minimizing components, and give a simultaneous minimizer for tailoring using a generically klt parameter.
On that chart the minimum \(a/e_*\) is \(\rho a'\), and likewise for \(h\) on the minimizing locus if used.
By linear optimization at cell vertices on the slice \(e_*=1\), both of these minima have bounded denominator.
Thus so do \(a'\) and \(h'\), uniformly in \(P\).

\emph{Uniformity of the chambers.}
If minimizers of \(a+t d\) leave the \(a\)-minimum for some \(t\) tending to zero, take the first (strictly positive rational) transition where minimizers of both types tie.
Connected interpolation gives a simultaneous minimizing place with \((a/e)-a'(P)\) arbitrarily small positive.
Tailor at that place for the transition discrepancy.
On the cell of its normalized weights (value \(e_*=1\)), the affine
function \(a-\rho a'(P)\) is nonnegative.  Its minimum is attained at
a vertex, because the cell is pointed.  Both the vertex values of
\(a\) and \(\rho a'(P)\) lie in a fixed rational grid.  A strictly
positive minimum would therefore be bounded below by a fixed positive
number.  Taking the interpolated excess smaller than this number forces
the face \(a=\rho a'(P)\) in that cell to be nonempty.
Let \(\varepsilon=a(w)-\rho a'(P)>0\).
The distance estimate in Claim~\ref{claim:A-quantized-cells}
gives \(w_0\) on that face with
\(\lvert w-w_0\rvert\le C_1\varepsilon\).
If \(C_2\) is the Lipschitz constant of \(d\), then
\[
 (a+td)(w_0)-(a+td)(w)
 \le-\varepsilon+tC_1C_2\varepsilon<0
\]
for a uniformly small \(t>0\).  This contradicts minimality;
continuity permits the comparison at real weights and rational
approximation recovers divisorial tests.
This proves uniform exact minimization on the \(a\)-face for \(a+t d\).
Identical reasoning with \(h+\theta c\) yields uniform small \(\lambda\) with \emph{all} minimizers still on the \(a\)-face for every \(0<\theta\le\theta_0\) in the extra-data case.
Fix a uniformly positive small such \(\lambda\).
If first positive \(\theta\)-transitions away from the \(h\)-minimum within the \(a\)-minimum approach zero, now interpolate with tiny positive excess of \(h/e\) above \(h'\); all simultaneous minimizers stay on the \(a\)-minimum.
Use joint \(a,h\) cells with the additional affine face \(a=\rho a'\); the grid and distance argument within this face applied to \(h\) contradicts minimality for \(h+\theta c\).
Within the \(a\)-minimum the minimizer condition no longer depends on the magnitude of \(\lambda>0\), only on \(\theta\).
Thus the full formulas hold uniformly on a smaller sector.

\emph{Induction on the base.}
For induction \(a'+t d'\) gives klt data for positive \(t\) by the minimum formula, and data at zero by the endpoint limit already explained.
In the extra-data case on \(a'=0\), \(d'\le c'\), \(h'\ge-\eta_i c'\), and testing at a \(d'\)-minimizer at fixed uniform small \(\theta>0\) gives \(h'+\theta c'\le C_3 d'\) by the exact formula.
Hence \(c'\le C_4 d'\) and \(h'\le C_3 d'\) there.
All hypotheses now hold at the next step.
Any negative \(h\) at \(a=0\) vertical in this step marks a negative \(h'\) at \(a'=0\), by restriction.
Induct to the base to obtain all assertions. \end{proof}

\section{Nonvanishing from signed phases}\label{sec:B}

We seek sections of a bounded positive multiple of a rational
polarization, even when its denominators are unbounded.  The hypothesis
will control the fractional orders of the polarization at primes whose
discrepancies tend to zero.  The sign of that control will determine
which corrections can be removed in the proof.

For a rational Weil divisor $H$ on a normal variety, set
\begin{equation}\label{eq:B-ramification-boundary}
 R(H)=\sum_P\left(1-\frac1{r_P}\right)P,
 \qquad
 r_P=\min\{r\in\Z_{>0}:r\operatorname{coeff}_P(H)\in\Z\}.
\end{equation}
Only finitely many terms are nonzero.  When $H$ is $\Q$-Cartier, its
order at a higher divisor means the order of its Cartier b-pullback,
rather than a rounded trace.  A \emph{signed phase} of this order is a lift of its image in
$\R/\Z$; whenever this lift tends to zero, it is eventually the unique
lift in $(-1/2,1/2)$.

We use the comparison results of Section~\ref{sec:A} and two boundedness
theorems.  Rationally connected projective varieties carrying effective
generalized $\delta$-lc Calabi--Yau data in bounded dimension are bounded
up to isomorphism in codimension one
\cite[Theorem~1.7]{BirkarCY}.  No bound on the Cartier index of the nef
b-divisor occurs in that theorem.  We also use effective birationality
for a nef and big integral Weil divisor $N$ on a projective
$\delta$-lc variety when $N-K$ is pseudo-effective
\cite[Theorem~1.1]{BirkarPolarised}.  Our applications of the latter
theorem take place on $\Q$-factorial models, so that $N$ is
$\Q$-Cartier, but need not be Cartier.

\begin{lemma}[Phase nonvanishing]\label{lem:B-phase-nonvanishing}
Let $(V_i,H_i)$ be a sequence in bounded dimension, where $V_i$ is
projective of Fano type and $H_i$ is a nef and big $\Q$-Cartier rational
Weil divisor.  Suppose that, for rational $t\in[0,L_i]$, where
$L_i\to\infty$, there are generalized data with discrepancies $a_i(t)$
such that
\begin{equation}\label{eq:B-family}
 B_i(t)\ge R(H_i),\qquad
 K_{V_i}+B_i(t)+\mathbf M_i(t)_{V_i}\sim_{\Q}tH_i,
 \qquad \mathbf M_i(t)\text{ is b-nef and $\Q$-Cartier}.
\end{equation}
Assume that $a_i(0)$ is klt and that $a_i(t)$ is nondecreasing in $t$
at every divisorial valuation.

Assume the following condition, also after restriction to any
subsequence.  For every choice of rational $t_i\to\infty$ in the
parameter range, one can pass to a further subsequence and choose a
fixed positive integer $q$ such that, for every sequence of primitive
divisorial valuations $T_i$ with $a_i(t_i,T_i)\to0$, the signed phases
$\theta_i(T_i)$ of $qH_i(T_i)$ satisfy
\begin{equation}\label{eq:B-phase-hypothesis}
 -o\bigl(a_i(t_i,T_i)\bigr)
 \le \theta_i(T_i)
 \le O\bigl(a_i(t_i,T_i)\bigr).
\end{equation}
Then there is an integer $M>0$ such that for every $i$ some
$m_i\in\{1,\ldots,M\}$ satisfies
$|\lfloor m_iH_i\rfloor|\ne\varnothing$.
Equivalently, after replacing $M$ by a common multiple, the same
positive degree works for all $i$.
\end{lemma}

All the constants in this lemma are sequence bounds.  In particular,
the phase condition is tested along arbitrary subsequences and with
arbitrarily chosen divisors.  This quantifier has two useful
consequences.  First, at a fixed growing parameter and after the
allowed scaling, there are $C>0$ and a sufficiently small fixed
$\eta>0$ such that all phases at $a_i(t_i,T)<\eta$ have their small
lifts bounded in absolute value by $Ca_i(t_i,T)$, eventually.
Otherwise one selects $a_i(t_i,T_i)\to0$ with an unbounded ratio.
Second, the condition also applies to integral multiples
$w_i=e_i\ord_{T_i}$ whenever $a_i(t_i,w_i)\to0$.  Indeed,
$a_i(t_i,T_i)\to0$, and
\begin{equation}\label{eq:B-multiple-phase}
 e_i\theta_i(T_i)=O\bigl(a_i(t_i,w_i)\bigr)\longrightarrow0,
 \qquad
 e_i\theta_i(T_i)\ge-o\bigl(a_i(t_i,w_i)\bigr).
\end{equation}
Thus multiplication introduces no wrap around the phase circle.  This
observation is needed each time a divisorial valuation restricts to a
multiple of an angular prime.

\subsection{Seifert bundles and two consequences of comparison}

We use the graded-algebra description of Seifert bundles and its
compactification \cite[Definition~5, Theorem~7 and Section~14]{KollarSeifert}.
The discrepancy and phase formulas needed here are proved below.

\begin{lemma}[Bundle construction and discrepancies]\label{lem:B-seifert}
Let $V$ be projective normal, and let $H$ be a $\Q$-Cartier rational
divisor.  There is a projective normal compactified Seifert bundle
$p:Y\to V$ with disjoint $\Q$-Cartier zero and infinity sections
$E_0,D$ and a rational fibre coordinate $z$ such that
\begin{equation}\label{eq:B-bundle-identities}
 \divisor(z)=p^*H+E_0-D,
 \qquad
 K_Y+E_0+D=p^*\bigl(K_V+R(H)\bigr).
\end{equation}
The second equality uses compatible canonical divisors.  The divisor
$D$ is relatively ample; if $H$ is nef, then $D$ is nef, and if $H$
is nef and big, then $D+p^*H$ is nef and big.

In these identities, pullback of a rational Weil divisor along the
equidimensional map $p$ is defined at the generic points of its prime
components, using the multiplicities of their inverse images.  This
agrees with Cartier pullback whenever the divisor is $\Q$-Cartier.
Orders of $H$, $E_0$, and $D$ on higher models use their Cartier
b-pullbacks.

Suppose $B\ge R(H)$ and $\mathbf M$ is a nef $\Q$-Cartier
b-divisor on $V$.  Give $Y$ boundary $p^*(B-R(H))$ and nef part
$p^*\mathbf M$.  Write $a$ for the generalized discrepancy on $V$
and $A$ for that on $Y$.  If a primitive divisorial valuation $v$ of
$k(Y)$ restricts nontrivially to $w$ on $k(V)$, then
\begin{equation}\label{eq:B-bundle-discrepancy}
 A(v)=a(w)+|h(v)|+\operatorname{def}(v),
 \qquad
 h(v)=v(E_0)-v(D)\equiv-w(H)\pmod\Z,
\end{equation}
where $\operatorname{def}(v)$ is a nonnegative integer.  It vanishes
precisely for the Gauss, or fibre-torus-invariant, valuations.
If $w=0$, the primitive valuations are the ordinary places of the
generic $\mathbb P^1$, and have discrepancy $1$.
Consequently klt base data give klt data on $Y$ without the two
sections; adding both sections gives lc data whose non-klt locus is
contained in the sections.
\end{lemma}

\begin{proof}
The two affine charts are the relative spectra of the algebras of
rational monomials $gz^j$ satisfying
$\divisor(g)+jH\ge0$, with $j\ge0$ and $j\le0$, respectively.
If $\ell H=\divisor(b)$ locally, these charts are finite cyclic
quotients of the ordinary line charts obtained using $b^{1/\ell}$
and the coordinate $z/b^{1/\ell}$.  They glue to the stated normal
projective bundle.  This description gives the two sections, their
$\Q$-Cartier property, and the first equality in
\eqref{eq:B-bundle-identities}.  The restriction of $D$ to itself is
numerically $H$, while its restriction to $E_0$ is trivial; this and
relative ampleness give the positivity assertions.

At a prime $P$ of $V$ the reduced vertical divisor has multiplicity
$r_P$.  Computing the relative canonical divisor by $d\log z$
therefore gives the second equality in
\eqref{eq:B-bundle-identities}, including exactly the coefficient
$1-1/r_P$.

For the discrepancy formula, write $w=e\ord_T$ on a smooth base
model and choose a uniformizer $x$ at $T$.  Wedge a base logarithmic
volume form with $d\log z$.  The additional logarithmic order is a
nonnegative integer; it is zero exactly when the residue of the
valuation-zero monomial $z^e/x^{v(z)}$ is transcendental over
$\kappa(T)$.  This is the Gauss condition.  Removing the two
sections from the boundary adds $v(E_0)+v(D)=|h(v)|$, since their
supports are disjoint.  Evaluating the first identity in
\eqref{eq:B-bundle-identities} gives the congruence.  The same
calculation on a determination of the nef part proves the generalized
formula.  For trivial restriction the assertion is the usual
discrepancy calculation on $\mathbb P^1_{k(V)}$.
\end{proof}

In the setting of \eqref{eq:B-family}, the data without the sections
have total logarithmic divisor
\begin{equation}\label{eq:B-bundle-total}
 K_Y+p^*(B(t)-R(H))+p^*\mathbf M(t)_Y
 \sim_{\Q}(t+1)p^*H-2D.
\end{equation}
At $t=0$, adding both sections instead gives effective generalized
lc Calabi--Yau data.

\begin{lemma}[Clearing the small denominators]\label{lem:B-denominator-clearing}
In bounded dimension, suppose $V$ is projective of Fano type and
carries effective generalized klt Calabi--Yau data with
$B\ge R(H)$.  After multiplying $H$ by a bounded positive integer,
there is an ordinary klt log Fano boundary $\Theta\ge R(H)$.
For this scaled divisor the Seifert bundle is of Fano type.
The assertion does not require a bound on the Cartier index of the
nef b-divisor.
\end{lemma}

\begin{proof}
Take a small $\Q$-factorialization.  The class $-(K_V+R(H))$ is
pseudo-effective, because it is the sum of $B-R(H)$ and the trace
of a b-nef divisor.  Moreover $(V,R(H))$ is ordinarily klt:
its discrepancies dominate the given generalized discrepancies,
by $B\ge R(H)$ and b-nef negativity.  Run its anticanonical MMP,
using Fano type to obtain a nef model.  Ordinary lc comparison
holds along this MMP by the negativity argument of
Section~\ref{sec:A}.  The coefficients of
$R(H)$ belong to the fixed hyperstandard set
$\{1-1/r:r\in\N\}$.  The bounded ordinary lc complement theorem
\cite[Theorems~1.7--1.8]{BirkarComplements}, followed by pullback
through this MMP, gives an ordinary lc boundary
$\Delta\ge R(H)$ with $n(K_V+\Delta)\sim0$, where $n$ is bounded
in terms of dimension.  The pullback dominates $R(H)$ because the
MMP has the anticanonical sign; this is the same complement
pullback used in Lemma~\ref{lem:A-lc-place-tools}.

If $r_P>n$, the $1/n$-quantization of $\Delta$ forces
$\operatorname{coeff}_P(\Delta)=1$.  Multiply $H$ by the least
common multiple of the integers at most the bound for $n$.
All smaller denominators disappear.  Every remaining component of
$R(H)$ is therefore contained in $\lfloor\Delta\rfloor$.
Let $\Theta_0$ be any klt log Fano boundary on $V$.  For a
sufficiently small positive rational $\rho$, chosen separately for
each variety,
\[
 \Theta=(1-\rho)\Delta+\rho\Theta_0
\]
is klt, dominates the new $R(H)$, and has ample negative log
canonical class.  No uniform lower bound for $\rho$ is needed.

On the bundle use boundary
$p^*(\Theta-R(H))+(1-\rho')(E_0+D)$, with $\rho'>0$ sufficiently
small.  Lemma~\ref{lem:B-seifert} proves klt.  If
$A=-(K_V+\Theta)$, its negative log canonical class is
$p^*(A-\rho'H)+2\rho'D$, which is ample for small $\rho'$ by
relative ampleness and base positivity.  Thus the bundle is of
Fano type.
\end{proof}

\begin{lemma}[An unbounded-index obstruction]\label{lem:B-unbounded-index}
Let $F_i$ be projective of positive bounded dimension and let $G_i$
be ample $\Q$-Cartier rational divisors such that $(F_i,R(G_i))$
is klt and
\[
 -(K_{F_i}+R(G_i))\sim_{\Q}I_iG_i,
 \qquad I_i>0,\quad I_i\longrightarrow\infty.
\]
Let $C_i$ be the compactified Seifert cone obtained by contracting
the zero section, and denote its infinity divisor by $D_i$.
Then $C_i$ is klt of Fano type, $D_i$ is ample, and
$-K_{C_i}\sim_{\Q}(I_i+1)D_i$.  It is impossible that
\begin{equation}\label{eq:B-cone-small-height}
 v_i(D_i)=o\bigl(A_{C_i}(v_i)\bigr)
\end{equation}
for every sequence of primitive divisorial valuations with
$A_{C_i}(v_i)\to0$.
\end{lemma}

\begin{proof}
The contraction exists by semiampleness of the infinity divisor.
One can also construct it by the finite fibre-power map to the
ordinary compactified bundle for a very ample Cartier multiple of
$G_i$, contracting its zero section and then Stein factorizing.
The bundle identity gives
\[
 K_Y+(1-I_i)E_0+(1+I_i)D\sim_{\Q}0.
\]
The pair with subboundary $(1-I_i)E_0$ is sub-klt by
\eqref{eq:B-bundle-discrepancy}; the discrepancy of the contracted
zero section is $I_i>0$.  This proves the assertions about the
cone and its canonical divisor.  Near infinity the discrepancy
formula reads, for $u=v(D)>0$,
\begin{equation}\label{eq:B-infinity-formula}
 A_C(v)=A_{F,R(G)}(w)+u+\operatorname{def}(v),
 \qquad u\equiv w(G)\pmod\Z.
\end{equation}

If \eqref{eq:B-cone-small-height} held, apply
Lemma~\ref{lem:A-comparison} over a point to the nonincreasing
family with discrepancies $A_C-tv(D)$, boundary $tD$, and nef part
$(I_i+1-t)D$ on a range tending to infinity.  The phase-free
ratio hypothesis is exactly \eqref{eq:B-cone-small-height}.
The comparison lemma would make this family klt on every fixed
bounded parameter range, whereas at $t=1$ the prime $D$ has
discrepancy zero.  This is a contradiction.
\end{proof}

\begin{lemma}[Bounding a principalization order]\label{lem:B-principal-order}
Let $F_i$ be projective of Fano type in bounded dimension, and let
$P_i\sim_{\Q}0$ be $\Q$-Cartier rational divisors.  Suppose that
$F_i$ carries effective generalized klt Calabi--Yau data with
discrepancies $a_i$ and boundary at least $R(P_i)$.  Suppose,
after a fixed integral scaling if necessary, that the signed phases
of $P_i$ at every sequence with $a_i(T_i)\to0$ lie in
$[-o(a_i(T_i)),O(a_i(T_i))]$.  Then the least positive integer $q_i$
for which $q_iP_i$ is principal is bounded.

More generally, the same cyclic-cover argument applies whenever
the induced isotrivial fibration below has
$u=o(A)$ at its small-discrepancy divisors over zero.  This
formulation permits the argument to be used after an independent
estimate of the height $u$.
\end{lemma}

\begin{proof}
Apply Lemma~\ref{lem:B-denominator-clearing} to obtain, after a
further bounded scaling, ordinary klt log Fano data dominating
$R(P_i)$.  This does not affect boundedness of the original
principalization orders.  A small $\Q$-factorialization preserves
all the data and these orders.  Suppress the index and suppose
for contradiction that $q\to\infty$.  Write
\[
 \divisor(\phi)=qP,
 \qquad
 \widetilde F=\operatorname{Norm}_F\bigl(k(F)(r)\bigr),
 \qquad r^q=\phi.
\]
The degree is $q$ by minimality of $q$ and the Kummer criterion
over a field containing all roots of unity.  The ramification in
codimension one is encoded by $R(P)$.  In particular the log
pullbacks of both the ordinary log Fano boundary and the
generalized boundary are effective and klt.

Form the diagonal quotient
\begin{equation}\label{eq:B-cyclic-fibration}
 X=(\widetilde F\times\mathbb P^1_y)/\mu_q
 \longrightarrow\mathbb P^1_{y^q},
\end{equation}
where $\mu_q$ acts inversely on $y$ and $r$, so that $yr$ is
invariant.  The map from the product to $X$ is quasi-\'etale:
a nonidentity element has a fixed locus of codimension at least
one on $\widetilde F$ and only the two fixed points on
$\mathbb P^1$.  The generic fibre of \eqref{eq:B-cyclic-fibration}
is connected.  Log pulling the ordinary log Fano data to the
product and descending proves that $X$ is of Fano type over the
curve.

Pull back the generalized data to the product and descend them.
They give effective generalized klt Calabi--Yau data over the
curve.  The invariant nef b-divisor descends on a sufficiently
high finite quotient; its $\Q$-Cartier property and relative
rational triviality descend by taking norms.  For a primitive
valuation over zero put
$u=v(y^q)/q>0$, and let $w$ be its restriction to $k(F)$.
The product discrepancy calculation gives
\begin{equation}\label{eq:B-cyclic-height}
 A(v)\ge a(w)+u,
 \qquad u\equiv-w(P)\pmod\Z,
\end{equation}
where $a(w)=0$ if $w=0$.  The congruence follows by evaluating
the invariant function $yr$.  These formulas can equally be
computed upstairs and divided by the ramification index, using
log pullback and homogeneous discrepancies.

If $A(v_i)\to0$, trivial restriction is impossible because it
would make $u_i$ a positive integer.  For nontrivial restriction,
\eqref{eq:B-multiple-phase} and \eqref{eq:B-cyclic-height} give
$u_i=o(A(v_i))$: all terms tend to zero, so there is no wrap,
and the lower bound for the signed phase bounds the positive
number $u_i$ from above.

Apply Lemma~\ref{lem:A-comparison} over the curve to
$A-t\,v(f^*[0])/q$.  These are effective generalized
Calabi--Yau data obtained by adding $tf^*[0]/q$ to the boundary;
the nef part is unchanged.  The parameter range may tend to
infinity.  The reduced zero fibre has discrepancy $1$ and
multiplicity $q$, hence its initial discriminant discrepancy
satisfies $0<b_0(0)\le1/q$.  Since the perturbation is a pure
base pullback,
\[
 b_t(0)=b_0(0)-t/q,
 \qquad b_1(0)\le0.
\]
This contradicts the ratio conclusion of
Lemma~\ref{lem:A-comparison}, because $b_0(0)\to0$.
\end{proof}

\subsection{Proof of phase nonvanishing}

We now prove the main nonvanishing lemma.  The first task is to subtract
small positive phase corrections from the polarization and show that the
remaining divisor is not big.  Its subsequent MMP will either contradict
the unbounded-index obstruction or produce a lower-dimensional base to
which the same phase hypotheses descend.

\begin{proof}[Proof of Lemma~\ref{lem:B-phase-nonvanishing}]
\emph{Choosing the phase corrections.}
We argue by dimension induction.  The assertion for a point is
immediate.  If the assertion fails in positive dimension, pass to
a sequence and choose integers $N_i\to\infty$ such that
\begin{equation}\label{eq:B-no-sections}
 |\lfloor mH_i\rfloor|=\varnothing
 \quad\text{for all integers }1\le m\le N_i.
\end{equation}
Choose rational $k_i\to\infty$ sufficiently slowly relative to
$N_i$ and $L_i$.  Use the phase hypothesis at $k_i$ and
Lemma~\ref{lem:B-denominator-clearing}, pass to a subsequence,
and scale $H_i$ by fixed positive integers.  Reparametrize the
given families by these same factors, and decrease $N_i$ by
bounded factors.  We retain \eqref{eq:B-no-sections}, obtain the
phase hypothesis without a displayed scaling at $k_i$, and have
Fano type Seifert bundles.  These bounded changes never replace
the given discrepancy families by different ones.

Fix a sufficiently small constant $\eta>0$.  Extract precisely
the exceptional primes $T$ with $a(k_i,T)<\eta$ whose small signed
phase $f_T$ of $H(T)$ is positive.  Include the original primes
with the same property in the list of corrections.  This is a
finite list: monotonicity gives $a(0,T)\le a(k_i,T)<\eta$, and
the extraction and finiteness argument of
Section~\ref{sec:A} applies to the klt data $a(0)$.
Write the resulting Fano type model as $U\to V$.  For small fixed
$\eta$ we have, eventually,
\begin{equation}\label{eq:B-positive-corrections}
 0<f_T\le C a(k_i,T)\le C\eta,
 \qquad
 C_U=\sum_T f_TT,
 \qquad P_U=H_U-C_U,
\end{equation}
where $H_U$ is the pullback of $H$.
We first prove that $P_U$ is not big, eventually.

\medskip
\noindent\emph{The threshold contradiction under bigness.}

Suppose instead that $P_U$ is big.  On the Seifert bundle let $A$
be the discrepancy of the data at $k_i$ without the sections.
Choose a small fixed $0<\delta<1/4$, sufficiently small relative
to $\eta$ and the phase bound.  Extract every exceptional prime
with $A<\delta$, and let $\mathcal S$ be the reduced sum of
\emph{all} primes with $A<\delta$ on this extraction, including
original primes.  Lemma~\ref{lem:B-seifert} shows that they are
Gauss valuations.  Their extraction is funded by the effective
zero-parameter generalized Calabi--Yau data with both sections,
mixed with a sufficiently small ordinary klt log Fano
contribution as in Section~\ref{sec:A}.  Thus the extraction is
of Fano type, and the trace boundary of $A$ is effective on it.

Put $\mathbf H=p^*H$, and denote by $\mathbf D$ the b-pullback
of $D$.  The b-divisor
\[
 \mathbf J=\mathbf D+\mathbf H
\]
is nef and big and has integral Weil trace on the original
bundle: the multiplicities $r_P$ clear the coefficients of
$H$ at all vertical primes.  Let $\tau$ be the pseudo-effective
threshold for subtracting $\mathcal S$ from its trace on the
extraction.  We claim that
\begin{equation}\label{eq:B-threshold-small}
 \mathcal S\ne0,
 \qquad k_i\tau_i\longrightarrow0.
\end{equation}

To prove the claim, if it fails pass to a subsequence and choose
integers
$k_i,\tau_i^{-1}\ll m_i\ll N_i$; omit $\tau_i^{-1}$ when
$\mathcal S=0$.  Round the coefficients on the extracted primes
to obtain
\[
 Q=\sum_{T\subset\mathcal S}q_TT,
 \quad 1\le q_T\le2,
 \qquad N=m\mathbf J_U-Q
\]
big and integral.  All nonintegral coefficients of $m\mathbf J_U$
are on the extracted primes, so this rounding suffices.  When
$\mathcal S=0$, use $N=mJ$.  Run the $N$-MMP to a nef and big
model.  Here and in what follows a prime denotes the output
trace, and a star denotes its Cartier b-pullback to a common
resolution.  Set
\[
 e_H=(H')^*-\mathbf H,\qquad
 e_D=(D')^*-\mathbf D,\qquad e_J=e_H+e_D.
\]
B-nef negativity and MMP negativity give
\begin{equation}\label{eq:B-rounded-errors}
 e_H,e_D\ge0,
 \qquad
 m e_J\le(Q')^*-Q^*,
 \qquad (Q')^*\le2(\mathcal S')^*.
\end{equation}
The total logarithmic divisor for $A$ has trace-pull error
$(k_i+1)e_H-2e_D$, by \eqref{eq:B-bundle-total}.  Its effective
trace boundary dominates $(1-\delta)\mathcal S'$.  Comparing
generalized and ordinary discrepancies on the output therefore
gives
\begin{align}
 A_{\mathrm{ord}}'
 &\ge A+(1-\delta)(\mathcal S')^*-(k_i+1)e_H\notag\\
 &\ge A+
 \left(1-\delta-\frac{2(k_i+1)}m\right)(\mathcal S')^*
 \ge A+\frac12(\mathcal S')^* .
 \label{eq:B-ordinary-gap}
\end{align}
The comparison also discards only a nonnegative trace-pull error
of the nef part.

This gives a gap at contracted primes as well as surviving ones.
Indeed, every primitive $T$ with $A(T)<\delta$ was put on the
extraction, and hence $Q^*(T)=q_T\ge1$, even if the subsequent
MMP contracts $T$.  Equation~\eqref{eq:B-rounded-errors} implies
\begin{equation}\label{eq:B-contracted-prime-bound}
 (Q')^*(T)\ge Q^*(T)\ge1,
 \qquad (\mathcal S')^*(T)\ge\frac12.
\end{equation}
Such a prime has $A_{\mathrm{ord}}'(T)\ge1/4$ by
\eqref{eq:B-ordinary-gap}; all other primitive primes have
$A_{\mathrm{ord}}'(T)\ge A(T)\ge\delta$.  Thus the output is
ordinarily $\min\{\delta,1/4\}$-lc.

Fano type makes $-K'$ pseudo-effective, so
$N'-K'$ is pseudo-effective.  Integral-Weil effective
birationality \cite[Theorem~1.1]{BirkarPolarised} supplies a
nonzero section of a bounded multiple of $N'$.  MMP negativity
pulls it back to the same multiple of $N$, and then to a bounded
multiple of $m\mathbf J$.
The fibre torus acts on the latter section space, so it contains
a weight vector $gz^j$.  A section of $\ell J$ satisfies
\begin{equation}\label{eq:B-angular-section}
 0\le j\le\ell,
 \qquad \divisor(g)+(j+\ell)H\ge0.
\end{equation}
Indeed the coefficients of its divisor along $E_0,D$ are
$j,\ell-j$, respectively, by \eqref{eq:B-bundle-identities}.
Thus it gives a section of $\lfloor (j+\ell)H\rfloor$ in positive
degree at most $2\ell=O(m_i)\ll N_i$, contrary to
\eqref{eq:B-no-sections}.  This proves
\eqref{eq:B-threshold-small}.

Run the MMP of $J-\tau\mathcal S$ to its semiample nonbig model
and let its Iitaka fibration be $Y'\to Z_\tau$.  The threshold
is rational because the extraction is of Fano type.  Negativity
now gives
\begin{equation}\label{eq:B-threshold-errors}
 0\le e_J\le\tau\bigl((\mathcal S')^*-\mathcal S^*\bigr).
\end{equation}
Choose a small fixed $c$ with $0<c<1-\delta$.  On this output
use the boundary $B_Y'-c\mathcal S'$ and nef b-divisor
\begin{equation}\label{eq:B-threshold-nef-part}
 \mathbf M^\#
 =\mathbf M(k_i)+
 \left(2+\frac c\tau\right)\mathbf D+
 \left(\frac c\tau-k_i-1\right)\mathbf H.
\end{equation}
This boundary is effective.  Both displayed coefficients in the
nef part are eventually nonnegative by
\eqref{eq:B-threshold-small}; in particular the nefness is
absolute, and uses no Cartier-index bound.  Its total
logarithmic divisor is $\Q$-linearly equivalent to
$(c/\tau)(J'-\tau\mathcal S')$.  The discrepancies are
\begin{equation}\label{eq:B-modified-gap}
 A^\#=A+c\left((\mathcal S')^*-\frac{e_J}{\tau}\right)
 \ge A+c\mathcal S^*.
\end{equation}
Every primitive prime with $A<\delta$ has $\mathcal S^*(T)=1$,
including a prime contracted along the MMP.  The remaining
primes have $A\ge\delta$.  Thus $A^\#$ has a uniform generalized
gap $\min\{c,\delta\}$.

On a geometric generic fibre $F$ of $Y'\to Z_\tau$ these are
effective generalized Calabi--Yau data.  Restriction can be
computed using horizontal divisors on all higher models over
the generic field, and then separable extension of constants;
it is crepant and preserves the displayed gap.  The fibre is
of Fano type and rationally connected.  Apply
\cite[Theorem~1.7]{BirkarCY} to $A^\#$ on $F$ and take a bounded
normal model isomorphic to it in codimension one.  Normalization
of the bounded models is still bounded after stratification and
spreading, and does not add divisors to a codimension-one
isomorphism.

The effective boundary $B_Y'-c\mathcal S'$ has coefficients at
least $1-\delta-c$ on $\mathcal S'$.  Intersecting its
Calabi--Yau equality with a fixed bounded polarization bounds
its degree: the trace of the nef b-part is pseudo-effective,
and the canonical degree is bounded in the bounded family.
Consequently
\begin{equation}\label{eq:B-small-degrees}
 \deg\mathcal S'=O(1),\qquad
 \deg J'=O(\tau),\qquad
 \deg\mathbf H_F,\ \deg\mathbf D_F=O(\tau).
\end{equation}
Here $J'\sim_{\Q}\tau\mathcal S'$ on the fibre, and the last
two degrees are nonnegative and sum to that of $J'$.

We now use the assumed bigness of $P_U$.  It gives an effective
rational divisor $\Sigma\sim_{\Q}H$ such that for every $v$
listed in $\mathcal S$, with restriction $w$, one has
\begin{equation}\label{eq:B-strict-improvement}
 w(\Sigma)>v(D)-v(E_0).
\end{equation}
For completeness, if $v(D)>0$, write $w=e\ord_T$.
The discrepancy formula gives $ea(k_i,T)<\delta$ and
$v(D)<\delta$.  Since $e|f_T|\le Cea(k_i,T)$ and $\delta$ is
sufficiently small, the congruence has no wrap and
$v(D)=ef_T$.  Thus $f_T>0$, and this prime was corrected in
$P_U$.  An effective representative of the big class $P_U$
may be chosen to vanish strictly at these primes and at any of
the remaining finitely many projected centres.  Every listed
valuation has nontrivial angular restriction, since trivial
restriction would give $A=1$; its centre on $U$ is therefore proper.
Reserve a tiny ample summand and choose a sufficiently divisible
member through all these centres.  Adding back $C_U$ gives
\eqref{eq:B-strict-improvement}.  Pushing down yields an
effective representative of $H$ on $V$; the representative on
$U$ is its pullback, since it differs from $H_U$ by a rational
principal divisor.

Every homogeneous normalized threshold section has divisor
\begin{equation}\label{eq:B-threshold-section}
 (1-b)D+bE_0+p^*\Xi,
 \qquad 0\le b\le1,
 \qquad \Xi\sim_{\Q}(1+b)H.
\end{equation}
If $b<1$, replace $b$ by $b+\alpha$ and $\Xi$ by
$\Xi+\alpha\Sigma$ for a sufficiently small positive rational
$\alpha$.  The divisor remains effective, while every
$\tau\mathcal S$ constraint improves strictly by
\eqref{eq:B-strict-improvement}.  The finite list of strict
improvements would increase the pseudo-effective threshold,
a contradiction.  Hence $b=1$ for every such section.
All these systems are torus invariant, because the listed
valuations and the original $J$ are torus invariant.  Equal-degree
ratios therefore lie in $k(V)$, by their weight decompositions
and $b=1$.  The valuation
$v_0=\ord_{E_0}$ is trivial on this field, so it is horizontal
over the Iitaka base $Z_\tau$.

We obtain a contradiction on the geometric generic fibre.
Tailor a bounded ordinary lc complement preserving $v_0$ as
an lc place, using the effective zero-parameter generalized
Calabi--Yau data with the two sections and the Fano type
funding of Lemma~\ref{lem:A-lc-place-tools}.  This may first
be done on a small $\Q$-factorialization of the fibre, using
the horizontal lift of $v_0$, and then transferred to the
bounded model.  Work on its homogeneous horizontal chart.
Proposition~\ref{prop:A-uniform-charts} bounds the Lipschitz
constant of $A^\#$.  The difference $A^\#-A$ is the order of
\[
 \frac c\tau\left(\mathbf J-(J'-\tau\mathcal S')^*\right).
\]
The parenthesized b-divisor has effective trace
$\tau\mathcal S'$ of degree $O(\tau)$ on this fibre and is
b-nef up to a rational principal divisor.  The shrinking-degree
estimate of Corollary~\ref{cor:A-horizontal-variants} therefore
gives an $O(1)$ Lipschitz bound for the difference, and hence
for $A$.

The same estimate gives an $O(\tau)$ Lipschitz bound for
$h_0(v)=v(E_0)$.  Indeed $E_0$ has effective trace and is
b-linearly equivalent to $\mathbf D-\mathbf H$, whose two nef
terms have degree $O(\tau)$ by \eqref{eq:B-small-degrees}.
In applying the estimate one replaces these nef terms by
effective representatives with arbitrarily small error,
avoiding the two centres in question.  The function-degree
argument on the bounded chart then applies: a rational
function with pole degree strictly below $1$ is constant.
This is why the Lipschitz bound tends to zero, rather than
merely remaining bounded.

Let $w_0$ be a chart vector representing $v_0$.  We have
$h_0(w_0)=A(w_0)>0$.  The shrinking Lipschitz bound implies
$|w_0/A(w_0)|\to\infty$.  The horizontal approximation in
Lemma~\ref{lem:A-lattice-approximation}, applied to this ray,
gives integral divisorial vectors $n$ and numbers $s\to0^+$
such that
\[
 A(n)=O(s)\longrightarrow0,
 \qquad h_0(n)\ge s/2.
\]
On the other hand, any such small-discrepancy valuation with
$h_0>0$ has $v(D)=0$.  Formula~\eqref{eq:B-bundle-discrepancy}
and the lower signed phase bound at $a(k_i)$, including
\eqref{eq:B-multiple-phase}, imply $h_0=o(A)$; trivial
restriction would instead make $h_0$ a positive integer and
is impossible.  This contradicts the two displayed estimates.
We have proved that $P_U$ is not big.

\medskip
\noindent\emph{Running the correction MMP.}

Diagonalize the preceding conclusion to choose
$\eta_i\to0^+$ sufficiently slowly so that $P_U$ is still
not big.  Run its MMP.  All parameters used from now on may
grow more slowly, remain below $k_i$, and be $o(1/\eta_i)$.
On the output put
$\mathcal T'$ equal to the reduced sum of surviving correction
primes.  With $\mathbf H$ now denoting the original angular
b-pullback, MMP negativity and b-nefness give
\begin{equation}\label{eq:B-correction-errors}
 0\le C_U^*\le
 E:=\mathbf H-(P')^*
 \le(C')^*,
 \qquad C'\le O(\eta_i)\mathcal T'.
\end{equation}
More explicitly, if $e_H=(H')^*-\mathbf H\ge0$, then
$E=(C')^*-e_H$, and $P$-negativity gives
$(C')^*-C_U^*\ge e_H$.

\medskip
\noindent\emph{Excluding a Mori fibre output.}
Suppose first that $P_U$ is not pseudo-effective.  On the
geometric generic fibre $F$ of the final anti-$P'$ Mori
contraction, put $G=-P'$, which is ample.  Choose a new slow
parameter $s_i\to\infty$, with $s_i\eta_i\to0$, and obtain
the phase condition at $a(s_i)$ after a fixed scaling.
Scale $H,P_U,C_U$ together.  The correction coefficients
remain $O(\eta_i)$, and $P_U$ is integral at correction
primes: subtracting their phases was precisely what made
those coefficients integral.  Thus $R(G)$ is supported
away from $\mathcal T'$, and the pushed boundary satisfies
\[
 B(s_i)'\ge R(G)+(1-o(1))\mathcal T'.
\]
Comparing with ordinary discrepancies on $F$, the only
adverse total-divisor error is $s_ie_H$, which is at most
$s_i(C')^*$.  B-nef negativity for the moduli part and
\eqref{eq:B-correction-errors} therefore give
\begin{equation}\label{eq:B-negative-discrepancy}
 A_{F,R(G)}
 \ge a(s_i)+\tfrac12(\mathcal T')^*.
\end{equation}
In particular $(F,R(G))$ is klt.  A uniform ordinary gap is
not asserted or needed in this case.

The relative Picard number of the Mori contraction is one.
The classes used here, including $R(G)$, are restrictions
from this model.  Horizontal prime coefficients restrict
unchanged over the generic field and split with unchanged
multiplicities after a separable constant extension, so the
same class comparisons hold on the geometric generic fibre.
There we have
\[
 G+H'=C'\le O(\eta_i)\mathcal T',
\]
where $H'$ is pseudo-effective.  The zero-parameter
generalized Calabi--Yau equality also shows that
$-(K_F+R(G))$ dominates a fixed positive multiple of
$\mathcal T'$ in this one-dimensional space of restricted
classes: its boundary contributes $(1-o(1))\mathcal T'$,
and the moduli trace is pseudo-effective.  Consequently
\begin{equation}\label{eq:B-large-index}
 -(K_F+R(G))\sim_{\Q}I_iG,
 \qquad I_i\longrightarrow\infty.
\end{equation}
Initially this is a numerical proportionality, with
$I_i\gg1/\eta_i$.  It is rational linear proportionality
after clearing denominators because $F$ is rationally
connected: pull numerically trivial Cartier multiples to a
smooth rationally connected resolution, where they are
torsion, and descend the resulting rational equivalence.

Consider the compactified cone for $(F,G)$ and a primitive
divisorial valuation with ordinary discrepancy $A\to0$ and
infinity height $u>0$.  Its angular restriction $w$ has
$A_{F,R(G)}(w)\le A$ and $u\le A$ by
\eqref{eq:B-infinity-formula}.  Moreover
\eqref{eq:B-correction-errors} and
\eqref{eq:B-negative-discrepancy} imply
\begin{equation}\label{eq:B-negative-error-small}
 0\le E(w)\le O(\eta_i)A_{F,R(G)}(w)\le O(\eta_i)A.
\end{equation}
Divisorial tests on the geometric fibre lift to homogeneous
horizontal tests on the original model.  Thus the signed
phase bound for $H$ applies, and
\[
 u\equiv w(G)\equiv-w(H)+E(w)\pmod\Z.
\]
All three quantities have small signed representatives.
There is no wrap, so the lower phase bound and
\eqref{eq:B-negative-error-small} give $u=o(A)$.
If $w=0$, the congruence makes $u$ a positive integer,
already impossible.  Valuations with $u=0$ satisfy the
same conclusion trivially.  This contradicts
Lemma~\ref{lem:B-unbounded-index} and
\eqref{eq:B-large-index}.

\medskip
\noindent\emph{Descent from the semiample output.}
We are left with a semiample nonbig $P'$ and its contraction
$g:U'\to W$.  For a newly trimmed parameter range tending
to infinity, with $t\eta_i\to0$ uniformly on that range,
use boundary and nef part
\begin{equation}\label{eq:B-new-family}
 \widetilde B(t)=B(t/2)'-tC',
 \qquad
 \widetilde{\mathbf M}(t)=\mathbf M(t/2)+(t/2)\mathbf H.
\end{equation}
Their total logarithmic divisor is $tP'$ up to rational
linear equivalence, and their discrepancies are exactly
\begin{equation}\label{eq:B-new-discrepancy}
 \widetilde a(t)=a(t/2)+tE.
\end{equation}
To see the equality, first pushing the old total divisor
contributes $-(t/2)e_H$; adding the nef term and subtracting
$tC'$ contributes another $-(t/2)e_H+t(C')^*$.
Their sum is $tE$.

The new boundary dominates $R(P')$.  At correction primes
$P'$ is integral, the old boundary has coefficient
$1-o(1)$, and the amount subtracted is $O(t\eta_i)=o(1)$.
Away from them the old denominator condition is unchanged.
The nef part is b-nef; \eqref{eq:B-new-discrepancy} and
$E\ge0$ show that the discrepancies are increasing and
klt, with effective Calabi--Yau data at zero.

This family has the required signed phase condition for
$P'$.  Indeed, at a new growing parameter, apply the
original condition to $a(t/2)$ and scale all the relevant
divisors by the resulting fixed integer.  If
$\widetilde a(t,T)\to0$, then
\[
 a(t/2,T)\le\widetilde a(t,T),
 \qquad E(T)\le\widetilde a(t,T)/t
      =o(\widetilde a(t,T)).
\]
Since $P'(T)=H(T)-E(T)$, subtracting this negligible
nonnegative order preserves the lower $-o(\widetilde a)$
bound and the upper $O(\widetilde a)$ bound for the phase.

On the geometric generic fibre of $g$, the divisor $P'$
is $\Q$-linearly trivial.  Its integral principalization
order is bounded by Lemma~\ref{lem:B-principal-order}:
along any putative unbounded-order subsequence choose one
growing parameter with the just-proved phase bound, use
the restricted effective generalized klt Calabi--Yau
data, and clear the remaining bounded small denominators
as in that lemma.  These operations use only fixed
multipliers.  A small $\Q$-factorialization of the
geometric fibre preserves all the divisor orders and
permits the cyclic-cover construction.

The principalization descends from the geometric generic
fibre to its field of definition.  Indeed being Cartier
descends faithfully flatly, and the geometrically trivial
line bundle is trivial by Hilbert 90; equivalently its
one-dimensional space of trivializing rational functions
descends.  Thus for a bounded positive integer $q$ and
some $\phi\in k(U')^*$, the divisor
$qP'-\divisor(\phi)$ has no horizontal terms.  There is
then an exact equality
\begin{equation}\label{eq:B-exact-descent}
 qP'-\divisor(\phi)=g^*H_W
\end{equation}
for an ample $\Q$-Cartier rational divisor $H_W$.
To justify exactness, semiampleness already gives
$P'\sim_{\Q}g^*A_W$ with $A_W$ ample.  After clearing
denominators the difference between the left side of
\eqref{eq:B-exact-descent} and $qg^*A_W$ is principal
and vertical.  A rational function with vertical divisor
restricts to a nowhere vanishing regular function on
the projective generic fibre, hence belongs to $k(W)^*$
because $g$ is a contraction.  Absorb this rational
principal divisor into $qA_W$ to obtain $H_W$.

If $W$ is a point, \eqref{eq:B-exact-descent} and $E\ge0$
already give a section of a bounded multiple of $H$.
Otherwise $W$ has strictly smaller dimension than $V$.
It is of Fano type.  Apply the generalized canonical
bundle formula to \eqref{eq:B-new-family}, and
replace the old parameter $s$ by $s=qt$.  The induced
data on $W$ are effective generalized data with total
$tH_W$, klt at zero and with nondecreasing discrepancies.
Their nef parts are b-nef and $\Q$-Cartier.

We check the denominator condition on $W$ explicitly.
For a prime $P$ on $W$, choose a nonexceptional vertical
component $T$ above it with multiplicity $e$.  Let $r$
be the denominator of $\operatorname{coeff}_T(P')$ and
let $r'$ be that of
$\operatorname{coeff}_T(qP'-\divisor(\phi))$.
Then $r'\mid r$.  The denominator $r_P$ of $H_W$
divides $er'$, by \eqref{eq:B-exact-descent}.
Since the upstairs boundary dominates $R(P')$, the
discriminant discrepancy at $P$ satisfies
\[
 b_t(P)\le\frac1{er}\le\frac1{er'}\le\frac1{r_P}.
\]
Thus the discriminant boundary dominates $R(H_W)$.

For each growing parameter, the pure pullback phase
condition descends by Lemma~\ref{lem:A-pullback-phases},
after its allowed fixed multiple.  Every hypothesis of
Lemma~\ref{lem:B-phase-nonvanishing} is now satisfied in
smaller dimension.  Induction gives bounded-degree
sections of $H_W$.  They pull back through
\eqref{eq:B-exact-descent} to bounded-degree sections
of $P'$, and the b-divisor inequality
$\mathbf H\ge(P')^*$ in
\eqref{eq:B-correction-errors} carries them to sections
of the same bounded multiples of $H$ on $V$.
This contradicts \eqref{eq:B-no-sections} and completes
the induction.

The contradiction excludes every sequence of failures
with $N_i\to\infty$, and hence proves boundedness for
the original sequence.  Any finite initial segment
omitted in an eventual assertion can also be included:
each of its big divisors has a section in some positive
degree.  Taking a common multiple of the bounded
degrees gives the common degree in the statement.
\end{proof}

\section{Cutoff fields, phase bounds, and section counts}\label{sec:C}

We continue to work with sequences of bounded dimension, and with the
subsequence convention of Section~\ref{sec:A}.  In particular, a statement
about all valuations with an order tending to zero is required on every
subsequence.  All constants below may depend on the dimension bound and
on the fixed constants explicitly present in the hypotheses.  They do
not depend on an individual member of a sequence.  Unless explicitly
stated otherwise, quantified divisorial tests are primitive prime
orders.  Homogeneous estimates extend to integral multiples; arbitrary
positive rescaling is used only for homogeneous inequalities or exact
zero conditions.

\subsection{Polarized discrepancy data and cutoff fields}
\label{subsec:C-cutoffs}

Let $V$ be a projective variety of Fano type.  Let $a$ be the discrepancy
function of an effective generalized rational lc Calabi--Yau pair over
the point, and suppose that
\[
                 a(T)\in\Z\qquad(T\text{ a primitive prime order}).
\]
Let $D$ be an effective nef and big rational $\Q$-Cartier divisor on
$V$.  Its Cartier b-pullback is denoted by $\mathbf D$.  Assume that
$b=a+\mathbf D$ is klt and has effective trace on $V$, or equivalently
$b(Q)\leq1$ for primes $Q$ on $V$.  This generalized pair is obtained
by adding $\mathbf D$ to the nef part and subtracting its trace from
the boundary; its total log canonical class is unchanged.  We call a
prime order $T$ with $a(T)=0$ a \emph{floor prime}.  Thus on $V$ every
positive coefficient of $D$ lies on the floor, and every nonfloor prime
has $a(Q)=1$ and $D(Q)=0$.

For a rational divisor $A$ on a normal projective model, put
\begin{equation}\label{eq:C-function-space}
 S(A)=\{0\}\cup\{f\in k(V)^*:\divisor(f)+A\geq0\}.
\end{equation}
Tests on this space always mean tests on its nonzero elements.  If
$A\geq0$, it contains $1$, so the field it generates is unambiguous.
We take the varying main scale $L\geq1$ to be rational.  Rational upper
approximations within a fixed factor give the same statements for real
scales, after changing constants.  A $\Q$-principal order is the order
function of a rational principal divisor.  Unless a new trace model is
specified, orders of rational divisors on $V$ mean their b-pullbacks.

We shall also use geometric generic fibres.  Divisors on such a fibre
spread over a finite extension of the generic field, and their closures
on resolutions give horizontal divisors on the total space.  Conversely,
horizontal divisors restrict divisorially.  Finite extensions of the
field of constants in characteristic zero are separable and are etale
on these generic-field models, so actual prime orders restrict to
primitive horizontal orders; integral multiples restrict to integral
multiples.  Canonical divisors from the base have horizontal order zero,
and discrepancies and Cartier b-orders therefore agree by homogeneity.
Nef b-parts restrict on a determination.  Relative numerical
proportionality for a Picard-number-one Mori contraction remains true on
a geometric generic fibre: spread any test curve over a finite extension
and specialize.  These fibres are of Fano type.  Consequently rational
numerical equivalence on them implies rational linear equivalence, by
rational connectedness and the Fano type basepoint-free consequences
recalled in Section~\ref{sec:A}.

Apply the extraction construction of Subsection~\ref{subsec:A-mori},
funded by $a+\mathbf D/2$, to obtain a projective $\Q$-factorial Fano
type model $U\to V$ containing every primitive prime with
$a=0$ and $\mathbf D\leq1$, and extracting no other prime.  The list
is finite: each selected prime has discrepancy at most $1/2$ for the
fixed generalized klt pair $a+\mathbf D/2$.  On $U$, precisely the
floor primes form $\Supp D_U$, and their coefficients lie in $(0,1]$.
For rational $0<\beta\leq1$, set
\begin{equation}\label{eq:C-cutoff}
 I_\beta=\{Q:a(Q)=0,\ D(Q)\leq\beta\},\qquad
 P_\beta=D_U-\sum_{Q\in I_\beta}D(Q)Q.
\end{equation}
All the primes in this sum occur on $U$.  Run a negative $P_\beta$-MMP
and let
\[
             p:U'\longrightarrow W_\beta
\]
be the contraction defined by the semiample output $P_\beta'$.  The
semiample divisor is the pullback of an ample rational divisor, allowing
a point as the base.  In fact there is an effective ample rational
divisor $D_W$ such that
\begin{equation}\label{eq:C-cutoff-pullback}
 P_\beta'=p^*D_W,\qquad
 D'=p^*D_W+\sum_{Q\in I_\beta}D(Q)Q'.
\end{equation}
To check the first equality as divisors, rather than just as classes,
$P_\beta'$ is effective and has no horizontal component.  A rational
principal difference between it and a pullback has divisor zero on the
projective generic fibre, hence is a function from the base.  Absorb
that function into the divisor downstairs.  Here the connected-fibre
property makes $k(W_\beta)$ relatively algebraically closed in $k(V)$.
Writing $\mathbf P$ for the nef Cartier b-divisor defined by
$P_\beta'$, MMP negativity gives
\begin{equation}\label{eq:C-cutoff-errors}
          0\leq\mathbf P\leq P_\beta^*\leq\mathbf D.
\end{equation}
Every prime in $I_\beta$ survives the MMP.  Its coefficient in
$P_\beta^*$ is zero, so the effective error
$P_\beta^*-\mathbf P$ has value zero there; a prime contracted by a
negative MMP would instead have strictly positive error.

Define
\[
                      K_\beta=k(W_\beta)\subset k(V).
\]
The field is generated by the spaces $S(mP_\beta)$, $m\in\N$, and is
relatively algebraically closed.  These fields decrease with $\beta$.
A strict inclusion of two nested relatively algebraically closed
subfields decreases transcendence degree: if their transcendence
degrees were equal, the larger would be algebraic over the smaller.
Thus there are at most $\dim V$ strict changes.  We have $K_1=k$, and
$K_{0^+}=k(V)$, since below the least coefficient in the finite support
of $D_U$ the divisor $P_\beta$ equals the big divisor $D_U$.

The following maximal-order condition, called \textup{(C-$\alpha$)}, will recur:
\begin{equation}\label{eq:C-alpha}
 T(\Gamma)\leq(1+\eta)D(T)
 \quad\text{for every effective }\Gamma\sim_\Q D,
 \qquad a(T)=0.
\end{equation}
The quantifier over $\Gamma$ is part of the condition.  When imposed on
a set of valuations in a sequence, the same $\eta=\eta_i$ works for
that entire set.

\begin{lemma}[Uniform cutoff models and order functions]
\label{lem:C-cutoff-claims}
Suppose $\dim W_\beta>0$.  After passage to a subsequence there is a
fixed $\delta>0$ and a projective $\Q$-factorial Fano type extraction
$\widehat W\to W_\beta$ such that $\widehat W$ has ordinary
$\delta$-lc singularities and every positive coefficient of the pullback
of $D_W/\beta$ is at least $\delta$.  A bounded constant $C$ can be
chosen so that $S(CD_W/\beta)$ generates $k(W_\beta)$.

There are algebraically independent functions
$\phi_1,\ldots,\phi_{\dim W_\beta}$ in this space and a b-nef
$\Q$-Cartier order function $H$ on $k(W_\beta)$, pulled back to
$k(V)$, with uniformly bounded denominator of its values on primitive
orders, such that:
\begin{enumerate}[label=\textup{(\roman*)}]
\item If $a(T)=0$, then $H(T)\leq CD(T)/\beta$.  If $T$ is nontrivial
on $K_\beta$ and all $T(\phi_j)\geq0$, then $H(T)\leq-1$.
\item If $D(Q)=0$ and $a(Q)\leq1$, then $H(Q)\leq0$.
\item If $a(T)=0$ and \eqref{eq:C-alpha} holds at $T$, then
$H(T)\geq-C\eta D(T)/\beta$.
\end{enumerate}
In particular, if \eqref{eq:C-alpha} holds on all of $I_\beta$ with
$\eta_i\to0$, those primes are eventually horizontal over $K_\beta$;
that is, their restrictions to $K_\beta$ are trivial.
\end{lemma}

\begin{proof}
On $U'\to W_\beta$ consider
\begin{equation}\label{eq:C-cutoff-direction}
 d=\frac{\mathbf D-\mathbf P}{\beta}+\mathbf P.
\end{equation}
It is nonnegative and strictly positive on $a=0$.  The discrepancies
$a+td$ have effective traces for $0\leq t\leq1$: at a surviving prime
in $I_\beta$ their trace increment is $tD(Q)/\beta\leq1$, and at an
unremoved floor prime it is $tD(Q)\leq1$.  Realize the family by a nef
increment $t\mathbf D/\beta$ and total log canonical divisor equal,
up to a rational principal divisor, to the pullback of
$t(1/\beta-1)D_W$.  Lemma~\ref{lem:A-linear-chamber} gives, on a fixed
smaller interval, discriminants
\[
                              a'+td'
\]
with $a'$ an effective generalized lc Calabi--Yau discrepancy over the
point and with uniformly bounded denominator.  These are klt for
positive $t$.

We first make explicit a lifting consequence needed here.  If
$a'(E_i)=0$, $d'(E_i)\to0$, and $k_i>0$ tends to zero, then, after a
subsequence, there are primitive lifts $T_i$ with
\begin{equation}\label{eq:C-small-lifts}
 T_i|_{k(W)}=e_i\ord_{E_i},\qquad
 a(T_i)=0,\qquad d(T_i)\to0,\qquad e_i k_i\to0.
\end{equation}
At one Mori step, choose a minimizer of $d$ on the $a$-minimum by
computing the discriminant at a sufficiently small positive parameter.
Its tailored vertical chart from Proposition~\ref{prop:A-uniform-charts}
has normalized values $0,\rho d'$ on the slice $e_*=1$.
Lemma~\ref{lem:A-lattice-approximation}, with a fixed small error
$\gamma$, gives an integral point with $e_*=lj$ bounded in terms of
$\gamma$, by rounding $ljU$ in the kernel lattice.  The error in $a$
is less than the gap of its value lattice; nonnegativity and quantization
therefore give $a=0$ exactly.  The slope bound in the same chart gives
\[
                   d\leq e_*\rho d'+C\gamma.
\]
First fix $\gamma$ and go sufficiently far in the sequence, then let
$\gamma\downarrow0$ diagonally slowly enough that also $e_*k_i\to0$.
Primitivizing the resulting order only decreases the orders and the
restriction multiplicity.  Iterate through the bounded chain of Mori
steps, at each step including the already accumulated multiplicity
times $k_i$ among the quantities required to be negligible.  Birational
steps preserve the discrepancy functions.  This proves
\eqref{eq:C-small-lifts} without asserting a uniform bound for $e_i$.

There cannot be a sequence of primes $E$ with, for one fixed sufficiently
small $t>0$,
\begin{equation}\label{eq:C-excluded-small-coefficients}
 (a'+td')(E)\longrightarrow0,
 \qquad 0<D_W(E)/\beta\longrightarrow0.
\end{equation}
The bounded denominator gives $a'(E)=0$ eventually.  Apply
\eqref{eq:C-small-lifts} with $k_i=D_W(E)/\beta$.  It gives
\[
 a(T)=0,\qquad
 \frac{D(T)}{\beta}\leq d(T)+\frac{eD_W(E)}{\beta}\longrightarrow0.
\]
Thus $T\in I_\beta$, where $\mathbf P(T)=0$, whereas
$\mathbf P(T)=eD_W(E)>0$, a contradiction.  Nor can an original prime
$E$ of $W$ have $0<D_W(E)/\beta\to0$.  A prime $Q$ above it on $U'$
satisfies
\[
 D(Q)=\mathbf P(Q)=eD_W(E)>0,\quad a(Q)=0,
 \quad a'(E)=0,\quad d'(E)\leq D_W(E).
\]
It would again give \eqref{eq:C-excluded-small-coefficients}.

Choose fixed rational $0<t\leq1/4$ in the chamber.  Extract every
exceptional prime of $A=a'+td'$ below a sufficiently small fixed
$\delta_*<1/4$.  This finite extraction is funded by mixing $a'$
slightly with an ordinary klt Calabi--Yau pair, as in
Subsection~\ref{subsec:A-mori}; the extracted primes have $a'=0$ and
positive boundary coefficient.  The preceding exclusions imply a
uniform positive lower bound on all nonzero coefficients of
\[
                 E_0=(D_W/\beta)|_{\widehat W}.
\]
Here and below the notation means the rational pullback.  The ordinary
discrepancies on the extraction dominate those of $A$, by b-nef
negativity and $\Q$-factoriality.  Every remaining exceptional primitive
prime has $A\geq\delta_*$, and every extracted prime has ordinary
discrepancy one.  This proves a uniform ordinary gap on $\widehat W$.

For completeness, we verify the additional ordinary gap needed to
apply polarized birationality to an integral Weil divisor.  Let $B_A$
be the trace boundary of $A$ on $\widehat W$.  On every component of
$E_0$ its coefficient is at least $1/2$.  For an original component the
covering-prime calculation above gives $a'=0$ and
$d'\leq D_W(\cdot)\leq1$; extracted components have $A<\delta_*$.
Choose a fixed integer $m>4/\delta$ and an integral divisor
\[
                     N=mE_0-Q_0,
\]
where the coefficients of $Q_0$ belong to $[1,2]$ on $\Supp E_0$ and
are zero elsewhere.  Such rounding is possible coefficient by
coefficient.  Since $mE_0-Q_0$ dominates a positive multiple of $E_0$,
$N$ is big.  Run an $N$-MMP to a nef and big integral Weil divisor
$N^+$ on a $\Q$-factorial model.  With stars denoting pullbacks to a
common determination, b-nef negativity and MMP negativity give
\begin{equation}\label{eq:C-rounded-errors}
 0\leq e_E=(E_0^+)^*-E_0^*,\qquad
             me_E\leq(Q_0^+)^*-Q_0^*.
\end{equation}
The total generalized log canonical divisor of $A$ is
$t(1-\beta)E_0$ up to a rational principal divisor.  Comparing its nef
part before and after the MMP, ordinary discrepancies on the output
are at least
\begin{equation}\label{eq:C-rounded-gap}
 A+(B_A^+)^*-t(1-\beta)e_E
 \ \geq\ A+\left(\frac14-\frac tm\right)(Q_0^+)^*.
\end{equation}
Indeed $B_A^+\geq Q_0^+/4$ and
$e_E\leq(Q_0^+)^*/m$ by \eqref{eq:C-rounded-errors}.
If a primitive prime with $A<\delta_*$ is in $\Supp E_0$, it has
$(Q_0^+)^*\geq Q_0^*\geq1$.  If it is not in this support, then it
is already a prime on $\widehat W$ and cannot be contracted by the
$N$-MMP: $N$ has coefficient zero there, the output pullback is
effective, and a contraction would require strictly positive MMP
error.  Its ordinary discrepancy on the output is one.  All other
primitive orders have $A\geq\delta_*$.  Formula
\eqref{eq:C-rounded-gap} therefore gives a uniform ordinary gap on the
output.  Also $N^+-K^+$ is pseudo-effective, since the output is of
Fano type.  Birkar's integral-Weil polarized birationality theorem
\cite[Theorem~1.1]{BirkarPolarised} now gives birational sections of a
bounded multiple of $N^+$.  Negativity identifies these with sections
of that multiple of $N$, hence with functions in $S(CD_W/\beta)$ for
bounded $C$.  This proves the field-generation assertion.

Choose algebraically independent $\phi_j$ in this space.  If
$\omega=\bigwedge_jd\phi_j$, define on the base
\begin{equation}\label{eq:C-logarithmic-order}
 H(E)=a'(E)-\nu_{\log,E}(\omega),\qquad
              \nu_{\log,E}(\omega)=1+\ord_E(\omega).
\end{equation}
The Calabi--Yau identity shows that $H$ is the moduli b-divisor after
a rational principal adjustment.  It is thus b-nef and $\Q$-Cartier,
and its values have bounded denominator because $a'$ does and the
log orders are integral.  If $T|_{k(W)}=e\ord_E$ and $a(T)=0$, the
discriminant minimum gives $a'(E)=0$.  Pole bounds for the $\phi_j$
on every model give
\[
                  \nu_{\log,E}(\omega)\geq-CD_W(E)/\beta.
\]
If all $E(\phi_j)\geq0$, the form is regular at that discrete
valuation ring and its logarithmic order is at least one.  These prove
(i).  For (ii), if restriction is nontrivial, then
$D_W(E)=0$ by $\mathbf P\leq\mathbf D$, so the functions are regular
and
\[
       H(Q)=e a'(E)-e\nu_{\log,E}(\omega)
                    \leq a(Q)-e\leq0.
\]
The trivial restriction gives $H(Q)=0$ directly.

On $W$ itself we also have $H_W\leq C D_W/\beta$: use $a'\leq1$,
the lower coefficient bound, and the same log-order estimate.  For any
arbitrarily small positive rational $\lambda$, choose on a determining
model an effective
\[
                  \Sigma\sim_\Q H+\lambda\mathbf D_W/\beta;
\]
this is possible because the displayed class is nef and big.  The
rational principal order
\[
           \chi=\mathbf\Sigma-H-\lambda\mathbf D_W/\beta
\]
has poles at most $C D_W/\beta$ on $W$, and therefore at most
$C\mathbf D/\beta$ upstairs.  Condition \eqref{eq:C-alpha}, applied to
$D+(\beta/C)\chi$, gives
$\chi(T)\leq C\eta D(T)/\beta$.  Since $\mathbf\Sigma\geq0$,
\[
 H(T)\geq-C\eta D(T)/\beta-\lambda\mathbf P(T)/\beta.
\]
Let $\lambda\downarrow0$ to prove (iii).  Finally, for $T\in I_\beta$,
$\mathbf P(T)=0$, so each $\phi_j$ has nonnegative order there.
If the restriction were nontrivial, (i) and (iii), together with
$D(T)\leq\beta$, would give $-1\geq H(T)\geq-C\eta_i$, which is
impossible eventually.
\end{proof}

\subsection{A phase estimate at a varying scale}

The cutoff construction identifies the part of the function field visible
at a prescribed pole cost.  We next use it to control the denominators of
rational principal orders at floor primes of very small polarization
order.  This is the phase estimate needed to regularize cone characters
in Section~\ref{sec:D}.
\label{subsec:C-scaled-phase}

\begin{lemma}[Scaled phase bound]\label{lem:C-scaled-phase}
Assume \eqref{eq:C-alpha} with $\eta_i\to0$ whenever
$a(T)=0$ and $LD(T)\to0$, along arbitrary subsequences.  Let $\ell$
be a rational principal order.  Suppose that for every nonfloor prime
$Q$ on $V$ there is an integer $p_Q$ with
\begin{equation}\label{eq:C-nonfloor-phase}
 |p_Q-\ell(Q)|\leq G(Q),\qquad
 G\geq0,\qquad G\sim_\Q c_gLD,
\end{equation}
where $c_g\geq0$ is rational and uniformly bounded.  There is a fixed
positive integer $q$ such that
\begin{equation}\label{eq:C-scaled-phase}
 \operatorname{dist}(q\ell(T),\Z)=O(LD(T))
 \quad\text{if }a(T)=0\text{ and }LD(T)\longrightarrow0.
\end{equation}
\end{lemma}

\begin{proof}
The hypothesis on nonfloor primes also holds on $U$, since every
additional extracted prime is on the floor.  If $c_g>0$, condition
\eqref{eq:C-alpha} applies to $G/(c_gL)$.  If $c_g=0$, the effective
rationally trivial divisor $G$ is zero.  Consequently, at every floor
order to which the smallness hypothesis applies,
\begin{equation}\label{eq:C-G-small}
                       G(T)\leq c_gL(1+\eta_i)D(T).
\end{equation}

Suppose the conclusion fails.  Factorial diagonalization gives a
subsequence and primitive targets $T$ with
\begin{equation}\label{eq:C-phase-failure}
 a(T)=0,\quad\beta=D(T),\quad L\beta\longrightarrow0,\qquad
 \frac{\operatorname{dist}(q\ell(T),\Z)}{L\beta}
                 \longrightarrow\infty
 \quad\text{for each fixed }q\in\N_{>0}.
\end{equation}
To justify the simultaneous last assertion, at stage $i$ choose a failure
for the multiplier $i!$ with ratio tending sufficiently quickly to
infinity.  For $q\mid i!$,
$\operatorname{dist}(i!x,\Z)\leq(i!/q)
\operatorname{dist}(qx,\Z)$; choosing the failure ratio larger than all
these factors gives the claimed diagonal sequence.
Among such sequences maximize the eventually constant
$\trdeg K_\beta$.  Lemma~\ref{lem:C-cutoff-claims} says that $T$ is
horizontal over this field.  Let $\tau$ be the first positive cutoff
for which $K_\tau=K_\beta$.  It exists because there are finitely many
coefficient cutoffs and the field $K_{0^+}=k(V)$ cannot have a nonzero
horizontal divisorial valuation.  In particular,
\begin{equation}\label{eq:C-first-change-scale}
                         0<\tau\leq\beta,
                    \qquad L\tau\longrightarrow0.
\end{equation}
After multiplying $\ell$ and $G$ by one fixed integer, maximality
implies a uniform bound
\begin{equation}\label{eq:C-before-threshold}
 \operatorname{dist}(\ell(Q),\Z)\leq CLD(Q)
                   \quad(a(Q)=0,\ D(Q)<\tau).
\end{equation}
Indeed, a failure of this assertion for every fixed multiple would,
by the same factorial diagonalization, produce a sequence at a cutoff
strictly before $\tau$, hence at strictly larger transcendence degree.

\paragraph{The first threshold.}
We first prove \eqref{eq:C-before-threshold} with $D(Q)\leq\tau$,
after another fixed multiplication.  If it fails, choose a primitive
$S$ with $D(S)=\tau$ such that every fixed multiple has phase mismatch
much larger than $L\tau$.  The prime $S$ lies on $U$.  Fix rational
$0<\gamma<1/4$ and put
\begin{equation}\label{eq:C-threshold-divisor}
 \mathcal S=S+
 \sum_{\substack{Q\text{ on }U\,,\ a(Q)=0\\Q\ne S}}
                 \gamma\min\{1,D(Q)/\tau\}\,Q.
\end{equation}
Let $t$ be the pseudo-effective threshold of $D_U-t\mathcal S$.
It is rational, since $U$ is of Fano type.  The divisor
$D_U-\tau\mathcal S$ is effective, while an effective representative
of $D_U-t\mathcal S$ would give an effective divisor rationally
equivalent to $D$ whose order at $S$ is at least $t$.
Rational pseudo-effectivity on a Fano type model is realized by such
an effective representative.  Hence \eqref{eq:C-alpha} at $S$ gives
\begin{equation}\label{eq:C-threshold-location}
                    \tau\leq t\leq(1+\eta_i)\tau,
                           \qquad\eta_i\longrightarrow0.
\end{equation}
Run the threshold MMP and take the semiample nonbig fibration.  On its
generic fibre,
\begin{equation}\label{eq:C-threshold-comparison}
 D'\sim_\Q t\mathcal S',\qquad
 \sum_{Q\ne S}D(Q)Q'\precsim2\eta_i\tau S'.
\end{equation}
Here $A\precsim B$ means that $B-A$ is rationally linearly equivalent
to an effective rational divisor.  To see the second inequality, the
coefficient of every $Q'\ne S'$ in $D'-t\mathcal S'$ is at least
$D(Q)/2$ for large $i$, whereas the coefficient of $S'$ is at least
$-\eta_i\tau$.  Restriction of the rationally trivial threshold
class proves the claim.  Enlarge $\eta_i$ rationally if necessary.
The restriction of $D'$ is big, so $S'$ is horizontal and big on this
generic fibre.  For the threshold MMP and every subsequent MMP over
its base, the pullback errors satisfy
\begin{equation}\label{eq:C-threshold-error}
 0\leq e_D=(D'')^*-\mathbf D
       \leq t\bigl((\mathcal S'')^*-\mathcal S^*\bigr).
\end{equation}
The first comparison uses that $\mathbf D$ is nef; the second is
MMP negativity for $D-t\mathcal S$, followed by its crepancy over the
threshold base.

We need a signed integral multiplier, allowed to grow slowly, for which
\begin{equation}\label{eq:C-phase-amplification}
 m=m_i\ne0,\quad\sigma=\{m\ell(S)\}\longrightarrow0^+,
                    \qquad\eta_i+|m|L\tau=o(\sigma).
\end{equation}
Pass first to a limiting phase.  If it is irrational, choose fixed
signed powers whose limiting residues approach zero from above, and
then let these powers vary sufficiently slowly.  For a rational limit,
first take a fixed multiple clearing that limit.  Its signed residual
$\zeta_i\to0$ satisfies $|\zeta_i|/(L\tau)\to\infty$ by the assumed
failure.  Choose its sign positive and multiply it by an integer so
that the resulting residual still tends to zero but dominates
$\eta_i$.  This is possible with no wrap: choose a target residual
$s_i\to0$ with $s_i\gg\eta_i+|\zeta_i|$, and a multiplier of size
$s_i/|\zeta_i|$.  Its product with $L\tau$, divided by the new
residual, is $O(L\tau/|\zeta_i|)\to0$.  This proves
\eqref{eq:C-phase-amplification}.

On the threshold model choose a rational correction $n$ so that
$N=m\ell'+n$ is an integral Weil divisor, as follows:
\begin{align}
 n(S')&=-\sigma,\label{eq:C-rounded-phase-data}\\
 |n(Q')|&\leq C|m|LD(Q)
                 &&(Q\ne S,\ a(Q)=0,\ D(Q)<\tau),\notag\\
 0\leq n(Q')&<1
                 &&(Q\ne S,\ a(Q)=0,\ D(Q)\geq\tau),\notag\\
 |n(Q')|&\leq |m|G'(Q')
                 &&(a(Q)>0).\notag
\end{align}
Use nearest integers in the second line, upward rounding in the third,
and $mp_Q$ in the last; outside the finite relevant support the
correction is zero.  Since $\ell'$ is rational principal,
\eqref{eq:C-threshold-comparison} bounds the positive errors in class
by
\[
 C|m|LD'+|m|G'+\tau^{-1}\sum_{Q\ne S}D(Q)Q'
          \precsim O(|m|L\tau+\eta_i)S'.
\]
Thus
\[
                    N\precsim-(\sigma-o(\sigma))S',
\]
and $N$ is not relatively pseudo-effective.  Run its Mori fibre MMP
over the threshold base.  On the geometric generic fibre $F$ of the
last Picard-number-one contraction,
\begin{equation}\label{eq:C-large-index-threshold}
 H_1=-N''\text{ is integral and ample},\qquad
                   H_1\precsim(\sigma+o(\sigma))S''.
\end{equation}
The upper comparison uses the lower bounds in
\eqref{eq:C-rounded-phase-data}; all the generic class comparisons
survive pushing and restricting.  The variety $F$ is of Fano type,
and $-K_F$ dominates $S''|_F$ numerically.  In fact the pushed crepant
boundary of $a$ is effective and contains $S''$ with coefficient one,
while the moduli trace on $F$ is pseudo-effective.  For this last
assertion restrict its nef determination first, add arbitrarily small
ample errors and a pullback of an ample divisor if needed, and push
effective approximations to $F$.  Letting the errors tend to zero
proves pseudo-effectivity.  Picard-number-one proportionality on the
total contraction, followed by its geometric generic restriction and
rational linear equivalence on $F$, therefore gives
\begin{equation}\label{eq:C-unbounded-index}
                         -K_F\sim_\Q I_iH_1,
                               \qquad I_i\longrightarrow+\infty.
\end{equation}

We verify the small-discrepancy hypothesis that contradicts
Lemma~\ref{lem:B-unbounded-index}.  Let $w$ be an integral divisorial
angular evaluation on $F$, allowing a nonprimitive one, with ordinary
discrepancy $A_F(w)\to0$.  Generalized-to-ordinary comparison gives
\begin{equation}\label{eq:C-small-angular}
                  A_F(w)\geq a(w)+w((\mathcal S'')^*).
\end{equation}
Indeed the pushed boundary has coefficient one at every surviving
floor prime, dominates $\mathcal S''$, and b-nef negativity controls
the moduli trace.  Hence integrality gives $a(w)=0$ eventually and,
by \eqref{eq:C-threshold-error},
$w(\mathcal S^*)\leq w((\mathcal S'')^*)\to0$.
The formula on $U$ gives $\mathbf D\leq C\mathcal S^*$, with fixed
$C$ because $D(Q)\leq1$ and $\tau\leq1$.  Thus the primitive
restriction of $w$ is eventually one of the primes extracted on $U$.
Since $w(\mathcal S^*)\to0$, that prime is neither $S$ nor a floor
prime of coefficient at least $\tau$: their $\mathcal S$-coefficients
are respectively $1$ and $\gamma$.  The remaining coefficient is
$\gamma D(Q)/\tau$, so
\begin{equation}\label{eq:C-small-D-threshold}
                   D(w)\leq C\tau A_F(w).
\end{equation}
The solved estimate \eqref{eq:C-before-threshold}, extended by integral
homogeneity, and \eqref{eq:C-phase-amplification} imply that
$m\ell(w)$ has a residual $o(A_F(w))$ modulo $\Z$.  Equations
\eqref{eq:C-G-small} and \eqref{eq:C-threshold-error} also give
\begin{equation}\label{eq:C-G-threshold-error}
 |m|w((G'')^*)
       =|m|\bigl(G(w)+c_gLe_D(w)\bigr)=o(A_F(w)).
\end{equation}
The equality holds because $G-c_gLD$ is rational principal and its
pullback is unchanged.  In the estimate, $e_D(w)\leq
 t\,w((\mathcal S'')^*)=O(\tau A_F(w))$.

The signs of the rounding corrections now matter.  Their negative
part is at most
\[
 \sigma(S'')^*+C|m|L(D'')^*+|m|(G'')^*;
\]
its value is $o(A_F(w))$ by the preceding bounds.  The positive
unit-size errors are supported on floor components, whose reduced
sum has order at most $A_F(w)$ by boundary comparison.  Thus
\begin{equation}\label{eq:C-signed-rounding}
                   -o(A_F(w))\leq w((n'')^*)\leq O(A_F(w)).
\end{equation}
All these are inequalities between effective divisors pulled back from
$\Q$-factorial models before restriction to $F$.

Consider the compact Seifert cone associated to the integral
polarization $H_1$.  For a divisorial evaluation over infinity with
positive height $u$ and ordinary discrepancy $A\to0$,
Lemma~\ref{lem:B-seifert} gives $A\geq u+A_F(w)$ and
\[
                    u\equiv-m\ell(w)-w((n'')^*)\pmod\Z.
\]
The trivial angular restriction is impossible when $A\to0$, since
then $u$ is a positive integer.  Otherwise all terms tend to zero,
so there is no wrap in this congruence.  The residual of $m\ell(w)$
is $o(A_F(w))$ and the negative part in
\eqref{eq:C-signed-rounding} is $o(A_F(w))$; positivity of $u$
therefore forces $u=o(A)$.  Orders with centre outside infinity have $v(D)=0$, so the same
small-height conclusion holds for every small-discrepancy prime.
This contradicts Lemma~\ref{lem:B-unbounded-index} and
\eqref{eq:C-unbounded-index}.  We have proved the phase estimate also
at $D(S)=\tau$, after a fixed multiplication.

\paragraph{Rational triviality on the cutoff fibre.}
Return to the cutoff model $U^\tau\to W_\tau$.  On its generic
fibre,
\[
                   D'=\sum_{Q\in I_\tau}D(Q)Q',
\]
and no other floor prime on this model is horizontal.  With the fixed
multiplications already made, choose an integral Weil divisor
$N=\ell'+n$ whose horizontal trace satisfies
\begin{equation}\label{eq:C-cutoff-rounding}
                         |n|\leq CLD'+G'.
\end{equation}
If either $N$ or $-N$ is not relatively pseudo-effective, run its Mori
fibre MMP.  On the final geometric generic Mori fibre, put
$R=\sum_{Q\in I_\tau}Q''$.  The same pullback and class comparisons
give
\begin{equation}\label{eq:C-second-large-index}
 H_1=\mp N''\precsim O(L\tau)R,\qquad
          0\leq\mathbf D\leq(D'')^*\leq\tau R^*.
\end{equation}
Again $-K$ dominates $R$ numerically, so the integral ample divisor
$H_1$ has anticanonical index tending to infinity.  If $A_F(w)\to0$,
ordinary comparison gives $a(w)=0$, $w(R^*)\leq A_F(w)$, and hence
\[
 D(w)=O(\tau A_F(w)),\qquad
                  w((G'')^*)=O(L\tau A_F(w)).
\]
For the latter estimate use \eqref{eq:C-G-small} at the primitive
restriction and the positive trace-pullback error bounded by
$(D'')^*$.  That primitive restriction is eventually in $I_\tau$,
since $D(w)\to0$, and all such primes survive on $U^\tau$.
Its solved phase gives a residual $o(A_F(w))$ for $\ell(w)$.
Equation \eqref{eq:C-cutoff-rounding}, pulled back after the MMP,
gives $w((n'')^*)=o(A_F(w))$.  The Seifert congruence consequently
forces $u=o(A)$ as before, contradicting
Lemma~\ref{lem:B-unbounded-index}.

Both $N$ and $-N$ are therefore relatively pseudo-effective.  On a
Fano type generic fibre they each have effective rationally equivalent
representatives, by semiample termination.  Their sum is an effective
rationally trivial divisor and is zero, so $N$ is rationally linearly
trivial on that fibre.

\paragraph{A bounded principalization multiple.}
Run an MMP over $W_\tau$ making $R=\sum_{I_\tau}Q''$ nef.  On the
geometric generic fibre the integral divisor $N''$ remains rationally
principal, and its pullback agrees with that of $N$, since a
rationally trivial divisor is crepant through these modifications.
Fix $0<c<1$.  The discrepancies
\begin{equation}\label{eq:C-torsion-CY}
                              a_*=a+cR^*
\end{equation}
are an effective generalized klt Calabi--Yau pair.  The trace is
effective because $R$ is reduced and lies on the floor; klt follows
from $\mathbf D\leq\tau R^*$ and positivity of $\mathbf D$ on
$a=0$.  The additional b-part is nef.  The same small-order argument
shows
\[
       \operatorname{dist}(N''(w),\Z)=o(a_*(w))
                          \quad\text{whenever }a_*(w)\to0.
\]
We record why this bounds its torsion order.  If its minimal
principalization order $l$ were unbounded, write
$lN''=\divisor(g)$ and take the normalization in a root of $g$.
The cover is connected by minimality of $l$ and is unramified in
codimension one because $N''$ is integral.  Thus ordinary Fano type
data pull back without an angular ramification boundary.  Form the
product with $\mathbf P^1$ and take the finite diagonal cyclic
quotient, as in Lemma~\ref{lem:B-principal-order}.  The product-to-quotient
map is quasi-etale; over zero the fibration has multiplicity $l$.
The log pullback of \eqref{eq:C-torsion-CY}, product with the line
factor, and descent provide effective generalized klt Calabi--Yau
data over the curve.

For a prime over zero, put
$u=\ord(\text{base coordinate})/l>0$.  Its discrepancy satisfies
$A\geq u+a_*(w)$, and the cyclic congruence is
$u\equiv-w(N'')\pmod\Z$.  A trivial angular restriction makes $u$
a positive integer and hence cannot occur when $A\to0$.  For a
nontrivial restriction the preceding phase estimate and the absence of
wrap give $u=o(A)$.  Add the rational divisor $f^*[0]/l$ with varying coefficient, extending
the height by zero away from the origin.  The discrepancies $A-su$ remain an
effective generalized Calabi--Yau family over the curve for arbitrarily
long nonnegative parameter intervals.  Lemma~\ref{lem:A-comparison}
applies because every $A\to0$ has $u=o(A)$.  But the base discrepancy
at zero is at most $1/l\to0$ and is nonpositive after the coefficient
$s=1$ is added.  This contradicts its conclusion.  Thus $l$ is bounded.

The bounded principalization descends to the original generic field.
Cartierness descends under the algebraically closed field extension for
the corresponding Weil divisor, and a geometrically trivial line
bundle and its inverse each have a nonzero section over the original
field by flat base change.  Their product is a nonzero constant, so
those sections give a trivialization.  This also explains explicitly
why no unbounded field extension is absorbed into the asserted bounded
multiple.

\paragraph{The final two-sided comparison.}
There is now a fixed positive integer $q$ and a rational function
$\chi$ on the generic fibre, hence in $k(V)$, such that the rational
principal order
\[
                         h=q\ell-\divisor(\chi)
\]
has horizontal trace on $U^\tau$ bounded by
\begin{equation}\label{eq:C-final-horizontal-bound}
                            |h'|\leq CLD'+qG'.
\end{equation}
For example, take $\divisor(\chi)=qN$ on that fibre, so
$h'=-qn$.  The original target $T$ is horizontal over $K_\tau=K_\beta$,
and subtracting an integral principal order does not remove its phase
failure:
\begin{equation}\label{eq:C-final-h-failure}
                          \frac{|h(T)|}{LD(T)}\longrightarrow\infty.
\end{equation}
On the geometric generic fibre consider
\begin{equation}\label{eq:C-final-family}
 a_0=a+c\mathbf D/\tau,\qquad
 a_s=a_0+\frac{s h-q|s|\mathbf G}{J_iL\tau}.
\end{equation}
Choose $J_i\to\infty$ sufficiently slowly compared with the ratio in
\eqref{eq:C-final-h-failure}, and then a two-sided parameter range
$|s|\leq R_i\to\infty$ with $R_i/J_i\to0$.  These are concave,
finite piecewise affine effective generalized Calabi--Yau data.
Indeed \eqref{eq:C-final-horizontal-bound} and
$D'\leq\tau\sum_{I_\tau}Q'$ give trace effectivity for this range;
off the floor the term $-q|s|\mathbf G$ absorbs the possible positive
trace of $sh$.  Up to rational principal terms the added nef part has
$\mathbf D$-coefficient
$c/\tau-q|s|c_g/(J_i\tau)>0$.  Thus b-nefness also holds.

The central member is klt.  If a primitive geometric generic order
satisfies
$a_0<\min\{1,c\}/2$, then $a=0$ and $D<\tau/2$.
Its primitive horizontal restriction is therefore a surviving prime
of $I_\tau$ on $U^\tau$.  Equations
\eqref{eq:C-G-small} and \eqref{eq:C-final-horizontal-bound} give
\[
                        |h|+q\mathbf G=O(L\mathbf D)
\]
there, uniformly for all such orders: these are primes already on the
model where the horizontal trace bound was proved.  Thus the fixed
discrepancy cutoff controls every order required by the two-sided
comparison lemma.  For every fixed signed $s$ the
relative difference $(a_s-a_0)/a_0$ is consequently $O(|s|/J_i)\to0$.
This is the strong two-sided hypothesis of
Lemma~\ref{lem:A-comparison}, with a fixed positive discrepancy
cutoff.  Its conclusion over a point says that $a_s$ is klt for fixed
$s$ eventually.  At the lift of $T$, however, choose $s=1$ or $s=-1$
so that $sh(T)=-|h(T)|$.  Equations
\eqref{eq:C-final-family} and \eqref{eq:C-final-h-failure}, with the
choice of $J_i$, give $a_s(T)<0$.  This contradiction proves the lemma.
\end{proof}

\subsection{Exact saturation of neutral fields}

For the pinning argument, a small error on a fixed list of functions is
not enough: later steps use rational expressions of unrestricted degree
in those functions.  The next lemma puts every bounded-cost function in
one relatively algebraically closed cutoff field and makes the relevant
valuations exactly trivial on that field.  Its count estimates use one
normalization for both the upper and lower bounds.
\label{subsec:C-saturation}

\begin{lemma}[Saturation and section counts]\label{lem:C-saturation-count}
Let $f_1,\ldots,f_s$ be a finite list of genuine rational functions on
$V$, whose length may vary with the sequence, and let
\begin{equation}\label{eq:C-neutral-polarization}
 J(T)=\max\{0,\max_j(-T(f_j))\},\qquad
 0\leq J\leq C L_0\mathbf D,\qquad1\leq L_0=O(L).
\end{equation}
Thus $J$ is an integral basepoint-free Cartier b-divisor on a high
model.  Suppose \eqref{eq:C-alpha} holds uniformly with $\eta_i\to0$
on
\[
                       \mathcal L=\{T:a(T)=0,\ J(T)=0\}.
\]
In addition, on every subsequence on which $L_0\to\infty$, assume
that all valuations in $\mathcal L$ eventually restrict trivially to
the \emph{field generated by} $S(L_0D)$.

After passage to a subsequence, there are cutoff fields $K_*$ of fixed
transcendence degree $p$, a fixed constant $B$, and positive real
numbers $v_i$ with the following properties.
\begin{enumerate}[label=\textup{(\roman*)}]
\item For every fixed $C_1>0$, eventually $S(C_1LD)\subset K_*$.
\item Eventually $S(BLD)$ generates $K_*$.  If $T\in\mathcal L$ is
nonnegative on this space, its restriction to $K_*$ is trivial.
\item There is $C'>0$, independent of the positive integer $n$, such
that for every fixed $n$, eventually
\begin{equation}\label{eq:C-uniform-counts}
 \dim S(nLD)\leq C'v_i n^p,\qquad
 \dim\operatorname{Span}\{g_1\cdots g_n:g_j\in S(BLD)\}
                                      \geq v_i n^p.
\end{equation}
\end{enumerate}
The eventual index in \textup{(i)} may depend on $C_1$, and that in
\textup{(iii)} may depend on $n$; $B$ and $C'$ are fixed before $n$.
All assertions concern divisorial evaluations.  Exact triviality also
holds on rational monomial subcones cut out by the indicated
conditions.
\end{lemma}

\begin{proof}[Proof of the field assertions]
Because $J$ is integral and $L_0=O(L)$, choose a fixed small $c>0$ so
that
\[
                  a(T)=0,\quad LD(T)\leq c
                            \quad\Longrightarrow\quad J(T)=0.
\]
There is, after a subsequence, a single field plateau containing two
cutoffs
\begin{equation}\label{eq:C-plateau}
 0<\beta_0\leq\beta_1\leq c/L,\qquad
 L\beta_0\longrightarrow0,\qquad L\beta_1\geq c'>0,\qquad
                         K_{\beta_0}=K_{\beta_1}=K_*.
\end{equation}
Here is a precise reason for this choice.  Pass to limits of the finitely
many change locations multiplied by $L$, adding dummy endpoints if the
number of changes varies.  The last change tending to zero lies below
a cutoff $\beta_0$ with $L\beta_0\to0$; choose it decreasing slowly
enough to lie beyond that change.  The next change either has a positive
limiting scaled location or lies beyond $c/L$.  Choose $\beta_1$
strictly before it, with $L\beta_1$ fixed positive and at most $c$.
This gives \eqref{eq:C-plateau}.  Passing again fixes $p=\trdeg K_*$.

All $I_{\beta_1}$ belong to $\mathcal L$, so
Lemma~\ref{lem:C-cutoff-claims} makes them horizontal over $K_*$.
It follows, for \emph{every} rational $m>0$, that
\begin{equation}\label{eq:C-exact-plateau-spaces}
                   S(mP_{\beta_0})=S(mP_{\beta_1}).
\end{equation}
Indeed a function in the left side belongs to $K_*$: enlarge $m$ to
an integer and use the section-field description.  Its order at every
removed prime in $I_{\beta_1}$ is zero by horizontality.  Removing
those pole allowances therefore leaves it in the right side.  The
reverse inclusion follows from $P_{\beta_1}\leq P_{\beta_0}$.
For fixed $C_1$, a function in $S(C_1LD)$ has no pole at
$I_{\beta_0}$ eventually, because its order is integral and
$C_1L\beta_0<1$.  Hence
\[
          S(C_1LD)\subset S(C_1LP_{\beta_0})
                       =S(C_1LP_{\beta_1})\subset K_*.
\]
This proves (i).  The cutoff generation statement at $\beta_1$, and
$\beta_1\geq c'/L$, place a birational field-generating space for
$K_*$ inside $S(B_0LD)$ for fixed $B_0$.  Enlarging $B$ will therefore
make $S(BLD)$ generate $K_*$; a basis contains $p$ algebraically
independent elements, by greedily choosing a transcendence basis from
any field-generating set.

We prove that some fixed $B$ forces the asserted exact triviality.
Otherwise choose $B_i\to\infty$ slowly and $T_i\in\mathcal L$ that
are nonnegative on $S(B_iLD)$ but nontrivial on $K_*$.  Since $K_1=k$,
there is a change at $\beta\geq\beta_1$ where $T$ is trivial on
$K_\beta$ but nontrivial on the preceding field $K_{\beta_-}$.
Take $\beta_-$ sufficiently close to $\beta$ from below that
$\beta_->\beta/2$.  Use $H,\phi_j$ from
Lemma~\ref{lem:C-cutoff-claims} at $\beta_-$.  Their pole bounds are
$C\mathbf D/\beta$, and $1/\beta=O(L)$; thus $T(\phi_j)\geq0$
eventually.  Its nontrivial restriction gives
\begin{equation}\label{eq:C-negative-cutoff-order}
                                  H(T)\leq-1.
\end{equation}

If $L_0\beta\to\infty$, then $L_0\to\infty$ because $\beta\leq1$.
The same bounded-cost field-generating space for $K_{\beta_-}$ is
contained in $S(L_0D)$ eventually.  The hypothesis of triviality on the
\emph{field} generated by the latter space contradicts the choice of
$T$.  It remains to treat a subsequence on which $L_0\beta$ is
bounded.

Restrict to the geometric generic fibre of the cutoff contraction
$U^\beta\to W_\beta$.  There $T$ is horizontal and
\[
                     D'\leq\beta\sum_{I_\beta}Q'.
\]
Apply the extra-data part of Lemma~\ref{lem:A-linear-chamber}, over a
point, with
\begin{equation}\label{eq:C-saturation-chamber}
 \widetilde a=a+\kappa J,\qquad
             d=c_* =\mathbf D/\beta,\qquad h=H,
\end{equation}
where $\kappa>0$ is fixed rational and sufficiently small; $c_*$
denotes the additional nonnegative direction of that lemma.
We verify all its hypotheses.  The orders $\widetilde a$ have bounded
denominator because $a$ and $J$ are integral.  The order $H$ has bounded
denominator by the cutoff lemma.  Both $J$ and $\mathbf D/\beta$ are
b-nef, and $H$ is pulled back from a b-nef divisor on $K_{\beta_-}$;
after a birational determination it restricts to a b-nef divisor on
this fibre.  Their Calabi--Yau realizations hold the total class
rationally trivial.  At floor primes on this fibre,
\[
                    H\leq C D/\beta,\qquad J\leq C D/\beta,
\]
where the second bound uses boundedness of $L_0\beta$.  At other
primes on the model, $a=1$, $D=0$, $J=0$, and $H\leq0$.
These inequalities give effective traces on uniform small positive
ranges in \eqref{eq:C-saturation-chamber}, including the family with
$\lambda H+sc_*$ for $0<s\leq\lambda\theta_0$.
The family $\widetilde a+td$ is klt for $t>0$ because $D$ is positive
on $a=0$.  Finally,
\[
 \widetilde a=0\quad\Longleftrightarrow\quad a=J=0,
 \qquad -C\eta_i c_*\leq H\leq Cc_*
                       \quad\text{on this zero set}.
\]
The lower bound is part (iii) of the cutoff lemma, using the assumed
uniform \eqref{eq:C-alpha}.  Thus the exact sign conclusion of
Lemma~\ref{lem:A-linear-chamber} over a point prohibits $H<0$ at the
lift of $T$, contradicting \eqref{eq:C-negative-cutoff-order}.
This proves (ii).

The proof of (iii), including its common normalization $v_i$, is given
below after establishing the adjoint lattice-point statement.  For the
extension to rational monomial subcones, take any finite list of tested
functions, resolve their divisors, and subdivide the cone so their
orders are linear on each piece.  Divisorial rational rays are dense
in each rational piece.  The exact zero assertions extend by continuity;
applying this separately to each function gives triviality on the
whole field.
\end{proof}

\subsection{Adjoint lattice points for integral Weil polarizations}
\label{subsec:C-adjoint}

We prove the precise lattice-point supply needed for the lower count.
The polarization will be nef and big and may be an integral Weil
divisor with unbounded Cartier index.  The rounding and pushdown steps
below keep those hypotheses throughout.

The weighted-blowup and adjoint-lifting construction below adapts the
``Adjoint lattice points'' and ``Orbifold vanishing'' lemmas in
\cite{OpenAIAdjointSemigroup}.  Here the argument is developed for a nef
and big integral Weil polarization with a klt Calabi--Yau boundary,
including the rounding and global regularity checks given below; the
earlier theorem is not used as a black box.

\begin{lemma}[Vanishing on a smooth orbifold]\label{lem:C-orbifold-vanishing}
Let $\mathcal X$ be a smooth proper Deligne--Mumford stack over an
algebraically closed field of characteristic zero, with trivial generic
stabilizer and projective coarse space.  Let $\Theta$ be an effective
rational divisor with simple normal crossing support and coefficients
less than one.  If $D$ is an integral Cartier divisor and
$D-(K_{\mathcal X}+\Theta)$ is nef and big, then
$H^q(\mathcal X,\mathcal O_{\mathcal X}(D))=0$ for $q>0$.
\end{lemma}
\begin{proof}
Write $q:\mathcal X\to X$ for the coarse morphism.  For each prime $B$
of $X$, let $P$ be its reduced inverse-image divisor and $e_B$ its
generic stabilizer order.  Thus $q^*B=e_BP$.  Rational divisors on
$\mathcal X$ descend uniquely to rational divisors on $X$; use a
subscript $c$ for the descended divisor.  The local normal-direction
equation $t=z^{e_B}$ gives
\[
 K_{\mathcal X}=q^*(K_X+R),\qquad
 R=\sum_B(1-e_B^{-1})B.
\]
Choose compatible canonical divisors using one rational top form.
Set $\Delta=R+\Theta_c-\{D_c\}$.  If the coefficients of $D$ and
$\Theta$ at $P$ are $d\in\Z$ and $\theta\in[0,1)$, write
$d=e_Bn+s$ with $0\le s<e_B$, also when $d$ is negative.  Then
\[
 \operatorname{coeff}_B\Delta
    =\frac{e_B-1+\theta-s}{e_B}\in[0,1).
\]
In particular $\Delta$ is effective.

The pair $(X,\Delta)$ is klt, including at higher-codimensional quotient
singularities.  Indeed, etale locally on $X$ there is a finite quotient
presentation $[U/G]$, with $U$ smooth and $f:U\to U/G$ finite.  Its
log pullback boundary is
$\Theta_U-f^*\{D_c\}$.  This possibly negative boundary is bounded
above by the simple normal crossing boundary $\Theta_U$.  Thus for
every divisorial valuation $F$ over $U$,
\[
 A(F;U,\Theta_U-f^*\{D_c\})
   =A(F;U,\Theta_U)+F(f^*\{D_c\})>0.
\]
The finite-map discrepancy formula multiplies downstairs log
discrepancies by positive ramification indices, so the downstairs
pair is klt.  This argument requires only local quotient charts.

The coarse pushforward is
\[
 q_*\mathcal O_{\mathcal X}(D)=\mathcal O_X(\lfloor D_c\rfloor).
\]
To check it, a rational section $h$ is regular along $P$ precisely when
$e_B\operatorname{ord}_B(h)+d\ge0$, equivalently
$\operatorname{ord}_B(h)+\lfloor d/e_B\rfloor\ge0$.  Both sides
extend these codimension-one conditions across codimension at least two:
on a smooth finite quotient chart this follows by normal extension and
then taking invariants.  The stack is tame in characteristic zero, so
coarse pushforward is exact and identifies the cohomology groups.

Quotient singularities are locally $\Q$-factorial.  Consequently
$M=\lfloor D_c\rfloor$ is an integral $\Q$-Cartier Weil divisor,
and
\[
 M-(K_X+\Delta)
   =D_c-(K_X+R+\Theta_c)
\]
is nef and big.  Kawamata--Viehweg vanishing for a klt pair and an
integral $\Q$-Cartier Weil divisor gives the result.  We use \cite[Theorem~3.2.4, manuscript version of 9 October 2023]{FujinoVanishing};
the additional log-bigness condition on log canonical centres is
vacuous for a klt pair.  The characteristic-zero statement also follows
from its complex form by descending the finite-type data to a finitely
generated field and using invariance of nefness, bigness and cohomology
under algebraically closed field extension.
\end{proof}

\begin{lemma}[Adjoint Taylor exponents]\label{lem:C-adjoint-lattice}
Let $X$ be a projective $\Q$-factorial variety of Fano type of
dimension $p>0$, and let $N$ be a nef and big integral Weil divisor.
Suppose that $H^0(X,\mathcal O_X(N))$ defines a birational map.
Choose an effective rational boundary $\Delta$ with $(X,\Delta)$ klt
and $K_X+\Delta\sim_{\Q}0$.  At a general smooth point outside
the supports of $N$ and $\Delta$, fix regular parameters and take
total-degree-then-lexicographic Taylor initial exponents, using the
compatible local trivializations of the powers of $N$.  Let
\[
 K=\overline{\operatorname{conv}}
 \left\{\frac{\nu(s)}j:
 j\in\Z_{>0},\quad
 0\ne s\in H^0(X,\mathcal O_X(jN))\right\}.
\]
If $t$ is a positive integer and $\alpha\in\Z_{\ge0}^p$ satisfies
$\alpha+\mathbf1\in\operatorname{int}(tK)$, then some section of
$tN$ has initial exponent $\alpha$.
\end{lemma}
\begin{proof}
The shift by $\mathbf1$ supplies the adjoint correction in this
argument.  A weighted blowup will convert it into the canonical
discrepancy, and the same shift will place the required monomial in a
multiplier ideal on the exceptional divisor.  Vanishing then lifts
that monomial to a global section.

\medskip
\noindent\emph{Preparing sections with weighted initials.}
Resolve $(X,\Delta)$ by $f:Y\to X$, isomorphically over a neighborhood
of the chosen point $x$.  Put $N_Y=f^*N$ and define the log pullback
$\Delta_Y$ by
\[
 K_Y+\Delta_Y=f^*(K_X+\Delta).
\]
The exceptional coefficients of $\Delta_Y$ may be negative.  This
will cause no difficulty: our vanishing boundary will be a fractional
rounding correction, not $\Delta_Y$.  Both $N_Y$ and $\Delta_Y$ have
zero divisor support near $x$ after compatible local trivialization.
Choose canonical divisors with no support near $x$ as well.

Let $C$ be the closed cone generated by the actual pairs $(j,\nu(s))$.
It has slice $K$ at degree one.  Choose rational $0<t_0<t$ with
$(t_0,\alpha+\mathbf1)\in\operatorname{int}C$.  There is a cone
generated by finitely many actual pairs that already contains this
point in its interior: approximate the vertices of a small simplex
around the point by finite nonnegative combinations of actual pairs.
Sufficiently close approximations still surround the point.

Choose positive primitive integral weights $u_1,\ldots,u_p$ that
replace the Taylor order by scalar weight on this finite list, each
listed section having its specified initial as a single weight initial.
Also require
\[
 \gamma\text{ precedes }\alpha\quad\Longrightarrow\quad
        u\cdot\gamma<u\cdot\alpha.
\]
Here is the finite-order justification.  Take an integer $d$ at least
the total degree of every listed initial and of $\alpha$.  Weights
sufficiently close to $(1,\ldots,1)$ preserve every strict total-degree
comparison through degree $d$, and satisfy
$(d+1)\min u_i>d\max u_i$, so all higher-degree monomials are larger
than every relevant initial.  Successive sufficiently small rational
perturbations realize lexicographic comparisons within the finitely
many monomials of degree at most $d$.  Clearing denominators and
dividing by the common divisor preserves all strict comparisons.
This also proves the assertion for formal Taylor expansions.

Put $p_0=u\cdot\alpha$ and $U=u\cdot(\alpha+\mathbf1)$.
The finite cone meets the plane with degree $t_0$ and weight $U$ in
a neighborhood of $(t_0,\alpha+\mathbf1)$ relative to that entire
plane.  Choose finitely many rational points surrounding it there.
Write these points as nonnegative rational combinations of the cone
generators and clear denominators.  Taking products, and then common
powers as necessary, gives a sufficiently divisible positive integer
$l$ and sections of the Cartier divisor $lt_0N_Y$ with single weight
initials $z^{\beta^{(1)}},\ldots,z^{\beta^{(r)}}$ such that
\[
 u\cdot\beta^{(a)}=lU,\qquad
 \alpha+\mathbf1\in
 \operatorname{relint}\operatorname{conv}
       \{\beta^{(1)}/l,\ldots,\beta^{(r)}/l\}.
\]
The relative interior contains a neighborhood in the whole plane
$u\cdot v=U$.  Pulling these sections to $Y$ is legitimate after
clearing denominators: the effective divisor of each product is
$\Q$-Cartier and its rational pullback is effective.

\medskip
\noindent\emph{The weighted modification and the vanishing divisor.}
Let $\pi:\mathcal Z\to Y$ be the smooth stack weighted blowup at $x$
with weights $u$, and let $E$ be its exceptional Cartier divisor.
On the coordinate model it is
\[
 \left[(\mathbb A^{p+1}_{x_1,\ldots,x_p,\tau}
       \setminus\{x_1=\cdots=x_p=0\})/\mathbb G_m\right],
 \qquad
 \rho\cdot(x_i,\tau)=(\rho^{u_i}x_i,\rho^{-1}\tau),
 \qquad z_i=\tau^{u_i}x_i.
\]
Its exceptional divisor is the weighted projective stack $\mathbb P(u)$.
This description pulls back along an etale parameter map near $x$
and glues to the identity away from $x$.  Equivalently use the weighted
ideal filtration generated by monomials of weight at least $k$; each
ideal is supported on a finite infinitesimal neighborhood of $x$ and
extends as the unit ideal outside $x$.  The finitely generated weighted
Rees algebra makes the coarse modification projective.  The stack is
smooth with finite stabilizers and trivial generic stabilizer.  On the
displayed atlas the Jacobian gives
\[
 K_{\mathcal Z}=\pi^*K_Y+(\textstyle\sum_i u_i-1)E.
\]
For $p=1$ this construction is the identity modification and $E=x$;
the same formulas and restriction sequence apply.

Write $\Delta_{\mathcal Z}=\pi^*\Delta_Y$, supported away from $E$,
and define
\[
 P_0=\pi^*(tN_Y)-p_0E,\qquad
 H_0=\pi^*(lt_0N_Y)-lUE.
\]
The divisor $P_0$ is rational globally but integral Cartier near $E$;
$H_0$ is Cartier globally.  The selected sections, after their common
weighted order is removed, are sections of $H_0$.  Their base ideal
$\mathfrak b$ does not vanish identically on $E$, and its restriction
to $E$ is generated by their monomial initials.  Take a projective
log resolution $\sigma:\mathcal Y\to\mathcal Z$ of these data and
all rational divisor supports used below, with strict transform
$\widetilde E$.  Arrange that the total support is simple normal
crossing, that $\sigma$ is an isomorphism at the generic point of
$E$, and that $\tau:\widetilde E\to E$ resolves $\mathfrak b|_E$.
First resolve the additional divisor supports, then principalize the
ideal with the resulting simple normal crossing boundary, including
the transform of $E$, on smooth charts
\cite[Theorems~2.4.1--2.4.3]{WlodarczykResolution}.  Smooth functoriality
identifies these constructions on chart overlaps, so they descend to
the required smooth stack; the sequence of blowups has projective
coarse space.  Write
\[
 \mathfrak b\mathcal O_{\mathcal Y}=\mathcal O_{\mathcal Y}(-G).
\]
Then $G\ge0$ does not contain $\widetilde E$, and
$\sigma^*H_0-G$ is globally generated.

Set
\[
 Q=P_0-K_{\mathcal Z}-\Delta_{\mathcal Z}-E,
 \quad R=\sigma^*Q-G/l,
 \quad A=K_{\mathcal Y}+\lceil R\rceil,
 \quad\Theta=\lceil R\rceil-R.
\]
Thus $A$ is an integral Cartier divisor, and $\Theta$ is effective
with simple normal crossing support and coefficients in $[0,1)$.
The weighted canonical identity, $U-p_0=\sum u_i$, and
$K_Y+\Delta_Y\sim_{\Q}0$ give
\[
 Q-H_0/l\sim_{\Q}(t-t_0)\pi^*N_Y.
\]
Consequently
\[
 A-(K_{\mathcal Y}+\Theta)=R
 \sim_{\Q}
 (t-t_0)\sigma^*\pi^*N_Y+(\sigma^*H_0-G)/l.
\]
The first summand is nef and big and the second is nef.  Lemma~\ref{lem:C-orbifold-vanishing} therefore gives $H^1(\mathcal Y,\mathcal O(A))=0$ and
the restriction surjection
\[
 H^0(\mathcal Y,A+\widetilde E)
 \longrightarrow
 H^0(\widetilde E,(A+\widetilde E)|_{\widetilde E}).
\]

\medskip
\noindent\emph{The boundary monomial.}
We now exhibit the required section on the exceptional divisor.
Near $E$ the original
boundary vanishes and $P_0$ is integral.  Thus, on a neighborhood of
$\widetilde E$,
\[
 A=K_{\mathcal Y}
   +\sigma^*(P_0-K_{\mathcal Z}-E)-\lfloor G/l\rfloor.
\]
Adjunction identifies the restricted bundle with
\[
 \tau^*(P_0|_E)+K_{\widetilde E}-\tau^*K_E
                   -\lfloor(G|_{\widetilde E})/l\rfloor.
\]
Round-down commutes with this restriction: simple normal crossings
prevent two components of $G$ from restricting to a common prime
divisor on $\widetilde E$, and each disconnected intersection retains
the same coefficient.  This is the log-resolution expression for
$P_0|_E\otimes\mathcal J_E((\mathfrak b|_E)^{1/l})$.

The weighted-homogeneous monomial $z^\alpha$ is a global section of
$P_0|_E=\mathcal O_E(p_0)$, after the fixed local frame.  On the chart
$x_k\ne0$, use the smooth finite atlas setting $x_k=1$, with residual
group $\mu_{u_k}$.  The ideal becomes the ordinary monomial ideal
generated by $x_{\mathrm{rem}}^{\beta^{(a)}_{\mathrm{rem}}}$ and the
proposed section becomes $x_{\mathrm{rem}}^{\alpha_{\mathrm{rem}}}$.
Dropping coordinate $k$ is an affine isomorphism from the plane
$u\cdot v=U$ to $\R^{p-1}$, because $u_k>0$.  Hence
$\alpha_{\mathrm{rem}}+\mathbf1$ lies in the ordinary interior of
the convex hull of the projected vectors $\beta^{(a)}_{\mathrm{rem}}/l$,
and therefore in the interior of the scaled Newton polyhedron.
Howald's monomial multiplier-ideal formula \cite[Main Theorem]{Howald}
gives the required membership
on this smooth atlas.  It is equivariant for $\mu_{u_k}$, so descends
on every chart.  No fractional coordinate exponents are introduced:
the root needed to set $x_k=1$ is already part of the atlas.

\medskip
\noindent\emph{Lifting and pushing down.}
The restriction surjection now lifts this monomial.  Near $E$ its
lifting bundle differs from $\sigma^*P_0$ by
\[
 K_{\mathcal Y}-\sigma^*K_{\mathcal Z}
  -\sigma^*E+\widetilde E-\lfloor G/l\rfloor.
\]
The part preceding the last term is exceptional over $\mathcal Z$;
indeed it is effective by the smooth blowup discrepancy calculation
for the smooth pair $(\mathcal Z,E)$.  At every nonexceptional prime
the displayed correction is nonpositive.  Therefore the lifted
rational section pushes down to a regular section of $P_0$ near $E$.
Its restriction to $E$ is $z^\alpha$, initially generically and then
everywhere by regularity.  Pushing further to the smooth neighborhood
of $x$ gives a regular section with weight initial exactly $z^\alpha$.

It remains to check global regularity downstairs, since $N_Y$ was
rational and $\Delta_Y$ could have negative coefficients.  At an
original prime $B$ of $X$ away from $x$, write
$n_B=\operatorname{coeff}_B(tN)\in\Z$,
$\delta_B=\operatorname{coeff}_B\Delta\ge0$, and let $g_B\ge0$ be
the corresponding coefficient of $G$.  Canonical coefficients cancel
inside the ceiling, and the allowed divisor coefficient of the
pushed lifting bundle is
\[
 n_B+\lceil-\delta_B-g_B/l\rceil\le n_B.
\]
Thus no forbidden pole appears at any original prime.  The point
$x$ was already handled by the local weighted calculation (also when
$p=1$).  Normality now gives a section of $\mathcal O_X(tN)$.
Its weight initial excludes every monomial preceding $\alpha$ in
the original Taylor order, by the choice of $u$, and the coefficient
of $z^\alpha$ is nonzero.  Its Taylor initial is therefore $\alpha$.
\end{proof}

\begin{lemma}[Uniform count from the jet body]\label{lem:C-jet-count}
In the setting of Lemma~\ref{lem:C-adjoint-lattice}, $K$ is bounded and contains the standard
simplex $\Sigma=\operatorname{conv}(0,e_1,\ldots,e_p)$.  For every
integer $n\ge1$,
\[
 h^0(X,nN)\le(n+p)^p\operatorname{vol}(K).
\]
Fix any integer $t>1+2p$.  The space spanned by products of $n$ sections
of $tN$ has dimension at least $n^p\operatorname{vol}(K)$.
\end{lemma}
\begin{proof}
On a nonempty open set, the map given by $|N|$ is an immersion and
has no base point, because it is birational in characteristic zero.
At the chosen general point its values and first derivatives thus
span all first jets.  Suitable linear combinations give initial
exponents $0,e_1,\ldots,e_p$, proving $\Sigma\subset K$.
The order of vanishing at this point of a section of $jN$ is at most
$Cj$: intersect its effective divisor with a general complete
intersection curve through the point, chosen not to lie in its
support.  Its degree is $jN\cdot H^{p-1}$ for a fixed sufficiently
ample $H$, and bounds the local intersection multiplicity.  The
constant may depend on this individual example, which suffices to
make $K$ bounded.

For a finite-dimensional space of germs, the number of distinct
leading exponents is its dimension, by successive elimination of
leading coefficients; alternatively use the filtration by the
well-ordered monomials, whose nonzero graded pieces are one-dimensional.
The exponent set at degree $n$ is contained in $nK\cap\Z^p$.
The disjoint half-open unit boxes based at these lattice points lie
in
\[
 nK+[0,1]^p\subset nK+p\Sigma\subset(n+p)K.
\]
Their total volume is their number, proving the upper bound.

For the lower bound let $y=nx\in nK$.  Write $x_i=a_i+f_i$ with
$a_i\in\Z_{\ge0}$ and $0\le f_i<1$.  Choose $n$ integers
in coordinate $i$, namely $a_i+1$ in exactly $\lfloor nf_i\rfloor$
of the positions and $a_i$ in the others.  Doing this independently
in every coordinate gives $\alpha^{(1)},\ldots,\alpha^{(n)}$ with
\[
 \sum_{r=1}^n\alpha^{(r)}=\lfloor y\rfloor,
 \qquad 0<\alpha_i^{(r)}+1-x_i<2.
\]
When $f_i=0$ all entries equal $a_i$, so the strict inequalities still
hold.  Therefore
$\alpha^{(r)}+\mathbf1-x\in\operatorname{int}((t-1)\Sigma)$:
all coordinates are positive and their sum is less than $2p<t-1$.
Since $\Sigma\subset K$ and $x\in K$, a neighborhood of
$\alpha^{(r)}+\mathbf1$ lies in
$K+(t-1)K=tK$.  Lemma~\ref{lem:C-adjoint-lattice} supplies a degree-$t$ section
with each exponent $\alpha^{(r)}$.  Their product has exponent
$\lfloor y\rfloor$.

Consequently all lattice points $\lfloor y\rfloor$, for $y\in nK$,
occur among product initials.  Their unit boxes cover $nK$, so their
number is at least $\operatorname{vol}(nK)=n^p\operatorname{vol}(K)$.
Sections with distinct initials are linearly independent, proving
the claimed lower bound.  The case $p=0$ is immediate separately.
\end{proof}

\subsection{Completion of the saturation count}

The field and exact-triviality assertions are already established.
We finish by placing both section counts on one nef integral Weil
polarization.  The preceding two lemmas then give the same jet-body
volume as the normalization in the upper and lower estimates.
\label{subsec:C-count-completion}

\begin{proof}[Proof of part \textup{(iii)} of Lemma~\ref{lem:C-saturation-count}]
If $p=0$, then $K_*=k$.  By part (i), each fixed-cost section space
is eventually the one-dimensional constant space.  Take $v_i=1$.
Assume $p>0$ and use the extraction $\widehat W$ from
Lemma~\ref{lem:C-cutoff-claims} at $\beta_1$ of
\eqref{eq:C-plateau}.  Let $D_*$ be the pullback of $D_W$ to
$\widehat W$.  For every rational $m>0$, the cutoff spaces identify
exactly as function spaces:
\begin{equation}\label{eq:C-cutoff-base-spaces}
                       S(mP_{\beta_1})=S(mD_*).
\end{equation}
The negative-MMP error is effective and exceptional, so pushing and
pulling the divisor inequality for an individual rational function
identifies the space with the one on the semiample output.  A function
in that space belongs to $k(W)$; the equality $P_{\beta_1}'=p^*D_W$
then identifies its inequalities with those downstairs.  Pullback to
$\widehat W$ does not change them.  This argument works for each
rational $m$, with no requirement that $mP_{\beta_1}$ be Cartier.

Every positive coefficient of $LD_*$ is at least $c'\delta$, because
$L\beta_1\geq c'$ and the coefficients of $D_*/\beta_1$ are at
least $\delta$.  Choose a fixed sufficiently large integer $M$ and put
\begin{equation}\label{eq:C-count-polarization}
                              N=\lfloor MLD_*\rfloor.
\end{equation}
The coefficient lower bound gives $N\geq LD_*$ once $M$ is large
enough.  Increase $M$ so that $N\geq C D_*/\beta_1$, with the
constant from the field-generation construction.  Then $S(N)$ defines
a birational map and $N$ is big and integral.  Run its MMP to a nef
and big integral Weil divisor $N'$ on a $\Q$-factorial projective
Fano type variety $X'$.  For every positive integer $j$, negativity
gives
\begin{equation}\label{eq:C-count-mmp-spaces}
                     S(jN)=S(jN'),\qquad
                              S(jN')\subset S(jMLD).
\end{equation}
Indeed the first equality follows by pushing an effective divisor and,
in the reverse direction, adding the effective exceptional MMP error;
the second follows from $N\leq MLD_*$ and
\eqref{eq:C-cutoff-base-spaces}.  For every fixed positive integer $n$,
part (i) and the exact plateau identities give eventually
\begin{equation}\label{eq:C-count-upper-inclusion}
               S(nLD)\subset S(nLD_*)\subset S(nN)=S(nN').
\end{equation}

Choose an effective ordinary rational klt Calabi--Yau boundary
$\Delta$ on $X'$.  Such a boundary is obtained from a klt Fano type
boundary by adding a general sufficiently divisible rational member
of its ample anti-log-canonical class.  At a general smooth point
outside the boundary and divisor supports, let $K_i$ be the jet body
of $N'$ in Lemma~\ref{lem:C-adjoint-lattice}.  Its volume is positive
and finite by Lemma~\ref{lem:C-jet-count}.  Set
\[
                              v_i=\vol(K_i).
\]
Equations \eqref{eq:C-count-upper-inclusion} and
Lemma~\ref{lem:C-jet-count} imply
\[
 \dim S(nLD)\leq(n+p)^p v_i
                    \leq(p+1)^p v_i n^p.
\]
Fix once and for all an integer $t>1+2p$.  The same lemma supplies a
space of products of $n$ sections of $tN'$ of dimension at least
$v_i n^p$.  By \eqref{eq:C-count-mmp-spaces},
$S(tN')\subset S(tMLD)$.  Enlarge the constant $B$ in part (ii) to
at least $tM$.  Those products then lie in the product space in
\eqref{eq:C-uniform-counts}, proving the lower bound with the same
$v_i$.  Both $t$ and $M$, and hence $B$, were fixed before $n$.
This completes the lemma.
\end{proof}

\begin{corollary}[Fixed generators and neutral algebraic extensions]
\label{cor:C-neutral-field}
In Lemma~\ref{lem:C-saturation-count}, let $g_{i,1},\ldots,g_{i,r_i}$
be any finite list contained in $S(C_1L_iD_i)$ for one fixed $C_1$.
For all sufficiently large $i$, the entire field
$F_i=k(g_{i,1},\ldots,g_{i,r_i})$, and its relative algebraic closure
inside $k(V_i)$, lie in $K_{*,i}$.  Thus every test from part
\textup{(ii)} is exactly zero on every nonzero element of these fields
at that same index, without a degree-dependent eventual qualification.
More generally, a valuation trivial on a subfield remains trivial on
any algebraic extension of it inside its valued overfield.
\end{corollary}

\begin{proof}
Part (i) of the lemma includes the whole space $S(C_1L_iD_i)$ at a
single eventual index; it does not supply separate indices for its
individual elements.  Field closure gives $F_i\subset K_{*,i}$.
The cutoff field is relatively algebraically closed, so it contains
the stated relative algebraic closure as well.  To see the last
assertion directly, let $x\ne0$ be algebraic and choose a polynomial
relation for it with nonzero constant term and coefficients in the
trivially valued subfield.  If the value of $x$ were positive, the
constant term would have uniquely least value; if it were negative,
the highest nonzero term would have uniquely least value.  Neither is
possible in a vanishing sum.  Hence its value is zero.  This applies
to every rational expression and every polynomial degree once the
finite generators have been included, rather than by a further limit
in that degree.
\end{proof}

\section{Characters and regularization on affine cones}\label{sec:D}

We translate the projective results of Sections~\ref{sec:B} and~\ref{sec:C}
into statements about homogeneous functions on affine cones.  The distinction
between a primitive prime order and a rescaled valuation is essential in this
translation.

\subsection{Angular fields and the discrepancy normalization}

Let $\breve R$ be a finitely generated normal domain over $k$, equipped with a
positive torus grading and a commuting diagonal action of a finite cyclic
group $G$.  Positivity means that there is a linear functional strictly
positive on every nonzero occurring weight and that the degree-zero ring is
$k$.  Set $R=\breve R^G$ and $X=\Spec R$, and assume that $X$ is
$\Q$-Gorenstein and klt.  We use the effective torus acting on $R$, with
character lattice $M$ of rank $r_T>0$.  Write
\[
 C=\operatorname{cone}\{m:R_m\ne0\}\subset M_{\R},\qquad
 \sigma=C^\vee\subset\operatorname{Hom}(M,\R),\qquad
 \mathcal K=\Frac(R)^{\mathbb T}.
\]
Thus $C$ is a full-dimensional pointed rational polyhedral cone.  Choose a
multiplicative splitting $m\mapsto z^m$ of the homogeneous rational
functions.  Such a splitting exists because $M$ is free abelian, and it gives
\begin{equation}\label{eq:D-angular-field}
 \Frac(R)=\mathcal K(z^M).
\end{equation}

Let $h\in\breve R$ be a nonzero homogeneous function of $G$-character $\chi$
and torus weight $\mathbf u\in M_{\Q}$.  If $r_G$ is the order of $G$, set
$H=r_G^{-1}\divisor(h^{r_G})$ on $X$.  Assume that there is a torus-invariant
rational top form $\theta$ on $X$ such that
\begin{equation}\label{eq:D-log-form}
 K_X(\theta)+H=0,\qquad (X,H)\text{ is lc}.
\end{equation}
Evaluation of a semi-invariant function always means evaluation of an
invariant power, divided by its exponent.  In particular, writing
$h^{r_G}=f_hz^{r_G\mathbf u}$ gives
\begin{equation}\label{eq:D-height-evaluation}
 (P,y)(h)=\gamma(P)+\langle y,\mathbf u\rangle,
 \qquad \gamma=\frac1{r_G}\divisor(f_h).
\end{equation}

For a primitive divisorial valuation $P$ of $\mathcal K$, define
\begin{equation}\label{eq:D-polyhedron}
 \Delta_P=\left\{y\in\operatorname{Hom}(M,\R):
 P(f)+\langle y,m\rangle\ge0\text{ whenever }fz^m\in R\right\}.
\end{equation}
The inequalities for finitely many homogeneous algebra generators suffice.
The polyhedron is nonempty: adding a sufficiently large multiple of a
strictly positive grading satisfies all these inequalities.  For rational
$y\in\Delta_P$, the Gauss rule
\[
 (P,y)\left(\sum_m f_mz^m\right)
   =\min_{f_m\ne0}\bigl(P(f_m)+\langle y,m\rangle\bigr)
\]
defines a valuation nonnegative on $R$.  After a positive rational
rescaling it is a primitive divisorial valuation.  For trivial restriction
to $\mathcal K$, the corresponding Gauss valuations are the nonzero rational
pure weights $y\in\sigma$, with value $\langle y,m\rangle$ on $fz^m$.

Choose a basis $z_1,\ldots,z_{r_T}$ of the weight monomials.  Torus
invariance gives a rational top form $\theta_0$ on $\mathcal K$ such that
\[
 \theta=\theta_0\wedge d\log z_1\wedge\cdots\wedge d\log z_{r_T}.
\]
Here $\theta_0$ is a nonzero constant when $\mathcal K=k$.  For a primitive
$P$ put $a(P)=\nu_{\log,P}(\theta_0)$, where log order is the order of a top
form plus one.  Extend $a$ homogeneously to positive multiples of prime
orders.

\begin{proposition}[Cone dictionary]\label{prop:D-cone-dictionary}
In the above setting the following statements hold.
\begin{enumerate}[label=\textup{(\roman*)}]
\item The discrepancy of $(X,H)$ at $(P,y)$ is $a(P)$, and its discrepancy
at a pure weight is zero.  Thus $a(P)$ is a nonnegative integer for every
primitive $P$, and $\mathbf u\in\operatorname{int}C$.
\item For rational $\mu\in\operatorname{int}C$, the ray quotient $V_\mu$ is
normal and projective, with function field $\mathcal K$.  The order function
\begin{equation}\label{eq:D-polarization}
 \mathbf L(\mu)(P)=\min_{y\in\Delta_P}\langle y,\mu\rangle
\end{equation}
is the b-pullback of an ample rational divisor on $V_\mu$.  Consequently
\begin{equation}\label{eq:D-height-divisor}
 \mathbf D(P)=\min_{y\in\Delta_P}(P,y)(h)
             =\gamma(P)+\mathbf L(\mathbf u)(P)
\end{equation}
is the b-pullback of an effective ample rational divisor $D$ on
$V_{\mathbf u}$.
\item At every actual prime $P$ of $V_\mu$, a minimizer $y$ of
\eqref{eq:D-polarization} can be chosen so that $e(P,y)$ is a primitive
prime order on $X$, for some positive integer $e$.  For this choice,
\begin{equation}\label{eq:D-actual-prime}
 a(P)+(P,y)(h)=\frac1e,
 \qquad e y\in\operatorname{Hom}(M,\Z).
\end{equation}
On each $V_\mu$, $a$ is the discrepancy function of an effective ordinary
lc Calabi--Yau pair, and $a+\mathbf D$ is that of an effective generalized
klt Calabi--Yau pair with nef part $\mathbf D$.  Each $V_\mu$ is of Fano
type.  At an actual prime outside the floor $\{a=0\}$, the same minimizer
satisfies
\begin{equation}\label{eq:D-off-floor}
 a(P)=1,\qquad e=1,\qquad (P,y)(h)=0.
\end{equation}
\end{enumerate}
\end{proposition}

\begin{proof}
Take a uniformizer $t$ at $P$ and separating residue coordinates
$s_1,\ldots,s_{p-1}$, where $p=\trdeg_k\mathcal K$.  Up to a unit the base
form is
$t^{a(P)}d\log t\wedge ds_1\wedge\cdots\wedge ds_{p-1}$.
For a rational Gauss extension, integral powers of the $z_j$ can be
adjusted by powers of $t$ to obtain angular units with algebraically
independent residues.  The resulting logarithmic wedge has log order zero.
Together with \eqref{eq:D-log-form}, this proves the discrepancy formula.
The calculation with trivial restriction to $\mathcal K$ is the same,
using a basis of the kernel of the pure weight.  Nonnegativity follows
from lc.  Integrality is the ordinary integral log order of $\theta_0$.
At a nonzero rational pure weight $y$, klt of $X$ gives
$A_X(y)=y(h)=\langle y,\mathbf u\rangle>0$.  Positivity on the nonzero
rays of $\sigma$ is equivalent to $\mathbf u\in\operatorname{int}C$.

Choose a positive integer $d$ with $d\mu\in M$ and set
\[
 V_\mu=\Proj\bigoplus_{j\ge0}R_{jd\mu}.
\]
The ray ring is finitely generated and normal: it is the appropriate
diagonalizable invariant subring, followed, if necessary, by a Veronese
subring.  Its degree-zero part is $k$, so its Proj is projective.  An
invariant rational function is a ratio of two homogeneous elements of the
same weight.  Multiplying their common weight by a suitable occurring
weight moves both elements onto a sufficiently divisible multiple of the
interior ray.  To see the latter assertion directly, express $\mu$ as a
strictly positive rational combination of all the finitely many generating
weights and clear denominators; sufficiently high multiples leave room to
subtract any prescribed occurring weight.  This proves
$k(V_\mu)=\mathcal K$.

Write the algebra generators as $f_jz^{m_j}$.  Linear programming duality
gives
\begin{equation}\label{eq:D-linear-program}
 \min_{y\in\Delta_P}\langle y,\mu\rangle
 =\max\left\{-\sum_j\lambda_jP(f_j):
       \lambda_j\ge0,\ \sum_j\lambda_jm_j=\mu\right\}.
\end{equation}
The polytope on the right is compact because the grading is positive.
Its rational vertices depend on the weights and $\mu$, not on $P$.
After clearing their denominators and taking a Veronese generated in
degree one, these vertices give actual monomials on the ray.  Conversely,
a sum of monomials has no worse pole order than its worst summand.
Therefore \eqref{eq:D-linear-program} is exactly the maximum allowed pole
order in the homogeneous linear system, expressed using $z^{d\mu}$ as
rational frame.  On the Proj this system is basepoint-free in sufficiently
divisible degree.  Its maximum-pole formula on every higher model is the
pullback of its ample tautological divisor.  This proves (ii), including
the b-divisor assertion.  Formula \eqref{eq:D-height-evaluation} identifies
$\mathbf D$ with a rational principal adjustment of this polarization.
Since $h$ is regular upstairs, its invariant powers have nonnegative
value at every Gauss extension, so $\mathbf D\ge0$.

For (iii), choose a homogeneous section $s$ on the defining ray whose
nonvanishing chart contains the generic point of $P$.  The map
$\Spec R_s\to\Spec(R_s)_0$ is the torus quotient, and the latter is the
corresponding chart of $V_\mu$.  Localize the base at $P$.  Its uniformizer
remains a nonunit upstairs: the invariant ring is a direct summand, and
an inverse upstairs would give an inverse downstairs by taking its
invariant part.  A component of its zero divisor therefore has height
one and dominates $P$, by the principal ideal theorem.  Such a component
is torus invariant, since a connected torus cannot permute its finitely
many irreducible components nontrivially.

The primitive order of a torus-invariant prime is Gauss.  Indeed its
valuation ideals are torus invariant, hence are direct sums of weight
spaces; there is no cancellation between different weights, and on each
weight space restriction to $\mathcal K$ supplies the indicated base
order.  Write the resulting valuation as $e(P,y)$.  Its restriction to
$\mathcal K$ is an integral multiple $eP$, and its values on $z^M$ are
integral.  Invertibility of $s$ implies that $y$ minimizes the ray
polarization.  A prime on a normal variety has ordinary log discrepancy
one.  Thus
$1=e\bigl(a(P)+(P,y)(h)\bigr)$, which proves
\eqref{eq:D-actual-prime}.  Notice that this argument concerns actual
primes on $V_\mu$; it asserts no such identity at an arbitrary exceptional
prime.

On $V_\mu$ set $\Omega=-K_{V_\mu}(\theta_0)$.  The log discrepancies of
$(V_\mu,\Omega)$ are $a$, and $K_{V_\mu}(\theta_0)+\Omega=0$.
At actual primes \eqref{eq:D-actual-prime} gives $0\le a\le1$; hence
$\Omega$ is effective.  Moreover
$a(P)+\mathbf D(P)\le1$ there.  At every divisorial $P$, a rational
minimizer for \eqref{eq:D-height-divisor} gives a Gauss extension, so klt
of $X$ implies $a(P)+\mathbf D(P)>0$.  Subtracting the trace of
$\mathbf D$ from $\Omega$ and taking $\mathbf D$ as nef part gives the
claimed generalized klt Calabi--Yau data.

For completeness these data imply Fano type without any uniform index
assumption.  On a common resolution on which $\mathbf D$ is semiample,
replace it by a general divided member of a sufficiently divisible free
system.  Added to the generalized boundary this gives an ordinary klt
sub-pair; its pushdown has effective boundary and is crepant, with total log
canonical divisor rationally trivial.  The added divisor is big.
Choose an effective ample divisor $A$ on $V_\mu$.  Bigness permits a second
effective representative of $\mathbf D$ containing a small positive
multiple of the pullback of $A$.  Mixing this representative with the
general free representative with sufficiently small positive weight
preserves klt.  The resulting effective boundary downstairs contains
$\varepsilon A$ for some $\varepsilon>0$.  Subtracting that summand leaves
an effective klt boundary whose negative log canonical class is
$\Q$-linearly equivalent to $\varepsilon A$, hence ample.  Finally, at
a non-floor prime the integer $a(P)>0$ in
\eqref{eq:D-actual-prime} forces \eqref{eq:D-off-floor}.
\end{proof}

\subsection{Which divisorial lc tests occur}

\begin{lemma}[All divisorial lc tests are Gauss]\label{lem:D-lc-tests}
Every nontrivial divisorial valuation nonnegative on $R$ and having
zero discrepancy for $(X,H)$ is Gauss.  Its restriction to $\mathcal K$
is either $eP$ with $a(P)=0$, or trivial, in which case it is a nonzero
pure weight in $\sigma$.  Conversely every rational Gauss extension
with $a(P)=0$, and every nonzero rational pure weight, has zero
discrepancy and a center whose closure contains the contracting vertex.

If $\Spec\breve R\to X$ is quasi-\'{e}tale and the pairs are related by
log pullback, every divisorial $h$-lc test upstairs restricts to one of
these tests, up to scaling.  Conversely all the rational tests just
listed lift to divisorial $h$-lc tests upstairs whose center closures
contain the upstairs vertex.  These statements suffice in particular
for testing invariant rational functions.
\end{lemma}

\begin{proof}
Let $v$ be a primitive divisorial valuation as in the first assertion.
By the divisorial restriction statement in Section~\ref{sec:A}, a
nontrivial restriction to $\mathcal K$ has the form $eP$, with $P$
primitive and $e$ a positive integer.  Use the uniformizer $t$ and
separating residue coordinates from the proof of
Proposition~\ref{prop:D-cone-dictionary}.  Every remaining logarithmic
differential has nonnegative log order at $v$.  Thus
\begin{equation}\label{eq:D-log-defect}
 A_{X,H}(v)\ge e a(P).
\end{equation}
More explicitly, pass to the finite-index sublattice of $M$ on which
$v(z^m)$ is divisible by $e$, and replace each of its basis monomials
by the unit $z^m t^{-v(z^m)/e}$.  If their residues are algebraically
dependent over the residue field of $P$, their leading differential
wedge vanishes, so the inequality in \eqref{eq:D-log-defect} is strict.
Zero discrepancy therefore forces $a(P)=0$ and independence of these
angular unit residues.  For a finite sum of tied terms, divide by a
term of least value.  All differences of its exponents lie in that
finite-index lattice; independence of the adjusted residues prevents
cancellation.  Thus $v$ is exactly the Gauss valuation determined by
its restrictions to $\mathcal K$ and $z^M$.

If the restriction to $\mathcal K$ is trivial, use a separating
transcendence basis of $\mathcal K$ in the same wedge computation.
If all angular weights were zero, the unit residues could have
transcendence degree at most $r_T-1$ over $\mathcal K$, and the leading
top wedge would vanish.  Zero log discrepancy therefore forces a
nonzero angular weight.  It further forces algebraic independence,
over $\mathcal K$, of the residues of a basis of the weight-zero
monomials.  The same tied-sum argument now proves that $v$ is a pure
Gauss weight.  Nonnegativity on $R$ places that weight in $\sigma$.

The converse discrepancy assertion was computed in
Proposition~\ref{prop:D-cone-dictionary}.  The center of a Gauss
valuation is torus invariant.  A positive grading contracts its
closure to the vertex, so the vertex lies in that closure; the
valuation need not itself be centered at the closed vertex.

For a finite quasi-\'{e}tale map with log-pulled pairs, the ramification
formula makes discrepancies homogeneous under restriction.  Hence
an upstairs lc valuation restricts to a downstairs lc valuation,
and every divisorial extension of a downstairs lc valuation is lc
upstairs.  Such extensions exist by taking the normalization of a
model carrying the downstairs prime.  The finite image of an upstairs
center closure is the downstairs center closure and contains the
vertex.  There is only one point upstairs over that vertex: every
positive homogeneous semi-invariant generator has an invariant
positive power, which vanishes there.  Therefore the upstairs center
closure contains the upstairs vertex as claimed.
\end{proof}

\subsection{Sections and a uniform sandwich argument}

Let $\xi=fz^k\in\Frac(R)$ be a nonzero rational eigenfunction, invariant
under the finite group, and let $m\in\Z$ satisfy
$\mu=k+m\mathbf u\in\operatorname{int}C$.  Define
\begin{equation}\label{eq:D-section-divisor}
 \mathbf J_m(P)=\min_{y\in\Delta_P}(P,y)(\xi h^m)
   =\divisor(f)(P)+m\gamma(P)+\mathbf L(k+m\mathbf u)(P).
\end{equation}
This is the b-pullback of an ample rational divisor $J_m$ on $V_\mu$;
the divisor itself may be signed.  At a non-floor prime of $V_\mu$ it has
integral coefficient, by \eqref{eq:D-off-floor} and the integral value of
the finite-group invariant function $\xi$ at a primitive cone prime.

For a rational b-divisor $\mathbf A$ that descends to a model, write
\[
 S(\mathbf A)=\{0\}\cup
 \{g\in\mathcal K^*: \divisor(g)+\mathbf A\ge0\}.
\]
Testing on its determining model is equivalent to testing all divisorial
orders.  This notation agrees with that of Section~\ref{sec:C} and also
applies to signed divisors.

\begin{lemma}[Sections as regular symbols]\label{lem:D-section-symbols}
For every positive integer $q$ and every $g\in\mathcal K$,
\begin{equation}\label{eq:D-symbol-equivalence}
 g\in S(q\mathbf J_m)
 \quad\Longleftrightarrow\quad
 g\xi^q h^{qm}\in\breve R.
\end{equation}
In particular the neutral functions of cost $b\in\N$ are precisely
\begin{equation}\label{eq:D-neutral-functions}
 \{g\in\mathcal K:g h^b\in\breve R\}=S(b\mathbf D).
\end{equation}
\end{lemma}

\begin{proof}
Regularity implies all Gauss inequalities, hence the implication to the
left.  Conversely take a sufficiently divisible power so that the
semi-invariant becomes invariant, its weight is on the integral ray, and
its degree is large enough for the ray ring to agree with the section ring
of the tautological polarization in that degree.  The divisor inequalities
and \eqref{eq:D-linear-program} then put that power in $R$.  The original
function belongs to $\Frac(\breve R)$ and is integral over $\breve R$, so
normality puts it in $\breve R$.  The second formula is the case $\xi=1$.
When $\mathcal K=k$, the same argument says that an interior weight becomes
effective in a divisible degree and then uses normality.
\end{proof}

All statements that follow are sequence statements in bounded dimension.
Constants implicit in $O(\cdot)$ are uniform in the sequence, and
hypotheses about small divisorial orders are imposed along every
subsequence, as in Sections~\ref{sec:A}--\ref{sec:C}.

\begin{proposition}[Sandwich regularization]\label{prop:D-sandwich}
Let $L\ge1$ vary with the sequence.  Choose integers $m>m_0$ with
$m-m_0\ge cL$ for a fixed $c>0$, and suppose both
$k+m_0\mathbf u$ and $k+m\mathbf u$ are interior.  Assume that whenever
primitive divisorial valuations $P$ satisfy $a(P)=0$ and
$L\mathbf D(P)\to0$, there are integers $j_P$ such that
\begin{equation}\label{eq:D-sandwich-hypothesis}
 j_P\le\mathbf J_{m_0}(P)\le\mathbf J_m(P)
       \le j_P+O(L\mathbf D(P)).
\end{equation}
Then there is a bounded positive integer $q$ and a nonzero
$g\in S(q\mathbf J_m)$.  Equivalently, $g\xi^qh^{qm}$ is regular on
$\Spec\breve R$.
\end{proposition}

\begin{proof}
For $\mathcal K=k$ the preceding lemma gives the assertion with $q=1$.
Otherwise work first on $V_\mu$, where $\mu=k+m\mathbf u$.  Use the
generalized klt Calabi--Yau discrepancy $a+\mathbf D$ to extract onto a
$\Q$-factorial Fano type model $U$ all exceptional primes below a fixed
small positive discrepancy cutoff, and no other exceptional primes.
The extraction statement is supplied by Section~\ref{sec:A}.
All extracted primes are floor primes: $a$ is integral and nonnegative.
On $U$, the positive coefficients of the trace $D_U$ are precisely the
floor primes, and they are at most one.  This holds for surviving primes
by \eqref{eq:D-actual-prime}, and for extracted primes by their small
$a+\mathbf D$ discrepancy.

For each floor prime $Q$ on $U$ put
$\rho_Q=\min\{1,L\mathbf D(Q)\}$.  Choose a small fixed rational
$c_1>0$, with $2c_1\le c$, and an effective correction $R_c$, supported
on these primes, such that
\begin{equation}\label{eq:D-rounding}
 c_1\rho_Q\le R_c(Q)\le2c_1\rho_Q,
 \qquad
 \operatorname{den}\bigl(J_m(Q)-R_c(Q)\bigr)\le C_1/\rho_Q.
\end{equation}
Here $\operatorname{den}$ is the reduced positive denominator.  Such a
choice is obtained by taking a rational grid of mesh at most
$c_1\rho_Q$ and denominator at most $C_1/\rho_Q$ in the interval
$[J_m(Q)-2c_1\rho_Q,J_m(Q)-c_1\rho_Q]$.  There is no correction outside
the floor: no such primes were extracted, and $J_m$ is integral there.

Minima of sums give the b-divisor inequality
\begin{equation}\label{eq:D-height-difference}
 \mathbf J_m-\mathbf J_{m_0}\ge(m-m_0)\mathbf D.
\end{equation}
Since $R_c\le2c_1 L D_U$, the divisor
$P_U=(J_m)_U-R_c$ dominates the trace of $\mathbf J_{m_0}$ and is big.
Run its MMP on $U$ to a nef big semiample divisor $P'$ on a
$\Q$-factorial Fano type model $U'$.  On a common higher model, denote
pullback from a displayed model by a star.  Then
\begin{equation}\label{eq:D-mmp-sandwich}
 \mathbf J_m\ge(P')^*\ge\mathbf J_{m_0},\qquad
 R_c^*\le \mathbf E:=\mathbf J_m-(P')^*\le(R_c')^*.
\end{equation}
Here is justification for both sides, including the comparison with the
possibly different ray $m_0$.  The divisor $J_m$ descends on $U$, and
MMP negativity gives $(P')^*\le P_U^*=\mathbf J_m-R_c^*$.
For the middle inequality, every divisible section of
$\mathbf J_{m_0}$ is a section of $P_U$, hence of $P'$ by the MMP.
The maximum-pole description of semiample divisors, applied on a common
determination, gives $(P')^*\ge\mathbf J_{m_0}$.  Finally b-nef
negativity gives $\mathbf J_m\le(J_m')^*$; since
$J_m'=P'+R_c'$, this is the last inequality in
\eqref{eq:D-mmp-sandwich}.  Also $\rho_Q\ge D_U(Q)$, so
\[
 \mathbf E\ge R_c^*\ge c_1(D_U)^*\ge c_1\mathbf D.
\]
The final comparison is again b-nef negativity for $\mathbf D$.

We apply Lemma~\ref{lem:B-phase-nonvanishing} to the polarization $P'$.
Choose fixed rational $\delta_2>0$ small enough for the denominator bound
below, and then $0<\delta_1\le\delta_2c_1$.  For
$0\le t\le\delta_2$ use the discrepancy family and nef part
\begin{align*}
 a_t&=a+(1-t/\delta_2)\delta_1\mathbf D+t\mathbf E,\\
 \mathbf M_t&=(1-t/\delta_2)\delta_1\mathbf D+t\mathbf J_m.
\end{align*}
Their total log canonical divisor is $\Q$-linearly equivalent to $tP'$.
For $t\ge\delta_2$ keep $a_t=a+\delta_2\mathbf E$ and add
$(t-\delta_2)(P')^*$ to the nef part.  This gives families on arbitrarily
long parameter intervals.  They are nondecreasing, klt, and have nef
$\Q$-Cartier b-parts.

We check effectivity and the denominator trace requirement of that lemma.
Every prime of $U'$ survives from $U$.  At a non-floor prime,
$a=1$, $\mathbf D=\mathbf E=0$, and $P'$ has integral coefficient.
At a surviving floor prime $Q$,
$\mathbf E(Q)=R_c(Q)\le2c_1\rho_Q$, while
\[
 \frac1{\operatorname{den}(P'(Q))}\ge\frac{\rho_Q}{C_1}.
\]
Take $2\delta_2c_1C_1\le1$.  Monotonicity gives
$0<a_t(Q)\le\delta_2R_c(Q)$ throughout the family, so
$a_t(Q)\le1/\operatorname{den}(P'(Q))$.  Equivalently, its effective
boundary dominates the denominator boundary $R(P')$ used in
Section~\ref{sec:B}.

It remains to check the signed phase hypothesis.  At a growing parameter,
small discrepancies mean $a(P)=0$ and $\mathbf E(P)\to0$.
Because $\mathbf E\ge c_1\mathbf D$, these valuations eventually belong
to the fixed-cutoff list already extracted on $U$.  On that list,
\eqref{eq:D-rounding} and \eqref{eq:D-mmp-sandwich} give
\[
 \mathbf E(P)\ge c_1\min\{1,L\mathbf D(P)\}.
\]
Thus $L\mathbf D(P)\to0$ and
$L\mathbf D(P)=O(\mathbf E(P))$.  Combining
\eqref{eq:D-sandwich-hypothesis} with
\eqref{eq:D-mmp-sandwich} yields
\[
 j_P\le(P')^*(P)\le j_P+O(\mathbf E(P)).
\]
The phase of $(P')^*(P)$ therefore has a nonnegative lift of size
$O(a_t(P))$, which is stronger than the signed phase hypothesis.
Lemma~\ref{lem:B-phase-nonvanishing} supplies a nonzero section in a
bounded positive degree of $P'$.  By
\eqref{eq:D-mmp-sandwich} it lies in $S(q\mathbf J_m)$, and
Lemma~\ref{lem:D-section-symbols} finishes the proof.
\end{proof}

\subsection{Bounded multiples, character counts, and tilted heights}

The sandwich proposition turns a phase interval into a regular symbol.
We use it in three forms: to supply symbols for bounded-width lattice
characters, to count symbols with one character and weight, and to
supply a symbol whose cost reflects a small tilted height.  The last form
will be used on newly optimized cones before any further rank refinement.

\begin{corollary}[Bounded-multiple character regularization]
\label{cor:D-regularization}
Suppose that the order condition \eqref{eq:C-alpha} on $V_{\mathbf u}$,
with error $\eta_i\to0$, holds whenever $a(P)=0$ and
$L\mathbf D(P)\to0$.  Fix $C_0>0$.  If $k\in M$ satisfies
\begin{equation}\label{eq:D-width}
 |\langle y,k\rangle|\le C_0L\langle y,\mathbf u\rangle
 \qquad(y\in\sigma),
\end{equation}
there are a bounded positive integer $q$, a nonzero invariant rational
eigenfunction $Y=fz^{qk}$, and a positive integer $b\le C_2L$ such that
$h^bY\in\breve R$.  The constants are uniform in $k$.  Both signs of
$k$ may be treated with the same bounded multiplier, and their costs may
be increased to a common cost $O(L)$.

The same bounds hold uniformly if $L$ is replaced by any $L'\ge L$ and
\eqref{eq:D-width} is imposed at scale $L'$.
\end{corollary}

\begin{proof}
Choose an integer $m_1\asymp L$ with $m_1>C_0L$, so that both
$\pm k+m_1\mathbf u$ lie in $\operatorname{int}C$.  Bigness of their
ray polarizations supplies rational-power functions $g_+,g_-$ on
$\mathcal K$ such that
\[
 (P,y)(z^{\pm k}h^{m_1})+P(g_\pm)\ge0
 \quad\text{for every }P\text{ and }y\in\Delta_P.
\]
Here a rational-power function means a positive integral root of an
element of $\mathcal K^*$ used only through its $\Q$-principal order;
it need not be an element of $\mathcal K$.  Put
\[
 \ell=-\divisor(g_+),\qquad
 G=2m_1D+\divisor(g_+g_-).
\]
Adding the two Gauss inequalities and taking the height minimum shows
that $G\ge0$ on $V_{\mathbf u}$, and
$G\sim_{\Q}2m_1D$.  At a non-floor prime $Q$ of this model, choose the
integral minimizing $y$ in \eqref{eq:D-off-floor}.  Since its height is
zero, the two inequalities give
\[
 0\le\langle y,k\rangle-\ell(Q)\le G(Q),
 \qquad \langle y,k\rangle\in\Z.
\]
Lemma~\ref{lem:C-scaled-phase} therefore gives a fixed positive integer
$q_0$ and a fixed $A>0$ such that, on all tiny sequences in the
statement, one can choose integers $j_P$ with
\[
 |q_0\ell(P)-j_P|\le A L\mathbf D(P).
\]

Apply Proposition~\ref{prop:D-sandwich} to $\xi=z^{q_0k}$.
For any sufficiently large integer $m_0=O(L)$ and $m=2m_0$, the first
Gauss inequality gives
\[
 \mathbf J_{m_0}\ge q_0\ell+(m_0-q_0m_1)\mathbf D.
\]
At a height-minimizing angular vector the second inequality gives
\[
 \mathbf J_m\le q_0\ell+(m-q_0m_1)\mathbf D+q_0\mathbf G
 \le q_0\ell+(m+q_0m_1)\mathbf D+q_0\mathbf G.
\]
The order condition \eqref{eq:C-alpha}, applied to
$G/(2m_1)\sim_{\Q}D$, bounds
$\mathbf G(P)\le2m_1(1+\eta_i)\mathbf D(P)$ on the tiny sequences.
Choose $m_0-q_0m_1\ge AL$.  These estimates give exactly
\eqref{eq:D-sandwich-hypothesis}.  A bounded further multiplier $q_1$
then supplies $g z^{q_1q_0k}h^{q_1m}\in\breve R$.  This is the desired
regularization.

Apply the argument to $-k$ as well, raise the two resulting functions
to bounded powers to align their multipliers, and multiply by powers of
$h$ to align their costs.  Uniformity in the weight follows by applying
the sequence argument to any proposed sequence of exceptions.  For an
enlarged scale $L'\ge L$, its tiny sequences satisfy
$L\mathbf D(P)\to0$ as well.  All the preceding hypotheses therefore
hold with $L'$ in place of $L$.  A sequence of exceptional pairs
$(L',k)$ would contradict the same proof; this proves the stated
uniformity, including for growing $L'/L$.
\end{proof}

\begin{corollary}[Count for a single character and weight]
\label{cor:D-character-count}
Assume the hypotheses of Lemma~\ref{lem:C-saturation-count} on
$V_{\mathbf u}$, and use its notation $p$ and $v_i$.  For a fixed
$C_0>0$, let $W_{t,k}$ be the space of regular symbols in $\breve R$
of finite-group character $\chi^t$ and torus weight $t\mathbf u+k$.
For every fixed positive integer $n$, eventually
\begin{equation}\label{eq:D-character-bound}
 \dim_k W_{t,k}\le C_3 v_i n^p
\end{equation}
simultaneously for all integers $0\le t\le C_0nL$ and all weights
satisfying
\[
 |\langle y,k\rangle|\le C_0nL\langle y,\mathbf u\rangle
 \qquad(y\in\sigma).
\]
The constant $C_3$ is independent of $n,t,k$.
\end{corollary}

\begin{proof}
If $W_{t,k}\ne0$, divide any nonzero member by $h^t$.  This is a
finite-group invariant rational eigenfunction, so $k\in M$.
Let $\mathbf J^{\mathrm{test}}$ denote the integral basepoint-free
test b-divisor called $J$ in Lemma~\ref{lem:C-saturation-count}.
That lemma has
$0\le\mathbf J^{\mathrm{test}}\le C L_0\mathbf D$ with
$L_0=O(L)$.  At a primitive $P$ with $a(P)=0$ and
$L\mathbf D(P)\to0$, its nonnegative integral value is eventually
zero.  Such $P$ belong to that lemma's test set $\mathcal L$ and hence
satisfy \eqref{eq:C-alpha}.  We may therefore apply
Corollary~\ref{cor:D-regularization} at the enlarged scale $nL$.

Choose an opposite invariant eigenfunction $Y$ of weight $-qk$, with
$q$ bounded and $h^bY$ regular for $b\le C_4 nL$.  For
$F\in W_{t,k}$ the neutral function
\[
 \Phi(F)=\frac{F^qY}{h^{qt}}
\]
satisfies $h^{qt+b}\Phi(F)=F^q h^bY\in\breve R$.  By
\eqref{eq:D-neutral-functions}, it belongs to
$S((qt+b)\mathbf D)\subseteq S(C_5nL\mathbf D)$.
Although $\Phi$ need not be linear, it is a homogeneous polynomial map
of degree $q$, nonzero away from the origin, and induces an injective
morphism
\[
 \mathbb P(W_{t,k})\longrightarrow
 \mathbb P\bigl(S(C_5nL\mathbf D)\bigr).
\]
Indeed proportional images imply that $(F_1/F_2)^q\in k^*$;
algebraic closedness of $k$ then implies $F_1/F_2\in k^*$.
An injective morphism of projective varieties preserves the dimension
of its source in its image, so
$\dim W_{t,k}\le\dim S(C_5nL\mathbf D)$.
The upper count in Lemma~\ref{lem:C-saturation-count}, used at an
integer at most $(C_5+1)n$, proves
\eqref{eq:D-character-bound}; $p$ is bounded by the ambient dimension.
The uniform enlarged-scale assertion in
Corollary~\ref{cor:D-regularization} makes $q,C_4,C_5$ independent of
$n,t,k$.  The eventual assertion retains the quantifier that $n$ is
fixed before passing sufficiently far along the sequence.
\end{proof}

\begin{corollary}[Characters for a tilted height]
\label{cor:D-tilted-characters}
Suppose there are boundedly many invariant rational eigenfunctions
$Y_j^\pm$ of weights $\pm k_j$, with
\[
 h^{m_j}Y_j^\pm\in\breve R,\qquad
 0<m_j\le A L_j,\qquad 1\le L_j\le L,
\]
where $A$ is fixed.  Fix bounded real numbers $u_j$ and $\delta>0$.
Assume that, at every nonzero pure weight and at all Gauss extensions
over primitive $P$ with $a(P)=0$, $L\mathbf D(P)\to0$, the evaluations
satisfy
\begin{equation}\label{eq:D-tilt-hypothesis}
 y_j=v(Y_j^+)=-v(Y_j^-),\qquad
 \sum_j\frac{u_jy_j}{L_j}\le\delta v(h).
\end{equation}
It suffices to check these inequalities on rational vectors.  For
fixed $M_0$ sufficiently large compared with $1/\delta$, choose integers
\[
 r_j=M_0u_jL/L_j+O(1),
\]
with uniformly bounded rounding errors.  Then there is a regular
nonzero symbol of finite-group character $\chi^{qm}$ and torus weight
\begin{equation}\label{eq:D-tilt-weight}
 q\left(m\mathbf u-\sum_jr_jk_j\right),
 \qquad 0<m\le C_6M_0\delta L,
\end{equation}
for a bounded positive integer $q$.  The constant $C_6$ depends only on
the fixed bounds for the number of functions, their costs, the $u_j$,
and rounding errors; it is independent of $\delta$ and $M_0$.
The bound for $q$ may depend on the fixed choices $\delta,M_0$.
No order condition \eqref{eq:C-alpha} is assumed here.
\end{corollary}

\begin{proof}
Set $\xi=\prod_j(Y_j^+)^{-r_j}$, which has weight
$-\sum_jr_jk_j$.  Regularity and opposite evaluations imply
$|y_j|\le m_jv(h)$ at every test under consideration.  Hence, with a
fixed constant $C_7$ controlling the rounding errors,
\begin{equation}\label{eq:D-tilt-lower}
 \sum_jr_jy_j
 \le(M_0\delta+C_7)Lv(h).
\end{equation}
Choose $M_0\delta\ge2C_7+2$ and an integer
$m_0\asymp M_0\delta L$ strictly larger than
$(M_0\delta+C_7)L$, and set $m=2m_0$.  The constants in these choices
are independent of $\delta,M_0$.  At pure weights
\eqref{eq:D-tilt-lower} implies strict positivity of the weights of
$\xi h^{m_0}$ and $\xi h^m$ on $\sigma\setminus\{0\}$; both rays
are therefore interior.  At the angular tests it gives
$\mathbf J_{m_0}(P)\ge0$.

At a height-minimizing angular vector, the elementary estimate
\[
 \left|\sum_jr_jy_j\right|
 \le\sum_j|r_j|m_j\mathbf D(P)
 \le C_8(M_0+1)L\mathbf D(P)
\]
also holds.  Consequently
$0\le\mathbf J_{m_0}(P)\le\mathbf J_m(P)
\le C_9(M_0,\delta)L\mathbf D(P)$ on the tiny sequences.
Proposition~\ref{prop:D-sandwich} applies with $j_P=0$ and produces
$g\xi^qh^{qm}\in\breve R$ for bounded $q$.  Its character, weight,
and cost are exactly those in \eqref{eq:D-tilt-weight}.
\end{proof}


\section{Optimization on common log canonical slices}\label{sec:E}

Throughout this section, $x\in S$ is a fixed normal complex affine germ of
positive dimension $n$, $S$ is $\Q$-Gorenstein and klt, and
$R=\mathcal O_{S,x}$, $\mathfrak m=\mathfrak m_x$, and $A=A_S$.
A valuation \emph{over $x$} is nonnegative on $R$; its centre may be a
proper generization of $x$. A valuation \emph{centred at $x$} has
$v(\mathfrak m)>0$. Unless otherwise stated, valuations are nontrivial
quasi-monomial real valuations. Discrepancy is extended linearly in the
weights on a log-smooth presentation. For every valuation $v$ over $x$,
put $v(0)=+\infty$ and, for $t\ge0$, set
\[
 \mathfrak a_t(v)=\{f\in R:v(f)\ge t\},\qquad
 \mathfrak a_{>t}(v)=\{f\in R:v(f)>t\}.
\]
The associated graded algebra and initial form are
\[
 \operatorname{gr}_v R=\bigoplus_{t\ge0}
       \mathfrak a_t(v)/\mathfrak a_{>t}(v),\qquad
 \operatorname{in}_v(f)=[f]_{v(f)}\quad(f\ne0).
\]
Degrees outside the value semigroup have zero summand. These conventions
also apply when the centre is a proper generization of $x$.
For a valuation centred at $x$, set
\[
 \vol(v)=\lim_{t\longrightarrow\infty}
       \frac{n!\,\operatorname{length}(R/\mathfrak a_t(v))}{t^n}.
\]
The limit exists by Cutkosky's theorem for graded families of
$\mathfrak m$-primary ideals \cite[Theorem~1.1 (=5.5)]{CutkoskyVolumes}.
Its hypothesis on the nilradical of the completion holds because $R$
is excellent and normal, hence has reduced completion.  The theorem
applies first to integral thresholds; squeezing between adjacent integral
thresholds gives the displayed real-parameter limit.
For a klt pair $(S,\Delta)$ and a centred valuation of finite log discrepancy,
the normalized volume is
\[
 \widehat{\vol}_{S,\Delta}(v)=A_{S,\Delta}(v)^n\vol(v).
\]
Its $n$th root is $A_{S,\Delta}(v)\vol(v)^{1/n}$, with the same minimizers.
Here and throughout the volume formulas, $n=\dim S$.

Fix a nonzero $h\in R$ such that $(S,H)$ is lc, where
$H=\divisor(h)$, and finitely many nonzero $F_j\in R$ and
$m_j\in\Z_{>0}$. The \emph{common lc slice} is
\begin{equation}\label{eq:E-slice}
 \mathcal Q=\{v:A(v)=v(h)=1,\quad v(F_j)\ge m_j\text{ for every }j\}.
\end{equation}
Assume that it contains a centred valuation. Its cone of rays is defined
by $A(v)=v(h)>0$ and $v(F_j)\ge m_jv(h)$.
An allowed height is
\begin{equation}\label{eq:E-height}
 b(v)=b_0v(h)-\sum_{\ell}b_{\ell}v(G_{\ell}),
 \qquad b_{\ell}\ge0,\quad G_{\ell}\in R\setminus\{0\},
 \qquad b|_{\mathcal Q}>0.
\end{equation}
On centred valuations for which $b(v)>0$, the degree-one homogeneous
height $b$ similarly gives the scale-invariant quantity $b(v)^n\vol(v)$ and its
$n$th root $b(v)\vol(v)^{1/n}$, with the same minimizers.
All constants in this section may depend on this individual germ and
its finite data. The constants in the strict volume gap below depend
only on dimension and the stated separation.

Our goal is to minimize this functional on the common lc rays and,
for rational heights, realize its minimizer as an ordinary
normalized-volume minimizer of an effective rational klt pair.
The realization must be exact: finite generation will be applied to
that pair and its actual minimizer. We first control retraction to
the lc locus and volume under interpolation; these estimates will
allow one rational perturbation to exclude every competing valuation
outside the common slice.

\subsection{Retractions, jets, and the topology of the slice}

Let $P\to S$ be a smooth model, $F$ a reduced snc divisor on $P$, and
$v$ a valuation with centre on $P$. Write $r_F(v)$ for its monomial
retraction using the components of $F$ through the centre, with their
$v$-weights. It is the trivial valuation if there are no such components.
Retraction decreases the value of every function regular on the chart
and preserves the values of its boundary monomials; see
\cite[Lemma~4.7, Corollary~4.8, and Proposition~5.1]{JonssonMustata}.
The following quantitative strengthening is needed here.

\begin{lemma}[Logarithmic jet estimate]\label{lem:E-log-jet}
Fix finitely many regular affine functions on $S$. On a fixed resolution
$P$, with $F$ reduced snc, there is a constant $C$ such that, for every
polynomial $g\ne0$ of degree at most $D$ in those functions and every
quasi-monomial $v$ over $x$,
\begin{equation}\label{eq:E-jet}
 r_F(v)(g)\le v(g)\le r_F(v)(g)+C(D+1)A_{P,F}(v).
\end{equation}
The same constant works when the centre is not closed and when some
weights disappear on passing to a face of the monomial cone.
\end{lemma}

\begin{proof}
Cover the relevant part of $P$ by finitely many smooth etale coordinate
charts. At the centre of $v$, write the coordinates of the components
of $F$ through that centre as $z_1,\ldots,z_s$ and the remaining
coordinates as $y_1,\ldots,y_{n-s}$. Thus $v(z_i)>0$.
Choose a closed specialization of the centre in this chart, avoiding
the other components of $F$, and translate the $y$-coordinates to
vanish there. Both $v$ and its retraction are nonnegative on this
closed-point local ring: the original centre generalizes this point.
A unit in that ring consequently has value zero for both valuations.
Put $a_i=v(z_i)$ and expand $g$ in the completed local ring.
The least $z$-weight is
\[
 d_0=\min\{\langle a,\alpha\rangle:
               [z^\alpha y^\beta]g\ne0\text{ for some }\beta\}.
\]
We first prove a degree bound independent of $a$.

The function field of the chart is a finite extension of
$\C(z,y)$. In a fixed basis, multiplication by each of the finitely
many affine functions has a matrix with fixed rational-function
entries. A polynomial of degree at most $D$ in these matrices has a
common denominator of degree $O(D+1)$ and numerator entries of the
same degree. Its characteristic equation, after clearing denominators,
gives
\begin{equation}\label{eq:E-algebraic-equation}
 \sum_{j=0}^{q}p_j(z,y)g^j=0,\qquad
 \deg p_j\le C_0(D+1),
\end{equation}
where $q$ is fixed and the coefficients are not all zero.
Translations of coordinates do not change this degree bound.
Give \emph{all} coordinates strictly positive rationally independent
weights. Let $\gamma$ be the leading exponent of $g$, and let $\delta_j$
be the leading exponent of a nonzero $p_j$. At least two distinct
powers, say $j$ and $k$, have the least weight in
\eqref{eq:E-algebraic-equation}. Independence of the weights gives
\[
 \delta_j+j\gamma=\delta_k+k\gamma.
\]
Consequently $|\gamma|\le C_1(D+1)$, after enlarging the constant.
This conclusion holds for every such choice of positive weights.

Now give the $z$-coordinates their fixed positive weights and give the
$y$-coordinates weight zero. There are finitely many $z$-exponents
of weight $d_0$, and only finitely many in a sufficiently small fixed
neighbourhood of this weight. Every vertex of the convex hull of the
tied $z$-exponents is exposed by a small generic perturbation of the
$z$-weights. Give the transverse coordinates sufficiently small positive
generic weights afterwards. The leading exponent projects to that
vertex: the positive gap to the other $z$-weights prevents an exponent
outside the least face from entering. The preceding bound therefore
bounds every vertex, and convexity bounds \emph{every} tied
$z$-exponent by $C_1(D+1)$. The same perturbation supplies at least one
pair $(\alpha,\beta)$ on the least face with
$|\alpha|+|\beta|\le C_1(D+1)$.

For this choice of $\alpha$, apply the product, over the $z$-coordinates,
of the operators
\[
 \prod_{\substack{0\le l\le C_1(D+1)\\ l\ne\alpha_i}}
             (z_i\partial_{z_i}-l).
\]
It kills every other tied $z$-exponent and multiplies the chosen
coefficient by a nonzero constant. Apply a transverse derivative
$\partial_y^\beta$ for a bounded exponent occurring in that coefficient.
Its new constant term is nonzero. These operators have total order
$O(D+1)$, and do not change any $z$-exponent. Their result $g'$ satisfies
\begin{equation}\label{eq:E-isolated-jet}
 g'=u z^\alpha\pmod{I_{>d_0}},\qquad u\text{ a local unit},
\end{equation}
where $I_{>d_0}$ is the monomial ideal generated by the $z$-monomials
of weight greater than $d_0$. This ideal is finitely generated.
The assertion descends from completion by faithful flatness. More
explicitly, faithful flatness first gives $g'=u z^\alpha+q$ with
$u$ regular and $q\in I_{>d_0}$. In the completion, $u$ differs from
the formal unit coefficient by an element of $(I_{>d_0}:z^\alpha)$.
This colon ideal is proper, because $z^\alpha$ has weight $d_0$,
and hence lies in the local maximal ideal. Thus $u$ is a unit.
Each monomial generator of $I_{>d_0}$ has value greater than $d_0$
for both $v$ and $r_F(v)$, and its regular coefficients have nonnegative
value. Thus
\[
 v(g')=r_F(v)(g')=d_0=r_F(v)(g).
\]
This also explains why a nonclosed centre causes no problem.

For completeness, a regular logarithmic vector field $L$ satisfies
\begin{equation}\label{eq:E-log-derivative}
 v(Lf)\ge v(f)-A_{P,F}(v).
\end{equation}
For a prime-order valuation, take a local logarithmic volume form
$\eta$ for $(P,F)$ and use
\[
 (Lf)\eta=f\,d\log(f)\wedge\iota_L\eta.
\]
At the prime on a higher smooth model, logarithmic coordinate forms
and $d\log(f)$ have nonnegative logarithmic order. Comparing the two
sides gives \eqref{eq:E-log-derivative}. Homogeneity gives it for rational
monomial weights, and rational approximation gives it for arbitrary
quasi-monomial weights. In this approximation the signs on a fixed
chart ring can be retained: the values of its finitely many generators
are rational piecewise linear functions on the monomial weight cone,
so one approximates in their common sign cells. Discrepancy and the
finitely many function values in question are continuous there.
The Euler operators and transverse derivatives used above are
logarithmic vector fields, and subtracting a scalar does not worsen
the estimate. Iterating \eqref{eq:E-log-derivative} and using
\eqref{eq:E-isolated-jet} proves the upper bound in \eqref{eq:E-jet}.
The lower bound is the retraction inequality.
\end{proof}

Suppose more generally that $\Psi$ is an effective lc rational sum of
Cartier divisors. Resolve its support and any further finite list of
ideals and divisors, and write
\[
 K_P+B_P=\pi^*(K_S+\Psi),\qquad F=B_P^{=1},\qquad
 e=A_{S,\Psi}(v),\qquad r=r_F(v).
\]
If $E$ is the reduced union of the support of $B_P$, the exceptional
locus, and the additional resolved divisors, then
\[
 e=A_{P,E}(v)+\sum_{D\subset E}(1-\operatorname{coeff}_D B_P)v(D).
\]
All summands are nonnegative. The positive numbers
$1-\operatorname{coeff}_D B_P$ for $D\not\subset F$ have a positive
minimum on this fixed resolution. Therefore
\begin{equation}\label{eq:E-retraction-errors}
 A_{P,F}(v)\le C e,\qquad
 |A_S(v)-A_S(r)|\le C e,
\end{equation}
and the evaluations of each resolved ideal or divisor differ at $v$
and $r$ by at most $Ce$. These statements follow also when $B_P$ has
negative coefficients: the displayed coefficients $1-\operatorname{coeff}_D B_P$
are still positive away from $F$. In particular, if $e=0$, then
Lemma~\ref{lem:E-log-jet} gives equality on every polynomial in the
fixed affine generators. Equality extends to the local ring and its
fraction field by localization. Hence $v=r$.

\begin{lemma}[Compactness and uniform centring]\label{lem:E-compactness}
The slice $\mathcal Q$ is a compact finite polyhedral complex of
monomial valuations on any log resolution adapted to $h$ and the
$F_j$. Evaluations of fixed functions are continuous on it.
For centred valuations with $A(v)=1$, there are constants $c,C>0$
depending on the germ such that
\begin{equation}\label{eq:E-centering}
 \vol(v)\ge c,\qquad
 \vol(v)\ge \frac{c}{v(\mathfrak m)}.
\end{equation}
Volume is continuous on the centred part of $\mathcal Q$.
\end{lemma}

\begin{proof}
Apply \eqref{eq:E-retraction-errors} to $\Psi=H$.
Every zero-discrepancy valuation is monomial on the coefficient-one
part of the fixed resolution. Conversely those monomial valuations
have zero discrepancy for $(S,H)$. Retain the strata whose closed
images contain $x$. This condition is preserved when a stratum is
replaced by a larger stratum, so the retained cones include their
faces. On a cone with primitive generators $E_i$ the equation
$A(v)=\sum_i a_iA(E_i)=1$ cuts out a compact simplex, because
$A(E_i)>0$. The additional inequalities are linear on an adapted
resolution. There are finitely many cones. Function evaluations are
monomial minima on them and agree on faces by localization to the
larger stratum, proving compactness and continuity.

We use Li's Izumi estimate \cite[Theorem~3.1]{LiIzumi}:
\begin{equation}\label{eq:E-izumi}
 v(f)\le C_I A(v)\ord_{\mathfrak m}(f)
 \quad(f\in R\setminus\{0\},\ v\text{ centred}).
\end{equation}
Thus, when $A(v)=1$, the valuation ideal $\mathfrak a_t(v)$ is contained
in $\mathfrak m^{\lceil t/C_I\rceil}$. Hilbert--Samuel growth gives the
first bound in \eqref{eq:E-centering}.
Choose fixed generators $z_1,\ldots,z_q$ of $\mathfrak m$, and let
$z=z_i$ have $v(z)=\delta=v(\mathfrak m)$. Since $R$ is a domain,
\[
 (\mathfrak a_t(v):z^j)=\mathfrak a_{t-j\delta}(v).
\]
Filtering $R/\mathfrak a_t(v)$ by the images of $z^jR$ gives
\[
 \operatorname{length}(R/\mathfrak a_t(v))
 =\sum_{0\le j<t/\delta}
   \operatorname{length}\bigl(R/(\mathfrak a_{t-j\delta}(v)+(z))\bigr).
\]
For $j\le t/(2\delta)$, Izumi bounds this summand below by
$\operatorname{length}(R/(\mathfrak m^{\lceil t/(2C_I)\rceil}+(z)))$.
The latter grows as a positive constant times $t^{n-1}$; this includes
$n=1$, when it is a positive constant. Since $z$ belongs to a fixed
finite list, the lower constant can be chosen uniformly. Divide by
$t^n/n!$ to obtain the second bound in \eqref{eq:E-centering}.

To prove continuity, consider valuations converging on one monomial
cone to a centred limit. Drop the components whose weights tend to
zero, and call the resulting retractions $r_i$. The retained positive
weights converge multiplicatively. On regular functions, the $r_i$
are therefore multiplicatively $1+o(1)$ equivalent to the limit.
Their maximal-ideal values, and those of the original valuations,
are bounded below. The reduced-log discrepancy measuring the dropped
weights tends to zero. Lemma~\ref{lem:E-log-jet} upgrades this fact to
a multiplicative comparison on \emph{all} regular functions as follows.

Write $q=\ord_{\mathfrak m}(f)\ge1$. With $c_0>0$ a common lower bound
for maximal-ideal values and $C_0q$ the common Izumi upper bound for
the values of $f$, fix $L$ with $Lc_0>2C_0$. Represent $f$ modulo
$\mathfrak m^{N}$, $N=\lceil Lq\rceil$, by a polynomial $p$ of degree
at most $N$ in fixed generators of $\mathfrak m$. Such a representative
exists also for localized regular functions by expanding their unit
denominators modulo $\mathfrak m^N$. The error has value greater than
both values of $f$, so $v_i(p)=v_i(f)$ and $r_i(p)=r_i(f)$.
The jet estimate bounds their difference by $o(1)q$, uniformly in $f$.
Since $r_i(f)\ge c_0q$, it follows that
$r_i\le v_i\le(1+o(1))r_i$ on all regular functions. Units have value
zero. Inclusion of valuation ideals now proves volume continuity.
Every convergent sequence has a subsequence in one of the finitely
many cones, which proves the assertion on the whole centred part.
\end{proof}

We record explicitly the uniform version of the truncation argument.
\begin{corollary}[Simultaneous approximation]\label{cor:E-all-functions}
In the situation of \eqref{eq:E-retraction-errors}, suppose
$A_S(v)$ and $A_S(r)$ are bounded above and
$v(\mathfrak m),r(\mathfrak m)\ge c_0>0$. Then a single constant $C$
works for every such $v$ and every $f\in R$:
\begin{equation}\label{eq:E-all-functions}
 r(f)\le v(f)\le(1+Ce)r(f).
\end{equation}
In particular,
$\vol(v)^{1/n}\ge(1+Ce)^{-1}\vol(r)^{1/n}$.
\end{corollary}
\begin{proof}
Use $N=\lceil L\ord_{\mathfrak m}(f)\rceil$ with the single constant
$L$ chosen in the preceding proof. Equations~\eqref{eq:E-jet} and
\eqref{eq:E-retraction-errors} give
$v(f)-r(f)\le C e\ord_{\mathfrak m}(f)\le Ce\,r(f)/c_0$.
Enlarge $C$. The ideal inclusions give the volume inequality.
\end{proof}

\subsection{Convex domination and strict volume comparison}

\begin{lemma}[Convex domination]\label{lem:E-convexity}
For $v_1,\ldots,v_k\in\mathcal Q$ and $p_s\ge0$ with
$\sum_s p_s=1$, there is $u\in\mathcal Q$ such that
\begin{equation}\label{eq:E-convexity}
 u(f)\ge\sum_s p_s v_s(f)\qquad(f\in R).
\end{equation}
If a positive-weight term is centred, then $u$ is centred. On the cone
of common lc rays, the same statement holds with
$A(u)=\sum_s p_sA(v_s)$, without requiring $A(v_s)=1$.
For an allowed height,
$b(u)\le\sum_s p_sb(v_s)$.
\end{lemma}

\begin{proof}
Fix finitely many functions $g_i$ whose averaged target values
$s_i=\sum_s p_sv_s(g_i)$ are positive, and choose positive rational
$s'_i\le s_i$. Include $h$ with the \emph{exact} target $1$ and each
$F_j$ with target $m_j$. Choose $N$ so that every $N/s'_i$ is integral,
and put $\mathfrak a=\sum_i(g_i^{N/s'_i})$. We claim that
$(S,\mathfrak a^{1/N})$ is not klt at $x$.
The multiplier-ideal summation formula on normal $\Q$-Gorenstein
complex varieties gives
\begin{equation}\label{eq:E-multiplier-sum}
 \mathcal J(\mathfrak a^{1/N})
 =\sum_{\substack{t_i\ge0\\\sum_i s'_it_i=1}}
       \mathcal J\Bigl(\prod_i(g_i)^{t_i}\Bigr).
\end{equation}
This is the mixed-ideal formula of
\cite[Theorem~3.2]{Takagi} and
\cite[Corollary~2 and Theorem~3.6]{JowMiller}; there is no smoothness
or Jacobian correction in this particular formula.
If the left side were the unit ideal at $x$, some summand would be the
unit ideal in the local ring. The resulting mixed pair would be klt,
so for every positive-weight $v_s$,
$\sum_i t_i v_s(g_i)<A(v_s)=1$. Averaging contradicts
\[
 \sum_i t_i\sum_s p_s v_s(g_i)=\sum_i t_is_i
       \ge\sum_i t_is'_i=1.
\]
The failure of klt is witnessed by a prime $E$ whose centre contains
$x$, and
$A(E)\le\min_i\{\ord_E(g_i)/s'_i\}$.
Thus $u=\ord_E/A(E)$ meets all the lower targets. In particular
$u(h)\ge1$, whereas lc of $(S,H)$ gives $u(h)\le A(u)=1$.
It is already in $\mathcal Q$.

Here is the geometric mechanism behind \eqref{eq:E-multiplier-sum}.
Resolve the summand ideals and their sum, with fixed divisors $D_i,D$,
and subdivide $\{p_i\ge0:\sum p_i=1/N\}$ by the rounding walls of
$\mathcal O_P(\lceil K_{P/S}-\sum_i p_iD_i\rceil)$.
The cellular complex indexed by chains of nonempty cells, with the
usual alternating face maps and inclusions on removing the first
cell, resolves
$\mathcal O_P(\lceil K_{P/S}-D/N\rceil)$.
In completed snc coordinates its exactness is coefficientwise:
the cells admitting a given Laurent monomial form a convex relatively
open subset of the exponent simplex, and their barycentric order
complex is contractible when nonempty. This subset is nonempty exactly
for monomials in the target. Indeed, locally one summand divided by
the invertible sum ideal is a unit, so its fixed divisor there equals
$D$. Local vanishing for each mixed multiplier-ideal term preserves
the surjective augmentation under pushforward. This gives the stated
sum formula in precisely the singular setting used above.

Let the rational lower targets increase to their averaged values.
Compactness from Lemma~\ref{lem:E-compactness} gives a point meeting
the finite list exactly as lower bounds. The closed sets imposing
these lower bounds, indexed by all regular functions, have the finite
intersection property. Compactness supplies a single $u$ satisfying
\eqref{eq:E-convexity} for every function. Positivity on generators of
$\mathfrak m$ proves the centring assertion. For unnormalized inputs,
put $a=\sum_s p_sA(v_s)$, apply the normalized assertion with weights
$p_sA(v_s)/a$ to $v_s/A(v_s)$, and multiply the result by $a$.
The formula for $b$ and the equality for $u(h)$ prove its height bound.
\end{proof}

\begin{corollary}[Opposite ratios]\label{cor:E-opposite-ratios}
Suppose rational functions $q,q'$ have regular numerators and
power-of-$h$ denominators, and $v(q')=-v(q)$ on $\mathcal Q$.
The interpolant in Lemma~\ref{lem:E-convexity} has
$u(q)=\sum_s p_sv_s(q)$ and $u(q')=\sum_s p_sv_s(q')$.
The assertion holds simultaneously for any number of such pairs.
\end{corollary}
\begin{proof}
Apply domination to both numerators and use the exact value of $h$.
The two lower inequalities have opposite sides and so must be equalities.
\end{proof}

\begin{lemma}[Volume interpolation and a strict gap]
\label{lem:E-volume-interpolation}
Let $w$ be centred, let $v$ be over $x$, and suppose
$u\ge tw+(1-t)v$ on $R$ for $0<t<1$. For $f\ne0$ with $w(f)>0$,
\begin{equation}\label{eq:E-volume-interpolation}
 \vol(u)\le t^{-n}
   \frac{tw(f)}{tw(f)+(1-t)v(f)}\vol(w).
\end{equation}
If $v$ is also centred, $t=1/2$, and
$v(f)\ge(1+\gamma)w(f)>0$ for $0<\gamma\le1$, then
\begin{equation}\label{eq:E-volume-gap}
 \vol(w)+\vol(v)-2\vol(u)\ge c(n,\gamma)\vol(w),
\end{equation}
where $c(n,\gamma)>0$ depends only on $n,\gamma$.
Consequently valuations whose volumes are within relative $o(1)$ of
the minimum on the same $A=1$ common slice are multiplicatively
$1+o(1)$ equal on \emph{all} regular functions. The conclusion is
uniform when the dimension is bounded. A sufficiently small fixed
relative error similarly gives any prescribed fixed multiplicative
closeness to one.
\end{lemma}

\begin{proof}
For a centred filtration, write its Laplace sum as the sum over its
jumps $a$ of
$\dim_\C(\mathfrak a_a/\mathfrak a_{>a})e^{-a/N}$.
Integration of the cumulative dimension function gives
\begin{equation}\label{eq:E-laplace}
 \lim_{N\longrightarrow\infty}N^{-n}
  \sum_a\dim_\C(\mathfrak a_a/\mathfrak a_{>a})e^{-a/N}
 =\vol(v).
\end{equation}
There are finitely many jumps below each bound, and polynomial growth
controls the tail.

Choose a homogeneous basis of
$\operatorname{gr}_wR/(\operatorname{in}_wf)$, lift it homogeneously,
and multiply the lifts by all nonnegative powers of
$\operatorname{in}_wf$. This is a basis of $\operatorname{gr}_wR$:
its independence follows from the domain property, and successive
division lowers the degree by $w(f)>0$, proving spanning.
Lift this basis to products $a_if^j$ in $R$. They span modulo arbitrary
high $w$-ideals, hence modulo high $u$-ideals since $u\ge tw$.
The $u$-value of $a_if^j$ is at least
$tw(a_i)+j(tw(f)+(1-t)v(f))$. Consequently its Laplace sum is at most
the $w$-Laplace sum at $N/t$, multiplied by
\[
 \frac{1-e^{-tw(f)/N}}
      {1-e^{-(tw(f)+(1-t)v(f))/N}}.
\]
Use \eqref{eq:E-laplace} to obtain
\eqref{eq:E-volume-interpolation}.

For the strict gap, fix $T>3$ and work in
$V_N=R/\mathfrak m^{\lceil CTN\rceil}$, where $C$ is large enough
for this fixed triple of centred valuations. The indicated power is
contained in all their threshold-$TN$ ideals. Cap both the $w$ and
$v$ flags at $TN$ and choose a basis simultaneously adapted to these
two flags; denote its capped values by $W_i,V_i$.
Lifts have $u$-value at least $(W_i+V_i)/2$. The identity
\[
 e^{-W_i/N}+e^{-V_i/N}-2e^{-(W_i+V_i)/(2N)}
     =\bigl(e^{-W_i/(2N)}-e^{-V_i/(2N)}\bigr)^2
\]
therefore bounds twice the capped $u$-Laplace sum above by the sum of
the other two, less the sum of these squares.

Choose a vector-space complement to
$\mathfrak a_{\gamma N/4}(w)$ in $R$ and multiply it by
$f^{\lfloor N/w(f)\rfloor}$. Every nonzero resulting element has
$w$-value less than $(1+\gamma/4)N$ and, for large $N$, has $v$-value
at least $(1+2\gamma/3)N$. Multiplication is injective, and this
subspace remains injective in $V_N$ by its upper $w$-value bound.
Its dimension is asymptotic to
\[
 \frac{\vol(w)}{n!}(\gamma N/4)^n.
\]
Its projection from the high-$v$ subspace modulo the high-$w$
subspace has this rank. Hence at least this many entries of a
simultaneously adapted basis satisfy
$W_i<(1+\gamma/4)N$ and $V_i\ge(1+2\gamma/3)N$.
Each contributes at least
\[
 \left(e^{-(1+\gamma/4)/2}-e^{-(1+2\gamma/3)/2}\right)^2
\]
to the preceding sum of squares. Divide by $N^n$, first let
$N\to\infty$, and then let $T\to\infty$.
The capped tail errors are bounded by polynomial functions of $T$
times $e^{-T}$ and disappear. This proves
\eqref{eq:E-volume-gap}, for example with any smaller positive
constant than
\[
 \frac{(\gamma/4)^n}{n!}
 \left(e^{-(1+\gamma/4)/2}-e^{-(1+2\gamma/3)/2}\right)^2.
\]

Finally interpolate two near-minimizers by
Lemma~\ref{lem:E-convexity}. Their interpolant has at least the
minimum volume. If their values on some regular function differed
by a fixed factor greater than one, apply \eqref{eq:E-volume-gap},
interchanging the valuations if necessary. It gives a fixed relative
loss, contrary to near-minimality. This proves the simultaneous
all-functions assertion and its fixed-error version.
\end{proof}

\subsection{Exact realization by a rational klt pair}

We have established compactness, a uniform comparison with monomial
retraction, and strict volume improvement when two valuations differ.
The remaining task is to encode the finite order constraints in an lc
boundary and then choose one rational perturbation for which every
ordinary minimizer lies exactly on its zero-discrepancy locus.  The
parameter must work for all competing valuations at once.

\begin{lemma}[A boundary encoding the common slice]\label{lem:E-slice-boundary}
Choose positive rational numbers $\theta_j$ with
$\sum_j\theta_j<1$. For sufficiently general scalars $a_j$, put
$E_j=\divisor(h^{m_j}+a_jF_j)$ and
\begin{equation}\label{eq:E-slice-boundary}
 \Psi=(1-\sum_j\theta_j)H+
                     \sum_j\frac{\theta_j}{m_j}E_j.
\end{equation}
Then $\Psi$ is effective and lc, and its nonzero zero-discrepancy
rays over the germ are exactly the common lc rays. On every such ray,
\begin{equation}\label{eq:E-pencil-order}
 v(h^{m_j}+a_jF_j)=m_jv(h).
\end{equation}
In particular, the coefficients $\theta_j$ may be chosen arbitrarily
small while retaining these conclusions.
\end{lemma}

\begin{proof}
Set $\mathfrak b_j=(h^{m_j},F_j)$. For any valuation,
\begin{align*}
 A(v)&-(1-\sum_j\theta_j)v(h)
               -\sum_j\frac{\theta_j}{m_j}v(\mathfrak b_j)\\
 &=A(v)-v(h)+\sum_j\theta_j
       \left(v(h)-\min\{v(h),v(F_j)/m_j\}\right)\ge0.
\end{align*}
Equality holds exactly on the common lc rays, together with the
trivial valuation. Resolve the ideals and choose the $E_j$ as
general members of their pencils. On the resolution the residual
free members meet the existing boundary with snc support and
contain none of its strata. Their coefficients $\theta_j/m_j$ are
less than one. They therefore introduce no new zero-discrepancy
strata, proving the assertion about $\Psi$.
On the coefficient-one monomial skeleton, each free residual member
has value zero because it contains no stratum. The value of its
pencil member is consequently the fixed-ideal value
$\min\{m_jv(h),v(F_j)\}=m_jv(h)$, proving
\eqref{eq:E-pencil-order}.
\end{proof}

\begin{theorem}[Exact common-slice minimizer]\label{thm:E-exact-minimizer}
For every allowed height $b$, the homogeneous functional
$b(v)\vol(v)^{1/n}$ on centred common lc rays has a unique minimizing
ray. Its $A=1$ representative depends continuously on the coefficients
of $b$ while the height stays strictly positive on the fixed slice.
It is fixed by every automorphism preserving the slice and the height.

If the coefficients of $b$ are rational, its minimizer $w$ is the
normalized-volume minimizer of an effective rational klt pair on
$x\in S$. More precisely, one can choose an effective rational lc
Cartier boundary $\Psi$ with exactly the common lc rays as its
zero-discrepancy locus, and one positive rational number $\lambda$,
fixed independently of the competing valuation, such that
\begin{equation}\label{eq:E-klt-perturbation}
 \Psi_\lambda=\Psi-\lambda b_0H
                       +\lambda\sum_\ell b_\ell\divisor(G_\ell)
\end{equation}
is effective and klt, and $w$ minimizes
$\widehat{\vol}_{S,\Psi_\lambda}$ among all centred valuations.
On the common lc rays its discrepancy is $\lambda b$.

For the same rational height, the algebra $\operatorname{gr}_wR$
is finitely generated and normal.
Its spectrum $S^0$ is a klt $\Q$-Gorenstein cone.
Write $h^0,F_j^0,G_\ell^0$ for initial forms, $w^0$ for the grading
valuation, and $\mathcal Q^0,b^0$ for the analogous central slice
and height. Then $(S^0,\divisor(h^0))$ is lc,
$A_{S^0}(w^0)=w^0(h^0)=1$, $w^0\in\mathcal Q^0$, and
$b^0$ is strictly positive on $\mathcal Q^0$.
The valuation $w^0$ minimizes $b^0\vol^{1/n}$ on its centred part.
\end{theorem}

\begin{proof}
\emph{Attainment, uniqueness, and continuity.}
Compactness gives $\min_{\mathcal Q}b>0$.
A minimizing sequence has bounded volume, so
\eqref{eq:E-centering} bounds its maximal-ideal values away from zero.
A subsequence converges in $\mathcal Q$ to a centred valuation.
Lemma~\ref{lem:E-compactness} gives volume continuity and hence
attainment. If two minimizing rays existed, normalize each by $b=1$.
Their midpoint has a dominating common lc interpolant $u$ with
$b(u)\le1$. If the minimum value of $b\vol^{1/n}$ is $M$, then
$\vol(u)\ge M^n/b(u)^n\ge M^n$.
The strict gap \eqref{eq:E-volume-gap} contradicts this unless the
two height-normalized valuations agree on every regular function.
Thus they are the same ray. For heights whose coefficients converge
to those of a positive height, the minimum heights stay bounded below
and one fixed centred competitor bounds the minimum values above.
The same compactness argument shows that every subsequential limit
of minimizers is the unique limiting minimizer. This proves continuity.
Automorphisms preserving the data preserve the functional, so
uniqueness proves equivariance.

\emph{Perturbing the boundary.}
Assume from now on that the height is rational. Choose $\Psi$ by
Lemma~\ref{lem:E-slice-boundary}, and resolve additionally all
$G_\ell$ and $\mathfrak m$.
For small rational $\lambda>0$, \eqref{eq:E-klt-perturbation} is
effective, since the coefficient of $H$ in \eqref{eq:E-slice-boundary}
is positive. It is klt near $x$: on each coefficient-one component
its new discrepancy is $\lambda b>0$, while the other finitely many
resolved discrepancies remain positive for sufficiently small
$\lambda$. The log-smooth criterion then tests every valuation.

\emph{Uniform control of every relevant competitor.}
Let $v$ be a centred quasi-monomial valuation with $A(v)=1$ and
\begin{equation}\label{eq:E-sublevel}
 A_{S,\Psi_\lambda}(v)\vol(v)^{1/n}\le\lambda M_0,
\end{equation}
where $M_0$ is the value of $b\vol^{1/n}$ at one fixed centred point
of $\mathcal Q$. This includes every normalized-volume minimizer of
the klt pair. Existence and quasi-monomiality for effective rational
klt pairs are supplied by \cite[Section~7]{Blum} and
\cite[Theorem~1.2]{XuQuasimonomial}.
The following constants are uniform over \emph{all} valuations in
\eqref{eq:E-sublevel} and all sufficiently small $\lambda$.

Put $e=A_{S,\Psi}(v)\ge0$ and retract $v$ to the coefficient-one
part of the fixed resolution, obtaining $r$.
Izumi bounds every fixed function value at $v$ when $A(v)=1$;
thus $|b(v)|\le C$. The first bound in
\eqref{eq:E-centering} and \eqref{eq:E-sublevel} give
\[
 e+\lambda b(v)\le C\lambda,\qquad e\le C\lambda.
\]
Equation~\eqref{eq:E-retraction-errors} yields
\begin{equation}\label{eq:E-uniform-retraction}
 A_S(r)=1+O(e),\qquad b(v)=b(r)+O(e).
\end{equation}
For small $\lambda$, $r$ is nontrivial and its normalized ray lies in
$\mathcal Q$. Strict positivity of $b$ on the compact slice therefore
gives uniform constants $\beta,B>0$ with
\[
 \beta\le b(r)\le B,\qquad b(v)\ge\beta/2.
\]
The sublevel inequality now bounds $\vol(v)$ above. The second bound
in \eqref{eq:E-centering} bounds $v(\mathfrak m)$ below; since
$\mathfrak m$ was resolved, its value at $r$ differs by $O(e)$.
Thus $r(\mathfrak m)$ is also bounded below and $r$ is an admissible
centred competitor. Corollary~\ref{cor:E-all-functions} supplies
\begin{equation}\label{eq:E-fixed-comparison}
 v(f)\le(1+Ce)r(f)\qquad\text{for every }f\in R,
\end{equation}
with a single $C$ independent of $f$, $v$, and small $\lambda$.

\emph{One fixed parameter forces exactness.}
Use \eqref{eq:E-uniform-retraction} and
\eqref{eq:E-fixed-comparison} to obtain
\begin{align}
 &(e+\lambda b(v))\vol(v)^{1/n}
       -\lambda b(r)\vol(r)^{1/n}\notag\\
 &\hspace{1cm}\ge
 \frac{e\,[1-\lambda C'-\lambda Cb(r)]}{1+Ce}
                 \vol(r)^{1/n}.\label{eq:E-exact-excess}
\end{align}
Choose \emph{one} positive rational $\lambda$ small enough for all
preceding bounds and for $\lambda(C'+CB)<1$.
This choice depends only on the fixed germ, boundary, and height;
it is independent of every competitor in \eqref{eq:E-sublevel} and
of every regular function. If such a competitor has $e>0$,
\eqref{eq:E-exact-excess} is strictly positive, so retraction strictly
improves its normalized-volume root functional.
An actual klt minimizer consequently has $e=0$, lies exactly on the
common lc slice, and is $w$ by uniqueness. Every competitor outside
\eqref{eq:E-sublevel} already has larger value than the fixed slice
competitor. This proves exact realization for this single rational
parameter; it is not merely a limiting assertion as $\lambda\to0$.
The argument applies for every smaller positive rational parameter
satisfying the same bounds.

\emph{Stable degeneration of this actual minimizer.}
We now apply stable degeneration to the effective rational klt pair
$(S,\Psi_\lambda)$ and its actual minimizer $w$.
By \cite[Theorems~1.1--1.2]{XuZhuangStable}, its associated graded is
finitely generated and defines a normal klt log Fano cone pair
$(S^0,\Psi_\lambda^0;w^0)$. The grading minimizes normalized volume
on the central pair; ordinary volume and pair discrepancy are
preserved, as in \cite[Lemma~4.10 and Theorem~3.5, arXiv version]{LiXuStable}.
The general uniqueness theorem for the klt-pair minimizer is
\cite[Theorem~1.1]{XuZhuangUnique}; the constrained uniqueness needed
here was proved before invoking it.

Each component used in \eqref{eq:E-klt-perturbation} is the divisor
of a regular function. For a principal ideal $(f)$ and a valuation
filtration, multiplicativity gives
$\operatorname{in}_w((f))=(\operatorname{in}_wf)$.
Equivalently its Rees homogenization, using its exact order, cuts a
Cartier divisor without a fibre component. Hence
$\Psi_\lambda^0$ is an effective rational sum of principal Cartier
divisors. Since $K_{S^0}+\Psi_\lambda^0$ is $\Q$-Cartier, so is
$K_{S^0}$, and removing the effective boundary shows that $S^0$ is klt.
Preservation of pair discrepancy and of all equation orders gives
\[
 A_{S^0}(w^0)=A_S(w)=1,\quad
 w^0(h^0)=1,\quad w^0(F_j^0)=w(F_j)\ge m_j.
\]

The choice of the $\theta_j$ can be made arbitrarily small. For each
such choice the exact minimizer is still the unique $w$, and thus
the graded algebra remains the same. Choose the corresponding
$\lambda$ smaller as necessary. The central effective klt boundaries
then contain $(1-o(1))\divisor(h^0)$.
For every prime over $S^0$, the resulting discrepancy inequalities
imply, on taking the coefficient limit,
$A_{S^0}(E)\ge\ord_E(h^0)$. This proves that
$(S^0,\divisor(h^0))$ is lc.

Finally take $v'\in\mathcal Q^0$. By \eqref{eq:E-pencil-order},
the initial equation of $h^{m_j}+a_jF_j$ is either $(h^0)^{m_j}$,
or $(h^0)^{m_j}+a_jF_j^0$ when $w(F_j)=m_j$.
In either case its $v'$-value is at least $m_jv'(h^0)$.
Therefore
\begin{equation}\label{eq:E-central-discrepancy}
 0<A_{S^0,\Psi_\lambda^0}(v')\le\lambda b^0(v'),
\end{equation}
with equality on the right at $w^0$.
The left inequality proves positivity of $b^0$ on the entire common
slice, including its noncentred faces. For centred $v'$, central
normalized-volume minimality and \eqref{eq:E-central-discrepancy}
give
\[
 \lambda b^0(v')\vol(v')^{1/n}
 \ge A_{S^0,\Psi_\lambda^0}(v')\vol(v')^{1/n}
 \ge\lambda b^0(w^0)\vol(w^0)^{1/n}.
\]
This is the required constrained minimality on the central slice.
\end{proof}

\begin{corollary}[Near-minimizers for a height]\label{cor:E-height-near-minimizers}
Fix an allowed height $b$ and normalize its common rays by $b=1$.
If $w_i,v_i$ have functional values within relative $o(1)$ of the
minimum, then $w_i=(1+o(1))v_i$ on all regular functions, with an error
independent of the function. The same conclusion, after $A=1$
normalization, holds for minimizers of heights whose coefficients
converge to those of $b$.
\end{corollary}
\begin{proof}
For the height-normalized pair, the midpoint interpolant has height
at most one. Its volume is therefore at least the minimum of volume
on the height-one slice. The proof using
\eqref{eq:E-volume-gap} applies verbatim: any fixed separation on one
function would force a fixed relative volume loss. For converging
heights $b_i$, positivity on the compact slice gives
$b_i/b=1+o(1)$ uniformly. Their minimizing rays, normalized by $b=1$,
are thus near-minimizers for $b$, and the preceding assertion applies.
Evaluation on $h$, which equals $A$ on the common rays, shows that
the factors needed to normalize by $A=1$ converge as well.
\end{proof}

\begin{corollary}[Domination by the minimizer]\label{cor:E-domination}
For the minimizing ray $w$ and every $v\in\mathcal Q$,
\begin{equation}\label{eq:E-domination}
 \frac{v(f)}{b(v)}\le 2^n\frac{w(f)}{b(w)}
                 \qquad(f\in R).
\end{equation}
For rational heights the same statement holds on the central slice
with all symbols carrying a superscript $0$.
\end{corollary}
\begin{proof}
Normalize both rays by $b=1$ and interpolate them with weights $1/2$.
The resulting centred $u$ has $b(u)\le1$, so minimality gives
$\vol(u)\ge\vol(w)$. For $w(f)>0$,
\eqref{eq:E-volume-interpolation} gives
\[
 1\le2^n\frac{w(f)}{w(f)+v(f)}.
\]
Thus $v(f)\le(2^n-1)w(f)$, which implies the stated bound.
If $f$ is a unit both values are zero. The central assertion follows
from the same argument and Theorem~\ref{thm:E-exact-minimizer}.
\end{proof}

\begin{lemma}[One-sided height tilt]\label{lem:E-one-sided-tilt}
Let $G\in R\setminus\{0\}$, $m\in\Z_{>0}$, and
$g(v)=v(G)/m-v(h)$. Suppose $g(v)/v(h)\ge2c$ on some common lc ray,
where $c>0$ is rational. Set
\[
 M=\max_{\mathcal Q}g\ge2c,\qquad
 \mu_*=(M-c)^{-1},\qquad
 b_\mu(v)=v(h)+\mu(cv(h)-g(v)).
\]
For $0\le\mu<\mu_*$ the height is positive. There is a rational
$\mu$ in this interval whose $A=1$ minimizing valuation $w$ satisfies
\[
 g(w)>0,\qquad b_\mu(w)\ge1.
\]
Moreover $b_\mu\le2+1/c$ on $\mathcal Q$, and the corresponding
central height obeys the same bound on $\mathcal Q^0$.
\end{lemma}

\begin{proof}
As $\mu\uparrow\mu_*$, the minimum of
$b_\mu\vol^{1/n}$ tends to zero. To prove this even when an endpoint
zero occurs at a noncentred valuation $v$, interpolate $v$ with one
fixed centred point $w_1\in\mathcal Q$, giving $w_1$ weight $t$.
For $0<\mu_*-\mu\le t$, the interpolant $u$ has
$b_\mu(u)=O(t)$. Since $v(h)=w_1(h)=1$, the improved volume estimate
\eqref{eq:E-volume-interpolation}, applied with $f=h$, gives
$\vol(u)\le t^{1-n}\vol(w_1)$.
Hence $b_\mu(u)\vol(u)^{1/n}=O(t^{1/n})$.
The lower volume bound in \eqref{eq:E-centering} then forces the
minimizing height to tend to zero. In particular its minimizing
valuation eventually has $g(w)>c$.

If the minimizer at $\mu=0$ has $g(w)>0$, take $\mu=0$.
Otherwise continuous dependence in
Theorem~\ref{thm:E-exact-minimizer} gives a nonempty open interval
of parameters for which $0<g(w)<c$. Choose a rational parameter
in that interval. Then $b_\mu(w)=1+\mu(c-g(w))\ge1$.
Finally $v(G)\ge0$ gives $g(v)\ge-1$ on the normalized slice, while
$\mu<\mu_*\le1/c$. Thus
$b_\mu(v)\le1+(c+1)/c=2+1/c$.
The same calculation applies centrally because $G^0$ is regular and
$v'(h^0)=1$ on $\mathcal Q^0$.
\end{proof}

\section{Degeneration, exact sign transfer, and rank refinement}\label{sec:F}

We work with a complex klt germ $x\in S$ of dimension $n$, a nonzero
regular function $h$ such that $(S,H)$ is lc, where $H=\divisor(h)$,
and a rational top form $\omega$ with $\divisor_S(\omega)=0$.
Thus the reflexive canonical sheaf is trivialized by $\omega$.
A finite cyclic group $G$ fixes $x$ and acts quasi-\'{e}tale on $S$;
we replace it by its effective image.  Suppose that $h$ and $\omega$
have the same character $\chi$.  In particular, $\omega/h$ descends
to the quotient.  Fix invariant ratios
\[
 f_j=F_j/h^{m_j},\qquad F_j\in\mathcal O_{S,x}\setminus\{0\},
 \quad m_j\in\N_{>0}.
\]
The common slice consists of the $h$-lc valuations satisfying
$q(f_j)\ge0$, with the convention on centres from Section~\ref{sec:E}:
the closed image of the centre contains $x$.  Unless another
normalization is displayed, $q(h)=A_S(q)=1$.

Let $w$ be a centred $G$-invariant quasi-monomial valuation for which
\[
 \breve R=\operatorname{gr}_w\mathcal O_{S,x}
\]
is finitely generated, normal, and klt.  These are hypotheses here;
they are supplied in our applications by the exact optimizer of
Section~\ref{sec:E}.  Write $S^0=\Spec\breve R$, with vertex $o$,
and denote initial symbols by a bar.  The fine grading is by the
finitely generated value group $\Lambda_w\subset\R$, so evaluation
at $w$ is injective on this lattice.  Its torus will be denoted by
$\mathbb T_w$.  Taking invariants in each filtered exact sequence gives
\begin{equation}\label{eq:F-invariant-graded}
 \operatorname{gr}_w(\mathcal O_{S,x}^{G})=\breve R^{G}.
\end{equation}
Indeed, the Reynolds operator makes taking invariants exact in
characteristic zero.  The weight lattices of $\breve R$ and
$\breve R^G$ have finite index in one another, and hence the same rank.

\subsection{Presentations and rational approximations}

\begin{lemma}[Weighted presentations]\label{lem:F-presentation}
Choose $x_1,\ldots,x_N\in\mathfrak m_x$ whose symbols generate
$\breve R$, and put $\alpha_i=w(x_i)>0$.  Any prescribed finite
list of functions of positive $w$-value can be included in this list.
Then
\begin{equation}\label{eq:F-presentation}
 \widehat{\mathcal O}_{S,x}=\C[[X_1,\ldots,X_N]]/I,
 \qquad
 \breve R=\C[X_1,\ldots,X_N]/\operatorname{in}_{\alpha}I,
\end{equation}
where initial forms have least weight.  The valuation ideals are
exactly the images of the ideals generated by monomials of weight
at least the prescribed threshold.  The same description holds
before completion by intersection.

For every sufficiently nearby linear functional
$L:\Lambda_w\otimes\R\longrightarrow\R$, positive on the listed
weights, the vector $\alpha_L=(L(\alpha_i))$ gives the same initial
ideal, with its grading composed with $L$.  Its filtration defines
a centred quasi-monomial valuation $w_L$.  The functionals can be
chosen rational, and the resulting valuations can be computed
monomially on a fixed smooth chart for $w$.
\end{lemma}

\begin{proof}
Subduction subtracts a polynomial in the $x_i$ with the prescribed
initial symbol, then repeats at a strictly larger value.  There are
only finitely many values in a bounded interval: each is a
nonnegative integral combination of the finitely many positive
$\alpha_i$.  The Izumi upper bound for the fixed valuation $w$
implies that remainders whose $w$-values tend to infinity tend to
zero in the $\mathfrak m_x$-adic topology.  In particular
$\mathfrak m_x\subset(x_1,\ldots,x_N)+\mathfrak m_x^2$, and
Nakayama's lemma shows that the $x_i$ generate $\mathfrak m_x$.
Completed subduction proves both assertions in
\eqref{eq:F-presentation} and the assertion about threshold ideals.
Monomial threshold ideals are finitely generated.  Their completed
extensions intersect the local ring in the original ideals, by
faithful flatness of completion.

We record why the dimension and primeness arguments used below are
valid also for other positive weight vectors.  For any such vector
$\gamma$, its monomial filtration is bounded above and below by
fixed linear reindexings of the $\mathfrak m_x$-adic filtration.
Consequently its cumulative quotient lengths grow with exponent
$n$.  The associated graded algebra is
$\C[X]/\operatorname{in}_{\gamma}I$; comparison of its positive
weighted filtration with ordinary polynomial degree shows that its
growth exponent is its Krull dimension.  Thus this dimension is
$n$.  If $\operatorname{in}_{\gamma}I$ is prime, the filtration is
multiplicative and separated, and defines a valuation.  If its
value group has rational rank $r$, the degree-zero homogeneous
fraction field is its residue field and has transcendence degree
$n-r$.  This proves Abhyankar equality.  Local monomialization of
Abhyankar valuations in characteristic zero
\cite[Proposition~3.7]{JonssonMustata} then makes the
valuation quasi-monomial.

Take finitely many series in $I$ whose $\alpha$-initials generate
$\operatorname{in}_{\alpha}I$.  Their leading faces persist after
composition with $L$: ties are preserved because their fine
weights are equal, and positivity leaves only finitely many other
terms that could compete in a fixed neighbourhood of $\alpha$.
Therefore
\[
 \operatorname{in}_{\alpha}I\subset
 \operatorname{in}_{\alpha_L}I.
\]
The first ideal is prime, and both quotients have dimension $n$.
A nonzero ideal in a finitely generated domain lowers its
dimension, so this inclusion is equality.

For the final assertion, use a smooth monomial chart for $w$ at its
stratum and expand the finitely many $x_i$ over the residue
coefficient field.  Their leading faces are finite polynomials in
the stratum parameters.  The weights of these parameters belong
to $\Lambda_w$.  Composing them with nearby rational $L$ preserves
the same leading polynomials.  The resulting monomial valuation
dominates the presentation filtration.  The induced map from the
presentation graded algebra to the graded ring of this monomial
valuation is injective: the polynomial relations among these
leading polynomials are exactly the old relations, unchanged by
regrading.  Hence the two valuations agree on every function.
\end{proof}

\subsection{The canonical calculation before lc equality}

Xu--Zhuang study specialization of non-degenerate reflexive
differential forms under stable degeneration of klt singularities
\cite[Theorem~1.1 and Section~3]{XuZhuangForms}.
The calculation below tracks a divisor-free canonical generator,
including its finite-group character and any additional grading,
under the finitely generated degeneration used here.  We give this
calculation before imposing lc equality because it will also be used
to prove that equality after two gradings are compressed.

\begin{proposition}[Canonical forms under degeneration]
\label{prop:F-canonical-degeneration}
Under the hypotheses preceding Lemma~\ref{lem:F-presentation},
$S^0$ has a divisor-free reflexive canonical generator $\omega^0$.
It has the $G$-character of $\omega$.  It is homogeneous for
$\mathbb T_w$, and its weight $\kappa$ satisfies
\begin{equation}\label{eq:F-canonical-weight}
 \langle w,\kappa\rangle=A_S(w).
\end{equation}
If an additional torus acts on the original germ, preserves $w$,
and makes $\omega$ homogeneous, its character is preserved
separately.  None of these assertions assumes $A_S(w)=w(h)$.
Moreover the action of $G$ on $S^0$ is quasi-\'{e}tale.
\end{proposition}

\begin{proof}
First track discrepancy under the rational approximations of
Lemma~\ref{lem:F-presentation}.  On a fixed monomial chart with
stratum parameters $z_1,\ldots,z_r$ and separating residue
coordinates $y_1,\ldots,y_{n-r}$, write
\[
 \omega=\psi\,
 \frac{dz_1}{z_1}\wedge\cdots\wedge\frac{dz_r}{z_r}
 \wedge dy_1\wedge\cdots\wedge dy_{n-r}.
\]
The quasi-monomial discrepancy formula gives
$A_S(w)=w(\psi)\in\Lambda_w$.  Include numerators and denominators
of the rational function $\psi$ among the finite functions whose
leading faces are tracked.  The same formula on the same chart
then gives the exact identity
\begin{equation}\label{eq:F-discrepancy-approximation}
 A_S(w_L)=L(A_S(w)).
\end{equation}
In particular this is an identity of linear weights, rather than
only a limiting assertion about discrepancies.

Choose rational $L$ and multiply it by a positive integer so that
the presentation weights are positive integers.  In this paragraph
write $v$ for this integral scaling of $w_L$, and $a_i=v(x_i)$.
After shrinking to a $G$-stable affine neighbourhood, the $x_i$
have common zero set $\{x\}$.  Define the extended Rees algebra
\begin{equation}\label{eq:F-rees}
 \mathcal R=\mathcal O(S)[t,x_1t^{-a_1},\ldots,x_Nt^{-a_N}]
       \subset\mathcal O(S)[t,t^{-1}].
\end{equation}
The threshold-ideal description proves that its special fibre is
the graded algebra in \eqref{eq:F-presentation}; positive-level
quotients are supported at $x$, so the affine and local
constructions have the same special fibre.  This algebra is of
finite type and torsion free over $\C[t]$, hence flat.  Let
$\mathcal S=\Spec\mathcal R$ and $E=(t=0)$.  The divisor $E$ is
normal, integral, reduced, and Cartier, and
$\mathcal S\setminus E=S\times\mathbb G_m$.

The total space is normal.  Away from $E$ this follows from the
normality of $S$.  At points on $E$, the nonzerodivisor $t$ and
the $S_2$ property of $\mathcal R/(t)$ imply $S_2$ for
$\mathcal R$.  At the generic point of $E$, the quotient by $t$
is a field, so the total local ring is a regular discrete
valuation ring.  These observations give $R_1$ as well.

The order along $E$ restricts to $v$ on $\C(S)$ and takes value
one on $t$.  Indeed, the residues of $f t^{-v(f)}$ are the
nonzero graded symbols of $f$, and the direct grading prevents
cancellation between distinct such leading terms.  It is the
Gauss extension of $v$.  The monomial log-form computation gives
\begin{equation}\label{eq:F-rees-log-order}
 \divisor_{\mathcal S}\left(\omega\wedge\frac{dt}{t}\right)
       =(A_S(v)-1)E.
\end{equation}
For example, on a model where the rational valuation $v$ is
$c\ord_P$, use a uniformizer of $P$ and separating residue
coordinates.  The extension has weights $(c,1)$ on that
uniformizer and $t$; the wedge of their logarithmic differentials
has log order zero, while the remaining coefficient has value
$cA_S(\ord_P)=A_S(v)$.  This proves the coefficient in
\eqref{eq:F-rees-log-order}.  There are no other divisor
coefficients, since off $E$ the form is the product of the
divisor-free $\omega$ and $dt/t$.

Since $K_S$ is Cartier and $v$ has integral values, $A_S(v)$ is an
integer.  Therefore
\[
 \eta=t^{-A_S(v)}\omega\wedge\frac{dt}{t}
 \qquad\text{satisfies}\qquad \divisor_{\mathcal S}(\eta)=-E.
\]
This establishes that the ordinary canonical divisor of the
total space is Cartier.  At every codimension-one point of $E$,
the fibre is regular and the flat morphism to the line is smooth.
The residue of $\eta$ is consequently a rational top form on $E$
with neither zeros nor poles at any codimension-one point.
It trivializes the reflexive canonical sheaf of $E$.
In particular this establishes ordinary canonical Cartierness,
not just Cartierness after adding a boundary.

Give $t$ weight $-1$ and $x_i t^{-a_i}$ weight $a_i$.  The residue
has weight $A_S(v)$ and the same $G$-character as $\omega$.
If an old torus is present, choose homogeneous generators;
$t$ is invariant for that torus, so its canonical character is
unchanged.  In the cone case such homogeneous generators have
common zero only at the vertex: their common zero set is
invariant and is isolated there, whereas the closure of every
positive-grading orbit contains the vertex.

Units of a positively graded affine domain with degree-zero part
$\C$ are constant.  Hence any two divisor-free canonical
generators on $E$ differ by a scalar.  The full fine grading
torus therefore acts on the generator by a character $\kappa$,
independent of the rational regrading used to construct it.
For all nearby rational $L$, its evaluation is
$L(A_S(w))$ by \eqref{eq:F-discrepancy-approximation} and the
residue computation.  Equality on this open set of rational
functionals identifies the two lattice weights and proves
\eqref{eq:F-canonical-weight}, including its equivariant version.

Finally suppose a nonidentity element of $G$ acquired a fixed
locus of codimension at most one on $E$.  At a general point of
that locus the fibre, and hence the morphism $\mathcal S\to\mathbb A^1$,
is smooth.  The fixed locus of the cyclic subgroup generated by
that element is smooth there.  Its tangent space is the invariant
tangent subspace, and it maps surjectively to the tangent space
of the line: averaging a lift of the base tangent vector gives
an invariant lift.  Thus this fixed locus is smooth over the
line and has the same fibre codimension on nearby nonzero
fibres.  This contradicts quasi-\'{e}taleness of the original
effective action.  The action on $S^0$ is therefore
quasi-\'{e}tale.
\end{proof}

\begin{corollary}[Lc equality and principal boundaries]
\label{cor:F-lc-degeneration}
If $w$ is $h$-lc, then $(S^0,\divisor(\bar h))$ is lc and
$\omega^0$ has the same fine torus weight and $G$-character as
$\bar h$.  Consequently $\omega^0/\bar h$ is invariant, and the
quotient cone has the log-form data of Section~\ref{sec:D}.
If $\Delta$ is any effective rational sum of principal divisors
near $x$, with $(S,\Delta)$ lc and
$A_S(w)=w(\Delta)$, its initial boundary $\Delta^0$ is lc on
$S^0$.  All these assertions also hold with specified additional
equivariant gradings.
\end{corollary}

\begin{proof}
Include the finitely many boundary equations in the presentation.
Their orders and discrepancy equality persist exactly under
rational regrading.  In \eqref{eq:F-rees}, each equation $f$ has
homogenization $f t^{-v(f)}$, which is a nonzerodivisor and has
no component $E$; it restricts to $\bar f$ on $E$.  In particular
the initial boundary is rationally Cartier.  The morphism
$\mathcal S\to S\times\mathbb A^1$ is birational.  Comparing the
coefficient at its only vertical prime by
\eqref{eq:F-rees-log-order} shows that
\[
 (\mathcal S,E+\Delta_{\mathcal S})
 \longrightarrow (S\times\mathbb A^1,
                       \Delta\times\mathbb A^1+S\times\{0\})
\]
is crepant precisely when $A_S(v)=v(\Delta)$.
Its source is therefore lc, and lc adjunction gives the stated
central pair.  At codimension-one points of the normal Cartier
fibre the morphism is smooth, so the different is exactly the
ordinary restriction of the homogenized boundary.
Apply this to $\Delta=H$.  Formula~\eqref{eq:F-canonical-weight}
and injectivity of evaluation on $\Lambda_w$ identify the weight
of $\omega^0$ with that of $\bar h$.  Quasi-\'{e}taleness and the
finite-map discrepancy formula descend klt, lc, and the invariant
log form to the quotient.  A suitable power of $\omega^0$
descends to a divisor-free pluricanonical form.
\end{proof}

\subsection{An exact bridge for inequalities}

\begin{proposition}[Sign transfer]\label{prop:F-bridge}
Suppose $w$ is common $h$-lc and $w(f_j)=0$ for every $j$.
Let $g=U/h^b$, where $0\ne U\in\mathcal O_{S,x}$,
$b\in\N_{>0}$, and $w(g)=0$.  Let $\mathcal Q$ be the actual
common slice and $\mathcal Q^0$ the common slice defined by
$\bar h$ and the $\bar f_j$ on $S^0$.  Then
\begin{equation}\label{eq:F-bridge}
 q(g)\le0\quad(q\in\mathcal Q)
 \quad\Longleftrightarrow\quad
 q^0(\bar g)\le0\quad(q^0\in\mathcal Q^0).
\end{equation}
One may add any finite list of tests of value zero at $w$.
There is no bound on the degree or the cost of $g$ in this statement.
\end{proposition}

\begin{proof}
Use Lemma~\ref{lem:E-slice-boundary} to encode the common slice.
Choose positive rational $\theta_j$ with $\sum_j\theta_j<1$ and
sufficiently general scalars $\lambda_j$, and set
\begin{equation}\label{eq:F-test-boundary}
 P_j=h^{m_j}+\lambda_jF_j,
 \qquad
 \Psi=(1-\textstyle\sum_j\theta_j)H+
                \sum_j\frac{\theta_j}{m_j}\divisor(P_j).
\end{equation}
The lemma says that $\Psi$ is effective and lc, with exactly the
common rays as its nonzero zero-discrepancy rays. It also gives
$q(P_j)=m_jq(h)$ on every common ray. These two properties let us
express a sign inequality as log canonicity of a small boundary
perturbation.

Assume the left side of \eqref{eq:F-bridge}.  For sufficiently
small positive rational $\delta$ the divisor
\begin{equation}\label{eq:F-perturbed-boundary}
 \Psi_\delta=\Psi+\delta(\divisor(U)-bH)
\end{equation}
is effective and lc near $x$.  Here is the exact reason for
smallness sufficing.  The coefficient of $H$ in
\eqref{eq:F-test-boundary} has a positive reserve, which preserves
effectivity.  On a fixed log resolution including $U$, every
coefficient below one stays below one for small $\delta$; for
a coefficient-one prime the discrepancy change is
$-\delta q(g)\ge0$.  Checking these finitely many components
checks lc on the entire resolution.  Restricting the germ if
necessary discards centres whose closed image does not contain
$x$.  No uniform choice of $\delta$ in $U$ is needed.

All equations in \eqref{eq:F-perturbed-boundary} have their
expected orders at $w$, and
$A_S(w)=w(\Psi_\delta)$ because $w(g)=w(f_j)=0$.
Corollary~\ref{cor:F-lc-degeneration} gives an lc initial boundary.
For every $q^0\in\mathcal Q^0$,
\[
 q^0(\overline{P_j})\ge m_jq^0(\bar h).
\]
Consequently
\[
 0\le A_{S^0,\Psi_\delta^0}(q^0)
       \le-\delta q^0(\bar g),
\]
which proves the right side of \eqref{eq:F-bridge}.

Conversely assume that right side.  Carry out exactly the same
boundary construction on $S^0$, choosing the $\lambda_j$ general
for both the actual and central constructions.  The intersection
of these finitely many nonempty open conditions is nonempty.
The assumed sign and a log resolution of the central equations
give a small rational $\delta>0$ for which $\Psi_\delta^0$ is
effective and lc near $o$.  Choose an integral rational regrading
that preserves all the equations involved.  Homogenize them in
the Rees family.  Their rational linear combination
$\Delta_{\mathcal S}$ is an effective rationally Cartier
horizontal boundary, disjoint in components from $E$, and its
restriction is $\Psi_\delta^0$.

The total space is normal, $E$ is normal and Cartier, and its
ordinary canonical divisor is Cartier by
Proposition~\ref{prop:F-canonical-degeneration}.  Thus
$K_{\mathcal S}+E+\Delta_{\mathcal S}$ is $\Q$-Cartier.
The different is the ordinary restriction, as checked above.
Inversion of adjunction for log canonicity
\citep[Theorem]{Kawakita} now makes
$(\mathcal S,E+\Delta_{\mathcal S})$ lc near $o$.
The section with all rescaled coordinates $x_i t^{-a_i}$ zero
specializes to $o$, so this neighbourhood meets it over a
nonzero parameter.  Off $E$ the family and its horizontal
boundary are the product with $(S,\Psi_\delta)$.  Hence the latter
pair is lc near $x$.
For a common actual ray, the initial boundary encoding has zero
discrepancy, and the last perturbation subtracts $\delta q(g)$.
It follows that $q(g)\le0$, proving the converse.

The resolution computations give inequalities on every divisorial
ray, and continuity on the monomial cones of the lc skeleton
gives them on the full quasi-monomial slices.  Additional finite
tests are treated in the same boundary and presentation.
\end{proof}

\begin{corollary}[Exact tests and pure weights]
\label{cor:F-exact-tests}
Suppose $w$ is common $h$-lc. If every common actual valuation
satisfies $q(f_j)=0$, then every common cone valuation satisfies
$q^0(\bar f_j)=0$.
The symbols $\bar f_j$ are invariant under $G$ and $\mathbb T_w$.
Every nonzero pure weight valuation of the cone is $h$-lc and,
after normalization by $\bar h=1$, is common.  Pure weights on
the boundary of the dual weight cone are allowed under the
closed-centre convention.
\end{corollary}

\begin{proof}
Apply Proposition~\ref{prop:F-bridge} to $g=f_j$, and combine its
upper bound with the defining lower bound.  The equality
$w(f_j)=0$ and injectivity of the fine grading imply that
$\bar f_j$ has weight zero.  The invariant log form
$\omega^0/\bar h$ has weight zero, so the pure-weight discrepancy
computation of Section~\ref{sec:D} gives lc equality for every
pure weight.  The test ratios then all have value zero.  Klt
implies that $\bar h$ is strictly positive on every nonzero
such weight, so the stated normalization is possible.
\end{proof}

\subsection{Compression of two gradings}

A second optimization takes place on the first cone.  To iterate the
construction on the original germ, we must realize these two successive
gradings by one actual valuation there.  The canonical-character
calculation will prove its lc equality, while the exact bridge keeps the
previous order constraints.

\begin{lemma}[Compression]\label{lem:F-compression}
Suppose $w$ is as above, is $h$-lc, and satisfies $w(f_j)=0$ for
every $j$. Let $v$ be a centred
$G\times\mathbb T_w$-invariant, $\bar h$-lc quasi-monomial
valuation on $S^0$, with normal finitely generated klt graded
algebra $R^{00}$.  Suppose also $v(\bar f_j)=0$.
The algebra $R^{00}$ has its joint grading by the old torus and
the $v$-grading.  For arbitrarily small positive $\delta$ outside
a countable exceptional set, it is the associated graded of an
actual centred quasi-monomial valuation $w_\delta$ on $S$, whose
grading evaluates a joint degree $(a,b)$ as $a+\delta b$.
The valuation $w_\delta$ is $h$-lc, is exact on all $f_j$, and,
after $h=1$ normalization, tends multiplicatively to $w$
uniformly on all nonzero regular functions.  If $v$ is nonzero
on an old degree-zero rational function, the rational rank of
$w_\delta$ is strictly larger than that of $w$.
\end{lemma}

\begin{proof}
Choose old-homogeneous elements of $\breve R$ whose $v$-symbols
generate $R^{00}$ and lift them to regular functions $x_i$ on
$S$.  Enlarge the list to generate the first graded algebra and
to include $h,F_j$ and every other specified equation.
Use $G$-eigenvectors throughout.  Put
$\alpha_i=w(x_i)$ and $\beta_i=v(\bar x_i)>0$.
Both successive subductions of
Lemma~\ref{lem:F-presentation} give
\begin{equation}\label{eq:F-iterated-initial}
 R^{00}=\C[X]/\operatorname{in}_\beta
                                  (\operatorname{in}_\alpha I).
\end{equation}
Completion of the first ideal does not change its initial ideal
for positive $\beta$.

Choose finitely many generators of the prime ideal on the right
that are iterated initials of elements of $I$.  This is possible
by first choosing $\alpha$-homogeneous elements of
$\operatorname{in}_\alpha I$ with the desired $\beta$-initials;
such elements lift with their prescribed $\alpha$-initial.
For each of these finitely many series, positivity and finiteness
of the relevant leading faces imply that its iterated initial
is its $(\alpha+\delta\beta)$-initial for all sufficiently small
$\delta>0$.  Hence
\[
 \operatorname{in}_\beta\operatorname{in}_\alpha I
           \subset\operatorname{in}_{\alpha+\delta\beta}I.
\]
The first ideal is prime and both quotients have dimension $n$.
They are equal.  Lemma~\ref{lem:F-presentation} supplies the
actual multiplicative valuation $w_\delta$ and its
quasi-monomiality.  Avoiding countably many values of $\delta$
makes $(a,b)\mapsto a+\delta b$ injective on the joint lattice.

We now prove the lc equality, after constructing the valuation.
The first degeneration has canonical character equal to the
weight of $\bar h$.  Apply
Proposition~\ref{prop:F-canonical-degeneration} to $v$ on that
cone, keeping the old torus action.  Since $v$ is $\bar h$-lc,
the resulting canonical generator has exactly the joint
character of $\bar{\bar h}$.  Independently apply the same
proposition to the actual valuation $w_\delta$, whose normal
finitely generated graded algebra has just been identified with
$R^{00}$.  Its discrepancy is the evaluation of this canonical
character under compression.  Therefore
\begin{equation}\label{eq:F-compressed-lc}
 A_S(w_\delta)
   =\langle w_\delta,\operatorname{wt}(\bar{\bar h})\rangle
   =w_\delta(h).
\end{equation}
This uses no prior lc assumption about $w_\delta$.
The included coordinates identify the specified symbols, and
the joint degree of each $\bar{\bar f_j}$ is zero; hence
$w_\delta(f_j)=0$.

For the uniform comparison, if
$(1-\varepsilon)\alpha_i\le\alpha_i+\delta\beta_i
\le(1+\varepsilon)\alpha_i$ for every $i$, the same inequalities
hold for every monomial weight.  The threshold-ideal description,
including its intersection with the actual local ring, gives
\[
 (1-\varepsilon)w(f)\le w_\delta(f)
                  \le(1+\varepsilon)w(f)
 \qquad(0\ne f\in\mathcal O_{S,x}).
\]
Normalization by $h$ changes the factors by another quantity
tending to one.  Finally the joint lattice surjects onto the old
lattice.  A degree-zero rational function with nonzero $v$-value
provides a nonzero vector in its kernel.  Injective compression
therefore strictly increases rational rank.
\end{proof}

\subsection{Balanced cones satisfying the angular order condition}

\begin{proposition}[Rank refinement]\label{prop:F-refinement}
Consider a sequence of the preceding data in dimensions at most
$n_0$.  Suppose that the common actual slices satisfy
$q(f_j)=0$ identically, that $w$ is a centred point of the slice
with the stipulated normal finitely generated klt graded algebra,
and that the cones satisfy the following balance condition with
one constant $C$ along the sequence:
\begin{equation}\label{eq:F-balance}
 q^0(s)\le Cw^0(s)
 \quad\begin{gathered}
 0\ne s\in\breve R\text{ fine-homogeneous},\\[-2pt]
 q^0\text{ common},\quad q^0(\bar h)=w(h)=1.
 \end{gathered}
\end{equation}
After passage to a subsequence, the $w$ can be replaced by $w'$
with all these properties, with a new fixed balance constant,
and with
\begin{equation}\label{eq:F-refinement-comparison}
 w'(f)=(1+o(1))w(f)
 \qquad(0\ne f\in\mathcal O_{S,x}),
\end{equation}
uniformly in $f$ in the normalization $w'(h)=w(h)=1$.
On the angular model $V_{\mathbf u}$ of the new cone, the
notation of Section~\ref{sec:D} then satisfies
\begin{equation}\label{eq:F-angular-order}
 T(\Gamma)\le(1+\eta_i)D(T),\qquad \eta_i\longrightarrow0,
\end{equation}
for every effective $\Q$-divisor $\Gamma\sim_\Q D$ and every
primitive divisorial $T$ such that
$a(T)=0$ and $T(\bar f_j)\ge0$ for all $j$.
This is condition \textup{(C-$\alpha$)} of Section~\ref{sec:C}.
The initial symbols here refer to $w'$.

The balance hypothesis holds when $w$ is the exact minimizer
on the common slice for the height $q(h)$.
\end{proposition}

\begin{proof}
The last assertion is the central domination estimate in
Corollary~\ref{cor:E-domination}; equivariance follows from the
uniqueness of the slice minimizer.  At every stage,
Corollary~\ref{cor:F-exact-tests} makes all common central test
values zero and makes every pure weight common.  The $\bar f_j$
are genuine functions of the angular field $\mathcal K$.
At a height-minimizing Gauss extension of $T$, regularity of
$\bar F_j$ gives
\begin{equation}\label{eq:F-angular-test-cost}
 T(\bar f_j)\ge-m_jD(T).
\end{equation}

Suppose \eqref{eq:F-angular-order} fails along a subsequence.
There are a fixed rational $c>0$, primitive witnesses $T$, and
effective divisors
\[
 \Gamma=D+\frac1m\divisor(\phi),\qquad m\in\N_{>0},
 \quad\phi\in\mathcal K^*,
\]
such that $a(T)=0$, $T(\bar f_j)\ge0$, and
\begin{equation}\label{eq:F-refinement-witness}
 T(\Gamma)-D(T)\ge2cD(T).
\end{equation}
Here $D(T)>0$ because $a+\mathbf D$ is klt.
After clearing denominators, the ray-ring description in
Section~\ref{sec:D} makes
$G_0=\bar h^m\phi$ a regular homogeneous function upstairs.
Choose a rational Gauss extension minimizing $\bar h$ at $T$.
It is $\bar h$-lc, and its lift to the finite quasi-\'{e}tale
cover cone remains lc.  Its centre has the vertex in its closed
image: this holds downstairs for a Gauss valuation, and finiteness
and the unique point over the vertex give it upstairs.
It is common, and \eqref{eq:F-refinement-witness} reads
\[
 \frac{g(q)}{q(\bar h)}\ge2c,
 \qquad g(q)=\frac1m q(G_0)-q(\bar h).
\]

Apply the one-sided tilt of Lemma~\ref{lem:E-one-sided-tilt} to
this cone.  Its exact minimizer $v$, normalized by $v(\bar h)=1$,
is centred, $G\times\mathbb T_w$-invariant, and has normal
finitely generated klt associated graded.  Localization at the
vertex does not change this graded algebra: a unit has constant
initial term.  The valuation $v$ is common, hence exact on the
old tests, and the tilt gives
\[
 v(\phi)>0,\qquad b(v)\ge1.
\]
The tilted height $b$ is bounded above by a fixed constant on
both the old and the new common slice, with constants depending
only on $n_0$ and the fixed $c$ at this stage.
Central domination therefore gives, on the new cone,
\begin{equation}\label{eq:F-tilt-balance}
 q^{00}(s)\le C_1v^0(s)
\end{equation}
for common $q^{00}$ normalized by $\bar{\bar h}=1$.
The joint canonical character equals that of $\bar{\bar h}$ by
Proposition~\ref{prop:F-canonical-degeneration}, including its
old-torus assertion.

Compress these two gradings by Lemma~\ref{lem:F-compression}.
Since $\phi$ has old weight zero and $v(\phi)>0$, rational rank
strictly increases.  The resulting actual valuation is $h$-lc,
exact on the tests, and can be chosen arbitrarily close to $w$
uniformly on every regular function.  We check balance, which
is necessary to repeat the argument.  A joint-homogeneous $s$
has an old-homogeneous lift $\widetilde s\in\breve R$; the
$v$-filtration splits into old weight spaces.  If its joint
degree is $(a,b)$, then
\[
 b=v(\widetilde s)\le Ca
\]
by \eqref{eq:F-balance}, since $v$ is common.  In the
normalization $w(h)=v(\bar h)=1$, the compressed degree is
$(a+\delta b)/(1+\delta)$.  Thus
\[
 v^0(s)=b\le C(1+\delta)w_\delta^0(s).
\]
Together with \eqref{eq:F-tilt-balance} this is the next
balance estimate, with a fixed constant.  A generic compression
makes every fine-homogeneous element joint-homogeneous, so this
checks precisely all elements required in \eqref{eq:F-balance}.

If the angular order condition still fails, pass to a subsequence
with a fixed positive defect and repeat.  Each repetition raises
rational rank, which is at most the dimension.  There can be at
most $n_0$ repetitions.  At a final stage the supremum of the
positive normalized defects must tend to zero.  This supremum is
finite: a height-minimizing Gauss witness and
\eqref{eq:F-balance}, applied to $G_0$ of grading value $m$,
bound $T(\Gamma)/D(T)$ by the balance constant.
Taking these suprema, or slightly larger positive numbers, gives
the single sequence $\eta_i$ in \eqref{eq:F-angular-order},
uniformly for all its displayed tests and divisors.
At each of the boundedly many compressions choose the
all-functions multiplicative error to tend to zero; their product
still tends to one and proves \eqref{eq:F-refinement-comparison}.
All exact test equalities, equivariance, canonical characters,
and quasi-\'{e}taleness have been preserved at every step.
\end{proof}

\section{Pinning the gradings}\label{sec:G}

We now combine the character supply of Section~\ref{sec:D} with the
optimization and transfer results of Sections~\ref{sec:E} and~\ref{sec:F}.
The discrepancy gap in this section is always a gap for primitive divisors
on the actual quotient germs.  No such gap is imposed on the auxiliary
central cones.

The local phase assertion of \cite[Theorem~3.1]{OpenAILocalPhase2026}
constructs centred divisorial valuations with bounded homogeneous discrepancy
and cyclic-character congruences, after passage to a subsequence, under its
hypotheses on Gorenstein klt germs of fixed dimension $n\ge2$.
It does not guarantee that these valuations lie on the prescribed $h$-lc
slice $A(T)=T(h)=1$ below, so it does not replace the pinning obstruction.

\begin{theorem}[Pinning obstruction]\label{thm:G-pinning}
Fix an integer $n>0$ and $\epsilon>0$.  There is no sequence of the following
data satisfying all the conditions below.  For each index $i$, let
$x_i\in S_i$ be an $n$-dimensional complex klt germ with an action of a
finite cyclic group $\Gamma_i$ fixing $x_i$, such that the quotient map is
quasi-\'etale and $S_i/\Gamma_i$ is $\epsilon$-lc.  Suppose $K_{S_i}$ has a
nowhere-vanishing local eigen-generator $\omega_i$, and let $h_i$ be a
nonzero regular eigenfunction with the same character $\chi_i$ as
$\omega_i$.  The character need not be faithful.  Assume that
$(S_i,\divisor(h_i))$ is lc and admits a centred quasi-monomial valuation
$T_i$ with
\[
 A_{S_i}(T_i)=T_i(h_i)=1.
\]
The excluded additional condition is: for every fixed integer $m>0$,
there is an index $i(m)$ such that, for all $i\ge i(m)$ and every regular
semi-invariant $s$ of character $\chi_i^m$,
\begin{equation}\label{eq:G-input}
                         T_i(s)\ge m.
\end{equation}
\end{theorem}

All constants in the proof are uniform along the sequence under
consideration; they may depend on the dimension and on the finitely many
previous stages of the construction.  An assertion made eventually also
holds after any further subsequence.  We suppress $i$ whenever this does
not obscure a limit.  We may successively pass to subsequences, since the
number of stages and of subsidiary row additions will be bounded.

\subsection{Costs and the induction data}

An invariant rational function $Y$ has \emph{cost at most $m$} if
$h^mY$ is regular, where $m$ is a nonnegative integer.  We increase costs
to positive integers and to common denominators when convenient.  Costs
are subadditive under multiplication.  The same terminology applies on
a graded cone, with $h$ replaced by its initial symbol.  Every homogeneous
ratio on such a cone of cost $m$ lifts to an actual invariant ratio of the
same cost: lift its regular numerator, and project to the character
$\chi^m$ by linear reductivity of the finite group.

We keep the test functions and accumulated variables on the original germ,
while replacing the working valuation and its cone.  Algebraic independence
of the accumulated symbols is the quantity that will force termination.
Their weights may nevertheless become rationally proportional: suitable
ratios of monomials then have weight zero and retain information in the
angular field.  Pinning the weights to a line therefore need not reduce
the transcendence degree of the accumulated symbols.

The following data are retained at the beginning of each stage.
\begin{enumerate}[label=\textup{(I\arabic*)}]
\item\label{it:G-tests}
There are finitely many actual invariant test ratios $f_j$, each of cost
$O(L_-)$, where initially $L_-=1$.  The common actual slice
\begin{equation}\label{eq:G-slice}
 \mathcal Q=\{q:A_S(q)=q(h)=1,\ q(f_j)\ge0\text{ for every }j\}
\end{equation}
contains a centred valuation, and every test has value exactly zero
throughout this slice.

\item\label{it:G-pairs}
There are accumulated pairs $Y_j^+,Y_j^-$ of invariant ratios of cost
$O(L_-)$, with opposite values on $\mathcal Q$.  We write
$y_j(q)=q(Y_j^+)=-q(Y_j^-)$.  The products $Y_j^+Y_j^-$ occur among the
exact tests.  Consequently the equality of the pair values also holds
on every later common cone slice, by Proposition~\ref{prop:F-bridge}.

\item\label{it:G-anchor}
Either all $y_j$ vanish on $\mathcal Q$, or their vector lies on a single
rational line.  In the latter case one accumulated variable, indexed by
$0$, is an \emph{anchor}.  It was chosen at an unbounded scale $L_1$;
its pair has cost $O(L_1)$, and $L_1\le L_-$.  For each other accumulated
variable there are integers $t_j,d_j$, with $d_j>0$ bounded, such that
\begin{equation}\label{eq:G-proportions}
 y_j=(t_j/d_j)y_0,
 \qquad |t_j|=O(L_-/L_1).
\end{equation}
These equations include the point where all values are zero and are
imposed by tests of both signs.  In the all-zero case there is no anchor.

\item\label{it:G-baseline}
The optimization height $\mathfrak B$ is rational and positive on
$\mathcal Q$, and in the normalization $q(h)=1$ satisfies
\begin{equation}\label{eq:G-baseline-bounds}
                  \rho\le\mathfrak B(q)\le C
\end{equation}
for fixed positive constants.  Initially $\mathfrak B(q)=q(h)$.
It can change once, when an anchor is first chosen, to
\begin{equation}\label{eq:G-baseline}
 \mathfrak B(q)=q(h)-\mu q(Y_0^{\mathrm{sign}})/L_1,
 \qquad \mu\ge0,
\end{equation}
where either member of the anchor pair can be used and $\mu$ is bounded.
This is an allowed height in Theorem~\ref{thm:E-exact-minimizer}.  On a common
central slice the same expression has a uniform upper bound after
normalizing $q'(h)=1$, because both anchor ratios have cost $O(L_1)$.

\item\label{it:G-working}
Take the exact $\mathfrak B$-minimizer and apply the rank refinement of
Proposition~\ref{prop:F-refinement}, choosing the approximation error to
tend to zero.  This produces an equivariant centred working valuation
$w$, with finitely generated normal graded algebra, whose values on all
regular functions are multiplicatively $1+o(1)$ times those of the exact
minimizer.  Its $\mathfrak B$-normalized volume is within $1+o(1)$ of the
minimum.  Its cone satisfies the balance estimate
\eqref{eq:F-balance} and the order condition \eqref{eq:C-alpha}, with error $\eta_i\to0$, on the angular tests
$\mathcal L$ specified by the current initial tests.  The symbols of all
accumulated $Y_j^+$ are algebraically independent.  If there is an anchor,
\begin{equation}\label{eq:G-anchor-size}
                        |w(Y_0^+)|\ge cL_1
\end{equation}
for a fixed $c>0$.
\end{enumerate}

The balance hypothesis needed for the refinement in
\ref{it:G-working} follows from the preceding height bounds.
Let $v$ be the exact $\mathfrak B$-minimizer, normalized by $v(h)=1$.
For every common central $q^0$ with $q^0(h)=1$,
Corollary~\ref{cor:E-domination} gives
\[
 q^0(s)\le 2^n\frac{\mathfrak B^0(q^0)}{\mathfrak B(v)}v^0(s)
          \le\frac{2^nC}{\rho}v^0(s)
 \qquad(0\ne s\text{ regular and homogeneous}).
\]
Here only the upper bound for $\mathfrak B^0$ is used on the central
slice; the denominator is bounded below at the actual minimizer.
Proposition~\ref{prop:F-refinement} then preserves balance with a fixed
constant.

Initially there are no tests and no pairs, so these conditions follow
from the centred point $T$ and Sections~\ref{sec:E}--\ref{sec:F}.
Corollary~\ref{cor:E-height-near-minimizers} applies to
$\mathfrak B$.  Indeed, normalize two valuations by $\mathfrak B=1$.
The radially scaled version of Lemma~\ref{lem:E-convexity} fixes
the averaged value of $h=A$ and gives exact interpolation of the opposite
anchor pair.  In particular its allowed height is at most $1$.
Its volume is therefore at least the minimum in height-one normalization.
The same gap argument gives multiplicative comparison on all functions.
The bounds \eqref{eq:G-baseline-bounds} then convert that comparison back
to $h=1$.  Applied to numerators and $h$-powers, it controls ratio values
up to the corresponding cost times the relative error.

Let $M,\mathbf u,\sigma$ be the invariant character lattice, the weight of
$h$, and the cone from Proposition~\ref{prop:D-cone-dictionary}, for the
working algebra.  Equip $M_\R$ with
\begin{equation}\label{eq:G-width}
 \|k\|=\max_{\substack{y\in\sigma\\\langle y,\mathbf u\rangle=1}}
                         |\langle y,k\rangle|.
\end{equation}
Fine weights evaluate injectively under $w$.  Thus the exact equations
in \eqref{eq:G-proportions} imply the same equations for the weights of
the symbols.  In the line case write
$\kappa_0=\operatorname{wt}(\operatorname{in}_wY_0^+)$.  Opposition on pure
common tests and the paired costs give $\|\kappa_0\|=O(L_1)$, whereas
\eqref{eq:G-anchor-size} gives the reverse bound.  Hence
\begin{equation}\label{eq:G-anchor-width}
                         \|\kappa_0\|\asymp L_1.
\end{equation}

After a stage with $L_-\to\infty$, we will additionally retain:
\begin{enumerate}[label=\textup{(H\arabic*)}]
\item\label{it:G-lattice}
Every lattice vector of width $O(L_-)$ belongs to $\Q\kappa_0$ if an
anchor is present, and is zero otherwise.  In the line case the index of
$\Z\kappa_0$ in $M\cap\Q\kappa_0$ is bounded.
\item\label{it:G-field}
The angular tests in $\mathcal L$, namely $a=0$ with nonnegative orders
on all initial tests, restrict trivially to the field generated by
$S(L_-D)$.
\end{enumerate}
Here $a,D,S(\cdot)$ have the angular meanings of
Sections~\ref{sec:C}--\ref{sec:D}.  No assertion
\ref{it:G-lattice} or~\ref{it:G-field} is required while $L_-$ remains
bounded, and there is then no anchor.  These two assertions will be proved
for the next stage after the high-scale pin, without being used to
construct its rank refinement.

\subsection{The next character block and the terminal case}

Suppose first that an anchor is present and $M$ has no direction
transverse to its line.  Write $\kappa_0=\ell k_0$, where $k_0$ is
primitive in $M$; assertion~\ref{it:G-lattice} bounds $\ell$.  The
restriction of $w$ to the actual quotient field has rational rank one
and is Abhyankar.  To justify that this is the actual value group, take
invariants in every level of the finite-group-stable filtration:
exactness of invariants identifies the invariant graded algebra with the
graded algebra of the quotient filtration.  Its homogeneous fraction
weights generate $M$, and evaluation is injective.  Thus the quotient
value group is generated by
\[
                         \gamma=|w(Y_0^+)|/\ell.
\]
Abhyankar equality is preserved under the finite field extension, so this
rank-one valuation is divisorial up to scaling.  Its primitive
normalization is $w|_{\Frac(S/\Gamma)}/\gamma$.  Quasi-\'etale log pullback
and $A_S(w)=1$ give its discrepancy as $1/\gamma$.  By
\eqref{eq:G-anchor-size} and $L_1\to\infty$, this tends to zero,
contradicting $\epsilon$-lc of the actual quotient.  Rank zero of $M$ is
impossible, since a finite extension of the quotient field cannot support
a nontrivial valuation with trivial restriction.  This disposes of the
terminal possibilities.

Otherwise choose a block of lattice directions as follows.  With an
anchor use the quotient norm modulo $\R\kappa_0$ on the projected
lattice; without one use \eqref{eq:G-width}.  After a previous high stage
its shortest nonzero length divided by $L_-$ tends to infinity.  For if a
projected vector had length $O(L_-)$, lift it and subtract the nearest
multiple of the primitive line vector.  Its remaining line contribution
is $O(L_1)$ by \eqref{eq:G-anchor-width} and the bounded index, producing
a transverse lattice vector of width $O(L_-)$, contrary to
assertion~\ref{it:G-lattice}.

Before the first high stage, if some successive minima remain bounded,
choose a nonempty maximal bounded block of independent directions,
of width $O(L)$ with $1\le L=O(1)$.  This is a \emph{bounded stage}.
In every other case choose an integral scale $L\to\infty$ comparable to
the first successive minimum, and take the whole first block of
successive minima comparable to $L$.  After a subsequence, the next
successive minimum, if any, divided by $L$ tends to infinity.  These are
\emph{high stages}.  Lift the selected directions $k_j$ to $M$ by rounding
along the anchor line when necessary.  They have width $O(L)$, and
$L_-=O(L)$.

In a high stage, let $s$ denote the block size plus one when an old anchor
is present, and just the block size otherwise.  Put
\begin{equation}\label{eq:G-Rstar}
 R_* = L/L_1\quad\text{with an old anchor},
 \qquad R_*=1\quad\text{without one}.
\end{equation}
For every fixed $N\ge1$, eventually
\begin{equation}\label{eq:G-lattice-count}
 \#\{k\in M:\|k\|\le CNL\}\le C'R_*N^s,
\end{equation}
with $C'$ independent of $N$.  Indeed, the gap to the next successive
minimum confines these vectors to the first-block span modulo the line.
Distinct projected lattice points have separation at least $cL$, so a
norm-ball packing in its bounded dimension gives $O(N^{s-1})$ points
when there is an anchor, and $O(N^s)$ otherwise.  Above each projected
point, the number of lifts in the ball is $O(NL/L_1)$, since the lattice
step on the anchor line has width comparable to $L_1$.  This proves
\eqref{eq:G-lattice-count}.

\subsection{Neutral tests and simultaneous non-cancellation}

Apply Lemma~\ref{lem:C-saturation-count} at scale $L$.
Its test-cost hypothesis uses $L_-=O(L)$; its additional field hypothesis,
when $L_-\to\infty$, is precisely assertion~\ref{it:G-field}.
Write its outputs as $K_*,p,v_i,B$.  If the angular field is the constant
field, set $p=0$ and $v_i=1$.  Increase the fixed constant $B$ if needed.
Lift a basis of $S(BLD)$ to actual ratios $G_a$ of cost $O(L)$ with the
prescribed neutral initial symbols.  This basis is finite for each germ,
but its cardinality may grow along the sequence; the cost constant is
uniform.  We include $1$ in the basis.  We may
also arrange that $p$ basis members have algebraically independent
initial symbols: the cutoff plateau has an endpoint $\beta_1\ge c/L$,
and its bounded birational system in $D_W/\beta_1$ pulls back into
$S(BLD)$ and generates $K_*$ as a field.  A transcendence basis can be
chosen from a field-generating set and extended to a vector-space basis.

For each selected direction $k_j$, Corollary~\ref{cor:D-regularization}
supplies homogeneous invariant ratios of weights $\pm q_jk_j$, with
bounded positive integers $q_j$, and cost $O(L)$.  Align the two
multipliers by taking bounded powers and lift the ratios to pairs
$Y_j^\pm$ on the actual germ.  Its small-height hypothesis holds: on a
primitive angular test with $a=0$ and $LD\to0$, an old test of cost
$O(L_-)$ has negative order bounded in absolute value by $O(L_-D)=o(1)$.
Its order is integral, hence is nonnegative eventually.  Such a test
therefore belongs to $\mathcal L$, where the required
\eqref{eq:C-alpha} holds.

Add the $G_a$ and the new products $Y_j^+Y_j^-$ as tests, before imposing
any individual new pair or proportionality constraint.  Call the result
the \emph{pre-pinning slice}.  All these tests have value zero at $w$,
so this slice contains $w$.  On the working cone, the old tests together
with the $G_a$ tests force trivial restriction on $K_*$ by saturation:
nonnegativity on a vector-space basis implies nonnegativity on its whole
linear system.  The assertion also holds for pure tests, which are
trivial on the angular field.  Lemma~\ref{lem:D-lc-tests} supplies the
homogeneous Gauss description and the finite-cover lifts for all other
common lc tests.  Rational monomial tests suffice, and the equalities
extend to their real subcones.  The $G_a$ and the new pair products have
initials in $K_*$, since their neutral costs are $O(L)$.  The bridge
therefore gives value at most zero on the actual pre-pinning slice.
Their test inequalities give the reverse sign, proving exactness and
opposition there.

\begin{lemma}[Non-cancellation in every degree]\label{lem:G-noncancellation}
At the preceding stage, consider any finite linear combination of
monomials in the accumulated and newly added $Y_j^+$ and the current
$G_a$, with nonnegative exponents.  Suppose its least-$w$ group has a
nonzero sum of initial symbols.  If all its terms have the same value at
an actual point of the pre-pinning slice, their sum has that value.
This assertion holds eventually simultaneously in all polynomial degrees.
Moreover, at a later working valuation on this slice where these terms
are tied, the same non-cancellation assertion transfers to its cone at
every common test where the terms remain tied and the pair values are
opposite.
\end{lemma}

\begin{proof}
Write $A_j=\operatorname{in}_wY_j^+$ and
$B_j=\operatorname{in}_wY_j^-$.  The pair products $A_jB_j$ lie in $K_*$.
With an old anchor, the finitely many old relations
\[
 \operatorname{wt}(A_j)=(t_j/d_j)\kappa_0
\]
give neutral rational functions $A_j^{d_j}A_0^{-t_j}$.  Each is a quotient
of two neutral functions of cost $O(L)$.  More explicitly, for $t_j\ge0$
write it as
\[
 \frac{A_j^{d_j}B_0^{t_j}}{(A_0B_0)^{t_j}};
\]
for $t_j<0$ its numerator is simply $A_j^{d_j}A_0^{|t_j|}$.  The required
costs are bounded by a fixed multiple of
$d_jL_-+|t_j|L_1=O(L_-)=O(L)$.  There are at most $n$ accumulated
variables, and their denominators $d_j$ are bounded.  Thus one fixed
constant works for this entire finite collection.  The inclusion
$S(CLD)\subset K_*$ puts their numerators and denominators in $K_*$ at
one eventual index.  Without an anchor all the old $A_j$ themselves are
neutral and of bounded cost, and belong to $K_*$ instead.

The new weights are independent modulo the old line, or modulo zero.
Consequently a monomial relation of total weight zero has zero net
exponent in each new variable.  Clearing the bounded old denominators
expresses a power of every remaining weight-kernel monomial in the
finite neutral generators just described.  Every such monomial is
therefore algebraic over $K_*$.  This is a simultaneous algebraic
statement: once the finite generators have entered $K_*$, there is no
further eventual index depending on the exponents.

Choose a term in the least-$w$ group, and multiply the whole sum by the
opposite factors $Y_j^-$ corresponding to its $Y_j^+$ factors.  Its
$G_a$ factors already have value zero.  The complemented sum has a
nonzero $w$-initial of weight zero.  Each term of that initial is
algebraic over $K_*$ by the preceding kernel argument, so the initial
itself is algebraic over $K_*$.  Every common cone test is trivial on
this element.  Here we use the elementary fact that a valuation trivial
on a field is trivial on its algebraic extension inside a valued
field: a polynomial equation with nonzero constant term would otherwise
have a unique term of smallest value.

The complemented sum is an invariant ratio with a regular numerator
after multiplication by some power of $h$.  Its cost is allowed to
depend on the chosen polynomial.  Proposition~\ref{prop:F-bridge} has
no uniform cost restriction, and yields an actual upper bound zero on
the pre-pinning slice.  At a point where the original terms tie,
opposition and the zero values of $G_a$ give value zero for every
complemented term.  The valuation inequality gives the lower bound
zero.  The complemented sum thus has exact value zero, which proves
non-cancellation.

At a later working valuation $w'$ where the terms tie, this establishes
exact value zero before any new bridge is invoked.  Its initial symbol
there is nonzero.  Apply the bridge in the other direction, using only
tests exact at $w'$.  At a common test on the new cone where the terms
remain tied, the transferred upper bound and the same termwise lower
bound again give equality.  This proves the onward assertion.
\end{proof}

The enlarged collection of positive symbols is algebraically independent
at $w$: the new weights are independent modulo the old line, and the
old symbols were independent.  It remains independent at every later
working valuation on the pre-pinning slice.  Indeed, a proposed
polynomial relation among the later symbols can be grouped by their
later grading.  Each nonempty group has a nonzero least-part initial at
$w$, and Lemma~\ref{lem:G-noncancellation} forbids cancellation of that
group at the later valuation.

At a bounded stage, add both new members $Y_j^\pm$ individually as tests.
Their values must then be zero.  The original centred point $T$ remains
in the slice by \eqref{eq:G-input}, because every cost used so far is
bounded.  This invocation is uniform over the finite interval of
possible integral costs, and applies to every numerator in each of the
corresponding semi-invariant spaces.  Thus the bounded stage is complete:
replace $L_-$ by $L$ and optimize and refine again.  No assertion of
\eqref{eq:G-input} at an unbounded cost has been used.

\subsection{The count at a high scale}

The neutral tests now prevent cancellation and preserve the independent
symbols when the working valuation changes.  At an unbounded scale,
the input saturation condition does not directly control all newly
chosen costs.  We instead bound the number of symbols available at that
scale.  The high-scale pin will force enough exact relations to avoid
violating this bound.

We now work at a high stage.  For each fixed integer $N\ge1$, each integer
$0\le t\le CNL$, and the character $\chi^t$, the span of regular
homogeneous symbols of that character of grading value at most $CNL$
has dimension at most
\begin{equation}\label{eq:G-count}
                            C'R_*v_iN^{p+s}.
\end{equation}
The constant $C'$ is independent of $N$ and $t$.  To see this, divide a
symbol by $h^t$.  It becomes an invariant homogeneous fraction with
weight in $M$.  Balance \eqref{eq:F-balance}, tested on rational pure
rays, bounds its width by $O(NL)$.  Those rays are common tests because
all current test weights are zero.  For each weight,
Corollary~\ref{cor:D-character-count} bounds its dimension by
$C''v_iN^p$.  Sum over \eqref{eq:G-lattice-count}.

The same upper bound, with only a fixed change of constant, holds on the
cone of any other actual equivariant centred valuation $v$ with
\begin{equation}\label{eq:G-domination}
                           v\ge c w
\end{equation}
on all regular functions, in $h=1$ normalization, for fixed $c>0$.
Indeed the order ideals satisfy
$\mathfrak a_{a/c}(w)\subseteq\mathfrak a_a(v)$.  The resulting surjection
of finite colength quotients is equivariant.  Taking each finite-group
eigenspace is exact, so the cumulative dimensions of its associated
graded pieces have the asserted comparison.  The strict or weak endpoint
of a grading cutoff can be handled by slightly increasing its fixed
constant.  In particular we transfer \eqref{eq:G-count} from the original
working cone; we never replace it by a count with a constant depending
on a later tilt.

We record the elementary product estimate used below.
\begin{lemma}[Uniform product growth]\label{lem:G-product-growth}
Let $V$ be a finite-dimensional subspace of a function field over an
algebraically closed field, containing $1$.  Put $g=\dim V$ and
$a=\trdeg k(V)\le n$.  If $V^{(N)}$ is the span of $N$-fold products of
members of $V$, there is $c_n>0$, depending only on $n$, such that
\[
                       \dim V^{(N)}\ge c_n gN^a
                       \qquad(N\ge1).
\]
\end{lemma}
\begin{proof}
The rational map defined by $V$ has a nondegenerate integral projective
image of dimension $a$ in $\mathbf P^{g-1}$.  Its degree-$N$ homogeneous
coordinate space is $V^{(N)}$.  Successively cut by general hyperplanes.
In positive dimension the reduced hyperplane section is nondegenerate
in its hyperplane.  One way to see the required nondegeneracy is to pass
to the normalization: a general hyperplane divisor is reduced in
codimension one, so another linear form vanishing on its reduced support
has a regular ratio with the cutting form, and the ratio is constant.
The reduced section is integral until the last curve section by Bertini;
the final reduced set of points spans its hyperplane and has Hilbert
function at least its dimension in degree one in every higher degree.
For the homogeneous coordinate ring, multiplication by a nonzero
cutting linear form is injective, and the quotient surjects onto the
coordinate ring of the reduced cut.  Thus successive differences of
Hilbert functions dominate the Hilbert functions of these cuts.
Summing these inequalities gives, for $a>0$, the sufficient bound
\[
        \dim V^{(N)}\ge(g-a)\binom{N+a-1}{a}.
\]
Since $g\ge a+1$, one has $g-a\ge g/(a+1)$, giving the asserted constant.
For $a=0$ the image is a point, $g=1$, and the assertion is immediate.
\end{proof}

\subsection{The high-scale pin}

At a high stage the new costs grow with $L$, so the original input
\eqref{eq:G-input} no longer supplies a centred point after individual
values are forced to zero.  We instead constrain their proportions while
retaining a centred point close enough to the current minimum.

\begin{proposition}[High-scale pin]\label{prop:G-high-pin}
Consider a high stage with pre-pinning slice, working valuation $w$,
baseline height $\mathfrak B$, scale $L$, and $s$ selected values as above.
After passage to a subsequence, finitely many invariant tests of cost
$O(L)$ can be added with the following properties.
\begin{enumerate}[label=\textup{(\roman*)}]
\item The resulting slice contains a centred valuation.  The added tests
are exact on that slice, all old tests and proportionalities remain in
force, and the accumulated positive symbols remain algebraically
independent at the exact minimizer and working valuation in
\textup{(iii)} below.
\item The selected values either vanish or lie on a single rational
line.  In the line case their proportions have the bounds
\eqref{eq:G-proportions}, with $L_-=L$.  An old anchor remains the anchor;
if a new anchor is chosen, its scale is $L_1=L$.
\item The baseline height is retained unless an anchor is first chosen,
in which case it may change to the rational height
\eqref{eq:G-baseline}.  It satisfies \eqref{eq:G-baseline-bounds} on the
new actual slice.  At its exact minimizer $v'$, and at an arbitrarily close
rank refinement $w'$, the anchor, when present, satisfies
$|v'(Y_0^+)|,|w'(Y_0^+)|\ge cL_1$.  Both valuations dominate a fixed positive
multiple of $w$ on every regular function.  In particular,
\[
                         w'\ge c w.
\]
\end{enumerate}
Consequently \ref{it:G-tests}--\ref{it:G-working} hold with $L_-=L$.
The next-scale assertions \ref{it:G-lattice}--\ref{it:G-field} are not
assumed in constructing the pin or its rank refinement.
\end{proposition}

Domination preserves the counting bound on the original working cone.
The rational proportions permit equations of cost $O(L)$, and the
nonzero anchor size will detect a small primitive quotient discrepancy
if only its line remains.  We first construct this pin, then prove that
it also gives the next-scale field and lattice assertions.

\subsection{Finding all necessary relations}\label{subsec:G-pin-rows}

We begin the proof of Proposition~\ref{prop:G-high-pin}.  Within this
stage, progress is measured by the rank of the integral relations imposed
on the $s$ selected values.  This row rank is separate from the
transcendence degree used to terminate the outer induction.

For the $s$ selected variables put $L_j=L_1$ for an old anchor, and
$L_j=L$ for every new variable, and set
\[
                     x(q)=(y_j(q)/L_j)_{j=1}^s.
\]
When there is an anchor we include its index $0$ in this notation.
Opposition and the paired costs bound all coordinates uniformly on the
actual pre-pinning slice.  Lemma~\ref{lem:E-convexity} and Corollary~\ref{cor:E-opposite-ratios}
give exact convex interpolation of the coordinates, so their image is compact and
convex.  Its points can be realized by centred valuations whenever the
interpolation assigns positive weight to a centred input.

We successively impose integral rows
\begin{equation}\label{eq:G-row}
 \sum_j k_jy_j=0,
 \qquad \widehat k=(k_jL_j/L)_j,
 \qquad \|\widehat k\|_{\mathrm{Eucl}}=O(1).
\end{equation}
A row is imposed by testing two opposite monomials.  For its positive
sign use
\[
 M_k=\prod_{k_j\ge0}(Y_j^+)^{k_j}
       \prod_{k_j<0}(Y_j^-)^{-k_j};
\]
exchange the signs to obtain the other monomial.
Each has cost $O(\sum_j|k_j|L_j)=O(L)$, and opposition makes the two
inequalities equivalent to the row equality.

Let $\mathcal R_i\subset\R^s$ denote the span of the normalized rows
already imposed.  We retain a comparator $v_0$ on the current slice
whose $\mathfrak B$-normalized volume is within $1+o(1)$ of the minimum
on the actual slice at the start of this high stage.  Initially use $w$.
The gap estimate then gives $v_0=(1+o(1))w$ on all regular functions.
With an old anchor, $|y_0(v_0)|/L_1$ stays bounded away from zero.
Pass to subsequences so the row spaces converge and the limits of the
chosen rows are independent.

For clarity, independence of their transverse integer rows is equivalent
to independence of these normalized row limits in the present situation.
The non-anchor coefficients $k_j$ are bounded integers, hence become
fixed on a subsequence.  If a rational combination of their transverse
rows vanished, the same combination of full rows would be supported
only on the anchor.  Its exact value at $v_0$ must vanish.  Since the
anchor coordinate is nonzero, its anchor coefficient vanishes as well.
Thus any dependence in the transverse rows is a dependence of the full
rows.  Conversely independent transverse rows keep the full normalized
rows independent in the limit.  With no anchor, all coefficients are
bounded integers and the statement is immediate.

Let $K_i$ be the projection image $\{x(q)\}$ of the current constrained
slice, and let
\[
                   C_i=\operatorname{conv}(\{0\}\cup K_i)
                       =[0,1]K_i.
\]
After a subsequence these bounded compact convex sets converge in
Hausdorff distance.  Their limit lies in
$\mathcal R_\infty^\perp=\lim\mathcal R_i^\perp$.
We first handle the case
\begin{equation}\label{eq:G-thin}
                      \operatorname{span}(\lim C_i)
                      \ne\mathcal R_\infty^\perp.
\end{equation}
Choose a unit vector $u\in\mathcal R_\infty^\perp$ perpendicular to the
left-hand span, and rational vectors $u_i\to u$.  Then
\begin{equation}\label{eq:G-projection-small}
                    u_i\cdot x(q)=o(1)
\end{equation}
uniformly on the current slice.  Fix a sufficiently small
$\delta>0$, with its size chosen below, and consider the rational height
\begin{equation}\label{eq:G-thin-height}
 b(q)=\delta\mathfrak B(q)-\sum_j u_{ij}y_j(q)/L_j.
\end{equation}
We take the scales integral, so all coefficients can be rational.
According to the sign of $u_{ij}$, use $Y_j^+$ or $Y_j^-$, and absorb its
$h$-denominator into the coefficient of $q(h)$.  This puts $b$ in the
allowed form with nonnegative coefficients on the subtracted regular
numerator values.  Equations \eqref{eq:G-baseline-bounds} and
\eqref{eq:G-projection-small} show that it is positive and equals
$(1+o(1))\delta\mathfrak B$ uniformly on this slice.

Let $v$ be its exact minimizer.  Comparison with $v_0$ shows that the
$\mathfrak B$-normalized volume of $v$ is within $1+o(1)$ of the original
stage minimum.  Consequently
\begin{equation}\label{eq:G-thin-close}
                           v=(1+o(1))w
\end{equation}
on all regular functions, and $v$ is exact on all present tests.
On its cone every nonzero pure test, and every Gauss test over a primitive
angular divisor with $a=0$ and $LD\to0$, is common.  Indeed all test
weights are zero, their costs are $O(L)$, and the integral angular orders
cannot be negative of size $o(1)$.  All pair equalities transfer to this
cone.  Positivity of the induced height, supplied by
Theorem~\ref{thm:E-exact-minimizer}, and the uniform central upper bound for $\mathfrak B$ imply
\begin{equation}\label{eq:G-central-tilt}
                \sum_j u_{ij}y_j'/L_j\le C\delta q'(h)
\end{equation}
on these tests.

Apply Corollary~\ref{cor:D-tilted-characters}.  For a fixed
$M_0\gg1/\delta$, choose integers
\begin{equation}\label{eq:G-rounding}
                r_j=M_0u_{ij}L/L_j+O(1).
\end{equation}
It gives a homogeneous invariant ratio $Z$ on the cone of $v$, of weight
\begin{equation}\label{eq:G-cheap-weight}
 -q\sum_jr_j\operatorname{wt}(\operatorname{in}_vY_j^+),
\end{equation}
where $q$ is a bounded positive integer, and of cost at most
\begin{equation}\label{eq:G-cheap-cost}
                              CqM_0\delta L.
\end{equation}
The constant $C$ is fixed before $\delta$ and $M_0$ are chosen.  The bound
on $q$ may depend on these fixed choices.  This distinction follows from
the proof of the tilted-character result: its height correction has
size $O(M_0\delta L)$; the $O(L)$ rounding error and the fixed constant in
\eqref{eq:G-central-tilt} are absorbed by first taking
$M_0\gg1/\delta$.  Only the later regularizing multiplier need depend
on these parameters.  The absolute grading value of $Z$ has the same
bound.  In fact \eqref{eq:G-cheap-weight} evaluates to
\[
 -qM_0L\,u_i\cdot x(v)+O(qL),
\]
which is bounded by \eqref{eq:G-cheap-cost} eventually, by
\eqref{eq:G-projection-small} and the chosen size of $M_0$.
No assumption \eqref{eq:C-alpha} on this newly tilted cone is used.

\begin{claim}\label{clm:G-new-row}
After a subsequence there is an additional row of the form
\eqref{eq:G-row}, exact at $v$, whose transverse part is independent of
the previously imposed transverse rows.
\end{claim}
\begin{proof}
Suppose otherwise.  Then, for each fixed bound on normalized coefficient
size, every exact integral tie at $v$ eventually belongs to the old row
span.  To verify this assertion, a tie with transverse part outside the
old span would give the desired row after fixing its bounded integer
transverse coefficients.  If its transverse part lies in the old span,
subtract the corresponding rational combination of the fixed old
transverse rows.  The remainder is a pure-anchor tie.  It is zero
because \eqref{eq:G-thin-close} preserves the nonzero anchor value at
$v$.  Without an anchor there is no remainder.  The claimed eventual
assertion for all ties with a given bound follows by choosing a failing
tie along a subsequence if necessary.

Let $\widetilde u_i$ be a unit vector in $\mathcal R_i^\perp$ tending to
$u$.  Choose a large fixed integer $N\gg qM_0$.  In the exponent box
\[
           0\le e_j\le\lfloor NL/L_j\rfloor,
           \qquad \widehat e=(e_jL_j/L)_j,
\]
there are at least $cR_*N^s$ integer points.  The range of
$\widetilde u_i\cdot\widehat e$ has length $O(N)$.  Partition this range
into intervals of width $qM_0/4$.  One slab contains at least
\begin{equation}\label{eq:G-slab-size}
                           cR_*N^{s-1}qM_0
\end{equation}
points.  For each exponent in this slab, consider on
$\operatorname{gr}_v$ its positive-variable monomial multiplied by $Z^l$,
where
\begin{equation}\label{eq:G-powers}
                     0\le l\le\left\lfloor
                                     N/(qM_0\delta)\right\rfloor,
\end{equation}
and by a basis of the products of $N$ of the $G_a$ whose $w$-initials
are independent.  Saturation and counting supply at least $v_iN^p$
such products.

Different powers $l$ have disjoint grading values in this construction.
A collision would give the exact integral tie
\[
                              e-e'-q(l-l')r.
\]
Its normalized size is bounded with $N,M_0,\delta,q$ fixed.  Under our
supposition it belongs to the old row span, so its normalized dot
product with $\widetilde u_i$ is zero.  The slab bounds the contribution
from $e-e'$ by $qM_0/4$.  But \eqref{eq:G-rounding} gives
\[
       \widehat r=M_0u_i+O(1),
       \qquad |\widetilde u_i\cdot\widehat r|\ge M_0/2
\]
eventually, after increasing the fixed $M_0$.  For $l\ne l'$ the other
contribution has magnitude at least $qM_0/2$, a contradiction.  It is
essential here that $\widetilde u_i$ is orthogonal to the entire old
normalized row span.

At any one power $l$, the slab size times $v_iN^p$ symbols are linearly
independent.  Before changing to $v$, the selected variable weights
separate the exponent vectors, and the chosen $G_a$ products have
independent initials at $w$.  Any possible relation after the change
can be grouped by its $v$-grading, and
Lemma~\ref{lem:G-noncancellation} rules out cancellation in each group.
Multiplication by the nonzero $Z^l$ preserves this independence.
Combining with disjoint grading values for different $l$ gives at least
\begin{equation}\label{eq:G-excess-count}
                     c\delta^{-1}R_*v_iN^{p+s}
\end{equation}
independent ratios.

Every ratio just counted has cost and absolute grading value at most
$C_0NL$, where $C_0$ is independent of $\delta,M_0,q$.  The slab
monomials cost $O(NL)$; the $G_a$ products do too; and
\eqref{eq:G-cheap-cost} multiplied by the exponent bound
\eqref{eq:G-powers} costs $O(NL)$.  A common $h$-denominator gives
regular numerators of one finite character and degree $O(NL)$.
By \eqref{eq:G-thin-close}, eventually $v\ge w/2$, so the transferred
upper bound \eqref{eq:G-count} has a constant independent of the tilt.
It contradicts \eqref{eq:G-excess-count} if $\delta$ is sufficiently
small.

The order of choices is as follows.  First fix the current stage
constants, the old rows, the saturation constant and the original
counting constant.  Then fix $\delta>0$ small enough for this last
contradiction.  Next fix $M_0$ sufficiently large compared with
$1/\delta$ and the rounding constants.  The character construction then
gives a bound $Q$ for $q$; pass to a subsequence fixing $q$, or choose
using $Q$.  Only now choose a sufficiently large fixed $N$, and finally
let $i$ be large enough for all of these fixed-parameter assertions.
Thus $\delta$ never tends to zero with $i$, and no eventual count is
used at a variable product degree.  This proves the claim.
\end{proof}

Impose the new row by both signs and retain $v$ as the new comparator.
It is centred and still within $1+o(1)$ of the original stage minimum.
The limits of the enlarged rows stay independent because their
transverse integer projections are fixed and independent.  Each time
\eqref{eq:G-thin} holds, Claim~\ref{clm:G-new-row} increases this rank.
There are at most $s$ such increases, and at most $s-1$ with a retained
nonzero anchor.  Therefore this procedure reaches the alternative
\begin{equation}\label{eq:G-full-span}
                    \operatorname{span}(\lim C_i)
                    =\mathcal R_\infty^\perp.
\end{equation}
This is strict finite progress within a stage.

\subsection{Rational directions and anchor size}\label{subsec:G-pin-directions}

Fix a small positive number $\lambda$.  Mix the comparator projection
with the full current projection set:
\[
                  K_i^\lambda=(1-\lambda)x(v_0)+\lambda K_i.
\]
These points have actual centred realizations $q$ with
$q\ge(1-\lambda)v_0$ on all regular functions, by Lemma~\ref{lem:E-convexity}.  Their baseline heights interpolate exactly because the
height depends only on $h$ and the anchor coordinate.  The bounds on
$\mathfrak B$ show, for sufficiently small fixed $\lambda$, that their
$\mathfrak B$-normalized volumes are within any prescribed small fixed
relative error of the original stage minimum.  For example domination
gives the volume factor $(1-\lambda)^{-n}$, and the height changes by
$O(\lambda)$ relative to its positive lower bound.

The radial hull $[0,1]K_i^\lambda$ still has a full-dimensional limit in
$\mathcal R_\infty^\perp$.  Indeed $K_i^\lambda$ contains $x(v_0)$,
and differences from that point span
$\lambda(K_i-x(v_0))$.  Its linear span is therefore the same as that of
$K_i$.  The same statement holds in the limit.  A compact convex set of
full linear span containing zero has nonempty relative interior, so by
Hausdorff convergence these radial hulls contain relative balls of
radius bounded below in $\mathcal R_i^\perp$, unless that space has
dimension zero.  If there is an old anchor, choose such a ball away from
anchor-coordinate zero: the comparator has nonzero anchor coordinate,
and a convex combination with an interior point gives an interior
point retaining that property.  Radial truncation is used only to choose
a direction; every nonzero direction chosen in it has an actual mixed
representative, with the centred domination just established.

With an old anchor, write an integral row matrix as $(b_i,A)$, where
$A$ is its transverse part.  After the subsequences already taken,
$A$ is a fixed integral matrix of full row rank, and the entries of
$b_i$ are $O(R_*)$.  The feasible ratios $z_j=y_j/y_0$ satisfy
\begin{equation}\label{eq:G-rational-affine}
                                Az=-b_i.
\end{equation}
Directions in the ball away from $x_0=0$ give a relative ball of feasible
ratios, in coordinates $z/R_*$.  We can choose a point there with bounded
common denominator and $z_j=O(R_*)$.  To make the denominator assertion
explicit, a fixed invertible minor of $A$ gives a particular solution of
\eqref{eq:G-rational-affine} with denominator dividing its determinant
and size $O(R_*)$.  The rational kernel of $A$ has a fixed lattice basis.
Round in a fixed sufficiently fine denominator refinement of that
lattice.  Since $R_*\ge1$, its mesh after division by $R_*$ is uniformly
small enough to enter the specified ball.  If the affine solution
space has dimension zero, the particular solution is already the
required exact ratio.  Thus the denominator is bounded and the numerator
bounds are $O(L/L_1)$.

Without an old anchor, the full row matrix is fixed integral.  If its
orthogonal space is nonzero, the ball contains a nonzero rational
direction of bounded complexity.  Choose a nonzero coordinate as the
new anchor and put $L_1=L$.  Choose the ball and direction so that this
coordinate has absolute value at least a fixed positive constant; its
actual mixed representative has that same lower bound or a larger one,
since the radial factor lies in $(0,1]$.  If the orthogonal space has
dimension zero, all selected values are already zero.

Impose the chosen proportions, or all zero values, by both-sign
monomial tests.  Their costs are $O(L)$: with an old anchor a row
$d_jy_j-t_jy_0=0$ has cost $O(d_jL+|t_j|L_1)$ with $d_j$ bounded and
$|t_j|=O(L/L_1)$; with a new anchor both coefficients are bounded.
The selected actual mixed representative survives all these equations,
so the new slice contains a centred point.  Old proportionalities are
retained by the old tests.  In the cases with no new anchor, the exact
baseline minimum on the pinned slice is as close to the preceding
working valuation $w$ on all functions as desired, by choosing
$\lambda$ sufficiently small and applying the gap estimate on the
original stage slice.  In particular an old anchor still has
$|y_0|\ge cL_1$, and domination by a fixed multiple of $w$ holds.

When a new anchor is chosen, a further step is needed: equations on
ratios alone do not prevent the minimizing vector from moving to zero.
The prior baseline is then $q(h)$.  Let $z$ denote the centred mixed
comparator on the pinned slice.  Choose the sign of the anchor so that
\[
 p_0(q)=q(Y_0^{\mathrm{sign}})/L_1,
 \qquad p_0(z)\ge c_1>0.
\]
Its ordinary volume is within a prescribed small fixed error of the
prior ordinary minimum $V_*$.  Let $M=\max p_0$ on this compact,
$h=1$-normalized slice.  Thus $M\ge c_1$ and $M=O(1)$.  For
$0\le\mu\le(1-\rho)/M$, consider
\begin{equation}\label{eq:G-anchor-height}
                              b_\mu(q)=q(h)-\mu p_0(q),
\end{equation}
where $\rho>0$ is a small fixed number.  These heights are bounded below
by $\rho$ and uniformly above.  The exact minimizing ray varies
continuously with $\mu$, by Theorem~\ref{thm:E-exact-minimizer}.

At the endpoint $\mu_*=(1-\rho)/M$ the minimizer $v_*$ has
$p_0(v_*)\ge M/2$ for small enough $\rho$.  Here are the estimates.
Interpolate a maximizing witness for $p_0$ with $z$, assigning weight
$\rho$ to $z$.  The interpolant $u$ is centred, has
$b_{\mu_*}(u)\le2\rho$, and \eqref{eq:E-volume-interpolation}, applied
with $f=h$ and both inputs normalized by $h=1$, gives
\[
                         \vol(u)\le\rho^{1-n}\vol(z).
\]
Minimality and $\vol(v_*)\ge V_*$ yield
\[
 b_{\mu_*}(v_*)
 \le2\rho^{1/n}\bigl(\vol(z)/V_*\bigr)^{1/n}<\tfrac12
\]
for sufficiently small fixed $\rho$.  Therefore
$p_0(v_*)>M/[2(1-\rho)]\ge M/2$.

If the minimizer at $\mu=0$ already has $p_0\ge c_1/4$, use $\mu=0$.
Otherwise continuity gives a nonempty open interval of parameters where
$c_1/4<p_0(v_\mu)<c_1/3$, and we choose a rational $\mu$ in that
interval.  Since $p_0(z)\ge c_1$, one has
$b_\mu(z)<b_\mu(v_\mu)$.  Tilted minimality consequently gives
\begin{equation}\label{eq:G-anchor-volume}
 \vol(v_\mu)
 \le\left(\frac{b_\mu(z)}{b_\mu(v_\mu)}\right)^n\vol(z)
 \le\vol(z).
\end{equation}
The same final bound is immediate at $\mu=0$.  Thus the new minimizer
still lies as near the original ordinary minimum as required, while
its anchor value is bounded below by $c_1L_1/4$.  The gap estimate gives
the required all-functions domination.  Adopt this rational height as
$\mathfrak B$ for all subsequent stages.  It has the form \eqref{eq:G-baseline} and bounds \eqref{eq:G-baseline-bounds}.

Finally take the rank refinement with error tending to zero, sufficiently
small to preserve the anchor lower bound.  This preserves the pinning
equalities, the exact pair and test identities, and the independent
symbols by Lemma~\ref{lem:G-noncancellation}.  It supplies balance and
\eqref{eq:C-alpha} and does not use the next-scale field and lattice
assertions.  Thus \ref{it:G-tests}--\ref{it:G-working} hold with $L_-=L$, proving
Proposition~\ref{prop:G-high-pin}.  It remains to verify the two
next-scale assertions.

\subsection{Field and lattice consequences of the pin}\label{subsec:G-pin-consequences}

\begin{lemma}\label{lem:G-pin-consequences}
The working valuation $w'$ furnished by Proposition~\ref{prop:G-high-pin}
satisfies \ref{it:G-lattice}--\ref{it:G-field} at the next scale $L_-=L$.
\end{lemma}
\begin{proof}
The proposition supplies the new exact tests, rational proportions and,
when present, the anchor lower bound.  It also gives $w'\ge c w$ on all
regular functions, so the transferred count \eqref{eq:G-count} remains
available with a fixed constant.

Use double bars for the initial symbols at $w'$.  Let $E$ be the field
generated by the $p$ chosen algebraically independent
$\overline{\overline{G_a}}$ and the $s$ selected variable symbols.
These have transcendence degree $p+s$, by
Lemma~\ref{lem:G-noncancellation}: they were independent together at $w$,
since the selected variable weights were independent and the chosen
$G_a$ were neutral.  Write $E^0$ for its subfield of weight zero.
Every common test on the new cone is trivial on $E^0$.  Indeed, the
both-sign pinning equations also hold on that cone.  If there is an
anchor, its nonzero value at $w'$ ensures that two selected monomials
with equal weight at $w'$ have equal value at every common test; in the
all-zero case this is immediate.  Lemma~\ref{lem:G-noncancellation} then
gives the predicted exact value for each nonzero homogeneous polynomial.
Ratios of homogeneous polynomials of equal weight generate $E^0$,
proving the assertion.

There is also a space $V_0$ of invariant ratios on this new cone, of
cost and absolute grading value $O(L)$, with
\begin{equation}\label{eq:G-large-space}
                       \dim V_0\ge cR_*v_i.
\end{equation}
Take the double-bar symbols of the whole $G_a$ basis, and, if an old
anchor was present, multiply by its powers with exponent from $0$ to
$\lfloor R_*\rfloor$.  Their costs remain $O(L)$.  They are linearly
independent at $w$ by their distinct anchor weights and the basis
property, and remain so by Lemma~\ref{lem:G-noncancellation}.

Choose an enlargement $V$ of $V_0$, with the same type of fixed cost and
grading bounds, containing $1$, the specified generators of $E$, and every
neutral function in $S(LD_{\mathrm{next}})$.  We may also include in $V$ a
nonzero homogeneous ratio of weight a bounded positive multiple of any
chosen new lattice vector $k$ of width $O(L)$.  For the latter use
Corollary~\ref{cor:D-regularization} at $w'$: the refined cone satisfies
\eqref{eq:C-alpha}, and tiny $LD_{\mathrm{next}}$ forces all new tests
by the same integrality and cost argument as before.  The space $V$
is spanned by homogeneous ratios.  Give all products of $N$ members of
$V$ one common denominator $h^t$, where $t=O(NL)$; their regular
numerators then have a single finite character and degree $O(NL)$.
The transferred upper bound \eqref{eq:G-count} applies.

If the field generated by $V$ had transcendence degree $a>p+s$,
Lemma~\ref{lem:G-product-growth} and \eqref{eq:G-large-space} would give
at least $cR_*v_iN^a$ independent products, against
$C'R_*v_iN^{p+s}$.  Choose one sufficiently large fixed $N$ and then let
$i$ be sufficiently large.  This contradiction shows that the field
generated by $V$ is algebraic over $E$.  Uniformity for the chosen $k$ follows by
applying this argument to any alleged sequence of failing choices.

A neutral element algebraic over $E$ is algebraic over $E^0$.  In fact,
clear denominators in an algebraic equation so that its coefficients
are polynomials in the homogeneous generators of $E$, take a nonzero
homogeneous part of the equation, and divide by a nonzero coefficient of
that weight.  All resulting coefficients have weight zero.  Common
tests are trivial on $E^0$, and hence on these algebraic elements.  The
Gauss description and finite-cover lifting translate this back to every
angular test in $\mathcal L$, proving assertion~\ref{it:G-field} for the
next scale.

The same algebraic equation argument shows that the weight of an
algebraic homogeneous element lies in the rational span of the weights
of $E$.  The latter is the anchor line or zero.  This proves the first
part of assertion~\ref{it:G-lattice}.  To bound its index in the line
case write $\kappa'_0=\ell k'_0$ with $k'_0$ primitive.  Its width is
$O(L_1)$ by the paired anchor costs.  If $\ell\to\infty$, apply the
bounded-multiple construction at the current scale $L$ to all $jk'_0$,
$1\le j\le\ell$.  Their widths are $O(L_1)$ and $L_1\le L$.
There are symbols of weights $q_jjk'_0$, with
$1\le q_j\le Q$ for fixed $Q$, and cost and absolute grading value
$O(L)$.  These bounds are uniform over all these choices: otherwise a
sequence of failing $j_i$ would contradict the arbitrary-sequence
version of Corollary~\ref{cor:D-regularization}.

Let $D_0$ be a fixed common multiple of the bounded denominators of the
finitely many pinned variable weights.  The subgroup generated by those
weights has primitive step $h_i\ge\ell/D_0$, and $h_i$ divides $\ell$.
Among the integers $q_jj$, an integer has at most $Q$ preimages, and each
residue modulo $h_i$ has at most $QD_0+1$ representatives in
$[1,Q\ell]$.  Hence at least
\[
                         \frac{\ell}{Q(QD_0+1)}
\]
distinct residues occur.  All weights of the original space $V_0$ in
\eqref{eq:G-large-space} lie in this subgroup.  With an old anchor its
basis consists of neutral symbols multiplied by powers of that anchor;
when the pin has just chosen a new anchor, $V_0$ consists only of neutral
symbols.
Multiplying this original $V_0$ by symbols representing distinct residues yields a
direct sum, since the weight supports are disjoint.  Its dimension is
an unbounded multiple of $R_*v_i$, while its costs and degrees are
$O(L)$.  The fixed-$N$ version of \eqref{eq:G-count} forbids this.
Thus $\ell$ is bounded, completing assertion~\ref{it:G-lattice}.

\end{proof}

\subsection{Termination}

\begin{proof}[Proof of Theorem~\ref{thm:G-pinning}]
Assume the excluded sequence exists.  The initial data satisfy
\ref{it:G-tests}--\ref{it:G-working}.  At every stage, either the sole
remaining anchor line gives the primitive-discrepancy contradiction
above, or a nonempty independent block of new directions can be chosen.
A bounded block is pinned directly by \eqref{eq:G-input}; a high block
is pinned by the finite row procedure, the rational-direction choice,
and, if required, the anchor-height adjustment.  The next-stage
invariants then hold by the field and lattice argument.  Each
nonterminal stage adds at least one new positive variable, and their
initial symbols remain algebraically independent at the next working
valuation.  A function field of dimension $n$ cannot contain more than
$n$ algebraically independent elements.  Thus the process cannot have
more than $n$ nonterminal stages, whereas all its terminal alternatives
have been excluded.  This contradiction proves the theorem.
\end{proof}

\section{From complement failure to the pinning input}\label{sec:H}

We prove Theorem~\ref{thm:birational} by translating a failure of bounded
klt complements into the sequence excluded by Theorem~\ref{thm:G-pinning}.
The construction uses a canonical cover when the birational morphism is an
isomorphism, and a signed anti-canonical torsor otherwise.  We then apply
Chen's equivalence and justify the passage between algebraically closed
fields of characteristic zero.

\subsection{A prime saturating the complete systems}

Fix a dimension $d>0$ and $0<\epsilon\le1$, and suppose the birational
complement assertion fails over $\C$.  Birkar's relative complement
theorem gives an integer $b=b(d)>0$ such that every example has an lc
$b$-complement near its specified point \cite[Theorem~1.8]{BirkarComplements}.
The hypotheses are precisely Fano type, the zero boundary, and nef
$-K_X$; in particular no generalized nef-part index is being imposed.
Choose a sequence with no klt $M_i$-complement, where $M_i>1$ tends to
infinity, $b\mid M_i$, and every fixed positive integer divides $M_i$
eventually.  For example, sufficiently large factorials have these
properties.  Failure of a uniform index allows this choice for each
prescribed $M_i$.

We may assume $-K_X$ is ample over $Z$.  On a variety of Fano type, the
relative Mori dream space consequences of the klt MMP imply that nef
$-K_X$ is semiample \cite{BCHM}.  Since $f$ is birational, this divisor
is relatively big.  Its ample model is thus a birational contraction
$g:X\to X^{\mathrm a}$ over $Z$.  A divisible multiple of $-K_X$ is the
pullback of its ample model divisor.  Pushing compatible canonical
divisors forward identifies the latter with the corresponding multiple
of $-K_{X^{\mathrm a}}$, so $K_X=g^*K_{X^{\mathrm a}}$ after compatible
choices.  Equivalently, a possible rational principal adjustment descends
because the function fields agree.  Thus the contraction is crepant,
$X^{\mathrm a}$ remains $\epsilon$-lc and of Fano type over $Z$, and
klt complements on it pull back to klt complements on $X$
\cite{ProkhorovShokurov}.  Applying the lc complement theorem on this
model and replacing $X$ by it preserves the failure under consideration.

Shrink $Z$ to an affine neighbourhood of the specified point.  For every
integer $k$, a section of $\mathcal O_X(-kK_X)$ is represented by a
rational function $s$ with
\begin{equation}\label{eq:H-section-divisor}
 \Delta_s=\divisor(s)-kK_X\ge0.
\end{equation}
The degree-$b$ complement is represented by $s_b$ and satisfies
$(X,\Delta_{s_b}/b)$ lc.  Here and below all section systems are complete
relative systems on this neighbourhood.

\begin{proposition}[Factorial saturation]\label{prop:H-saturation}
For each example there is a primitive divisorial order $w$ over $X$ whose
image centre on $Z$ contains the specified point and for which
\begin{equation}\label{eq:H-saturation}
 w(\Delta_s)\ge k A_X(w)\quad
       (k>0,\ k\mid M_i),\qquad
 w(\Delta_{s_b})=bA_X(w).
\end{equation}
After replacing the example by a geometric generic slice, one may assume
that the image centre is a closed point $z$, that
$1\le\dim X=\dim Z\le d$, and that all the inequalities hold for the
complete local systems at $z$.
\end{proposition}
\begin{proof}
The coherent sheaf $f_*\mathcal O_X(-M_iK_X)$ has finitely many
generators on the affine base.  Resolve their fractional ideal of ratios
to a fixed section, together with the relevant divisors.  On the resulting
smooth model $p:Y\to X$, their pulled-back divisor system has a common
fixed $\Q$-divisor $F$ and a basepoint-free residual system.  This
construction is valid even if $M_iK_X$ is not Cartier: each section divisor
is $\Q$-Cartier, their differences are principal, and principalizing the
fractional ratio ideal makes the moving differences Cartier.  The fixed
divisor has order equal to the minimum of the section orders at every
prime on the model.

The section $s_b^{M_i/b}$ belongs to this system.  Therefore
$F/M_i\le p^*(\Delta_{s_b}/b)$, and the fixed-order discrepancy
\[
 A_X(E)-\frac{\ord_E(F)}{M_i}
\]
is nonnegative for every prime $E$ over $X$.  Take the support on $Y$ to
be simple normal crossings.  A general residual member, with coefficient
$1/M_i<1$, is transverse to this support.  If every fixed-order
discrepancy were strictly positive over a neighbourhood of the specified
point, this general member would give a klt $M_i$-complement there.
The only possible zero discrepancies of the fixed log smooth data come
from strata containing a coefficient-one component.  There are finitely
many such components, and their images in $Z$ are closed because
$Y\to Z$ is proper.  Thus the absence of a klt complement on every
neighbourhood implies that one component $E$ with image containing the
specified point has
$\ord_E(F)=M_i A_X(E)$.  Set $w=\ord_E$.

For $k\mid M_i$ and any section $s$ in degree $k$, its power
$s^{M_i/k}$ is in degree $M_i$.  Hence
$(M_i/k)w(\Delta_s)\ge\ord_E(F)=M_i A_X(w)$.
The same inequality for $s_b$, together with lc of the chosen complement,
gives equality in degree $b$.  These statements also apply to local
sections: coefficients in the local base ring have nonnegative $w$-order,
and shrinking the base only inverts functions of order zero whenever
the centre still contains the specified point.

Write $W$ for the image centre of $w$ on $Z$.  It is a proper closed
subset: $w(\Delta_{s_b})>0$ and a divisorial centre cannot dominate the
birational base.  If $q=\dim W>0$, choose a general linear projection
$Z\to\mathbf A^q$ that is dominant and generically finite on $W$.
Pass to the geometric generic fibre and follow a component of $E$.
The image of that component in the new base is a closed point, and the
new dimension is $\dim X-q\ge1$.

Here are the required preservation details.  The prime $E$ is horizontal
over $\mathbf A^q$.  After shrinking its base, the chosen resolution and
its finitely many relevant strata are smooth wherever they dominate this
base.  Dividing a top form by a base volume form therefore restricts
compatible canonical divisors to the generic fibre, preserving the
log discrepancies and section-divisor orders.  Extending the generic
constant field separably does not ramify these prime orders: the
nonzero discriminant of a finite separable constant extension is a unit
in their valuation rings.  Geometric integrality can be obtained by
choosing the corresponding components after this extension.

Divisorial sheaves restrict to the corresponding divisorial sheaves by
localization and separable scalar extension, as is seen on the smooth
locus and then in codimension one.  Flat base change for the proper morphism
shows that the original generators generate the new complete systems;
only the finitely many degrees dividing $M_i$ and the degree $b$ are
needed here.  Their inequalities hence hold for every local section at
the new image point.  Normality, birationality, the discrepancy bound,
relative ampleness of $-K_X$, and a Fano type witness are preserved by
this operation.  Thus all the conclusions survive the slice.  The
algebraically closed generic field of the complex base is isomorphic
to $\C$; alternatively the finite data can again be embedded into $\C$.
We retain the notation $X,Z,z,w$ for these new data.
\end{proof}

After a subsequence the new dimension and which of the following two
constructions applies are fixed.  We keep the original bounded integer
$b$, even if the dimension has decreased.

\subsection{The canonical cover and the signed torsor}

\begin{lemma}[The isomorphic case]\label{lem:H-canonical-cover}
Suppose $f$ is an isomorphism near $z$.  There is a normal integral
cyclic cover $S_0\to X$, with a unique point $o$ over $z$, a
quasi-\'{e}tale cyclic action fixing $o$, and a divisor-free canonical
generator of character one.  It has a regular eigenfunction
$f_b=s_bt^b$ of character $b$ and a divisorial valuation $w_0$ centred
at $o$ such that, with $A_0=A_{S_0}(w_0)>0$,
\begin{equation}\label{eq:H-upstairs-saturation}
 (S_0,\divisor(f_b)/b)\text{ is lc},\qquad
 w_0(f_b)=bA_0,\qquad
 w_0(a)\ge kA_0
\end{equation}
for every regular eigenfunction $a$ of character $k$ when
$k>0$ and $k\mid M_i$.  The variety $S_0$ is klt.
\end{lemma}
\begin{proof}
Let $r$ be the local Cartier index of $K_X$ and shrink so that
$rK_X=\divisor(c)$.  Normalize $X$ in
\[
 t^r=c^{-1},
\]
with the cyclic group acting on $t$ by its standard character.  The cover
is integral of degree $r$.  Indeed a smaller Kummer degree would make
$c^{-1}$ a $p$th power for a prime $p\mid r$, and then
$(r/p)K_X$ would be principal, contrary to minimality.  At a prime of
$X$, $K_X$ is Cartier and the exponent of $c$ is divisible by $r$.
The codimension-one Kummer calculation shows that the cover is
quasi-\'{e}tale.  If $\divisor(\omega_X)=K_X$, the form
$t\omega_X$ has divisor zero upstairs.  Finite crepant pullback proves
that $S_0$ is klt.

A fibre of this finite quotient is a single orbit: invariant functions
separate distinct finite orbits.  If the fibre over $z$ had more than one
point, regular functions separating its points would have a nontrivial
eigencomponent nonzero at one, hence at every, point of that orbit.
After shrinking by its invariant norm, it would be an eigenunit.  Write
it $s t^j$ with $0<j<r$.  The codimension-one calculation gives
$\divisor(s)=jK_X$, contradicting the minimal index.  Thus the fibre is
one fixed point.

For every integer $k$, the regular functions of character $k$ are
precisely $s t^k$ with $s\in\mathcal O_X(-kK_X)$; changing $k$ by
$r$ is accommodated by $t^r=c^{-1}$.  Their divisors are the pullbacks
of $\Delta_s$.  Extend $w$ to a primitive divisorial valuation $w_0$
on the cover.  Its centre is $o$.  If its restriction is $e w$, finite
crepant pullback gives $A_{S_0}(w_0)=eA_X(w)$ and
$w_0(s t^k)=e w(\Delta_s)$.  Equation~\eqref{eq:H-saturation} proves
all the assertions.
\end{proof}

For the second construction, retain the grading variable $t$ as an
indeterminate, and write $K=\C(X)=\C(Z)$.
Choose a rational top form $\omega_X$ with $\divisor(\omega_X)=K_X$.
The positive degrees record the anti-canonical section inequalities.
We also keep negative degrees: after restriction to a finite scalar
group, one character space contains degrees of both signs.  Bounds on
both signs will therefore be needed to control every section of a
fixed character.

\begin{lemma}[The signed anti-canonical torsor]\label{lem:H-signed-torsor}
Suppose $f$ is not an isomorphism near $z$.  The signed algebra
\begin{equation}\label{eq:H-signed-algebra}
 R=\bigoplus_{k\in\Z}
       H^0(X,\mathcal O_X(-kK_X))t^k,\qquad S_0=\Spec R,
\end{equation}
over the affine base is a finitely generated normal domain of dimension
$\dim X+1$. Write $R_k=H^0(X,\mathcal O_X(-kK_X))t^k$ for its
homogeneous part of degree $k$. Its \emph{positive chamber} is the open
\[
 S_0^+=\bigcup_{k>0}\ \bigcup_{0\ne a\in R_k}D_{S_0}(a),
\]
where $D_{S_0}(a)$ is the principal open on which $a$ does not vanish.
This chamber has complement of codimension at least two. The form
\begin{equation}\label{eq:H-torsor-form}
 \Omega_0=t\omega_X\wedge d\log t
\end{equation}
is a divisor-free canonical generator of grading weight one, and every
finite scalar grading subgroup acts quasi-\'{e}talely.

The variety $S_0$ is klt, and
$(S_0,\divisor(s_bt^b)/b)$ is lc.  There is a grading-fixed closed point
$o$ over $z$ and a rational Gauss valuation $w_0$ centred there which
satisfies \eqref{eq:H-upstairs-saturation} for every actual positive
degree $k\mid M_i$.  Moreover there are numbers $\alpha>0$ and
$F_w>\alpha$, depending on the example, for which every nonzero
homogeneous section satisfies
\begin{equation}\label{eq:H-two-slopes}
 w_0(s t^k)\ge
 \begin{cases}
 k\alpha,&k>0,\\
 |k|(F_w-\alpha),&k<0.
 \end{cases}
\end{equation}
\end{lemma}
\begin{proof}
\textit{Finite generation and normality.}
The positive algebra is finitely generated by relative ampleness of
$-K_X$.  For the negative algebra, take a small $\Q$-factorialization
and a $K$-MMP over $Z$.  Fano type gives termination and a semiample
output, since every divisor is relatively big for a birational
contraction \cite{BCHM}.  The integral-degree section spaces are
preserved in these steps.  To check this without assuming a Cartier
index in degree one, on a common model write
$p^*K_{X_1}=q^*K_{X_2}+E$ for a $K$-MMP step, where $E$ is effective
and $q$-exceptional.  An effective section divisor
$\divisor(s)+mK_{X_1}$ pushes to an effective section divisor on $X_2$.
Conversely, the pullback of an effective section divisor on $X_2$,
plus $mE$, is effective, and pushes to its section divisor on $X_1$.
These are rational pullbacks, valid for every integral $m>0$ because
the divisors involved are $\Q$-Cartier; one can equivalently verify
the inequalities after taking a Cartier power.  The initial small
modification identifies the divisorial sheaves in codimension one.
On the semiample output $Y$, let $g:Y\to Y^{\mathrm c}$ be the
semiample contraction and choose $rK_Y=g^*L$ with $L$ relatively ample.
For each $0\le j<r$, the projection formula identifies the residue
module with
\[
 \bigoplus_{a\ge0}
 H^0\bigl(Y^{\mathrm c},g_*\mathcal O_Y(jK_Y)\otimes
                           \mathcal O_{Y^{\mathrm c}}(aL)\bigr).
\]
The pushforward is coherent by properness, so this is finite over the
ample section ring.  The Cartier Veronese is finitely generated and
the finitely many residue modules are finite over it.  The negative
algebra is therefore finitely generated.  The union of generators of
the two half-algebras generates $R$.

The ring $R$ is the intersection of $K[t,t^{-1}]$ with the Gauss
valuation rings specified, for every prime $D$ of $X$, by
\[
 \ord_D(s)-k\operatorname{coeff}_D(K_X)\ge0
 \quad\text{on }s t^k.
\]
This follows from the codimension-one description of a divisorial
sheaf, and proves normality.  Its fraction field is $K(t)$: the
fraction field of degree zero is $K$, and the generic fibre contains
every integral degree.  In particular $\dim S_0=\dim X+1$.

\medskip
\noindent\textit{The positive chamber and the canonical form.}
We prove that every prime divisor meets the positive chamber, where
the canonical form can be computed on a Laurent bundle.  Suppose
that a prime divisor $P$ lies outside it.  The connected grading torus
preserves every irreducible component of this complement, so $P$ is
homogeneous.  It cannot contain all nonzero graded pieces.  Indeed the
locus where they all vanish is a closed subscheme of $Z$ missing its
generic point, where the signed algebra is $K[t,t^{-1}]$.  That locus
has dimension at most $\dim Z-1=\dim S_0-2$.

Choose a negative homogeneous element $q$ of degree $-m$ not in $P$.
Set
\[
 A=R[q^{-1}]=R_{\le0}[q^{-1}],\qquad C=A_0.
\]
The equality follows by multiplying any positive-degree element by
sufficiently many powers of $q$.  The ring
$C[q,q^{-1}]$ is a Laurent polynomial ring: its nonzero powers have
different degrees.  For a homogeneous $a\in A_k$ the identity
\begin{equation}\label{eq:H-laurent-integrality}
 a^m q^k\in C,\qquad
 a^m=(a^m q^k)q^{-k}
\end{equation}
proves integrality over this Laurent ring.  Both formulas are valid for
negative $k$, because $q$ is a unit on this chart.  Finite generation
then makes $A$ finite over $C[q,q^{-1}]$.
The ring $C$ contains the affine base ring $R_0$ and lies in $K$, so
$\Frac C=K$.

The height-one homogeneous prime induced by $P$ contracts to a
height-one homogeneous prime of $C[q,q^{-1}]$.  Since $q$ is invertible,
that contracted prime is generated by its intersection with $C$ and
lies over a prime divisor of $\Spec C$, a chart of the negative Proj.
This also excludes trivial restriction to $K$: there is no homogeneous
prime divisor solely in the invertible Laurent direction.  The negative
Proj is the relatively ample $K$-model obtained by the small modification,
the $K$-MMP and its semiample contraction.  None of these operations
extracts divisors.  The prime in question therefore has a strict
transform $D$ on $X$; its restriction has the form
$w'=e\ord_D$, $e>0$.

There are sections of both sufficiently divisible signs whose
$\Delta$-orders at $D$ are zero.  For the positive sign this is relative
generation of a multiple of $-K_X$.  For the negative sign use a
generated multiple of $K$ on its ample model avoiding the surviving
prime.  The birational map identifies the two generic DVRs and their
compatible canonical coefficients.  Its section lifts through the
preserved section space, and the effective MMP correction has zero
coefficient at this strict divisor.  Thus its $\Delta$-order on $X$ is
also zero.  Relative generation supplies both sections globally over
the affine base with the indicated degree signs.

Let $v=\ord_P$, and put
$\beta=v(t)+w'(K_X)$, where $w'(K_X)$ means rational pullback order.
For a homogeneous element,
\begin{equation}\label{eq:H-gauss-order}
 v(s t^k)=w'(\Delta_s)+k\beta.
\end{equation}
Nonnegativity of $v$ on the two zero-multiplicity sections forces
$\beta\ge0$ and $\beta\le0$.  Hence $\beta=0$, and the positive
section just chosen has value zero, contrary to $P$ being outside the
positive chamber.  This proves the codimension assertion.

On the positive chamber the signed algebra is the usual Laurent
divisorial bundle over $X$.  To see this on a Proj chart, localize at a
divisible positive section.  Poles along its complementary Cartier
divisor can be cleared by its powers, identifying every degree piece
with the corresponding localized divisorial sheaf on $X$.  A Cartier
Veronese is a line bundle with its zero section removed, and the full
chart is finite over that bundle by taking powers.  In particular the
inverse image of a codimension-at-least-two subset of $X$ has codimension
at least two.  At a regular codimension-one point write
$K_X=\divisor(c)$; then $ct$ is an invertible Laurent coordinate and
\[
 t\omega_X\wedge d\log t
   =(ct)(\omega_X/c)\wedge d\log(ct).
\]
This is a nowhere-vanishing top form, and every nonidentity finite
scalar acts freely there.  The established codimension-two complements
extend these conclusions to the assertions about $\Omega_0$ and
quasi-\'{e}taleness on $S_0$.

\medskip
\noindent\textit{Discrepancies and the lc boundary.}
The form also controls valuations whose centres lie outside the positive
chamber.  We will use its discrepancy estimate both for the boundary
defined by $s_bt^b$ and for the finite quotients in
Lemma~\ref{lem:H-padding}.  If a divisorial multiple $v$ centred on $S_0$ has
nontrivial restriction $w'$ to $K$, the relative Abhyankar inequality
makes $w'$ a divisorial multiple: its value group is discrete, and
adjoining $t$ increases residue transcendence degree by at most one,
forcing its residue field to have transcendence degree $\dim X-1$ over
$\C$.
Its centre exists on $X$ by properness.
Relative generation of a positive Cartier multiple gives
$\beta=v(t)+w'(K_X)\ge0$.  The top form gives
\begin{equation}\label{eq:H-torsor-discrepancy}
 A_{S_0}(v)\ge A_X(w')+\beta,
\end{equation}
with equality for a Gauss extension.  One can compute this on a smooth
model at the prime defining $w'$.  Choose a uniformizer $x$ and
separating residue coordinates $y_1,\ldots,y_{\dim X-1}$.  The ratio of
$\omega_X$ to $d\log x\wedge dy_1\wedge\cdots\wedge dy_{\dim X-1}$
has $w'$-order $A_X(w')+w'(K_X)$.  After wedging with $d\log t$, the remaining
logarithmic differential has nonnegative log order at $v$.  Indeed, at
an upstairs uniformizer $\rho$, every logarithmic differential has at
most a simple $d\rho/\rho$ pole, at most one such factor survives a
wedge, and the differentials of the residue-coordinate units are
regular.  At a Gauss
valuation its residues are independent and this log order is zero.
Multiplication by $t$ yields \eqref{eq:H-torsor-discrepancy}.

If $w'$ is trivial, work on the smooth generic fibre
$\Spec K[t,t^{-1}]$; a centred valuation there has $v(t)=0$ and positive
discrepancy.  Thus $S_0$ is klt.  For a nontrivial restriction,
\eqref{eq:H-gauss-order}, \eqref{eq:H-torsor-discrepancy}, and lc of the
actual degree-$b$ complement give
\[
 A_{S_0}(v)-\frac{v(s_bt^b)}b
 \ge A_X(w')-\frac{w'(\Delta_{s_b})}b\ge0.
\]
For trivial restriction, $v(s_bt^b)=0$.  Consequently the displayed
boundary pair is lc.  It has the lc place constructed next and is not
being asserted to be klt.

\medskip
\noindent\textit{A centred order with two slopes.}
It remains to construct the closed point and its saturated valuation.
The positive slope will come from a choice of Gauss weight; the negative
slope comes from vanishing on the fibre.
Every irreducible component of the reduced fibre over $z$ has positive
dimension.  An isolated zero-dimensional component would disconnect
the fibre unless the whole fibre were finite; a proper birational
morphism to a normal base with a finite fibre is an isomorphism after
shrinking, the excluded case.  A negative-degree section, raised to a
Cartier power, restricts on each such projective component to a section
of an antiample line bundle.  It is therefore zero on every component.
Its effective divisor contains the entire reduced fibre, so its order
at the specified $w$ is strictly positive.  This argument only concerns
reduced supports and also covers nonreduced fibres.

Choose finitely many negative-ring generators $s_jt^{-m_j}$ over the
degree-zero ring and set
\[
 F_w=\min_j\frac{w(\Delta_{s_j})}{m_j}>0.
\]
Every negative homogeneous section is a homogeneous polynomial in these
generators with base coefficients of nonnegative $w$-order.  It follows
that $w(\Delta_s)\ge |k|F_w$ for $k<0$.
Choose a rational $0<\alpha<F_w$, and let $w_0$ be the Gauss extension
of $w$ with $w_0(t)+w(K_X)=\alpha$.  Formula
\eqref{eq:H-gauss-order} defines it, and gives
\eqref{eq:H-two-slopes}; for the positive sign use effectivity of
$\Delta_s$.  All nonzero degrees have strictly positive value, while
the degree-zero maximal ideal is $\mathfrak m_z$.  Thus $w_0$ is centred
at the fixed closed point with ideal generated by these elements.  A
rational Gauss extension of a divisorial valuation is itself divisorial
up to scaling.  Equality in \eqref{eq:H-torsor-discrepancy} gives
$A_0=A_X(w)+\alpha$, and \eqref{eq:H-saturation} now yields
\eqref{eq:H-upstairs-saturation} in actual positive degrees.
\end{proof}

\subsection{The quotient gap and matching the canonical character}

In the canonical-cover case retain its cyclic group.  In the torsor
case choose any finite scalar subgroup $\mu_r$; its order will be fixed
below, after the individual slopes in \eqref{eq:H-two-slopes} are known.
All characters in the next lemma are expressed relative to the standard
character of this cyclic group.

\begin{lemma}[Padding and the actual quotient discrepancy]\label{lem:H-padding}
In either construction, set
\[
 S=S_0\times\mathbf A^{b-1},
 \qquad (o,0)\in S,
\]
and give every added coordinate $u_j$ character one.  Then $S$ is klt,
the finite quotient $Y=S/\mu_r$ is $\epsilon$-lc, and the canonical
generator on $S$ has character $b$.  There is a regular function $h$ of
character $b$ such that $(S,\divisor(h))$ is lc and a centred
quasi-monomial valuation $T_0$ such that
\begin{equation}\label{eq:H-padding-equality}
 A_S(T_0)=T_0(h)=bA_0.
\end{equation}
\end{lemma}
\begin{proof}
The product is klt.  The action remains quasi-\'{e}tale: the fixed locus
of every nonidentity element has codimension at least two already in
$S_0$.  Wedging its canonical generator with
$du_1\wedge\cdots\wedge du_{b-1}$ gives a divisor-free generator of
character $b$.  A suitable power descends through the finite quotient,
so $K_Y$ is $\Q$-Cartier and the quotient map is crepant.

Quasi-\'{e}taleness alone would not preserve a prescribed discrepancy
gap.  Let $E$ be an actual prime divisor over $Y$, and normalize an
extension $v$ to $\C(S)$ so that its restriction to $\C(Y)$ is
$\ord_E$.  Finite crepant pullback, with homogeneous discrepancies,
gives $A_Y(E)=A_S(v)$.  Although $v$ may take fractional values upstairs,
the field $K=\C(X)$ is pointwise invariant and is contained in
$\C(Y)$.  If $w'=v|_K$ is nontrivial, its values are integral.  The
relative Abhyankar inequality then gives
\[
 w'=e\ord_F\quad\text{for a prime }F\text{ over }X
 \quad\text{and an integer }e\ge1.
\]
For the torsor, the log-form calculation above, with the additional
coordinate forms, yields
\[
 A_S(v)\ge A_X(w')+\beta+\sum_jv(u_j)
          \ge e\epsilon\ge\epsilon.
\]
Here $\beta\ge0$ by positive generation, each $u_j$ is regular, and
the non-Gauss logarithmic differential defect is nonnegative.  For the
canonical cover, first restrict to $\C(S_0)$ and use the product
log-form estimate; finite crepant pullback from $S_0$ to $X$ then gives
the same lower bound $A_X(w')$.

If $w'$ is trivial, the valuation lies over the common generic point
of the birational varieties $X$ and $Z$.  The generic quotient is smooth: the canonical cover there is
a cyclic torsor, and the signed construction is a Laurent line with a
free scalar action; their products with the coordinate representation
have smooth quotients.  Generic restriction of the ordinary prime
$E$ is still a primitive prime, since localization at the trivially
valued base field does not change its value group or uniformizer.
Its discrepancy on this smooth generic quotient is at least one, hence
at least $\epsilon$.  This proves the required gap for every actual
quotient prime.

Put $f_b=s_bt^b$ and take a general combination
\begin{equation}\label{eq:H-h-function}
 h=a_0 f_b+a_1u_1^b+\cdots+a_{b-1}u_{b-1}^b.
\end{equation}
We check lc using the ideal $I=(f_b,u_1^b,\ldots,u_{b-1}^b)$.
For any divisorial valuation $v$ with nontrivial restriction $v_0$ to
$\C(S_0)$, that restriction is divisorial up to scaling and the
product differential calculation gives
\[
 A_S(v)\ge A_{S_0}(v_0)+\sum_jv(u_j)
 \ge \frac{v(f_b)}b+\sum_jv(u_j)
 \ge \min\{v(f_b),bv(u_1),\ldots,bv(u_{b-1})\}.
\]
If the restriction is trivial, $v(f_b)=0$ and the same ideal inequality
holds because $S$ is klt.  Therefore $(S,I)$ is lc.  On a resolution of
the ideal, a general residual member is smooth and transverse to the
fixed log smooth boundary; with coefficient one this proves
$(S,\divisor(h))$ lc.  All summands of $h$ have character $b$.

Finally extend $w_0$ by the joint Gauss valuation giving every $u_j$
weight $A_0$.  It is centred at $(o,0)$ and quasi-monomial, and its
product discrepancy is $A_0+(b-1)A_0=bA_0$.  Every summand of $h$ has
that value.  Their initials do not cancel for a general choice of the
$a_j$ (the added coordinates are independent), giving
\eqref{eq:H-padding-equality}.  When $b=1$ there are no added
coordinates and $h$ is simply a nonzero scalar multiple of $f_b$.
\end{proof}

\begin{proposition}[The pinning input]\label{prop:H-g-input}
The failure sequence produces, after a subsequence, a fixed-dimensional
sequence satisfying every hypothesis of Theorem~\ref{thm:G-pinning}.
\end{proposition}
\begin{proof}
Let $\chi$ be character $b$ of the chosen cyclic group; it need not be
faithful.  Normalize $T=T_0/(bA_0)$.  The preceding lemma gives
\[
 A_S(T)=T(h)=1,
 \qquad (S,\divisor(h))\text{ lc},
 \qquad S/\mu_r\text{ is }\epsilon\text{-lc},
\]
and $h$ and the divisor-free canonical generator both have character
$\chi$.  It remains to check, for each fixed positive integer $m$, that
eventually
\begin{equation}\label{eq:H-final-saturation}
 T(a)\ge m
 \quad\text{for every regular eigenfunction }a\text{ of character }\chi^m.
\end{equation}

Expand a polynomial eigenfunction in monomials in the added coordinates.
A monomial of total degree $j$ has $T_0$-value $jA_0$.  If $j\ge bm$,
this already suffices, since its coefficient is regular and has
nonnegative value.  Otherwise its coefficient on $S_0$ has character
\[
 k=bm-j>0.
\]
In the canonical-cover case, that entire eigenspace is the degree-$k$
section space described in Lemma~\ref{lem:H-canonical-cover}.  Since
every fixed $k$ divides $M_i$ eventually, its coefficients have value
at least $kA_0$.  The monomial therefore has value at least
$(j+k)A_0=bmA_0$.

In the signed torsor, coefficients decompose into actual degrees
$k+\ell r$, $\ell\in\Z$.  For $\ell=0$ use the same factorial
saturation.  For $\ell\ne0$ use both strict slopes in
\eqref{eq:H-two-slopes}.  At example $i$, choose $r_i>i$ so large that
\begin{equation}\label{eq:H-cyclic-order}
 (r_i-i)\min\{\alpha,F_w-\alpha\}\ge iA_0.
\end{equation}
For $1\le k\le i$, every nonzero residue shift has
$|k+\ell r_i|\ge r_i-i$, so its value is at least $iA_0\ge kA_0$.
This also gives the desired monomial bound.  The choice is made after
$\alpha,F_w,A_0$ are known; no uniform lower bound on either slope is
needed.  For fixed $m$, eventually $bm\le i$, and every required
unshifted degree divides $M_i$.

A finite sum cannot have value smaller than the minimum of its terms.
The argument therefore proves \eqref{eq:H-final-saturation} in the
affine coordinate ring.  It also proves it in the local ring at the
fixed point.  Indeed any local function can be written with an invariant
unit denominator: multiply an initial denominator by its other group
translates and use its invariant norm.  For an eigenfunction the resulting
numerator has the same character.  The invariant denominator has value
zero, so the inequality is unchanged.

The dimensions satisfy
\[
 \dim S=
 \begin{cases}
 \dim X+b-1,&\text{canonical-cover case},\\
 \dim X+b,&\text{signed-torsor case},
 \end{cases}
 \qquad\text{and hence}\qquad \dim S\le d+b.
\]
They are thus bounded independently of $i$, and a subsequence has fixed
dimension.  The valuations are centred and quasi-monomial, the covers
are finite cyclic and quasi-\'{e}tale, and every remaining requirement
of Theorem~\ref{thm:G-pinning} was checked above.
\end{proof}

\subsection{Chen's reduction and change of field}

We use the following form of Chen's reduction.
For every positive integer $\ell$, Chen's Theorem~1.6, assertions (2) and (3),
identifies the following two assertions \cite{Chen2025}:
\begin{enumerate}[label=(\roman*)]
\item bounded monotonic klt complements for birational Fano type
contractions with zero boundary and $\dim Z\le\ell$;
\item the Cartier-section assertion for Fano type contractions with
$\dim Z\le\ell$, with a constant depending on the total-space dimension
and the discrepancy gap.
\end{enumerate}
In (i), the total-space and base dimensions are equal.  Consequently,
Theorem~\ref{thm:birational} in each of the finitely many positive
dimensions at most $\ell$, for every positive discrepancy parameter,
supplies (i); its dimension-zero case is
trivial.  One may take a common multiple of the finitely many complement
indices for each fixed discrepancy parameter.  Thus the birational
inputs have dimension at most $\ell$, whereas assertion (ii) also allows
total spaces of higher dimension.

The real-pair conventions in Chen's statement cause no extra hypothesis
in (i).  If an integral Weil divisor $K_X$ is $\R$-Cartier, express it
locally as a finite real combination of Cartier divisors.  Comparing Weil
coefficients gives a finite linear system with rational coefficients and
rational right-hand side.  A real solution implies a rational solution,
so $K_X$ is $\Q$-Cartier.  Similarly, a real effective log Fano witness
can be replaced by a rational one.  Keep its finite boundary support
and its zero coefficients.  On a finite open cover, choose Cartier
representations of its log canonical class, using the same boundary
coefficient variables on every chart.  These impose finitely many
rational affine equations.  Rational points are dense in their real
solution space.  Every rational solution defines a global rational
divisor whose log canonical divisor is locally $\Q$-Cartier; clearing
denominators on the finite cover makes it $\Q$-Cartier on $X$.  A rational
affine basis of the solution space therefore places all its log canonical
divisors in a fixed finite-dimensional real span of global $\Q$-Cartier
divisors.  Relative ampleness is open in that span.  On a fixed log
resolution, the klt condition and positivity of the remaining boundary
coefficients also persist under a sufficiently small perturbation.
Hence the rational birational
statement proves the full zero-boundary assertion needed in Chen's
real-pair formulation.

\begin{lemma}[Point ideals and change of field]\label{lem:H-field-reduction}
It suffices to prove Theorem~\ref{thm:main} over $\C$.  More precisely, a
uniform positive lower bound for the log canonical threshold of
$\mathfrak m_z\mathcal O_X$ implies the Cartier-section assertion after
replacing the bound by a smaller number less than one.  Both the hypotheses
and this ideal threshold are preserved by algebraically closed field
extensions of characteristic zero.
\end{lemma}
\begin{proof}
Shrink $Z$ to an affine neighbourhood of $z$, choose generators
$g_1,\ldots,g_N$ of $\mathfrak m_z$, and resolve both $(X,B)$ and the
pulled-back ideal.  If $p:Y\to X$ is the resulting log resolution, write
\[
 K_Y+B_Y=p^*(K_X+B),\qquad
 \mathfrak m_z\mathcal O_Y=\mathcal O_Y(-F).
\]
The ideal threshold is the minimum of
$A_{X,B}(E)/\ord_E(F)$ over the finitely many components of $F$ having
positive order.  Its interpretation is lc on the inverse image of a
neighbourhood of $z$.  Indeed the ideal is the unit ideal away from
$f^{-1}(z)$, and the original pair is already lc there.

For a general linear combination $g=\sum a_jg_j$, the pullback of
$\divisor(g)$ is $F+H$, where $H$ is a general member of a basepoint-free
system.  On a further restriction of the base, $H$ is smooth and transverse
to the fixed simple normal crossings support.  If $0<t<1$ is at most the
ideal threshold, then $(Y,B_Y+tF+tH)$ is lc, hence so is
$(X,B+t f^*\divisor(g))$.  Choose $g\ne0$.  Since $g\in\mathfrak m_z$,
$D=\divisor(g)$ is a nonzero effective Cartier divisor containing $z$.
Conversely, a Cartier divisor through $z$ has a local equation
$g\in\mathfrak m_z$, so that $v(\mathfrak m_z\mathcal O_X)\le v(f^*g)$
for every valuation.  Thus an admissible coefficient for that divisor is
also a lower bound for the ideal threshold.

All the varieties, morphisms, rational boundaries, a rational Fano type
witness, generators, and the chosen resolutions descend to an
algebraically closed subfield $k_0$ of finite transcendence degree over
$\overline{\Q}$.  Such a field embeds in $\C$.  Normality, geometric
integrality and the identity $f_*\mathcal O_X=\mathcal O_Z$ are unchanged
by algebraically closed extension; the latter also follows from flat base
change for the proper morphism.  Projectivity and relative ampleness are
geometric properties.  Relative nefness is geometric as well: any curve
on an extension can be spread over a finite-type parameter scheme and
specialized to a curve cycle over the smaller algebraically closed field,
with the same intersection number.  Thus a negative intersection after
extension would already yield a negative intersection before extension.

For discrepancies one can use the same log resolution.  In the
nonvacuous range $0<\epsilon\le1$, a log smooth sub-boundary with all
coefficients at most $1-\epsilon$ is $\epsilon$-lc.  At a stratum this
follows by comparing with its reduced simple normal crossings support:
a primitive monomial order has at least one positive integral weight,
and each boundary weight contributes at least $\epsilon$; a
nonmonomial prime has a further nonnegative discrepancy, with an integral
contribution at least one when needed away from the strata.  This also
shows preservation of the discrepancy bound after extension.  The
coefficients and orders in the displayed resolution formula are unchanged,
so the ideal threshold is unchanged.  The complex assertion therefore
gives the same lower bound over $k_0$ and over the original field, where
the preceding general-member construction finishes the proof.
\end{proof}

\subsection{Conclusion of the two theorems}

\begin{proof}[Proof of Theorem~\ref{thm:birational}]
Over $\C$, a failure in any fixed dimension and with any fixed
$0<\epsilon\le1$ would, by Proposition~\ref{prop:H-g-input}, give the
sequence prohibited by Theorem~\ref{thm:G-pinning}.  Therefore a uniform
positive integral complement index exists.  The range $\epsilon>1$
in positive dimension is vacuous, since a prime divisor on the smooth
locus has log discrepancy one.

The birational conclusion also descends to every algebraically closed
field of characteristic zero.  Increase the complex uniform index to a
multiple greater than one, which preserves the complement assertion.
Descend a given example, its Fano type
witness, and a resolution to an algebraically closed field embeddable
in $\C$.  For the complex uniform index $n$, the complete relative
system $f_*\mathcal O_X(-nK_X)$ commutes with these field extensions.
Resolve its finitely generated fractional base ideal as in
Proposition~\ref{prop:H-saturation}.  A klt member over $\C$ implies
that every fixed-order discrepancy relevant to the fibre is strictly
positive.  Those discrepancies are unchanged by extension.  A general
member over the original infinite field consequently gives a klt
$n$-complement after shrinking, since the general residual coefficient
$1/n$ is less than one.  Its
section divisor $\Delta_s$ defines
$\Delta=\Delta_s/n\ge0$, and
$n(K_X+\Delta)=\divisor(s)$ is Cartier and linearly trivial on the
chosen neighbourhood, as required.
\end{proof}

\begin{proof}[Proof of Theorem~\ref{thm:main}]
First work over $\C$ and fix the total-space dimension $d$ and gap
$\epsilon>0$.  Theorem~\ref{thm:birational} in every dimension
$1,\ldots,d$ and for every positive discrepancy parameter, together
with its trivial dimension-zero case and the
rationalization above, verifies Chen's assertion (2) for base dimension
bound $\ell=d$.  His Theorem~1.6 therefore gives assertion (3), with a
positive constant depending only on $d$ and $\epsilon$
\cite{Chen2025}.  Every base of a contraction from a $d$-dimensional
variety has dimension at most $d$, so this applies to the given $f$.
It supplies lc on the full inverse image of a neighbourhood of $z$.
Lemma~\ref{lem:H-field-reduction}, decreasing the coefficient below one
if necessary, transports the uniform constant to an arbitrary
algebraically closed characteristic-zero field and produces a nonzero
effective Cartier divisor through $z$ there.  This proves the stated
neighbourhood assertion.
\end{proof}

\begin{proof}[Proof of Corollary~\ref{cor:real-boundaries}]
The same application of Chen's implication (2)$\Rightarrow$(3) over
$\C$ gives his full real-boundary assertion.  Only the zero-boundary
birational inputs to (2) were rationalized above.  The resulting
assertion (3) already allows arbitrary effective real $B$ with
$K_X+B$ real Cartier, and supplies the same dependence on $d$ and
$\epsilon$.  No perturbation of the nef real boundary is involved.
\end{proof}

\begingroup
\small
\bibliographystyle{plainnat}
\bibliography{references}
\endgroup
\end{document}
